% Standalone mathematical supplement. Compile twice with pdfLaTeX.
\documentclass[11pt,a4paper]{article}
\usepackage[T1]{fontenc}
\usepackage[utf8]{inputenc}
\usepackage[english]{babel}
\usepackage{lmodern}
\usepackage[margin=20mm]{geometry}
\usepackage{microtype}
\usepackage{amsmath,amssymb,mathtools}
\usepackage{graphicx,float,longtable,array,booktabs}
\usepackage[numbers,sort&compress]{natbib}
\usepackage{xurl}
\usepackage[hidelinks,unicode,pdfencoding=auto]{hyperref}
\usepackage{bookmark}
\hypersetup{pdftitle={Mathematical supplement to A minimal local covariant action for holographic dark energy: constraints, perturbations, and nonlinear dynamics},pdfauthor={Danylo Yerokhin},
  pdfsubject={Mathematical catalogue},
  pdfkeywords={holographic dark energy, covariant action, constraints, perturbations, strong coupling}}
\urlstyle{same}
\setcounter{tocdepth}{2}
\setcounter{secnumdepth}{3}
\allowdisplaybreaks[2]
\setlength{\emergencystretch}{2em}
\title{Mathematical supplement to\\
A minimal local covariant action for holographic dark energy:\\
constraints, perturbations, and nonlinear dynamics}
\author{Danylo Yerokhin\thanks{Email: \href{mailto:denyerokhin@gmail.com}{denyerokhin@gmail.com}.
ORCID: \href{https://orcid.org/0009-0003-8645-327X}{0009-0003-8645-327X}.}\\[0.5em]
\normalsize V.\,N. Karazin Kharkiv National University\\
\normalsize 4 Svobody Sq., Kharkiv 61022, Ukraine}
\date{}
\begin{document}
\maketitle
\tableofcontents
\clearpage

\section{Mathematical catalogue}
\label{app:mathematical-catalogue}
\subsection{Scope and indexing}
\label{sec:cat-scope}
This standalone mathematical supplement accompanies the main article~\cite{YerokhinArticle}. References explicitly identified as belonging to the main article, including Appendices~A--E, use its section and equation numbers. All other section and equation numbers refer to this supplement.

Subsections~\ref{sec:cat-contacts}--\ref{sec:cat-angular} and~\ref{sec:cat-rationals}, together with the kinematics and normalizations given below, apply to the radiation-like branch $ U=0 $ in the absence of an independent matter sector. The independent parameters used are the ratio of the magnitudes of the comoving momenta $ r=k'/k>0 $, the cosine of the angle between them $ -1<z<1 $, and the variable $ x=k\eta $. Subsection~\ref{sec:cat-local-vertices} contains a separate local result for a general potential $ U\in C^6 $ at $ \chi<0 $ and $ \Delta_\Phi\ne0 $. The condition $ U=0 $ and the kinematics of the radiation-like branch do not apply to that subsection. The catalogue does not give the general gravitational amplitude on an arbitrary background. The tabulated coefficients do not include the external frequency factors, the overall scale, or the separate volume normalization of exactly homogeneous lines. These factors are retained in the formulas of the main article.

The position indices $0,1,2,3$ coincide with those in Subsection~C.3 of the main article. In its main text, the description of the channel $13|24$ preceding Eq.~(195) labels the same positions $1,2,3,4$. The field indices $0,1,2$ denote the scalar and the first and second linear tensor polarizations. The inward-directed momenta in units of $k$ are
\begin{equation}
\shortstack[l]{$\displaystyle \mathbf p_0=(0,0,1),\quad \mathbf p_1=(0,0,-1),$\\
$\displaystyle \mathbf p_2=(-r\sqrt{1-z^2},0,-rz),\quad \mathbf p_3=(r\sqrt{1-z^2},0,rz).$}
\label{eq:cat-momenta}
\end{equation}
The signed frequencies are $(c_{A_0},c_{A_1},-rc_{A_2},-rc_{A_3})$, where $c_0=1/\sqrt3$, $c_1=c_2=1$. For each external momentum $\mathbf p$, set $\mathbf n=\mathbf p/|\mathbf p|$. We use the vectors $\mathbf e=(0,1,0)$, $\mathbf b=(-n_z,0,n_x)$ and the tensors $E_1=\mathbf e\mathbf e-\mathbf b\mathbf b$, $E_2=\mathbf e\mathbf b+\mathbf b\mathbf e$. Both tensors have squared norm two. The parameter $t_\theta$ of the original algebraic calculation is related to $z$ by $t_\theta^2=(1-z)/(1+z)$.

For an internal momentum $\mathbf p$ of nonzero norm, we use the unnormalized vector $\mathbf b=(-p_z,0,p_x)$: $E_1=\mathbf e\mathbf e-\mathbf b\mathbf b/p^2$, $E_2=\mathbf e\mathbf b+\mathbf b\mathbf e$. The squared norms of the scalar and the two tensors are $1,2,2p^2$. Reversing $\mathbf p$ to $-\mathbf p$ between the ends of the internal line gives the factors $\tau_I=(1,1,-1)$.

At $\mathbf p=0$, we use five symmetric trace-free matrices. The two diagonal matrices are $\mathrm{diag}(1,-1,0)$ and $\mathrm{diag}(1,1,-2)$. The other three have unit entries at the pairs of positions $xy,yx$, $xz,zx$, and $yz,zy$, respectively. Their squared norms are $2,6,2,2,2$. The kinetic sign of the homogeneous scalar coordinate is $-1$. The covariances of the geometric coordinates are divided by the positive squared geometric norm. The kinetic sign of the scalar coordinate is accounted for in the free equation of motion without introducing a negative state probability.

\subsection{Thirteen representative contact components and reconstruction of 81 components}
\label{sec:cat-contacts}
For each contact component, four sets of coefficients of powers of $x$ are given below: $h_p$ for the full Hamiltonian contact term; $b_p$ for the original Lagrangian contact term taken with the opposite sign; $u_p$ for the sum of contact terms from eliminating auxiliary variables; and $v_p$ for the sum of Legendre-transform contact terms. In all cases,
\begin{equation}
\frac{H_4}{g_{\mathrm R}^2k^4}=\sum_{p=-4}^{0}h_px^p,\qquad h_p=b_p+u_p+v_p.
\label{eq:cat-contact-parts}
\end{equation}
The individual partitions needed for the angular reconstruction are specified in Subsection~\ref{sec:cat-contact-partitions}. The sums $u_p,v_p$ alone do not determine them. Coefficients of powers not listed are zero. The symbols $R_s$ in the following formulas denote the explicit algebraic functions in Subsection~\ref{sec:cat-rationals}, not external electronic records. When expressed in terms of $z$, some sources for the global Cartesian tensors contain $\sqrt{1-z^2}$; after changing to $t_\theta$, all functions are rational.

Let $T$ exchange the incoming and outgoing pairs, and let $S$ exchange the two positions within one pair. The exact rules are
\begin{equation}
 h_p^{B A C D}(r,z)=h_p^{A B D C}(r,z)=h_p^{A B C D}(r,-z),\qquad
 h_p^{C D A B}(r,z)=r^{4+p}\overline{h_p^{A B C D}(1/r,z)}.
\label{eq:cat-contact-symmetries}
\end{equation}
The indices within each pair are ordered first, and the pairs are then ordered lexicographically. An odd number of within-pair exchanges gives $z\mapsto-z$. Exchanging the pairs corresponds to the substitution $r\mapsto1/r$, complex conjugation, and multiplication by $r^{4+p}$. These identities have been checked separately for each of the 81 components. The same identities also hold separately for $ b_p,u_p,v_p $ defined in \eqref{eq:cat-contact-parts}. The starting point is the total Lagrangian density $ \mathcal L_\psi+\mathcal L_g $ on the radiation-like background, given in Eqs. (342)--(343) of the main article. Substituting the solutions of the linear constraints into its quadrilinear part gives $ b_p $ with the opposite sign. Eliminating the auxiliary variables at second order in the perturbations gives $ u_p $, and the Legendre transform in Eq. (191) of the main article gives $ v_p $. The corresponding decomposition of the full contact term is given in Eq. (351) of the main article.
% Main article: eq:v34r-exact-scalar-density; eq:v34r-exact-gravity-density; eq:v33-legendre-contact; eq:radiation-contact-decomposition.

In the quadrilinear coefficients of these parts of the contact term, all external-wave data are permuted: field types, momenta, signed frequencies, and polarizations. For $ u_p $ and $ v_p $, each partition is first transformed according to \eqref{eq:cat-partition-permutations}, including the replacement $ +\leftrightarrow- $ when the two positions within one external pair are exchanged. The three partitions $ 0,+,- $ are then summed according to \eqref{eq:cat-partition-contact-sum}. The resulting sums satisfy \eqref{eq:cat-contact-symmetries}. For a fixed partition label, the rules \eqref{eq:cat-partition-permutations} must be used. The original part $ b_p $ requires no such summation.

Components containing an odd number of fields with the second linear tensor polarization vanish identically by reflection symmetry under $y\mapsto-y$. There are 40 such components; the remaining 41 are reconstructed from 13 representatives.
\subsubsection{Component 0000}
\begin{equation}
\shortstack[l]{$\displaystyle h^{0000}_{-4}=R_{1}$\\
$\displaystyle h^{0000}_{-3}=R_{2}$\\
$\displaystyle h^{0000}_{-2}=R_{3}$\\
$\displaystyle h^{0000}_{-1}=R_{4}$\\
$\displaystyle h^{0000}_{0}=R_{5}$}
\label{eq:cat-map-2}
\end{equation}
\begin{equation}
\shortstack[l]{$\displaystyle b^{0000}_{-4}=R_{6}$\\
$\displaystyle b^{0000}_{-3}=R_{7}$\\
$\displaystyle b^{0000}_{-2}=R_{8}$\\
$\displaystyle b^{0000}_{-1}=R_{9}$\\
$\displaystyle b^{0000}_{0}=R_{10}$}
\label{eq:cat-map-3}
\end{equation}
\begin{equation}
\shortstack[l]{$\displaystyle u^{0000}_{-4}=R_{11}$\\
$\displaystyle u^{0000}_{-3}=R_{12}$\\
$\displaystyle u^{0000}_{-2}=R_{13}$\\
$\displaystyle u^{0000}_{-1}=R_{14}$\\
$\displaystyle u^{0000}_{0}=R_{15}$}
\label{eq:cat-map-4}
\end{equation}
\begin{equation}
\shortstack[l]{$\displaystyle v^{0000}_{-4}=R_{16}$\\
$\displaystyle v^{0000}_{-3}=R_{17}$\\
$\displaystyle v^{0000}_{-2}=R_{18}$\\
$\displaystyle v^{0000}_{-1}=R_{19}$\\
$\displaystyle v^{0000}_{0}=R_{20}$}
\label{eq:cat-map-5}
\end{equation}
\subsubsection{Component 0001}
\begin{equation}
\shortstack[l]{$\displaystyle h^{0001}_{-4}=R_{21}$\\
$\displaystyle h^{0001}_{-3}=R_{22}$\\
$\displaystyle h^{0001}_{-2}=R_{23}$\\
$\displaystyle h^{0001}_{-1}=R_{24}$}
\label{eq:cat-map-6}
\end{equation}
\begin{equation}
\shortstack[l]{$\displaystyle b^{0001}_{-4}=R_{25}$\\
$\displaystyle b^{0001}_{-3}=R_{26}$\\
$\displaystyle b^{0001}_{-2}=R_{27}$\\
$\displaystyle b^{0001}_{-1}=R_{28}$}
\label{eq:cat-map-7}
\end{equation}
\begin{equation}
\shortstack[l]{$\displaystyle u^{0001}_{-4}=R_{29}$\\
$\displaystyle u^{0001}_{-3}=R_{30}$\\
$\displaystyle u^{0001}_{-2}=R_{31}$\\
$\displaystyle u^{0001}_{-1}=R_{32}$}
\label{eq:cat-map-8}
\end{equation}
\begin{equation}
\shortstack[l]{$\displaystyle v^{0001}_{-4}=R_{33}$\\
$\displaystyle v^{0001}_{-3}=R_{34}$\\
$\displaystyle v^{0001}_{-2}=R_{35}$\\
$\displaystyle v^{0001}_{-1}=R_{36}$}
\label{eq:cat-map-9}
\end{equation}
\subsubsection{Component 0011}
\begin{equation}
\shortstack[l]{$\displaystyle h^{0011}_{-4}=R_{37}$\\
$\displaystyle h^{0011}_{-3}=R_{38}$\\
$\displaystyle h^{0011}_{-2}=R_{39}$\\
$\displaystyle h^{0011}_{-1}=R_{40}$}
\label{eq:cat-map-10}
\end{equation}
\begin{equation}
\shortstack[l]{$\displaystyle b^{0011}_{-4}=R_{41}$\\
$\displaystyle b^{0011}_{-3}=R_{42}$\\
$\displaystyle b^{0011}_{-2}=R_{43}$}
\label{eq:cat-map-11}
\end{equation}
\begin{equation}
\shortstack[l]{$\displaystyle u^{0011}_{-4}=R_{44}$\\
$\displaystyle u^{0011}_{-3}=R_{45}$\\
$\displaystyle u^{0011}_{-2}=R_{46}$\\
$\displaystyle u^{0011}_{-1}=R_{47}$}
\label{eq:cat-map-12}
\end{equation}
\begin{equation}
\shortstack[l]{$\displaystyle v^{0011}_{-4}=R_{48}$\\
$\displaystyle v^{0011}_{-3}=R_{49}$\\
$\displaystyle v^{0011}_{-2}=R_{50}$\\
$\displaystyle v^{0011}_{-1}=R_{51}$}
\label{eq:cat-map-13}
\end{equation}
\subsubsection{Component 0022}
\begin{equation}
\shortstack[l]{$\displaystyle h^{0022}_{-4}=R_{52}$\\
$\displaystyle h^{0022}_{-3}=R_{53}$\\
$\displaystyle h^{0022}_{-2}=R_{54}$}
\label{eq:cat-map-14}
\end{equation}
\begin{equation}
\shortstack[l]{$\displaystyle b^{0022}_{-4}=R_{55}$\\
$\displaystyle b^{0022}_{-3}=R_{56}$\\
$\displaystyle b^{0022}_{-2}=R_{57}$}
\label{eq:cat-map-15}
\end{equation}
\begin{equation}
\shortstack[l]{$\displaystyle u^{0022}_{-4}=R_{58}$\\
$\displaystyle u^{0022}_{-3}=R_{59}$\\
$\displaystyle u^{0022}_{-2}=R_{60}$\\
$\displaystyle u^{0022}_{-1}=R_{51}$}
\label{eq:cat-map-16}
\end{equation}
\begin{equation}
\shortstack[l]{$\displaystyle v^{0022}_{-4}=R_{61}$\\
$\displaystyle v^{0022}_{-3}=R_{62}$\\
$\displaystyle v^{0022}_{-2}=R_{63}$\\
$\displaystyle v^{0022}_{-1}=R_{64}$}
\label{eq:cat-map-17}
\end{equation}
\subsubsection{Component 0101}
\begin{equation}
\shortstack[l]{$\displaystyle h^{0101}_{-4}=R_{65}$\\
$\displaystyle h^{0101}_{-3}=R_{66}$\\
$\displaystyle h^{0101}_{-2}=R_{67}$\\
$\displaystyle h^{0101}_{-1}=R_{68}$}
\label{eq:cat-map-18}
\end{equation}
\begin{equation}
\shortstack[l]{$\displaystyle b^{0101}_{-4}=R_{69}$\\
$\displaystyle b^{0101}_{-3}=R_{70}$\\
$\displaystyle b^{0101}_{-2}=R_{71}$}
\label{eq:cat-map-19}
\end{equation}
\begin{equation}
\shortstack[l]{$\displaystyle u^{0101}_{-4}=R_{72}$\\
$\displaystyle u^{0101}_{-3}=R_{73}$\\
$\displaystyle u^{0101}_{-2}=R_{74}$\\
$\displaystyle u^{0101}_{-1}=R_{75}$}
\label{eq:cat-map-20}
\end{equation}
\begin{equation}
\shortstack[l]{$\displaystyle v^{0101}_{-4}=R_{76}$\\
$\displaystyle v^{0101}_{-3}=R_{77}$\\
$\displaystyle v^{0101}_{-2}=R_{78}$\\
$\displaystyle v^{0101}_{-1}=R_{79}$}
\label{eq:cat-map-21}
\end{equation}
\subsubsection{Component 0111}
\begin{equation}
\shortstack[l]{$\displaystyle h^{0111}_{-4}=R_{80}$\\
$\displaystyle h^{0111}_{-3}=R_{81}$\\
$\displaystyle h^{0111}_{-2}=R_{82}$\\
$\displaystyle h^{0111}_{-1}=R_{83}$}
\label{eq:cat-map-22}
\end{equation}
\begin{equation}
\shortstack[l]{$\displaystyle b^{0111}_{-3}=R_{84}$\\
$\displaystyle b^{0111}_{-2}=R_{85}$}
\label{eq:cat-map-23}
\end{equation}
\begin{equation}
\shortstack[l]{$\displaystyle u^{0111}_{-4}=R_{86}$\\
$\displaystyle u^{0111}_{-3}=R_{87}$\\
$\displaystyle u^{0111}_{-2}=R_{88}$\\
$\displaystyle u^{0111}_{-1}=R_{83}$}
\label{eq:cat-map-24}
\end{equation}
\begin{equation}
\shortstack[l]{$\displaystyle v^{0111}_{-4}=R_{89}$\\
$\displaystyle v^{0111}_{-3}=R_{90}$\\
$\displaystyle v^{0111}_{-2}=R_{91}$}
\label{eq:cat-map-25}
\end{equation}
\subsubsection{Component 0122}
\begin{equation}
\shortstack[l]{$\displaystyle h^{0122}_{-3}=R_{92}$\\
$\displaystyle h^{0122}_{-2}=R_{93}$}
\label{eq:cat-map-26}
\end{equation}
\begin{equation}
\shortstack[l]{$\displaystyle b^{0122}_{-3}=R_{94}$\\
$\displaystyle b^{0122}_{-2}=R_{95}$}
\label{eq:cat-map-27}
\end{equation}
\begin{equation}
\shortstack[l]{$\displaystyle u^{0122}_{-3}=R_{96}$\\
$\displaystyle u^{0122}_{-2}=R_{97}$}
\label{eq:cat-map-28}
\end{equation}
\begin{equation}
v^{0122}=0.
\label{eq:cat-map-29}
\end{equation}
\subsubsection{Component 0202}
\begin{equation}
\shortstack[l]{$\displaystyle h^{0202}_{-4}=R_{98}$\\
$\displaystyle h^{0202}_{-3}=R_{99}$\\
$\displaystyle h^{0202}_{-2}=R_{100}$\\
$\displaystyle h^{0202}_{-1}=R_{101}$}
\label{eq:cat-map-30}
\end{equation}
\begin{equation}
\shortstack[l]{$\displaystyle b^{0202}_{-4}=R_{102}$\\
$\displaystyle b^{0202}_{-3}=R_{103}$\\
$\displaystyle b^{0202}_{-2}=R_{104}$}
\label{eq:cat-map-31}
\end{equation}
\begin{equation}
\shortstack[l]{$\displaystyle u^{0202}_{-4}=R_{105}$\\
$\displaystyle u^{0202}_{-3}=R_{106}$\\
$\displaystyle u^{0202}_{-2}=R_{107}$\\
$\displaystyle u^{0202}_{-1}=R_{108}$}
\label{eq:cat-map-32}
\end{equation}
\begin{equation}
\shortstack[l]{$\displaystyle v^{0202}_{-4}=R_{109}$\\
$\displaystyle v^{0202}_{-3}=R_{110}$\\
$\displaystyle v^{0202}_{-2}=R_{111}$\\
$\displaystyle v^{0202}_{-1}=R_{112}$}
\label{eq:cat-map-33}
\end{equation}
\subsubsection{Component 0212}
\begin{equation}
\shortstack[l]{$\displaystyle h^{0212}_{-4}=R_{113}$\\
$\displaystyle h^{0212}_{-3}=R_{114}$\\
$\displaystyle h^{0212}_{-2}=R_{115}$\\
$\displaystyle h^{0212}_{-1}=R_{116}$}
\label{eq:cat-map-34}
\end{equation}
\begin{equation}
\shortstack[l]{$\displaystyle b^{0212}_{-3}=R_{117}$\\
$\displaystyle b^{0212}_{-2}=R_{118}$}
\label{eq:cat-map-35}
\end{equation}
\begin{equation}
\shortstack[l]{$\displaystyle u^{0212}_{-4}=R_{119}$\\
$\displaystyle u^{0212}_{-3}=R_{120}$\\
$\displaystyle u^{0212}_{-2}=R_{121}$\\
$\displaystyle u^{0212}_{-1}=R_{116}$}
\label{eq:cat-map-36}
\end{equation}
\begin{equation}
\shortstack[l]{$\displaystyle v^{0212}_{-4}=R_{122}$\\
$\displaystyle v^{0212}_{-3}=R_{123}$\\
$\displaystyle v^{0212}_{-2}=R_{124}$}
\label{eq:cat-map-37}
\end{equation}
\subsubsection{Component 1111}
\begin{equation}
\shortstack[l]{$\displaystyle h^{1111}_{-4}=R_{125}$\\
$\displaystyle h^{1111}_{-3}=R_{126}$\\
$\displaystyle h^{1111}_{-2}=R_{127}$\\
$\displaystyle h^{1111}_{-1}=R_{128}$}
\label{eq:cat-map-38}
\end{equation}
\begin{equation}
\shortstack[l]{$\displaystyle b^{1111}_{-2}=R_{129}$}
\label{eq:cat-map-39}
\end{equation}
\begin{equation}
\shortstack[l]{$\displaystyle u^{1111}_{-4}=R_{130}$\\
$\displaystyle u^{1111}_{-3}=R_{131}$\\
$\displaystyle u^{1111}_{-2}=R_{132}$\\
$\displaystyle u^{1111}_{-1}=R_{128}$}
\label{eq:cat-map-40}
\end{equation}
\begin{equation}
\shortstack[l]{$\displaystyle v^{1111}_{-4}=R_{133}$\\
$\displaystyle v^{1111}_{-3}=R_{134}$\\
$\displaystyle v^{1111}_{-2}=R_{135}$}
\label{eq:cat-map-41}
\end{equation}
\subsubsection{Component 1122}
\begin{equation}
\shortstack[l]{$\displaystyle h^{1122}_{-4}=R_{136}$\\
$\displaystyle h^{1122}_{-3}=R_{137}$\\
$\displaystyle h^{1122}_{-2}=R_{138}$}
\label{eq:cat-map-42}
\end{equation}
\begin{equation}
\shortstack[l]{$\displaystyle b^{1122}_{-2}=R_{139}$}
\label{eq:cat-map-43}
\end{equation}
\begin{equation}
\shortstack[l]{$\displaystyle u^{1122}_{-4}=R_{140}$\\
$\displaystyle u^{1122}_{-3}=R_{141}$\\
$\displaystyle u^{1122}_{-2}=R_{142}$}
\label{eq:cat-map-44}
\end{equation}
\begin{equation}
\shortstack[l]{$\displaystyle v^{1122}_{-4}=R_{143}$\\
$\displaystyle v^{1122}_{-3}=R_{144}$\\
$\displaystyle v^{1122}_{-2}=R_{145}$}
\label{eq:cat-map-45}
\end{equation}
\subsubsection{Component 1212}
\begin{equation}
\shortstack[l]{$\displaystyle h^{1212}_{-4}=R_{146}$\\
$\displaystyle h^{1212}_{-3}=R_{147}$\\
$\displaystyle h^{1212}_{-2}=R_{148}$\\
$\displaystyle h^{1212}_{-1}=R_{149}$}
\label{eq:cat-map-46}
\end{equation}
\begin{equation}
\shortstack[l]{$\displaystyle b^{1212}_{-2}=R_{150}$}
\label{eq:cat-map-47}
\end{equation}
\begin{equation}
\shortstack[l]{$\displaystyle u^{1212}_{-4}=R_{151}$\\
$\displaystyle u^{1212}_{-3}=R_{152}$\\
$\displaystyle u^{1212}_{-2}=R_{153}$\\
$\displaystyle u^{1212}_{-1}=R_{149}$}
\label{eq:cat-map-48}
\end{equation}
\begin{equation}
\shortstack[l]{$\displaystyle v^{1212}_{-4}=R_{154}$\\
$\displaystyle v^{1212}_{-3}=R_{155}$\\
$\displaystyle v^{1212}_{-2}=R_{156}$}
\label{eq:cat-map-49}
\end{equation}
\subsubsection{Component 2222}
\begin{equation}
\shortstack[l]{$\displaystyle h^{2222}_{-4}=R_{157}$\\
$\displaystyle h^{2222}_{-3}=R_{158}$\\
$\displaystyle h^{2222}_{-2}=R_{159}$\\
$\displaystyle h^{2222}_{-1}=R_{160}$}
\label{eq:cat-map-50}
\end{equation}
\begin{equation}
\shortstack[l]{$\displaystyle b^{2222}_{-2}=R_{161}$}
\label{eq:cat-map-51}
\end{equation}
\begin{equation}
\shortstack[l]{$\displaystyle u^{2222}_{-4}=R_{162}$\\
$\displaystyle u^{2222}_{-3}=R_{163}$\\
$\displaystyle u^{2222}_{-2}=R_{164}$\\
$\displaystyle u^{2222}_{-1}=R_{160}$}
\label{eq:cat-map-52}
\end{equation}
\begin{equation}
\shortstack[l]{$\displaystyle v^{2222}_{-4}=R_{165}$\\
$\displaystyle v^{2222}_{-3}=R_{166}$\\
$\displaystyle v^{2222}_{-2}=R_{167}$}
\label{eq:cat-map-53}
\end{equation}
\subsubsection{Complete table}
\label{sec:cat-component-table}
A value of one in column $S$ denotes $z\mapsto-z$. In column $T$, it denotes the substitution $r\mapsto1/r$ with complex conjugation and multiplication by $r^{4+p}$ according to \eqref{eq:cat-contact-symmetries}. No transformations are required for zero rows.
\begin{longtable}{ccccc}
Component & Representative & $S$ & $T$ & Nonzero\\\hline
0000 & 0000 & 0 & 0 & 1\\
0001 & 0001 & 0 & 0 & 1\\
0002 & 0002 & 0 & 0 & 0\\
0010 & 0001 & 1 & 0 & 1\\
0011 & 0011 & 0 & 0 & 1\\
0012 & 0012 & 0 & 0 & 0\\
0020 & 0002 & 1 & 0 & 0\\
0021 & 0012 & 1 & 0 & 0\\
0022 & 0022 & 0 & 0 & 1\\
0100 & 0001 & 0 & 1 & 1\\
0101 & 0101 & 0 & 0 & 1\\
0102 & 0102 & 0 & 0 & 0\\
0110 & 0101 & 1 & 0 & 1\\
0111 & 0111 & 0 & 0 & 1\\
0112 & 0112 & 0 & 0 & 0\\
0120 & 0102 & 1 & 0 & 0\\
0121 & 0112 & 1 & 0 & 0\\
0122 & 0122 & 0 & 0 & 1\\
0200 & 0002 & 0 & 1 & 0\\
0201 & 0102 & 0 & 1 & 0\\
0202 & 0202 & 0 & 0 & 1\\
0210 & 0102 & 1 & 1 & 0\\
0211 & 0211 & 0 & 0 & 0\\
0212 & 0212 & 0 & 0 & 1\\
0220 & 0202 & 1 & 0 & 1\\
0221 & 0212 & 1 & 0 & 1\\
0222 & 0222 & 0 & 0 & 0\\
1000 & 0001 & 1 & 1 & 1\\
1001 & 0101 & 1 & 0 & 1\\
1002 & 0102 & 1 & 0 & 0\\
1010 & 0101 & 0 & 0 & 1\\
1011 & 0111 & 1 & 0 & 1\\
1012 & 0112 & 1 & 0 & 0\\
1020 & 0102 & 0 & 0 & 0\\
1021 & 0112 & 0 & 0 & 0\\
1022 & 0122 & 1 & 0 & 1\\
1100 & 0011 & 0 & 1 & 1\\
1101 & 0111 & 0 & 1 & 1\\
1102 & 0211 & 0 & 1 & 0\\
1110 & 0111 & 1 & 1 & 1\\
1111 & 1111 & 0 & 0 & 1\\
1112 & 1112 & 0 & 0 & 0\\
1120 & 0211 & 1 & 1 & 0\\
1121 & 1112 & 1 & 0 & 0\\
1122 & 1122 & 0 & 0 & 1\\
1200 & 0012 & 0 & 1 & 0\\
1201 & 0112 & 0 & 1 & 0\\
1202 & 0212 & 0 & 1 & 1\\
1210 & 0112 & 1 & 1 & 0\\
1211 & 1112 & 0 & 1 & 0\\
1212 & 1212 & 0 & 0 & 1\\
1220 & 0212 & 1 & 1 & 1\\
1221 & 1212 & 1 & 0 & 1\\
1222 & 1222 & 0 & 0 & 0\\
2000 & 0002 & 1 & 1 & 0\\
2001 & 0102 & 1 & 1 & 0\\
2002 & 0202 & 1 & 0 & 1\\
2010 & 0102 & 0 & 1 & 0\\
2011 & 0211 & 1 & 0 & 0\\
2012 & 0212 & 1 & 0 & 1\\
2020 & 0202 & 0 & 0 & 1\\
2021 & 0212 & 0 & 0 & 1\\
2022 & 0222 & 1 & 0 & 0\\
2100 & 0012 & 1 & 1 & 0\\
2101 & 0112 & 1 & 1 & 0\\
2102 & 0212 & 1 & 1 & 1\\
2110 & 0112 & 0 & 1 & 0\\
2111 & 1112 & 1 & 1 & 0\\
2112 & 1212 & 1 & 0 & 1\\
2120 & 0212 & 0 & 1 & 1\\
2121 & 1212 & 0 & 0 & 1\\
2122 & 1222 & 1 & 0 & 0\\
2200 & 0022 & 0 & 1 & 1\\
2201 & 0122 & 0 & 1 & 1\\
2202 & 0222 & 0 & 1 & 0\\
2210 & 0122 & 1 & 1 & 1\\
2211 & 1122 & 0 & 1 & 1\\
2212 & 1222 & 0 & 1 & 0\\
2220 & 0222 & 1 & 1 & 0\\
2221 & 1222 & 1 & 1 & 0\\
2222 & 2222 & 0 & 0 & 1\\
\end{longtable}
\subsection{Coefficients of the internal field and its \texorpdfstring{$x$}{x}-derivative for all 54 pairs}
\label{sec:cat-pairs}
In the heading $ij:ab$, the first two indices denote the momentum positions, and the second two denote the types of the two external fields. The coefficients are defined by
\begin{equation}
\frac{[L_3]_{ij,I}}{g_{\mathrm R}k^3}=e^{-i\Omega_{ij}x}\big[A_{ij,I}(x)X_I+B_{ij,I}(x)\partial_xX_I\big],\qquad
A_{ij,I}=\sum_{p=-3}^{0}A_{ij,I,p}x^p,\quad B_{ij,I}=\sum_{p=-2}^{0}B_{ij,I,p}x^p.
\label{eq:cat-pair-polynomials}
\end{equation}
The overall factor $g_{\mathrm R}k^3$ has been factored out in accordance with the normalization in the article. The derivative $\partial_x$ is taken with respect to $x=k\eta$. The spatial momentum of the trial internal line is $-(\mathbf p_i+\mathbf p_j)$, but its polarization basis is constructed from $\mathbf p_i+\mathbf p_j$. The index $I=0,1,2$ refers to a line with nonzero momentum; for the pairs $01$ and $23$, the index $I=0,1,\ldots,5$ refers to six homogeneous directions. All coefficients not listed are zero.
\subsubsection{Pair 01:00}
\begin{equation}
A^{0100;0}=0.
\label{eq:cat-map-54}
\end{equation}
\begin{equation}
A^{0100;1}=0.
\label{eq:cat-map-55}
\end{equation}
\begin{equation}
\shortstack[l]{$\displaystyle A^{0100;2}_{-3}=R_{168}$\\
$\displaystyle A^{0100;2}_{-2}=R_{169}$\\
$\displaystyle A^{0100;2}_{-1}=R_{170}$}
\label{eq:cat-map-56}
\end{equation}
\begin{equation}
A^{0100;3}=0.
\label{eq:cat-map-57}
\end{equation}
\begin{equation}
A^{0100;4}=0.
\label{eq:cat-map-58}
\end{equation}
\begin{equation}
A^{0100;5}=0.
\label{eq:cat-map-59}
\end{equation}
\begin{equation}
\shortstack[l]{$\displaystyle B^{0100;0}_{-2}=R_{140}$\\
$\displaystyle B^{0100;0}_{-1}=R_{171}$\\
$\displaystyle B^{0100;0}_{0}=R_{172}$}
\label{eq:cat-map-60}
\end{equation}
\begin{equation}
B^{0100;1}=0.
\label{eq:cat-map-61}
\end{equation}
\begin{equation}
\shortstack[l]{$\displaystyle B^{0100;2}_{-2}=R_{173}$\\
$\displaystyle B^{0100;2}_{-1}=R_{174}$}
\label{eq:cat-map-62}
\end{equation}
\begin{equation}
B^{0100;3}=0.
\label{eq:cat-map-63}
\end{equation}
\begin{equation}
B^{0100;4}=0.
\label{eq:cat-map-64}
\end{equation}
\begin{equation}
B^{0100;5}=0.
\label{eq:cat-map-65}
\end{equation}
\subsubsection{Pair 01:01}
\begin{equation}
A^{0101;0}=0.
\label{eq:cat-map-66}
\end{equation}
\begin{equation}
\shortstack[l]{$\displaystyle A^{0101;1}_{-3}=R_{143}$\\
$\displaystyle A^{0101;1}_{-2}=R_{175}$}
\label{eq:cat-map-67}
\end{equation}
\begin{equation}
A^{0101;2}=0.
\label{eq:cat-map-68}
\end{equation}
\begin{equation}
A^{0101;3}=0.
\label{eq:cat-map-69}
\end{equation}
\begin{equation}
A^{0101;4}=0.
\label{eq:cat-map-70}
\end{equation}
\begin{equation}
A^{0101;5}=0.
\label{eq:cat-map-71}
\end{equation}
\begin{equation}
B^{0101;0}=0.
\label{eq:cat-map-72}
\end{equation}
\begin{equation}
\shortstack[l]{$\displaystyle B^{0101;1}_{-2}=R_{176}$\\
$\displaystyle B^{0101;1}_{-1}=R_{177}$}
\label{eq:cat-map-73}
\end{equation}
\begin{equation}
B^{0101;2}=0.
\label{eq:cat-map-74}
\end{equation}
\begin{equation}
B^{0101;3}=0.
\label{eq:cat-map-75}
\end{equation}
\begin{equation}
B^{0101;4}=0.
\label{eq:cat-map-76}
\end{equation}
\begin{equation}
B^{0101;5}=0.
\label{eq:cat-map-77}
\end{equation}
\subsubsection{Pair 01:02}
\begin{equation}
A^{0102;0}=0.
\label{eq:cat-map-78}
\end{equation}
\begin{equation}
A^{0102;1}=0.
\label{eq:cat-map-79}
\end{equation}
\begin{equation}
A^{0102;2}=0.
\label{eq:cat-map-80}
\end{equation}
\begin{equation}
\shortstack[l]{$\displaystyle A^{0102;3}_{-3}=R_{176}$\\
$\displaystyle A^{0102;3}_{-2}=R_{177}$}
\label{eq:cat-map-81}
\end{equation}
\begin{equation}
A^{0102;4}=0.
\label{eq:cat-map-82}
\end{equation}
\begin{equation}
A^{0102;5}=0.
\label{eq:cat-map-83}
\end{equation}
\begin{equation}
B^{0102;0}=0.
\label{eq:cat-map-84}
\end{equation}
\begin{equation}
B^{0102;1}=0.
\label{eq:cat-map-85}
\end{equation}
\begin{equation}
B^{0102;2}=0.
\label{eq:cat-map-86}
\end{equation}
\begin{equation}
\shortstack[l]{$\displaystyle B^{0102;3}_{-2}=R_{143}$\\
$\displaystyle B^{0102;3}_{-1}=R_{175}$}
\label{eq:cat-map-87}
\end{equation}
\begin{equation}
B^{0102;4}=0.
\label{eq:cat-map-88}
\end{equation}
\begin{equation}
B^{0102;5}=0.
\label{eq:cat-map-89}
\end{equation}
\subsubsection{Pair 01:10}
\begin{equation}
A^{0110;0}=0.
\label{eq:cat-map-90}
\end{equation}
\begin{equation}
\shortstack[l]{$\displaystyle A^{0110;1}_{-3}=R_{143}$\\
$\displaystyle A^{0110;1}_{-2}=R_{175}$}
\label{eq:cat-map-91}
\end{equation}
\begin{equation}
A^{0110;2}=0.
\label{eq:cat-map-92}
\end{equation}
\begin{equation}
A^{0110;3}=0.
\label{eq:cat-map-93}
\end{equation}
\begin{equation}
A^{0110;4}=0.
\label{eq:cat-map-94}
\end{equation}
\begin{equation}
A^{0110;5}=0.
\label{eq:cat-map-95}
\end{equation}
\begin{equation}
B^{0110;0}=0.
\label{eq:cat-map-96}
\end{equation}
\begin{equation}
\shortstack[l]{$\displaystyle B^{0110;1}_{-2}=R_{176}$\\
$\displaystyle B^{0110;1}_{-1}=R_{177}$}
\label{eq:cat-map-97}
\end{equation}
\begin{equation}
B^{0110;2}=0.
\label{eq:cat-map-98}
\end{equation}
\begin{equation}
B^{0110;3}=0.
\label{eq:cat-map-99}
\end{equation}
\begin{equation}
B^{0110;4}=0.
\label{eq:cat-map-100}
\end{equation}
\begin{equation}
B^{0110;5}=0.
\label{eq:cat-map-101}
\end{equation}
\subsubsection{Pair 01:11}
\begin{equation}
A^{0111;0}=0.
\label{eq:cat-map-102}
\end{equation}
\begin{equation}
A^{0111;1}=0.
\label{eq:cat-map-103}
\end{equation}
\begin{equation}
\shortstack[l]{$\displaystyle A^{0111;2}_{-1}=R_{178}$}
\label{eq:cat-map-104}
\end{equation}
\begin{equation}
A^{0111;3}=0.
\label{eq:cat-map-105}
\end{equation}
\begin{equation}
A^{0111;4}=0.
\label{eq:cat-map-106}
\end{equation}
\begin{equation}
A^{0111;5}=0.
\label{eq:cat-map-107}
\end{equation}
\begin{equation}
\shortstack[l]{$\displaystyle B^{0111;0}_{-2}=R_{179}$\\
$\displaystyle B^{0111;0}_{-1}=R_{174}$}
\label{eq:cat-map-108}
\end{equation}
\begin{equation}
B^{0111;1}=0.
\label{eq:cat-map-109}
\end{equation}
\begin{equation}
B^{0111;2}=0.
\label{eq:cat-map-110}
\end{equation}
\begin{equation}
B^{0111;3}=0.
\label{eq:cat-map-111}
\end{equation}
\begin{equation}
B^{0111;4}=0.
\label{eq:cat-map-112}
\end{equation}
\begin{equation}
B^{0111;5}=0.
\label{eq:cat-map-113}
\end{equation}
\subsubsection{Pair 01:12}
\begin{equation}
A^{0112;0}=0.
\label{eq:cat-map-114}
\end{equation}
\begin{equation}
A^{0112;1}=0.
\label{eq:cat-map-115}
\end{equation}
\begin{equation}
A^{0112;2}=0.
\label{eq:cat-map-116}
\end{equation}
\begin{equation}
A^{0112;3}=0.
\label{eq:cat-map-117}
\end{equation}
\begin{equation}
A^{0112;4}=0.
\label{eq:cat-map-118}
\end{equation}
\begin{equation}
A^{0112;5}=0.
\label{eq:cat-map-119}
\end{equation}
\begin{equation}
B^{0112;0}=0.
\label{eq:cat-map-120}
\end{equation}
\begin{equation}
B^{0112;1}=0.
\label{eq:cat-map-121}
\end{equation}
\begin{equation}
B^{0112;2}=0.
\label{eq:cat-map-122}
\end{equation}
\begin{equation}
B^{0112;3}=0.
\label{eq:cat-map-123}
\end{equation}
\begin{equation}
B^{0112;4}=0.
\label{eq:cat-map-124}
\end{equation}
\begin{equation}
B^{0112;5}=0.
\label{eq:cat-map-125}
\end{equation}
\subsubsection{Pair 01:20}
\begin{equation}
A^{0120;0}=0.
\label{eq:cat-map-126}
\end{equation}
\begin{equation}
A^{0120;1}=0.
\label{eq:cat-map-127}
\end{equation}
\begin{equation}
A^{0120;2}=0.
\label{eq:cat-map-128}
\end{equation}
\begin{equation}
\shortstack[l]{$\displaystyle A^{0120;3}_{-3}=R_{143}$\\
$\displaystyle A^{0120;3}_{-2}=R_{175}$}
\label{eq:cat-map-129}
\end{equation}
\begin{equation}
A^{0120;4}=0.
\label{eq:cat-map-130}
\end{equation}
\begin{equation}
A^{0120;5}=0.
\label{eq:cat-map-131}
\end{equation}
\begin{equation}
B^{0120;0}=0.
\label{eq:cat-map-132}
\end{equation}
\begin{equation}
B^{0120;1}=0.
\label{eq:cat-map-133}
\end{equation}
\begin{equation}
B^{0120;2}=0.
\label{eq:cat-map-134}
\end{equation}
\begin{equation}
\shortstack[l]{$\displaystyle B^{0120;3}_{-2}=R_{176}$\\
$\displaystyle B^{0120;3}_{-1}=R_{177}$}
\label{eq:cat-map-135}
\end{equation}
\begin{equation}
B^{0120;4}=0.
\label{eq:cat-map-136}
\end{equation}
\begin{equation}
B^{0120;5}=0.
\label{eq:cat-map-137}
\end{equation}
\subsubsection{Pair 01:21}
\begin{equation}
A^{0121;0}=0.
\label{eq:cat-map-138}
\end{equation}
\begin{equation}
A^{0121;1}=0.
\label{eq:cat-map-139}
\end{equation}
\begin{equation}
A^{0121;2}=0.
\label{eq:cat-map-140}
\end{equation}
\begin{equation}
A^{0121;3}=0.
\label{eq:cat-map-141}
\end{equation}
\begin{equation}
A^{0121;4}=0.
\label{eq:cat-map-142}
\end{equation}
\begin{equation}
A^{0121;5}=0.
\label{eq:cat-map-143}
\end{equation}
\begin{equation}
B^{0121;0}=0.
\label{eq:cat-map-144}
\end{equation}
\begin{equation}
B^{0121;1}=0.
\label{eq:cat-map-145}
\end{equation}
\begin{equation}
B^{0121;2}=0.
\label{eq:cat-map-146}
\end{equation}
\begin{equation}
B^{0121;3}=0.
\label{eq:cat-map-147}
\end{equation}
\begin{equation}
B^{0121;4}=0.
\label{eq:cat-map-148}
\end{equation}
\begin{equation}
B^{0121;5}=0.
\label{eq:cat-map-149}
\end{equation}
\subsubsection{Pair 01:22}
\begin{equation}
A^{0122;0}=0.
\label{eq:cat-map-150}
\end{equation}
\begin{equation}
A^{0122;1}=0.
\label{eq:cat-map-151}
\end{equation}
\begin{equation}
\shortstack[l]{$\displaystyle A^{0122;2}_{-1}=R_{180}$}
\label{eq:cat-map-152}
\end{equation}
\begin{equation}
A^{0122;3}=0.
\label{eq:cat-map-153}
\end{equation}
\begin{equation}
A^{0122;4}=0.
\label{eq:cat-map-154}
\end{equation}
\begin{equation}
A^{0122;5}=0.
\label{eq:cat-map-155}
\end{equation}
\begin{equation}
\shortstack[l]{$\displaystyle B^{0122;0}_{-2}=R_{181}$\\
$\displaystyle B^{0122;0}_{-1}=R_{169}$}
\label{eq:cat-map-156}
\end{equation}
\begin{equation}
B^{0122;1}=0.
\label{eq:cat-map-157}
\end{equation}
\begin{equation}
B^{0122;2}=0.
\label{eq:cat-map-158}
\end{equation}
\begin{equation}
B^{0122;3}=0.
\label{eq:cat-map-159}
\end{equation}
\begin{equation}
B^{0122;4}=0.
\label{eq:cat-map-160}
\end{equation}
\begin{equation}
B^{0122;5}=0.
\label{eq:cat-map-161}
\end{equation}
\subsubsection{Pair 02:00}
\begin{equation}
\shortstack[l]{$\displaystyle A^{0200;0}_{-3}=R_{182}$\\
$\displaystyle A^{0200;0}_{-2}=R_{183}$\\
$\displaystyle A^{0200;0}_{-1}=R_{184}$\\
$\displaystyle A^{0200;0}_{0}=R_{185}$}
\label{eq:cat-map-162}
\end{equation}
\begin{equation}
\shortstack[l]{$\displaystyle A^{0200;1}_{-3}=R_{186}$\\
$\displaystyle A^{0200;1}_{-2}=R_{187}$\\
$\displaystyle A^{0200;1}_{-1}=R_{188}$}
\label{eq:cat-map-163}
\end{equation}
\begin{equation}
A^{0200;2}=0.
\label{eq:cat-map-164}
\end{equation}
\begin{equation}
\shortstack[l]{$\displaystyle B^{0200;0}_{-2}=R_{189}$\\
$\displaystyle B^{0200;0}_{-1}=R_{190}$\\
$\displaystyle B^{0200;0}_{0}=R_{191}$}
\label{eq:cat-map-165}
\end{equation}
\begin{equation}
\shortstack[l]{$\displaystyle B^{0200;1}_{-2}=R_{192}$\\
$\displaystyle B^{0200;1}_{-1}=R_{193}$}
\label{eq:cat-map-166}
\end{equation}
\begin{equation}
B^{0200;2}=0.
\label{eq:cat-map-167}
\end{equation}
\subsubsection{Pair 02:01}
\begin{equation}
\shortstack[l]{$\displaystyle A^{0201;0}_{-3}=R_{194}$\\
$\displaystyle A^{0201;0}_{-2}=R_{195}$\\
$\displaystyle A^{0201;0}_{-1}=R_{196}$}
\label{eq:cat-map-168}
\end{equation}
\begin{equation}
\shortstack[l]{$\displaystyle A^{0201;1}_{-3}=R_{197}$\\
$\displaystyle A^{0201;1}_{-2}=R_{198}$\\
$\displaystyle A^{0201;1}_{-1}=R_{199}$}
\label{eq:cat-map-169}
\end{equation}
\begin{equation}
A^{0201;2}=0.
\label{eq:cat-map-170}
\end{equation}
\begin{equation}
\shortstack[l]{$\displaystyle B^{0201;0}_{-2}=R_{200}$\\
$\displaystyle B^{0201;0}_{-1}=R_{201}$}
\label{eq:cat-map-171}
\end{equation}
\begin{equation}
\shortstack[l]{$\displaystyle B^{0201;1}_{-2}=R_{202}$\\
$\displaystyle B^{0201;1}_{-1}=R_{203}$}
\label{eq:cat-map-172}
\end{equation}
\begin{equation}
B^{0201;2}=0.
\label{eq:cat-map-173}
\end{equation}
\subsubsection{Pair 02:02}
\begin{equation}
A^{0202;0}=0.
\label{eq:cat-map-174}
\end{equation}
\begin{equation}
A^{0202;1}=0.
\label{eq:cat-map-175}
\end{equation}
\begin{equation}
\shortstack[l]{$\displaystyle A^{0202;2}_{-3}=R_{204}$\\
$\displaystyle A^{0202;2}_{-2}=R_{205}$\\
$\displaystyle A^{0202;2}_{-1}=R_{206}$}
\label{eq:cat-map-176}
\end{equation}
\begin{equation}
B^{0202;0}=0.
\label{eq:cat-map-177}
\end{equation}
\begin{equation}
B^{0202;1}=0.
\label{eq:cat-map-178}
\end{equation}
\begin{equation}
\shortstack[l]{$\displaystyle B^{0202;2}_{-2}=R_{207}$\\
$\displaystyle B^{0202;2}_{-1}=R_{208}$}
\label{eq:cat-map-179}
\end{equation}
\subsubsection{Pair 02:10}
\begin{equation}
\shortstack[l]{$\displaystyle A^{0210;0}_{-3}=R_{209}$\\
$\displaystyle A^{0210;0}_{-2}=R_{210}$\\
$\displaystyle A^{0210;0}_{-1}=R_{211}$}
\label{eq:cat-map-180}
\end{equation}
\begin{equation}
\shortstack[l]{$\displaystyle A^{0210;1}_{-3}=R_{212}$\\
$\displaystyle A^{0210;1}_{-2}=R_{213}$\\
$\displaystyle A^{0210;1}_{-1}=R_{214}$}
\label{eq:cat-map-181}
\end{equation}
\begin{equation}
A^{0210;2}=0.
\label{eq:cat-map-182}
\end{equation}
\begin{equation}
\shortstack[l]{$\displaystyle B^{0210;0}_{-2}=R_{215}$\\
$\displaystyle B^{0210;0}_{-1}=R_{216}$}
\label{eq:cat-map-183}
\end{equation}
\begin{equation}
\shortstack[l]{$\displaystyle B^{0210;1}_{-2}=R_{217}$\\
$\displaystyle B^{0210;1}_{-1}=R_{218}$}
\label{eq:cat-map-184}
\end{equation}
\begin{equation}
B^{0210;2}=0.
\label{eq:cat-map-185}
\end{equation}
\subsubsection{Pair 02:11}
\begin{equation}
\shortstack[l]{$\displaystyle A^{0211;0}_{-3}=R_{219}$\\
$\displaystyle A^{0211;0}_{-2}=R_{220}$\\
$\displaystyle A^{0211;0}_{-1}=R_{221}$}
\label{eq:cat-map-186}
\end{equation}
\begin{equation}
\shortstack[l]{$\displaystyle A^{0211;1}_{-1}=R_{222}$}
\label{eq:cat-map-187}
\end{equation}
\begin{equation}
A^{0211;2}=0.
\label{eq:cat-map-188}
\end{equation}
\begin{equation}
\shortstack[l]{$\displaystyle B^{0211;0}_{-2}=R_{223}$\\
$\displaystyle B^{0211;0}_{-1}=R_{224}$}
\label{eq:cat-map-189}
\end{equation}
\begin{equation}
B^{0211;1}=0.
\label{eq:cat-map-190}
\end{equation}
\begin{equation}
B^{0211;2}=0.
\label{eq:cat-map-191}
\end{equation}
\subsubsection{Pair 02:12}
\begin{equation}
A^{0212;0}=0.
\label{eq:cat-map-192}
\end{equation}
\begin{equation}
A^{0212;1}=0.
\label{eq:cat-map-193}
\end{equation}
\begin{equation}
\shortstack[l]{$\displaystyle A^{0212;2}_{-1}=R_{225}$}
\label{eq:cat-map-194}
\end{equation}
\begin{equation}
B^{0212;0}=0.
\label{eq:cat-map-195}
\end{equation}
\begin{equation}
B^{0212;1}=0.
\label{eq:cat-map-196}
\end{equation}
\begin{equation}
B^{0212;2}=0.
\label{eq:cat-map-197}
\end{equation}
\subsubsection{Pair 02:20}
\begin{equation}
A^{0220;0}=0.
\label{eq:cat-map-198}
\end{equation}
\begin{equation}
A^{0220;1}=0.
\label{eq:cat-map-199}
\end{equation}
\begin{equation}
\shortstack[l]{$\displaystyle A^{0220;2}_{-3}=R_{226}$\\
$\displaystyle A^{0220;2}_{-2}=R_{227}$\\
$\displaystyle A^{0220;2}_{-1}=R_{228}$}
\label{eq:cat-map-200}
\end{equation}
\begin{equation}
B^{0220;0}=0.
\label{eq:cat-map-201}
\end{equation}
\begin{equation}
B^{0220;1}=0.
\label{eq:cat-map-202}
\end{equation}
\begin{equation}
\shortstack[l]{$\displaystyle B^{0220;2}_{-2}=R_{229}$\\
$\displaystyle B^{0220;2}_{-1}=R_{230}$}
\label{eq:cat-map-203}
\end{equation}
\subsubsection{Pair 02:21}
\begin{equation}
A^{0221;0}=0.
\label{eq:cat-map-204}
\end{equation}
\begin{equation}
A^{0221;1}=0.
\label{eq:cat-map-205}
\end{equation}
\begin{equation}
\shortstack[l]{$\displaystyle A^{0221;2}_{-1}=R_{231}$}
\label{eq:cat-map-206}
\end{equation}
\begin{equation}
B^{0221;0}=0.
\label{eq:cat-map-207}
\end{equation}
\begin{equation}
B^{0221;1}=0.
\label{eq:cat-map-208}
\end{equation}
\begin{equation}
B^{0221;2}=0.
\label{eq:cat-map-209}
\end{equation}
\subsubsection{Pair 02:22}
\begin{equation}
\shortstack[l]{$\displaystyle A^{0222;0}_{-3}=R_{232}$\\
$\displaystyle A^{0222;0}_{-2}=R_{233}$\\
$\displaystyle A^{0222;0}_{-1}=R_{234}$}
\label{eq:cat-map-210}
\end{equation}
\begin{equation}
\shortstack[l]{$\displaystyle A^{0222;1}_{-1}=R_{235}$}
\label{eq:cat-map-211}
\end{equation}
\begin{equation}
A^{0222;2}=0.
\label{eq:cat-map-212}
\end{equation}
\begin{equation}
\shortstack[l]{$\displaystyle B^{0222;0}_{-2}=R_{236}$\\
$\displaystyle B^{0222;0}_{-1}=R_{237}$}
\label{eq:cat-map-213}
\end{equation}
\begin{equation}
B^{0222;1}=0.
\label{eq:cat-map-214}
\end{equation}
\begin{equation}
B^{0222;2}=0.
\label{eq:cat-map-215}
\end{equation}
\subsubsection{Pair 03:00}
\begin{equation}
\shortstack[l]{$\displaystyle A^{0300;0}_{-3}=R_{238}$\\
$\displaystyle A^{0300;0}_{-2}=R_{239}$\\
$\displaystyle A^{0300;0}_{-1}=R_{240}$\\
$\displaystyle A^{0300;0}_{0}=R_{241}$}
\label{eq:cat-map-216}
\end{equation}
\begin{equation}
\shortstack[l]{$\displaystyle A^{0300;1}_{-3}=R_{242}$\\
$\displaystyle A^{0300;1}_{-2}=R_{243}$\\
$\displaystyle A^{0300;1}_{-1}=R_{244}$}
\label{eq:cat-map-217}
\end{equation}
\begin{equation}
A^{0300;2}=0.
\label{eq:cat-map-218}
\end{equation}
\begin{equation}
\shortstack[l]{$\displaystyle B^{0300;0}_{-2}=R_{245}$\\
$\displaystyle B^{0300;0}_{-1}=R_{246}$\\
$\displaystyle B^{0300;0}_{0}=R_{247}$}
\label{eq:cat-map-219}
\end{equation}
\begin{equation}
\shortstack[l]{$\displaystyle B^{0300;1}_{-2}=R_{248}$\\
$\displaystyle B^{0300;1}_{-1}=R_{249}$}
\label{eq:cat-map-220}
\end{equation}
\begin{equation}
B^{0300;2}=0.
\label{eq:cat-map-221}
\end{equation}
\subsubsection{Pair 03:01}
\begin{equation}
\shortstack[l]{$\displaystyle A^{0301;0}_{-3}=R_{250}$\\
$\displaystyle A^{0301;0}_{-2}=R_{251}$\\
$\displaystyle A^{0301;0}_{-1}=R_{196}$}
\label{eq:cat-map-222}
\end{equation}
\begin{equation}
\shortstack[l]{$\displaystyle A^{0301;1}_{-3}=R_{252}$\\
$\displaystyle A^{0301;1}_{-2}=R_{253}$\\
$\displaystyle A^{0301;1}_{-1}=R_{254}$}
\label{eq:cat-map-223}
\end{equation}
\begin{equation}
A^{0301;2}=0.
\label{eq:cat-map-224}
\end{equation}
\begin{equation}
\shortstack[l]{$\displaystyle B^{0301;0}_{-2}=R_{255}$\\
$\displaystyle B^{0301;0}_{-1}=R_{256}$}
\label{eq:cat-map-225}
\end{equation}
\begin{equation}
\shortstack[l]{$\displaystyle B^{0301;1}_{-2}=R_{257}$\\
$\displaystyle B^{0301;1}_{-1}=R_{258}$}
\label{eq:cat-map-226}
\end{equation}
\begin{equation}
B^{0301;2}=0.
\label{eq:cat-map-227}
\end{equation}
\subsubsection{Pair 03:02}
\begin{equation}
A^{0302;0}=0.
\label{eq:cat-map-228}
\end{equation}
\begin{equation}
A^{0302;1}=0.
\label{eq:cat-map-229}
\end{equation}
\begin{equation}
\shortstack[l]{$\displaystyle A^{0302;2}_{-3}=R_{259}$\\
$\displaystyle A^{0302;2}_{-2}=R_{260}$\\
$\displaystyle A^{0302;2}_{-1}=R_{261}$}
\label{eq:cat-map-230}
\end{equation}
\begin{equation}
B^{0302;0}=0.
\label{eq:cat-map-231}
\end{equation}
\begin{equation}
B^{0302;1}=0.
\label{eq:cat-map-232}
\end{equation}
\begin{equation}
\shortstack[l]{$\displaystyle B^{0302;2}_{-2}=R_{262}$\\
$\displaystyle B^{0302;2}_{-1}=R_{263}$}
\label{eq:cat-map-233}
\end{equation}
\subsubsection{Pair 03:10}
\begin{equation}
\shortstack[l]{$\displaystyle A^{0310;0}_{-3}=R_{264}$\\
$\displaystyle A^{0310;0}_{-2}=R_{265}$\\
$\displaystyle A^{0310;0}_{-1}=R_{211}$}
\label{eq:cat-map-234}
\end{equation}
\begin{equation}
\shortstack[l]{$\displaystyle A^{0310;1}_{-3}=R_{266}$\\
$\displaystyle A^{0310;1}_{-2}=R_{267}$\\
$\displaystyle A^{0310;1}_{-1}=R_{268}$}
\label{eq:cat-map-235}
\end{equation}
\begin{equation}
A^{0310;2}=0.
\label{eq:cat-map-236}
\end{equation}
\begin{equation}
\shortstack[l]{$\displaystyle B^{0310;0}_{-2}=R_{269}$\\
$\displaystyle B^{0310;0}_{-1}=R_{270}$}
\label{eq:cat-map-237}
\end{equation}
\begin{equation}
\shortstack[l]{$\displaystyle B^{0310;1}_{-2}=R_{271}$\\
$\displaystyle B^{0310;1}_{-1}=R_{272}$}
\label{eq:cat-map-238}
\end{equation}
\begin{equation}
B^{0310;2}=0.
\label{eq:cat-map-239}
\end{equation}
\subsubsection{Pair 03:11}
\begin{equation}
\shortstack[l]{$\displaystyle A^{0311;0}_{-3}=R_{219}$\\
$\displaystyle A^{0311;0}_{-2}=R_{273}$\\
$\displaystyle A^{0311;0}_{-1}=R_{274}$}
\label{eq:cat-map-240}
\end{equation}
\begin{equation}
\shortstack[l]{$\displaystyle A^{0311;1}_{-1}=R_{275}$}
\label{eq:cat-map-241}
\end{equation}
\begin{equation}
A^{0311;2}=0.
\label{eq:cat-map-242}
\end{equation}
\begin{equation}
\shortstack[l]{$\displaystyle B^{0311;0}_{-2}=R_{223}$\\
$\displaystyle B^{0311;0}_{-1}=R_{276}$}
\label{eq:cat-map-243}
\end{equation}
\begin{equation}
B^{0311;1}=0.
\label{eq:cat-map-244}
\end{equation}
\begin{equation}
B^{0311;2}=0.
\label{eq:cat-map-245}
\end{equation}
\subsubsection{Pair 03:12}
\begin{equation}
A^{0312;0}=0.
\label{eq:cat-map-246}
\end{equation}
\begin{equation}
A^{0312;1}=0.
\label{eq:cat-map-247}
\end{equation}
\begin{equation}
\shortstack[l]{$\displaystyle A^{0312;2}_{-1}=R_{277}$}
\label{eq:cat-map-248}
\end{equation}
\begin{equation}
B^{0312;0}=0.
\label{eq:cat-map-249}
\end{equation}
\begin{equation}
B^{0312;1}=0.
\label{eq:cat-map-250}
\end{equation}
\begin{equation}
B^{0312;2}=0.
\label{eq:cat-map-251}
\end{equation}
\subsubsection{Pair 03:20}
\begin{equation}
A^{0320;0}=0.
\label{eq:cat-map-252}
\end{equation}
\begin{equation}
A^{0320;1}=0.
\label{eq:cat-map-253}
\end{equation}
\begin{equation}
\shortstack[l]{$\displaystyle A^{0320;2}_{-3}=R_{278}$\\
$\displaystyle A^{0320;2}_{-2}=R_{279}$\\
$\displaystyle A^{0320;2}_{-1}=R_{280}$}
\label{eq:cat-map-254}
\end{equation}
\begin{equation}
B^{0320;0}=0.
\label{eq:cat-map-255}
\end{equation}
\begin{equation}
B^{0320;1}=0.
\label{eq:cat-map-256}
\end{equation}
\begin{equation}
\shortstack[l]{$\displaystyle B^{0320;2}_{-2}=R_{281}$\\
$\displaystyle B^{0320;2}_{-1}=R_{282}$}
\label{eq:cat-map-257}
\end{equation}
\subsubsection{Pair 03:21}
\begin{equation}
A^{0321;0}=0.
\label{eq:cat-map-258}
\end{equation}
\begin{equation}
A^{0321;1}=0.
\label{eq:cat-map-259}
\end{equation}
\begin{equation}
\shortstack[l]{$\displaystyle A^{0321;2}_{-1}=R_{283}$}
\label{eq:cat-map-260}
\end{equation}
\begin{equation}
B^{0321;0}=0.
\label{eq:cat-map-261}
\end{equation}
\begin{equation}
B^{0321;1}=0.
\label{eq:cat-map-262}
\end{equation}
\begin{equation}
B^{0321;2}=0.
\label{eq:cat-map-263}
\end{equation}
\subsubsection{Pair 03:22}
\begin{equation}
\shortstack[l]{$\displaystyle A^{0322;0}_{-3}=R_{284}$\\
$\displaystyle A^{0322;0}_{-2}=R_{285}$\\
$\displaystyle A^{0322;0}_{-1}=R_{286}$}
\label{eq:cat-map-264}
\end{equation}
\begin{equation}
\shortstack[l]{$\displaystyle A^{0322;1}_{-1}=R_{287}$}
\label{eq:cat-map-265}
\end{equation}
\begin{equation}
A^{0322;2}=0.
\label{eq:cat-map-266}
\end{equation}
\begin{equation}
\shortstack[l]{$\displaystyle B^{0322;0}_{-2}=R_{288}$\\
$\displaystyle B^{0322;0}_{-1}=R_{289}$}
\label{eq:cat-map-267}
\end{equation}
\begin{equation}
B^{0322;1}=0.
\label{eq:cat-map-268}
\end{equation}
\begin{equation}
B^{0322;2}=0.
\label{eq:cat-map-269}
\end{equation}
\subsubsection{Pair 12:00}
\begin{equation}
\shortstack[l]{$\displaystyle A^{1200;0}_{-3}=R_{238}$\\
$\displaystyle A^{1200;0}_{-2}=R_{239}$\\
$\displaystyle A^{1200;0}_{-1}=R_{240}$\\
$\displaystyle A^{1200;0}_{0}=R_{241}$}
\label{eq:cat-map-270}
\end{equation}
\begin{equation}
\shortstack[l]{$\displaystyle A^{1200;1}_{-3}=R_{242}$\\
$\displaystyle A^{1200;1}_{-2}=R_{243}$\\
$\displaystyle A^{1200;1}_{-1}=R_{244}$}
\label{eq:cat-map-271}
\end{equation}
\begin{equation}
A^{1200;2}=0.
\label{eq:cat-map-272}
\end{equation}
\begin{equation}
\shortstack[l]{$\displaystyle B^{1200;0}_{-2}=R_{245}$\\
$\displaystyle B^{1200;0}_{-1}=R_{246}$\\
$\displaystyle B^{1200;0}_{0}=R_{247}$}
\label{eq:cat-map-273}
\end{equation}
\begin{equation}
\shortstack[l]{$\displaystyle B^{1200;1}_{-2}=R_{248}$\\
$\displaystyle B^{1200;1}_{-1}=R_{249}$}
\label{eq:cat-map-274}
\end{equation}
\begin{equation}
B^{1200;2}=0.
\label{eq:cat-map-275}
\end{equation}
\subsubsection{Pair 12:01}
\begin{equation}
\shortstack[l]{$\displaystyle A^{1201;0}_{-3}=R_{250}$\\
$\displaystyle A^{1201;0}_{-2}=R_{251}$\\
$\displaystyle A^{1201;0}_{-1}=R_{196}$}
\label{eq:cat-map-276}
\end{equation}
\begin{equation}
\shortstack[l]{$\displaystyle A^{1201;1}_{-3}=R_{252}$\\
$\displaystyle A^{1201;1}_{-2}=R_{253}$\\
$\displaystyle A^{1201;1}_{-1}=R_{254}$}
\label{eq:cat-map-277}
\end{equation}
\begin{equation}
A^{1201;2}=0.
\label{eq:cat-map-278}
\end{equation}
\begin{equation}
\shortstack[l]{$\displaystyle B^{1201;0}_{-2}=R_{255}$\\
$\displaystyle B^{1201;0}_{-1}=R_{256}$}
\label{eq:cat-map-279}
\end{equation}
\begin{equation}
\shortstack[l]{$\displaystyle B^{1201;1}_{-2}=R_{257}$\\
$\displaystyle B^{1201;1}_{-1}=R_{258}$}
\label{eq:cat-map-280}
\end{equation}
\begin{equation}
B^{1201;2}=0.
\label{eq:cat-map-281}
\end{equation}
\subsubsection{Pair 12:02}
\begin{equation}
A^{1202;0}=0.
\label{eq:cat-map-282}
\end{equation}
\begin{equation}
A^{1202;1}=0.
\label{eq:cat-map-283}
\end{equation}
\begin{equation}
\shortstack[l]{$\displaystyle A^{1202;2}_{-3}=R_{259}$\\
$\displaystyle A^{1202;2}_{-2}=R_{260}$\\
$\displaystyle A^{1202;2}_{-1}=R_{261}$}
\label{eq:cat-map-284}
\end{equation}
\begin{equation}
B^{1202;0}=0.
\label{eq:cat-map-285}
\end{equation}
\begin{equation}
B^{1202;1}=0.
\label{eq:cat-map-286}
\end{equation}
\begin{equation}
\shortstack[l]{$\displaystyle B^{1202;2}_{-2}=R_{262}$\\
$\displaystyle B^{1202;2}_{-1}=R_{263}$}
\label{eq:cat-map-287}
\end{equation}
\subsubsection{Pair 12:10}
\begin{equation}
\shortstack[l]{$\displaystyle A^{1210;0}_{-3}=R_{264}$\\
$\displaystyle A^{1210;0}_{-2}=R_{265}$\\
$\displaystyle A^{1210;0}_{-1}=R_{211}$}
\label{eq:cat-map-288}
\end{equation}
\begin{equation}
\shortstack[l]{$\displaystyle A^{1210;1}_{-3}=R_{266}$\\
$\displaystyle A^{1210;1}_{-2}=R_{267}$\\
$\displaystyle A^{1210;1}_{-1}=R_{268}$}
\label{eq:cat-map-289}
\end{equation}
\begin{equation}
A^{1210;2}=0.
\label{eq:cat-map-290}
\end{equation}
\begin{equation}
\shortstack[l]{$\displaystyle B^{1210;0}_{-2}=R_{269}$\\
$\displaystyle B^{1210;0}_{-1}=R_{270}$}
\label{eq:cat-map-291}
\end{equation}
\begin{equation}
\shortstack[l]{$\displaystyle B^{1210;1}_{-2}=R_{271}$\\
$\displaystyle B^{1210;1}_{-1}=R_{272}$}
\label{eq:cat-map-292}
\end{equation}
\begin{equation}
B^{1210;2}=0.
\label{eq:cat-map-293}
\end{equation}
\subsubsection{Pair 12:11}
\begin{equation}
\shortstack[l]{$\displaystyle A^{1211;0}_{-3}=R_{219}$\\
$\displaystyle A^{1211;0}_{-2}=R_{273}$\\
$\displaystyle A^{1211;0}_{-1}=R_{274}$}
\label{eq:cat-map-294}
\end{equation}
\begin{equation}
\shortstack[l]{$\displaystyle A^{1211;1}_{-1}=R_{275}$}
\label{eq:cat-map-295}
\end{equation}
\begin{equation}
A^{1211;2}=0.
\label{eq:cat-map-296}
\end{equation}
\begin{equation}
\shortstack[l]{$\displaystyle B^{1211;0}_{-2}=R_{223}$\\
$\displaystyle B^{1211;0}_{-1}=R_{276}$}
\label{eq:cat-map-297}
\end{equation}
\begin{equation}
B^{1211;1}=0.
\label{eq:cat-map-298}
\end{equation}
\begin{equation}
B^{1211;2}=0.
\label{eq:cat-map-299}
\end{equation}
\subsubsection{Pair 12:12}
\begin{equation}
A^{1212;0}=0.
\label{eq:cat-map-300}
\end{equation}
\begin{equation}
A^{1212;1}=0.
\label{eq:cat-map-301}
\end{equation}
\begin{equation}
\shortstack[l]{$\displaystyle A^{1212;2}_{-1}=R_{277}$}
\label{eq:cat-map-302}
\end{equation}
\begin{equation}
B^{1212;0}=0.
\label{eq:cat-map-303}
\end{equation}
\begin{equation}
B^{1212;1}=0.
\label{eq:cat-map-304}
\end{equation}
\begin{equation}
B^{1212;2}=0.
\label{eq:cat-map-305}
\end{equation}
\subsubsection{Pair 12:20}
\begin{equation}
A^{1220;0}=0.
\label{eq:cat-map-306}
\end{equation}
\begin{equation}
A^{1220;1}=0.
\label{eq:cat-map-307}
\end{equation}
\begin{equation}
\shortstack[l]{$\displaystyle A^{1220;2}_{-3}=R_{278}$\\
$\displaystyle A^{1220;2}_{-2}=R_{279}$\\
$\displaystyle A^{1220;2}_{-1}=R_{280}$}
\label{eq:cat-map-308}
\end{equation}
\begin{equation}
B^{1220;0}=0.
\label{eq:cat-map-309}
\end{equation}
\begin{equation}
B^{1220;1}=0.
\label{eq:cat-map-310}
\end{equation}
\begin{equation}
\shortstack[l]{$\displaystyle B^{1220;2}_{-2}=R_{281}$\\
$\displaystyle B^{1220;2}_{-1}=R_{282}$}
\label{eq:cat-map-311}
\end{equation}
\subsubsection{Pair 12:21}
\begin{equation}
A^{1221;0}=0.
\label{eq:cat-map-312}
\end{equation}
\begin{equation}
A^{1221;1}=0.
\label{eq:cat-map-313}
\end{equation}
\begin{equation}
\shortstack[l]{$\displaystyle A^{1221;2}_{-1}=R_{283}$}
\label{eq:cat-map-314}
\end{equation}
\begin{equation}
B^{1221;0}=0.
\label{eq:cat-map-315}
\end{equation}
\begin{equation}
B^{1221;1}=0.
\label{eq:cat-map-316}
\end{equation}
\begin{equation}
B^{1221;2}=0.
\label{eq:cat-map-317}
\end{equation}
\subsubsection{Pair 12:22}
\begin{equation}
\shortstack[l]{$\displaystyle A^{1222;0}_{-3}=R_{284}$\\
$\displaystyle A^{1222;0}_{-2}=R_{285}$\\
$\displaystyle A^{1222;0}_{-1}=R_{286}$}
\label{eq:cat-map-318}
\end{equation}
\begin{equation}
\shortstack[l]{$\displaystyle A^{1222;1}_{-1}=R_{287}$}
\label{eq:cat-map-319}
\end{equation}
\begin{equation}
A^{1222;2}=0.
\label{eq:cat-map-320}
\end{equation}
\begin{equation}
\shortstack[l]{$\displaystyle B^{1222;0}_{-2}=R_{288}$\\
$\displaystyle B^{1222;0}_{-1}=R_{289}$}
\label{eq:cat-map-321}
\end{equation}
\begin{equation}
B^{1222;1}=0.
\label{eq:cat-map-322}
\end{equation}
\begin{equation}
B^{1222;2}=0.
\label{eq:cat-map-323}
\end{equation}
\subsubsection{Pair 13:00}
\begin{equation}
\shortstack[l]{$\displaystyle A^{1300;0}_{-3}=R_{182}$\\
$\displaystyle A^{1300;0}_{-2}=R_{183}$\\
$\displaystyle A^{1300;0}_{-1}=R_{184}$\\
$\displaystyle A^{1300;0}_{0}=R_{185}$}
\label{eq:cat-map-324}
\end{equation}
\begin{equation}
\shortstack[l]{$\displaystyle A^{1300;1}_{-3}=R_{186}$\\
$\displaystyle A^{1300;1}_{-2}=R_{187}$\\
$\displaystyle A^{1300;1}_{-1}=R_{188}$}
\label{eq:cat-map-325}
\end{equation}
\begin{equation}
A^{1300;2}=0.
\label{eq:cat-map-326}
\end{equation}
\begin{equation}
\shortstack[l]{$\displaystyle B^{1300;0}_{-2}=R_{189}$\\
$\displaystyle B^{1300;0}_{-1}=R_{190}$\\
$\displaystyle B^{1300;0}_{0}=R_{191}$}
\label{eq:cat-map-327}
\end{equation}
\begin{equation}
\shortstack[l]{$\displaystyle B^{1300;1}_{-2}=R_{192}$\\
$\displaystyle B^{1300;1}_{-1}=R_{193}$}
\label{eq:cat-map-328}
\end{equation}
\begin{equation}
B^{1300;2}=0.
\label{eq:cat-map-329}
\end{equation}
\subsubsection{Pair 13:01}
\begin{equation}
\shortstack[l]{$\displaystyle A^{1301;0}_{-3}=R_{194}$\\
$\displaystyle A^{1301;0}_{-2}=R_{195}$\\
$\displaystyle A^{1301;0}_{-1}=R_{196}$}
\label{eq:cat-map-330}
\end{equation}
\begin{equation}
\shortstack[l]{$\displaystyle A^{1301;1}_{-3}=R_{197}$\\
$\displaystyle A^{1301;1}_{-2}=R_{198}$\\
$\displaystyle A^{1301;1}_{-1}=R_{199}$}
\label{eq:cat-map-331}
\end{equation}
\begin{equation}
A^{1301;2}=0.
\label{eq:cat-map-332}
\end{equation}
\begin{equation}
\shortstack[l]{$\displaystyle B^{1301;0}_{-2}=R_{200}$\\
$\displaystyle B^{1301;0}_{-1}=R_{201}$}
\label{eq:cat-map-333}
\end{equation}
\begin{equation}
\shortstack[l]{$\displaystyle B^{1301;1}_{-2}=R_{202}$\\
$\displaystyle B^{1301;1}_{-1}=R_{203}$}
\label{eq:cat-map-334}
\end{equation}
\begin{equation}
B^{1301;2}=0.
\label{eq:cat-map-335}
\end{equation}
\subsubsection{Pair 13:02}
\begin{equation}
A^{1302;0}=0.
\label{eq:cat-map-336}
\end{equation}
\begin{equation}
A^{1302;1}=0.
\label{eq:cat-map-337}
\end{equation}
\begin{equation}
\shortstack[l]{$\displaystyle A^{1302;2}_{-3}=R_{204}$\\
$\displaystyle A^{1302;2}_{-2}=R_{205}$\\
$\displaystyle A^{1302;2}_{-1}=R_{206}$}
\label{eq:cat-map-338}
\end{equation}
\begin{equation}
B^{1302;0}=0.
\label{eq:cat-map-339}
\end{equation}
\begin{equation}
B^{1302;1}=0.
\label{eq:cat-map-340}
\end{equation}
\begin{equation}
\shortstack[l]{$\displaystyle B^{1302;2}_{-2}=R_{207}$\\
$\displaystyle B^{1302;2}_{-1}=R_{208}$}
\label{eq:cat-map-341}
\end{equation}
\subsubsection{Pair 13:10}
\begin{equation}
\shortstack[l]{$\displaystyle A^{1310;0}_{-3}=R_{209}$\\
$\displaystyle A^{1310;0}_{-2}=R_{210}$\\
$\displaystyle A^{1310;0}_{-1}=R_{211}$}
\label{eq:cat-map-342}
\end{equation}
\begin{equation}
\shortstack[l]{$\displaystyle A^{1310;1}_{-3}=R_{212}$\\
$\displaystyle A^{1310;1}_{-2}=R_{213}$\\
$\displaystyle A^{1310;1}_{-1}=R_{214}$}
\label{eq:cat-map-343}
\end{equation}
\begin{equation}
A^{1310;2}=0.
\label{eq:cat-map-344}
\end{equation}
\begin{equation}
\shortstack[l]{$\displaystyle B^{1310;0}_{-2}=R_{215}$\\
$\displaystyle B^{1310;0}_{-1}=R_{216}$}
\label{eq:cat-map-345}
\end{equation}
\begin{equation}
\shortstack[l]{$\displaystyle B^{1310;1}_{-2}=R_{217}$\\
$\displaystyle B^{1310;1}_{-1}=R_{218}$}
\label{eq:cat-map-346}
\end{equation}
\begin{equation}
B^{1310;2}=0.
\label{eq:cat-map-347}
\end{equation}
\subsubsection{Pair 13:11}
\begin{equation}
\shortstack[l]{$\displaystyle A^{1311;0}_{-3}=R_{219}$\\
$\displaystyle A^{1311;0}_{-2}=R_{220}$\\
$\displaystyle A^{1311;0}_{-1}=R_{221}$}
\label{eq:cat-map-348}
\end{equation}
\begin{equation}
\shortstack[l]{$\displaystyle A^{1311;1}_{-1}=R_{222}$}
\label{eq:cat-map-349}
\end{equation}
\begin{equation}
A^{1311;2}=0.
\label{eq:cat-map-350}
\end{equation}
\begin{equation}
\shortstack[l]{$\displaystyle B^{1311;0}_{-2}=R_{223}$\\
$\displaystyle B^{1311;0}_{-1}=R_{224}$}
\label{eq:cat-map-351}
\end{equation}
\begin{equation}
B^{1311;1}=0.
\label{eq:cat-map-352}
\end{equation}
\begin{equation}
B^{1311;2}=0.
\label{eq:cat-map-353}
\end{equation}
\subsubsection{Pair 13:12}
\begin{equation}
A^{1312;0}=0.
\label{eq:cat-map-354}
\end{equation}
\begin{equation}
A^{1312;1}=0.
\label{eq:cat-map-355}
\end{equation}
\begin{equation}
\shortstack[l]{$\displaystyle A^{1312;2}_{-1}=R_{225}$}
\label{eq:cat-map-356}
\end{equation}
\begin{equation}
B^{1312;0}=0.
\label{eq:cat-map-357}
\end{equation}
\begin{equation}
B^{1312;1}=0.
\label{eq:cat-map-358}
\end{equation}
\begin{equation}
B^{1312;2}=0.
\label{eq:cat-map-359}
\end{equation}
\subsubsection{Pair 13:20}
\begin{equation}
A^{1320;0}=0.
\label{eq:cat-map-360}
\end{equation}
\begin{equation}
A^{1320;1}=0.
\label{eq:cat-map-361}
\end{equation}
\begin{equation}
\shortstack[l]{$\displaystyle A^{1320;2}_{-3}=R_{226}$\\
$\displaystyle A^{1320;2}_{-2}=R_{227}$\\
$\displaystyle A^{1320;2}_{-1}=R_{228}$}
\label{eq:cat-map-362}
\end{equation}
\begin{equation}
B^{1320;0}=0.
\label{eq:cat-map-363}
\end{equation}
\begin{equation}
B^{1320;1}=0.
\label{eq:cat-map-364}
\end{equation}
\begin{equation}
\shortstack[l]{$\displaystyle B^{1320;2}_{-2}=R_{229}$\\
$\displaystyle B^{1320;2}_{-1}=R_{230}$}
\label{eq:cat-map-365}
\end{equation}
\subsubsection{Pair 13:21}
\begin{equation}
A^{1321;0}=0.
\label{eq:cat-map-366}
\end{equation}
\begin{equation}
A^{1321;1}=0.
\label{eq:cat-map-367}
\end{equation}
\begin{equation}
\shortstack[l]{$\displaystyle A^{1321;2}_{-1}=R_{231}$}
\label{eq:cat-map-368}
\end{equation}
\begin{equation}
B^{1321;0}=0.
\label{eq:cat-map-369}
\end{equation}
\begin{equation}
B^{1321;1}=0.
\label{eq:cat-map-370}
\end{equation}
\begin{equation}
B^{1321;2}=0.
\label{eq:cat-map-371}
\end{equation}
\subsubsection{Pair 13:22}
\begin{equation}
\shortstack[l]{$\displaystyle A^{1322;0}_{-3}=R_{232}$\\
$\displaystyle A^{1322;0}_{-2}=R_{233}$\\
$\displaystyle A^{1322;0}_{-1}=R_{234}$}
\label{eq:cat-map-372}
\end{equation}
\begin{equation}
\shortstack[l]{$\displaystyle A^{1322;1}_{-1}=R_{235}$}
\label{eq:cat-map-373}
\end{equation}
\begin{equation}
A^{1322;2}=0.
\label{eq:cat-map-374}
\end{equation}
\begin{equation}
\shortstack[l]{$\displaystyle B^{1322;0}_{-2}=R_{236}$\\
$\displaystyle B^{1322;0}_{-1}=R_{237}$}
\label{eq:cat-map-375}
\end{equation}
\begin{equation}
B^{1322;1}=0.
\label{eq:cat-map-376}
\end{equation}
\begin{equation}
B^{1322;2}=0.
\label{eq:cat-map-377}
\end{equation}
\subsubsection{Pair 23:00}
\begin{equation}
A^{2300;0}=0.
\label{eq:cat-map-378}
\end{equation}
\begin{equation}
\shortstack[l]{$\displaystyle A^{2300;1}_{-3}=R_{290}$\\
$\displaystyle A^{2300;1}_{-2}=R_{291}$\\
$\displaystyle A^{2300;1}_{-1}=R_{292}$}
\label{eq:cat-map-379}
\end{equation}
\begin{equation}
\shortstack[l]{$\displaystyle A^{2300;2}_{-3}=R_{293}$\\
$\displaystyle A^{2300;2}_{-2}=R_{294}$\\
$\displaystyle A^{2300;2}_{-1}=R_{295}$}
\label{eq:cat-map-380}
\end{equation}
\begin{equation}
A^{2300;3}=0.
\label{eq:cat-map-381}
\end{equation}
\begin{equation}
\shortstack[l]{$\displaystyle A^{2300;4}_{-3}=R_{296}$\\
$\displaystyle A^{2300;4}_{-2}=R_{297}$\\
$\displaystyle A^{2300;4}_{-1}=R_{298}$}
\label{eq:cat-map-382}
\end{equation}
\begin{equation}
A^{2300;5}=0.
\label{eq:cat-map-383}
\end{equation}
\begin{equation}
\shortstack[l]{$\displaystyle B^{2300;0}_{-2}=R_{140}$\\
$\displaystyle B^{2300;0}_{-1}=R_{299}$\\
$\displaystyle B^{2300;0}_{0}=R_{300}$}
\label{eq:cat-map-384}
\end{equation}
\begin{equation}
\shortstack[l]{$\displaystyle B^{2300;1}_{-2}=R_{301}$\\
$\displaystyle B^{2300;1}_{-1}=R_{302}$}
\label{eq:cat-map-385}
\end{equation}
\begin{equation}
\shortstack[l]{$\displaystyle B^{2300;2}_{-2}=R_{303}$\\
$\displaystyle B^{2300;2}_{-1}=R_{304}$}
\label{eq:cat-map-386}
\end{equation}
\begin{equation}
B^{2300;3}=0.
\label{eq:cat-map-387}
\end{equation}
\begin{equation}
\shortstack[l]{$\displaystyle B^{2300;4}_{-2}=R_{305}$\\
$\displaystyle B^{2300;4}_{-1}=R_{306}$}
\label{eq:cat-map-388}
\end{equation}
\begin{equation}
B^{2300;5}=0.
\label{eq:cat-map-389}
\end{equation}
\subsubsection{Pair 23:01}
\begin{equation}
A^{2301;0}=0.
\label{eq:cat-map-390}
\end{equation}
\begin{equation}
\shortstack[l]{$\displaystyle A^{2301;1}_{-3}=R_{307}$\\
$\displaystyle A^{2301;1}_{-2}=R_{308}$}
\label{eq:cat-map-391}
\end{equation}
\begin{equation}
\shortstack[l]{$\displaystyle A^{2301;2}_{-3}=R_{309}$\\
$\displaystyle A^{2301;2}_{-2}=R_{310}$}
\label{eq:cat-map-392}
\end{equation}
\begin{equation}
A^{2301;3}=0.
\label{eq:cat-map-393}
\end{equation}
\begin{equation}
\shortstack[l]{$\displaystyle A^{2301;4}_{-3}=R_{311}$\\
$\displaystyle A^{2301;4}_{-2}=R_{312}$}
\label{eq:cat-map-394}
\end{equation}
\begin{equation}
A^{2301;5}=0.
\label{eq:cat-map-395}
\end{equation}
\begin{equation}
B^{2301;0}=0.
\label{eq:cat-map-396}
\end{equation}
\begin{equation}
\shortstack[l]{$\displaystyle B^{2301;1}_{-2}=R_{219}$\\
$\displaystyle B^{2301;1}_{-1}=R_{313}$}
\label{eq:cat-map-397}
\end{equation}
\begin{equation}
\shortstack[l]{$\displaystyle B^{2301;2}_{-2}=R_{314}$\\
$\displaystyle B^{2301;2}_{-1}=R_{315}$}
\label{eq:cat-map-398}
\end{equation}
\begin{equation}
B^{2301;3}=0.
\label{eq:cat-map-399}
\end{equation}
\begin{equation}
\shortstack[l]{$\displaystyle B^{2301;4}_{-2}=R_{316}$\\
$\displaystyle B^{2301;4}_{-1}=R_{317}$}
\label{eq:cat-map-400}
\end{equation}
\begin{equation}
B^{2301;5}=0.
\label{eq:cat-map-401}
\end{equation}
\subsubsection{Pair 23:02}
\begin{equation}
A^{2302;0}=0.
\label{eq:cat-map-402}
\end{equation}
\begin{equation}
A^{2302;1}=0.
\label{eq:cat-map-403}
\end{equation}
\begin{equation}
A^{2302;2}=0.
\label{eq:cat-map-404}
\end{equation}
\begin{equation}
\shortstack[l]{$\displaystyle A^{2302;3}_{-3}=R_{232}$\\
$\displaystyle A^{2302;3}_{-2}=R_{318}$}
\label{eq:cat-map-405}
\end{equation}
\begin{equation}
A^{2302;4}=0.
\label{eq:cat-map-406}
\end{equation}
\begin{equation}
\shortstack[l]{$\displaystyle A^{2302;5}_{-3}=R_{319}$\\
$\displaystyle A^{2302;5}_{-2}=R_{320}$}
\label{eq:cat-map-407}
\end{equation}
\begin{equation}
B^{2302;0}=0.
\label{eq:cat-map-408}
\end{equation}
\begin{equation}
B^{2302;1}=0.
\label{eq:cat-map-409}
\end{equation}
\begin{equation}
B^{2302;2}=0.
\label{eq:cat-map-410}
\end{equation}
\begin{equation}
\shortstack[l]{$\displaystyle B^{2302;3}_{-2}=R_{284}$\\
$\displaystyle B^{2302;3}_{-1}=R_{321}$}
\label{eq:cat-map-411}
\end{equation}
\begin{equation}
B^{2302;4}=0.
\label{eq:cat-map-412}
\end{equation}
\begin{equation}
\shortstack[l]{$\displaystyle B^{2302;5}_{-2}=R_{322}$\\
$\displaystyle B^{2302;5}_{-1}=R_{323}$}
\label{eq:cat-map-413}
\end{equation}
\subsubsection{Pair 23:10}
\begin{equation}
A^{2310;0}=0.
\label{eq:cat-map-414}
\end{equation}
\begin{equation}
\shortstack[l]{$\displaystyle A^{2310;1}_{-3}=R_{307}$\\
$\displaystyle A^{2310;1}_{-2}=R_{308}$}
\label{eq:cat-map-415}
\end{equation}
\begin{equation}
\shortstack[l]{$\displaystyle A^{2310;2}_{-3}=R_{309}$\\
$\displaystyle A^{2310;2}_{-2}=R_{310}$}
\label{eq:cat-map-416}
\end{equation}
\begin{equation}
A^{2310;3}=0.
\label{eq:cat-map-417}
\end{equation}
\begin{equation}
\shortstack[l]{$\displaystyle A^{2310;4}_{-3}=R_{311}$\\
$\displaystyle A^{2310;4}_{-2}=R_{312}$}
\label{eq:cat-map-418}
\end{equation}
\begin{equation}
A^{2310;5}=0.
\label{eq:cat-map-419}
\end{equation}
\begin{equation}
B^{2310;0}=0.
\label{eq:cat-map-420}
\end{equation}
\begin{equation}
\shortstack[l]{$\displaystyle B^{2310;1}_{-2}=R_{219}$\\
$\displaystyle B^{2310;1}_{-1}=R_{313}$}
\label{eq:cat-map-421}
\end{equation}
\begin{equation}
\shortstack[l]{$\displaystyle B^{2310;2}_{-2}=R_{314}$\\
$\displaystyle B^{2310;2}_{-1}=R_{315}$}
\label{eq:cat-map-422}
\end{equation}
\begin{equation}
B^{2310;3}=0.
\label{eq:cat-map-423}
\end{equation}
\begin{equation}
\shortstack[l]{$\displaystyle B^{2310;4}_{-2}=R_{316}$\\
$\displaystyle B^{2310;4}_{-1}=R_{317}$}
\label{eq:cat-map-424}
\end{equation}
\begin{equation}
B^{2310;5}=0.
\label{eq:cat-map-425}
\end{equation}
\subsubsection{Pair 23:11}
\begin{equation}
A^{2311;0}=0.
\label{eq:cat-map-426}
\end{equation}
\begin{equation}
\shortstack[l]{$\displaystyle A^{2311;1}_{-1}=R_{324}$}
\label{eq:cat-map-427}
\end{equation}
\begin{equation}
\shortstack[l]{$\displaystyle A^{2311;2}_{-1}=R_{325}$}
\label{eq:cat-map-428}
\end{equation}
\begin{equation}
A^{2311;3}=0.
\label{eq:cat-map-429}
\end{equation}
\begin{equation}
\shortstack[l]{$\displaystyle A^{2311;4}_{-1}=R_{326}$}
\label{eq:cat-map-430}
\end{equation}
\begin{equation}
A^{2311;5}=0.
\label{eq:cat-map-431}
\end{equation}
\begin{equation}
\shortstack[l]{$\displaystyle B^{2311;0}_{-2}=R_{179}$\\
$\displaystyle B^{2311;0}_{-1}=R_{327}$}
\label{eq:cat-map-432}
\end{equation}
\begin{equation}
B^{2311;1}=0.
\label{eq:cat-map-433}
\end{equation}
\begin{equation}
B^{2311;2}=0.
\label{eq:cat-map-434}
\end{equation}
\begin{equation}
B^{2311;3}=0.
\label{eq:cat-map-435}
\end{equation}
\begin{equation}
B^{2311;4}=0.
\label{eq:cat-map-436}
\end{equation}
\begin{equation}
B^{2311;5}=0.
\label{eq:cat-map-437}
\end{equation}
\subsubsection{Pair 23:12}
\begin{equation}
A^{2312;0}=0.
\label{eq:cat-map-438}
\end{equation}
\begin{equation}
A^{2312;1}=0.
\label{eq:cat-map-439}
\end{equation}
\begin{equation}
A^{2312;2}=0.
\label{eq:cat-map-440}
\end{equation}
\begin{equation}
A^{2312;3}=0.
\label{eq:cat-map-441}
\end{equation}
\begin{equation}
A^{2312;4}=0.
\label{eq:cat-map-442}
\end{equation}
\begin{equation}
A^{2312;5}=0.
\label{eq:cat-map-443}
\end{equation}
\begin{equation}
B^{2312;0}=0.
\label{eq:cat-map-444}
\end{equation}
\begin{equation}
B^{2312;1}=0.
\label{eq:cat-map-445}
\end{equation}
\begin{equation}
B^{2312;2}=0.
\label{eq:cat-map-446}
\end{equation}
\begin{equation}
B^{2312;3}=0.
\label{eq:cat-map-447}
\end{equation}
\begin{equation}
B^{2312;4}=0.
\label{eq:cat-map-448}
\end{equation}
\begin{equation}
B^{2312;5}=0.
\label{eq:cat-map-449}
\end{equation}
\subsubsection{Pair 23:20}
\begin{equation}
A^{2320;0}=0.
\label{eq:cat-map-450}
\end{equation}
\begin{equation}
A^{2320;1}=0.
\label{eq:cat-map-451}
\end{equation}
\begin{equation}
A^{2320;2}=0.
\label{eq:cat-map-452}
\end{equation}
\begin{equation}
\shortstack[l]{$\displaystyle A^{2320;3}_{-3}=R_{284}$\\
$\displaystyle A^{2320;3}_{-2}=R_{321}$}
\label{eq:cat-map-453}
\end{equation}
\begin{equation}
A^{2320;4}=0.
\label{eq:cat-map-454}
\end{equation}
\begin{equation}
\shortstack[l]{$\displaystyle A^{2320;5}_{-3}=R_{322}$\\
$\displaystyle A^{2320;5}_{-2}=R_{323}$}
\label{eq:cat-map-455}
\end{equation}
\begin{equation}
B^{2320;0}=0.
\label{eq:cat-map-456}
\end{equation}
\begin{equation}
B^{2320;1}=0.
\label{eq:cat-map-457}
\end{equation}
\begin{equation}
B^{2320;2}=0.
\label{eq:cat-map-458}
\end{equation}
\begin{equation}
\shortstack[l]{$\displaystyle B^{2320;3}_{-2}=R_{232}$\\
$\displaystyle B^{2320;3}_{-1}=R_{318}$}
\label{eq:cat-map-459}
\end{equation}
\begin{equation}
B^{2320;4}=0.
\label{eq:cat-map-460}
\end{equation}
\begin{equation}
\shortstack[l]{$\displaystyle B^{2320;5}_{-2}=R_{319}$\\
$\displaystyle B^{2320;5}_{-1}=R_{320}$}
\label{eq:cat-map-461}
\end{equation}
\subsubsection{Pair 23:21}
\begin{equation}
A^{2321;0}=0.
\label{eq:cat-map-462}
\end{equation}
\begin{equation}
A^{2321;1}=0.
\label{eq:cat-map-463}
\end{equation}
\begin{equation}
A^{2321;2}=0.
\label{eq:cat-map-464}
\end{equation}
\begin{equation}
A^{2321;3}=0.
\label{eq:cat-map-465}
\end{equation}
\begin{equation}
A^{2321;4}=0.
\label{eq:cat-map-466}
\end{equation}
\begin{equation}
A^{2321;5}=0.
\label{eq:cat-map-467}
\end{equation}
\begin{equation}
B^{2321;0}=0.
\label{eq:cat-map-468}
\end{equation}
\begin{equation}
B^{2321;1}=0.
\label{eq:cat-map-469}
\end{equation}
\begin{equation}
B^{2321;2}=0.
\label{eq:cat-map-470}
\end{equation}
\begin{equation}
B^{2321;3}=0.
\label{eq:cat-map-471}
\end{equation}
\begin{equation}
B^{2321;4}=0.
\label{eq:cat-map-472}
\end{equation}
\begin{equation}
B^{2321;5}=0.
\label{eq:cat-map-473}
\end{equation}
\subsubsection{Pair 23:22}
\begin{equation}
A^{2322;0}=0.
\label{eq:cat-map-474}
\end{equation}
\begin{equation}
\shortstack[l]{$\displaystyle A^{2322;1}_{-1}=R_{328}$}
\label{eq:cat-map-475}
\end{equation}
\begin{equation}
\shortstack[l]{$\displaystyle A^{2322;2}_{-1}=R_{329}$}
\label{eq:cat-map-476}
\end{equation}
\begin{equation}
A^{2322;3}=0.
\label{eq:cat-map-477}
\end{equation}
\begin{equation}
\shortstack[l]{$\displaystyle A^{2322;4}_{-1}=R_{330}$}
\label{eq:cat-map-478}
\end{equation}
\begin{equation}
A^{2322;5}=0.
\label{eq:cat-map-479}
\end{equation}
\begin{equation}
\shortstack[l]{$\displaystyle B^{2322;0}_{-2}=R_{181}$\\
$\displaystyle B^{2322;0}_{-1}=R_{331}$}
\label{eq:cat-map-480}
\end{equation}
\begin{equation}
B^{2322;1}=0.
\label{eq:cat-map-481}
\end{equation}
\begin{equation}
B^{2322;2}=0.
\label{eq:cat-map-482}
\end{equation}
\begin{equation}
B^{2322;3}=0.
\label{eq:cat-map-483}
\end{equation}
\begin{equation}
B^{2322;4}=0.
\label{eq:cat-map-484}
\end{equation}
\begin{equation}
B^{2322;5}=0.
\label{eq:cat-map-485}
\end{equation}
\subsection{Twenty-four time-ordered exchange entries}
\label{sec:cat-exchanges}
For each set of external fields, the table below specifies all entries, including those with zero coefficients.
\begin{longtable}{cccc}
Partition & Later pair & Earlier pair & Internal indices\\\hline
$02|13$ & $02$ & $13$ & $0,1,2$\\
$02|13$ & $13$ & $02$ & $0,1,2$\\
$03|12$ & $03$ & $12$ & $0,1,2$\\
$03|12$ & $12$ & $03$ & $0,1,2$\\
$01|23$ & $01$ & $23$ & $0,1,2,3,4,5$\\
$01|23$ & $23$ & $01$ & $0,1,2,3,4,5$\\
\end{longtable}
For a line with nonzero momentum, $\nu_I=c_I|\mathbf p_i+\mathbf p_j|$. The exchange coefficients at two times are obtained by expanding the product
\begin{equation}
 \frac{\tau_I}{2\nu_I\mathcal N_I}
 [A_{a,I}(x)-i\nu_I B_{a,I}(x)]
 [A_{b,I}(y)+i\nu_I B_{b,I}(y)].
\label{eq:cat-ordered-exchange}
\end{equation}
Each coefficient of $x^py^s$ is multiplied by $-D_{ps}(\Omega_a+\nu_I,\Omega_b-\nu_I)$. The corresponding row of the table is used for the reverse time ordering.

For a homogeneous coordinate, set $X_a=A_a$, $Y_a=\sigma_I[(x-x_i)A_a+B_a]$. The covariance bilinear form (402) of the main article is used instead of \eqref{eq:cat-ordered-exchange}.
% Main article: eq:v34-global-contraction.
For the scalar, $\sigma_I=-1$, and there is no division by a negative norm; for the tensors, $\sigma_I=1$, and division by $2,6,2,2,2$ is retained. The physical covariances $C_{QQ},C_{PP},C_{QP}$ are fixed before choosing $k$, whereas the dimensionless covariances transform as $C_{qq}=kC_{QQ}$, $C_{pp}=C_{PP}/k$, $C_{qp}=C_{QP}$. The additional volume factor and the two spatial delta functions are not included in $A,B$.
\subsection{Finite angular reconstruction}
\label{sec:cat-angular}
The coefficients in Subsections~\ref{sec:cat-contacts}--\ref{sec:cat-pairs} provide the starting data for the angular reconstruction. The individual contact partitions, boundary coefficients, and finite convolutions required to complete it are given below. The following finite procedure computes each coefficient without requiring data in machine-readable form. All operations are performed before taking the limit $r\to1$.

For the partition $\varsigma=\pm1$, substitute $z_\varsigma=(1+r^2-q^2)\varsigma/(2r)$. After cancellation of common factors in the rational expression, its denominator is a product of a power of $q$ and a factor depending only on $r$. Multiplying out and collecting terms with equal powers of $q$ therefore gives the finite series
\begin{equation}
 A_{a,I,p}(r,z_\varsigma)=\sum_u a_{a,I,p,u}(r)q^u,\qquad
 B_{a,I,p}(r,z_\varsigma)=\sum_u b_{a,I,p,u}(r)q^u.
\label{eq:cat-angular-source-transfer}
\end{equation}
The left-hand sides are given explicitly in terms of the functions $R_s$. Obtaining these series requires neither differentiating the action nor solving new constraints. The absence of a nonmonomial denominator in $q$ is checked separately by the program that computes the angular coefficients.

Let $m=h_0-h_1$, $n=h_2-h_3$, $Z=(1+r^2)/(2r)$, and
\begin{equation}
 d^j_{mn}(\theta)=p_{jmn}W_{jmn}(z),\qquad
 p_{jmn}=\sqrt{(j+m)!(j-m)!(j+n)!(j-n)!}.
\label{eq:cat-angular-wigner-factor}
\end{equation}
For each admissible $b$, set $\alpha_b=j+(n-m)/2-b$, $\beta_b=(m-n)/2+b$. The finite set of coefficients $w_{2d}$ of the polynomial $W_{jmn}(z_\varsigma)=\sum_d w_{2d}q^{2d}$ is given by
\begin{equation}
\shortstack[l]{
$\displaystyle w_{2d}=\sum_b\frac{(-1)^{m-n+b}}{2^j(j+n-b)!b!(m-n+b)!(j-m-b)!}$\\
$\displaystyle \hspace{8mm}\times\sum_{u+v=d}\binom{\alpha_b}{u}\binom{\beta_b}{v}
 (1+\varsigma Z)^{\alpha_b-u}(1-\varsigma Z)^{\beta_b-v}
 \left(-\frac{\varsigma}{2r}\right)^u
 \left(\frac{\varsigma}{2r}\right)^v.$}
\label{eq:cat-angular-explicit-wigner-coefficients}
\end{equation}
The outer sum is restricted by the nonnegativity of the factorial arguments, and the inner sum by the conditions $0\leq u\leq \alpha_b$, $0\leq v\leq \beta_b$. All coefficients $w$ with odd indices are zero. Set $\mathcal N_I=n_Iq^{d_I}$, where $(n_I,d_I)=(1,0),(2,0),(2,2)$, and
\begin{equation}
 v_{a,I,p,u}^{\pm}=a_{a,I,p,u}\mp ic_I b_{a,I,p,u-1}.
\label{eq:cat-angular-velocity-shift}
\end{equation}
For each set of external fields, partition, and time ordering, the exact exchange coefficients in the linear polarization basis are given by the finite convolution
\begin{equation}
 C_{ps\ell}^{a|b,I}=\frac{\tau_I}{2c_In_I}
 \sum_{t+u+v=\ell+d_I}w_t
 v_{a,I,p,u}^{+}v_{b,I,s,v}^{-}.
\label{eq:cat-angular-exchange-convolution}
\end{equation}
The individual parts $h_p^{(d)}$ of the contact term are first reconstructed using Subsection~\ref{sec:cat-contact-partitions}. The corresponding momentum transfer is chosen for each crossed partition before adding the contributions of the two partitions. If $h_p^{(d)}(r,z_\varsigma)=\sum_u h_{p,u}^{(d)}q^u$, then
\begin{equation}
 H_{p\ell}^{(d)}=\sum_{t+u=\ell-1}w_t h_{p,u}^{(d)}.
\label{eq:cat-angular-contact-convolution}
\end{equation}
The original contact term and the homogeneous part have no soft momentum-transfer denominator. They can be integrated with respect to $z$ using the rule $\int_{-1}^{1}z^v dz=2/(v+1)$ for even $v$; for odd $v$, the integral is zero. For homogeneous exchange, this rule applies to the products $A_aA_b$, $A_aB_b$, $B_aA_b$, $B_aB_b$ with the appropriate normalization. The fixed covariance bilinear form of Subsection~\ref{sec:cat-exchanges} is then reconstructed using $X,Y$.

The transformation to helicities is performed by a finite sum over the four linear polarization indices with $C_{00}=1$, $C_{h1}=1/2$, $C_{h2}=i\operatorname{sgn}(h)/2$. Coefficients not listed are zero. The overall factors $p_{jmn}/r$, $2\pi/\sqrt{(1+\delta_{h_0h_1})(1+\delta_{h_2h_3})}$, the frequency factors in the normalization of the external one-particle states, and $g_{\mathrm R}^2k^3$ are restored according to the article. No additional tensor normalization is applied after this.

For the boundary contributions, set $V_a^\pm(x)=A_{a,I}(x)\mp ic_IqB_{a,I}(x)$ at $I=0$. Their coefficients $V_{a,I,p}^\pm=A_{a,I,p}\mp ic_IqB_{a,I,p}$ have expansions in $q$ with the coefficients given by \eqref{eq:cat-angular-velocity-shift}. The same convolutions use $f_a V_b^-$, $V_a^+f_b$, and $f_af_b$. The nonzero coefficients $f$ in the linear polarization basis at $\zeta=\varsigma z$ are
\begin{equation}
\shortstack[l]{
$\displaystyle f_{00}=\frac{6ic_0r(r-1)(1+\zeta)}{q^2},\quad
 f_{11}=\frac{6ir(r-1)(1+\zeta)(1+\zeta^2)}{q^2},$\\
$\displaystyle f_{22}=-\frac{12ir(r-1)\zeta(1+\zeta)}{q^2},\qquad f_{ab}=0\quad(a\ne b).$}
\label{eq:cat-boundary-explicit-basis}
\end{equation}
The coefficients of their finite series in $q$ replace one or both coefficients $v$ in \eqref{eq:cat-angular-exchange-convolution} at $I=0$. The coefficients are multiplied by exactly the time factors in Eqs. (390)--(391) of the main article. The factors $x_i^{-1},x_f^{-1}$ are not included in the coefficients a second time. The contributions from both time orderings are summed. The five boundary terms are given in Subsection~C.6 of the main article, Eqs. (393)--(394). That subsection also shows the exact cancellation of the leading contribution proportional to $q^{-3}$ when the boundary terms are combined with the bulk exchange.
% Main article: eq:v34r-boundary-table-mixed-dictionary; eq:v34r-boundary-table-square-dictionary; eq:boundary-five-groups-leading; eq:boundary-scalar-five-coefficients.

The coefficients, individual contact partitions, and finite convolutions given here determine all bulk and boundary entries of the following table. Equation \eqref{eq:cat-angular-explicit-wigner-coefficients}, the convolutions \eqref{eq:cat-angular-exchange-convolution}--\eqref{eq:cat-angular-contact-convolution}, and the boundary coefficients \eqref{eq:cat-boundary-explicit-basis} define an analytical procedure for any integer $ j\geq\max(|m|,|n|) $ allowed for the chosen external channels. If identical particles have equal helicities in the incoming or outgoing pair, $ j $ must be even, as stated after Eq. (369) of the main article. The four checked values $ j=0,2,3,4 $ determine only the contents of the following table, not the domain of applicability of this procedure.
% Main article: eq:v34-bose-radial-factor; eq:angular-block-complete-order.

\subsubsection{Table of 74 bulk and 74 boundary entries}
\label{sec:cat-angular-table}
In the table, $a$ is the incoming channel number and $c$ is the outgoing channel number, with numbering starting from zero. The columns give the helicity values, so the external indices and the indices $m,n$ of the Wigner polynomial are uniquely determined. Zero entries are retained.
\begin{longtable}{rrrrrrrr}
$j$ & $c$ & $a$ & $h_0,h_1$ & $h_2,h_3$ & $m$ & $n$ & Number of exchanges\\\hline
0 & 0 & 0 & $0,0$ & $0,0$ & 0 & 0 & 24\\
0 & 1 & 0 & $0,0$ & $2,2$ & 0 & 0 & 24\\
0 & 2 & 0 & $0,0$ & $-2,-2$ & 0 & 0 & 24\\
0 & 0 & 1 & $2,2$ & $0,0$ & 0 & 0 & 24\\
0 & 1 & 1 & $2,2$ & $2,2$ & 0 & 0 & 24\\
0 & 2 & 1 & $2,2$ & $-2,-2$ & 0 & 0 & 24\\
0 & 0 & 2 & $-2,-2$ & $0,0$ & 0 & 0 & 24\\
0 & 1 & 2 & $-2,-2$ & $2,2$ & 0 & 0 & 24\\
0 & 2 & 2 & $-2,-2$ & $-2,-2$ & 0 & 0 & 24\\
2 & 0 & 0 & $0,0$ & $0,0$ & 0 & 0 & 24\\
2 & 1 & 0 & $0,0$ & $0,2$ & 0 & -2 & 24\\
2 & 2 & 0 & $0,0$ & $0,-2$ & 0 & 2 & 24\\
2 & 3 & 0 & $0,0$ & $2,2$ & 0 & 0 & 24\\
2 & 4 & 0 & $0,0$ & $-2,-2$ & 0 & 0 & 24\\
2 & 0 & 1 & $0,2$ & $0,0$ & -2 & 0 & 24\\
2 & 1 & 1 & $0,2$ & $0,2$ & -2 & -2 & 24\\
2 & 2 & 1 & $0,2$ & $0,-2$ & -2 & 2 & 24\\
2 & 3 & 1 & $0,2$ & $2,2$ & -2 & 0 & 24\\
2 & 4 & 1 & $0,2$ & $-2,-2$ & -2 & 0 & 24\\
2 & 0 & 2 & $0,-2$ & $0,0$ & 2 & 0 & 24\\
2 & 1 & 2 & $0,-2$ & $0,2$ & 2 & -2 & 24\\
2 & 2 & 2 & $0,-2$ & $0,-2$ & 2 & 2 & 24\\
2 & 3 & 2 & $0,-2$ & $2,2$ & 2 & 0 & 24\\
2 & 4 & 2 & $0,-2$ & $-2,-2$ & 2 & 0 & 24\\
2 & 0 & 3 & $2,2$ & $0,0$ & 0 & 0 & 24\\
2 & 1 & 3 & $2,2$ & $0,2$ & 0 & -2 & 24\\
2 & 2 & 3 & $2,2$ & $0,-2$ & 0 & 2 & 24\\
2 & 3 & 3 & $2,2$ & $2,2$ & 0 & 0 & 24\\
2 & 4 & 3 & $2,2$ & $-2,-2$ & 0 & 0 & 24\\
2 & 0 & 4 & $-2,-2$ & $0,0$ & 0 & 0 & 24\\
2 & 1 & 4 & $-2,-2$ & $0,2$ & 0 & -2 & 24\\
2 & 2 & 4 & $-2,-2$ & $0,-2$ & 0 & 2 & 24\\
2 & 3 & 4 & $-2,-2$ & $2,2$ & 0 & 0 & 24\\
2 & 4 & 4 & $-2,-2$ & $-2,-2$ & 0 & 0 & 24\\
3 & 0 & 0 & $0,2$ & $0,2$ & -2 & -2 & 24\\
3 & 1 & 0 & $0,2$ & $0,-2$ & -2 & 2 & 24\\
3 & 0 & 1 & $0,-2$ & $0,2$ & 2 & -2 & 24\\
3 & 1 & 1 & $0,-2$ & $0,-2$ & 2 & 2 & 24\\
4 & 0 & 0 & $0,0$ & $0,0$ & 0 & 0 & 24\\
4 & 1 & 0 & $0,0$ & $0,2$ & 0 & -2 & 24\\
4 & 2 & 0 & $0,0$ & $0,-2$ & 0 & 2 & 24\\
4 & 3 & 0 & $0,0$ & $2,2$ & 0 & 0 & 24\\
4 & 4 & 0 & $0,0$ & $2,-2$ & 0 & 4 & 24\\
4 & 5 & 0 & $0,0$ & $-2,-2$ & 0 & 0 & 24\\
4 & 0 & 1 & $0,2$ & $0,0$ & -2 & 0 & 24\\
4 & 1 & 1 & $0,2$ & $0,2$ & -2 & -2 & 24\\
4 & 2 & 1 & $0,2$ & $0,-2$ & -2 & 2 & 24\\
4 & 3 & 1 & $0,2$ & $2,2$ & -2 & 0 & 24\\
4 & 4 & 1 & $0,2$ & $2,-2$ & -2 & 4 & 24\\
4 & 5 & 1 & $0,2$ & $-2,-2$ & -2 & 0 & 24\\
4 & 0 & 2 & $0,-2$ & $0,0$ & 2 & 0 & 24\\
4 & 1 & 2 & $0,-2$ & $0,2$ & 2 & -2 & 24\\
4 & 2 & 2 & $0,-2$ & $0,-2$ & 2 & 2 & 24\\
4 & 3 & 2 & $0,-2$ & $2,2$ & 2 & 0 & 24\\
4 & 4 & 2 & $0,-2$ & $2,-2$ & 2 & 4 & 24\\
4 & 5 & 2 & $0,-2$ & $-2,-2$ & 2 & 0 & 24\\
4 & 0 & 3 & $2,2$ & $0,0$ & 0 & 0 & 24\\
4 & 1 & 3 & $2,2$ & $0,2$ & 0 & -2 & 24\\
4 & 2 & 3 & $2,2$ & $0,-2$ & 0 & 2 & 24\\
4 & 3 & 3 & $2,2$ & $2,2$ & 0 & 0 & 24\\
4 & 4 & 3 & $2,2$ & $2,-2$ & 0 & 4 & 24\\
4 & 5 & 3 & $2,2$ & $-2,-2$ & 0 & 0 & 24\\
4 & 0 & 4 & $2,-2$ & $0,0$ & 4 & 0 & 24\\
4 & 1 & 4 & $2,-2$ & $0,2$ & 4 & -2 & 24\\
4 & 2 & 4 & $2,-2$ & $0,-2$ & 4 & 2 & 24\\
4 & 3 & 4 & $2,-2$ & $2,2$ & 4 & 0 & 24\\
4 & 4 & 4 & $2,-2$ & $2,-2$ & 4 & 4 & 24\\
4 & 5 & 4 & $2,-2$ & $-2,-2$ & 4 & 0 & 24\\
4 & 0 & 5 & $-2,-2$ & $0,0$ & 0 & 0 & 24\\
4 & 1 & 5 & $-2,-2$ & $0,2$ & 0 & -2 & 24\\
4 & 2 & 5 & $-2,-2$ & $0,-2$ & 0 & 2 & 24\\
4 & 3 & 5 & $-2,-2$ & $2,2$ & 0 & 0 & 24\\
4 & 4 & 5 & $-2,-2$ & $2,-2$ & 0 & 4 & 24\\
4 & 5 & 5 & $-2,-2$ & $-2,-2$ & 0 & 0 & 24\\
\end{longtable}
\subsubsection{Reconstruction of the individual contact partitions}
\label{sec:cat-contact-partitions}
The sums $u_p$ and $v_p$ in Subsection~\ref{sec:cat-contacts} do not determine the distribution of regular terms among the partitions. This distribution is specified explicitly below. The labels $0,+,-$ refer, respectively, to the partitions $01|23$, $02|13$, $03|12$ and to the momentum transfers
\begin{equation}
 q_0=0,\qquad q_+^2=1+r^2-2rz,\qquad q_-^2=1+r^2+2rz.
\label{eq:cat-partition-transfers}
\end{equation}
The coefficients decompose as
\begin{equation}
\shortstack[l]{$\displaystyle u_p=u_p^{(0)}+u_p^{(+)}+u_p^{(-)},\qquad
 v_p=v_p^{(0)}+v_p^{(+)}+v_p^{(-)},$\\
$\displaystyle h_p=b_p+u_p^{(0)}+u_p^{(+)}+u_p^{(-)}
                   +v_p^{(0)}+v_p^{(+)}+v_p^{(-)}.$}
\label{eq:cat-partition-contact-sum}
\end{equation}
The external indices $ABCD$ are omitted in these equalities. The seven terms in the second line are the individual parts of the contact term denoted by $h_p^{(d)}$ in \eqref{eq:cat-angular-contact-convolution} and in Eq. (378) of the main article.
% Main article: eq:angular-contact-finite-reconstruction.

\paragraph{Legendre-transform contact terms.}
For $a|b=02|13$ or $03|12$, use the coefficients $B_{a,I,s}$ from Subsection~\ref{sec:cat-pairs}, calculated with the internal momentum direction of the corresponding pair taken into account. Their finite convolution gives
\begin{equation}
 v_p^{(a|b)}=\sum_{I=0}^{2}\frac{\tau_I}{\mathcal N_I}
 \sum_{s+t=p}B_{a,I,s}B_{b,I,t},\qquad
 (\tau_I)=(1,1,-1),\quad (\mathcal N_I)=(1,2,2q_{a|b}^{2}).
\label{eq:cat-partition-velocity-nonzero}
\end{equation}
Here $q_{02|13}=q_+$, $q_{03|12}=q_-$. The homogeneous partition uses six coordinates and positive squared geometric norms:
\begin{equation}
\shortstack[l]{$\displaystyle v_p^{(0)}=\sum_{I=0}^{5}\frac{\sigma_I}{\mathcal N_I^{(0)}}
 \sum_{s+t=p}B_{01,I,s}B_{23,I,t},$\\
$\displaystyle (\sigma_I)=(-1,1,1,1,1,1),\qquad
 (\mathcal N_I^{(0)})=(1,2,6,2,2,2).$}
\label{eq:cat-partition-velocity-homogeneous}
\end{equation}
In both convolutions, $s,t\in\{-2,-1,0\}$; coefficients not listed are zero. The minus sign at $I=0$ in \eqref{eq:cat-partition-velocity-homogeneous} corresponds to the inverse kinetic coefficient of the homogeneous scalar coordinate. No additional factor of $1/2$ is introduced: the two cross terms of the quadratic correction arising from the Legendre transform have already been combined within each unordered partition. Nor is a second reversal of the polarizations at the other end of the internal line required, since it is already accounted for by the factor $\tau_I$.

\paragraph{Homogeneous contact term from auxiliary-variable elimination.}
Define the dimensionless Laurent polynomials $\mathcal J_{01}^{AB}(x)$ and $\mathcal J_{23}^{CD}(x)$ by
\begin{equation}
\shortstack[l]{$\displaystyle \mathcal J_{01}^{00}=-18x^{-2}-\frac{20\sqrt3\,i}{3}x^{-1}+\frac43,$\\
$\displaystyle \mathcal J_{23}^{00}=-18x^{-2}+\frac{20\sqrt3\,ir}{3}x^{-1}+\frac{4r^2}{3},$\\
$\displaystyle \mathcal J_{01}^{11}=-\mathcal J_{01}^{22}=-2x^{-2}-4ix^{-1},$\\
$\displaystyle \mathcal J_{23}^{11}=-\mathcal J_{23}^{22}=-2x^{-2}+4irx^{-1}.$}
\label{eq:cat-partition-homogeneous-polynomials}
\end{equation}
All polynomials with mixed external indices are zero. Denoting by $[x^p]$ the operation of extracting the coefficient of $x^p$, we obtain
\begin{equation}
 u_p^{(0),ABCD}=2[x^p]\bigl(\mathcal J_{01}^{AB}(x)\mathcal J_{23}^{CD}(x)\bigr).
\label{eq:cat-partition-auxiliary-homogeneous}
\end{equation}
This finite rule determines the homogeneous part independently of the behavior of the nonzero-transfer contributions as $r\to1$ and $z\to\pm1$.

\paragraph{Contact terms from auxiliary-variable elimination in the crossed partitions.}
The 46 nonzero coefficients $u_p^{(+),ABCD}$ for the 13 representatives in Subsection~\ref{sec:cat-contacts} are given below. Coefficients of all other powers, including $p=0$, are zero. The formulas retain the full rational functions, including the regular terms. For each representative, the coefficients of the second crossed partition are determined from the sum in Subsection~\ref{sec:cat-contacts}:
\begin{equation}
 u_p^{(-),ABCD}=u_p^{ABCD}-u_p^{(0),ABCD}-u_p^{(+),ABCD}.
\label{eq:cat-partition-auxiliary-minus}
\end{equation}
Thus, the coefficients below, together with Eqs. \eqref{eq:cat-partition-auxiliary-homogeneous} and \eqref{eq:cat-partition-auxiliary-minus}, specify both crossed parts without repeating the auxiliary-variable elimination.

\paragraph{Permutations and the order of the angular reconstruction.}
For $w=u$ or $w=v$, permutations of the external indices also act on the partition label. Let $\iota(0)=0$, $\iota(+)=-$, $\iota(-)=+$. Then
\begin{equation}
\shortstack[l]{$\displaystyle
 w_p^{(d),BACD}(r,z)=w_p^{(d),ABDC}(r,z)
                  =w_p^{(\iota(d)),ABCD}(r,-z),$\\
$\displaystyle
 w_p^{(d),CDAB}(r,z)=r^{4+p}\overline{w_p^{(d),ABCD}(1/r,z)}.$}
\label{eq:cat-partition-permutations}
\end{equation}
The table in Subsection~\ref{sec:cat-component-table} applies to the entire triple of partitions: an odd number of within-pair exchanges requires the simultaneous replacements $z\mapsto-z$ and $+\leftrightarrow-$. Exchanging the incoming and outgoing pairs does not change the partition label. All 40 components containing an odd number of fields with the second linear tensor polarization are zero separately in each partition.

For the part labeled $+$, substitute $z=(1+r^2-q^2)/(2r)$ into \eqref{eq:cat-angular-contact-convolution}; for the part labeled $-$, substitute $z=-(1+r^2-q^2)/(2r)$. The coefficients in $q$ are extracted separately and only then summed. If the polynomial part is also separated out, division with remainder is performed in the variable $ q $ at fixed $ r $: the numerator of the reduced fraction $ u_p^{(\pm)}(r,z_\varsigma) $ or $ v_p^{(\pm)}(r,z_\varsigma) $ is divided by its denominator after the stated substitution with the corresponding sign $ \varsigma $. As in \eqref{eq:cat-angular-source-transfer}, the denominator contains only a nonnegative integer power of $ q $ and a nonzero factor depending on $ r $. The polynomial quotient therefore consists of terms with nonnegative powers of $ q $, while the remainder divided by the denominator consists of terms with negative powers. This is equivalent to separating the finite Laurent series already obtained according to the sign of the power. The division is performed separately within each reconstructed function, and the polynomial quotient retains the same partition label. If the denominator is independent of $ q $, the entire function belongs to the quotient. The sum of negative powers alone does not replace the full function: the regular part is neither discarded nor redistributed between $ + $ and $ - $.

The terms $u_p^{(0)}$ and $v_p^{(0)}$ are retained separately from $b_p$ and from the regular parts at nonzero momentum transfer. Their volume normalization remains that of the isolated zero mode, specified by Eq. (403) of the main article. This normalization is not obtained by substituting $q=0$ into the expression for a line with nonzero momentum. The absence of a soft momentum-transfer pole in a regular term does not by itself make that term an exactly homogeneous contribution.
% Main article: eq:v34-isolated-zero-distribution.

\paragraph{Coefficients $u_p^{(+),0000}$.}
\begin{equation}
u_{-4}^{(+),0000}=\frac{1}{q_{+}^{4} r}\left[\vcenter{\hbox{\shortstack[l]{$\displaystyle  - 72 z^{3} + 10 r + 10 r^{5} + 20 r^{3} + 24 z$\\
$\displaystyle  - 264 r^{2} z^{3} - 264 r^{4} z^{3} - 104 r^{3} z^{2}$\\
$\displaystyle  - 72 r^{6} z^{3} + 24 r^{6} z + 60 r z^{2}$\\
$\displaystyle  + 60 r^{5} z^{2} + 66 r z^{4} + 66 r^{5} z^{4}$\\
$\displaystyle  + 72 r^{2} z^{5} + 72 r^{4} z^{5} + 80 r^{2} z$\\
$\displaystyle  + 80 r^{4} z + 132 r^{3} z^{4}$}}}\right]
\label{eq:cat-auxiliary-plus-0000-m4}
\end{equation}
\begin{equation}
u_{-3}^{(+),0000}=\frac{- 4 \sqrt{3} i r + 4 \sqrt{3} i}{3 q_{+}^{4}}\left[\vcenter{\hbox{\shortstack[l]{$\displaystyle 1 + r^{4} - 9 z^{4} - 6 z + 2 r^{2} + 18 z^{2} + 36 z^{3}$\\
$\displaystyle  - 42 r z^{4} - 42 r^{3} z^{4} - 36 r^{2} z^{4}$\\
$\displaystyle  - 26 r^{2} z^{2} - 18 r z^{2} - 18 r^{2} z^{5}$\\
$\displaystyle  - 18 r^{3} z^{2} - 12 r^{2} z - 9 r^{4} z^{4} - 6 r z^{3}$\\
$\displaystyle  - 6 r^{4} z - 6 r^{3} z^{3} + 8 r z + 8 r^{3} z$\\
$\displaystyle  + 18 r z^{5} + 18 r^{3} z^{5} + 18 r^{4} z^{2}$\\
$\displaystyle  + 36 r^{4} z^{3} + 90 r^{2} z^{3}$}}}\right]
\label{eq:cat-auxiliary-plus-0000-m3}
\end{equation}
\begin{equation}
u_{-2}^{(+),0000}=\frac{2 r z}{3 q_{+}^{4}}\left[\vcenter{\hbox{\shortstack[l]{$\displaystyle  - 8 - 36 z^{3} - 16 r^{2} - 8 r^{4} + 12 r + 12 r^{3}$\\
$\displaystyle  + 24 z^{2} + 36 z - 72 r^{2} z^{4} - 66 r z^{2}$\\
$\displaystyle  - 66 r^{3} z^{2} - 55 r z - 55 r^{3} z - 42 r^{2} z^{3}$\\
$\displaystyle  - 36 r^{4} z^{3} + 6 r^{2} z + 9 r z^{3} + 9 r^{3} z^{3}$\\
$\displaystyle  + 24 r^{4} z^{2} + 36 r^{4} z + 72 r z^{4}$\\
$\displaystyle  + 72 r^{3} z^{4} + 148 r^{2} z^{2}$}}}\right]
\label{eq:cat-auxiliary-plus-0000-m2}
\end{equation}
\begin{equation}
u_{-1}^{(+),0000}=- \frac{4 \sqrt{3} i r^{2} z^{2} \left(r - 1\right) \left(z + 1\right) \left(3 z - 1\right)}{3 q_{+}^{2}}
\label{eq:cat-auxiliary-plus-0000-m1}
\end{equation}
\paragraph{Coefficients $u_p^{(+),0001}$.}
\begin{equation}
u_{-4}^{(+),0001}=\frac{2 \sqrt{3} z^{2} - 2 \sqrt{3}}{q_{+}^{4} r}\left[\vcenter{\hbox{\shortstack[l]{$\displaystyle  - r - 6 z - 6 r^{3} - 5 r^{5} - 26 r^{2} z - 18 r^{6} z$\\
$\displaystyle  - 12 r^{2} z^{3} - 9 r^{5} z^{2} - 2 r^{4} z + 11 r z^{2}$\\
$\displaystyle  + 26 r^{3} z^{2} + 48 r^{4} z^{3}$}}}\right]
\label{eq:cat-auxiliary-plus-0001-m4}
\end{equation}
\begin{equation}
u_{-3}^{(+),0001}=\frac{- 2 i z^{2} + 2 i}{q_{+}^{4}}\left[\vcenter{\hbox{\shortstack[l]{$\displaystyle  - r - 6 z - 6 z^{2} - 4 r^{3} - 3 r^{5} + 8 r^{2}$\\
$\displaystyle  + 8 r^{4} + \sqrt{3} r^{5} - \sqrt{3} r - 96 r^{3} z^{3}$\\
$\displaystyle  - 24 r^{2} z - 12 r^{2} z^{3} - 10 r^{3} z$\\
$\displaystyle  - 9 r^{5} z^{2} - 6 \sqrt{3} z + 4 r^{2} z^{2} + 11 r z^{2}$\\
$\displaystyle  + 12 r z^{3} + 18 r^{4} z^{2} + 18 r^{4} z^{3} + 26 r z$\\
$\displaystyle  + 30 r^{3} z^{2} + 36 r^{5} z - \sqrt{3} r z^{2}$\\
$\displaystyle  - 8 \sqrt{3} r^{4} z - 4 \sqrt{3} r^{3} z^{2}$\\
$\displaystyle  - 3 \sqrt{3} r^{5} z^{2} + 4 \sqrt{3} r^{2} z$\\
$\displaystyle  + 6 \sqrt{3} r^{4} z^{3} + 12 \sqrt{3} r^{2} z^{3}$}}}\right]
\label{eq:cat-auxiliary-plus-0001-m3}
\end{equation}
\begin{equation}
u_{-2}^{(+),0001}=\frac{2 r z^{2} - 2 r}{3 q_{+}^{4}}\left[\vcenter{\hbox{\shortstack[l]{$\displaystyle  - 18 z^{2} + 3 r + 3 r^{2} + 3 r^{3} + 3 r^{4} + 18 z$\\
$\displaystyle  - 39 r^{2} z^{2} - 36 r^{2} z^{3} - 18 r^{2} z$\\
$\displaystyle  - 18 r^{4} z^{2} - 15 r^{3} z - 9 r^{4} z - 3 \sqrt{3} r^{2}$\\
$\displaystyle  - 3 \sqrt{3} r^{4} - 2 \sqrt{3} z + 3 \sqrt{3} r + 3 r z^{2}$\\
$\displaystyle  + 3 \sqrt{3} r^{3} + 6 \sqrt{3} z^{2} + 12 r z + 36 r z^{3}$\\
$\displaystyle  + 36 r^{3} z^{2} + 36 r^{3} z^{3} - 48 \sqrt{3} r^{2} z^{3}$\\
$\displaystyle  - 27 \sqrt{3} r z - 18 \sqrt{3} r^{4} z^{2} - 12 \sqrt{3} r z^{3}$\\
$\displaystyle  - 8 \sqrt{3} r^{2} z + 3 \sqrt{3} r^{3} z^{2} + 21 \sqrt{3} r^{4} z$\\
$\displaystyle  + 21 \sqrt{3} r^{2} z^{2} + 28 \sqrt{3} r z^{2}$\\
$\displaystyle  + 36 \sqrt{3} r^{3} z^{3}$}}}\right]
\label{eq:cat-auxiliary-plus-0001-m2}
\end{equation}
\begin{equation}
u_{-1}^{(+),0001}=\frac{- 2 i r^{2} z^{3} + 2 i r^{2} z}{3 q_{+}^{2}}\left[\vcenter{\hbox{\shortstack[l]{$\displaystyle  - 3 r - 2 \sqrt{3} - 6 \sqrt{3} z + 3 \sqrt{3} r + 9 r z + 3 \sqrt{3} r z$}}}\right]
\label{eq:cat-auxiliary-plus-0001-m1}
\end{equation}
\paragraph{Coefficients $u_p^{(+),0011}$.}
\begin{equation}
u_{-4}^{(+),0011}=\frac{- 2 z^{2} + 2}{q_{+}^{4}}\left[\vcenter{\hbox{\shortstack[l]{$\displaystyle  - 9 - 45 r^{4} - 26 r^{2} + 25 z^{2} - 324 r^{3} z^{3}$\\
$\displaystyle  - 108 r z^{3} - 60 r^{3} z - 20 r z + 144 r^{2} z^{4}$\\
$\displaystyle  + 189 r^{4} z^{2} + 234 r^{2} z^{2}$}}}\right]
\label{eq:cat-auxiliary-plus-0011-m4}
\end{equation}
\begin{equation}
u_{-3}^{(+),0011}=\frac{- 4 i z^{2} + 4 i}{q_{+}^{4}}\left[\vcenter{\hbox{\shortstack[l]{$\displaystyle  - \sqrt{3} + 9 r + 9 r^{5} + 32 r^{3} - r z^{2} - 42 r^{4} z$\\
$\displaystyle  - 28 r^{2} z - 15 \sqrt{3} r^{4} - 12 r^{2} z^{3}$\\
$\displaystyle  - 9 r^{5} z^{2} - 2 \sqrt{3} r^{2} + 9 \sqrt{3} z^{2}$\\
$\displaystyle  + 18 r^{4} z^{3} + 24 r^{3} z^{2} - 108 \sqrt{3} r^{3} z^{3}$\\
$\displaystyle  - 42 \sqrt{3} r z^{3} - 28 \sqrt{3} r^{3} z - 14 \sqrt{3} r z$\\
$\displaystyle  + 48 \sqrt{3} r^{2} z^{4} + 63 \sqrt{3} r^{4} z^{2}$\\
$\displaystyle  + 90 \sqrt{3} r^{2} z^{2}$}}}\right]
\label{eq:cat-auxiliary-plus-0011-m3}
\end{equation}
\begin{equation}
u_{-2}^{(+),0011}=\frac{2 z^{2} - 2}{q_{+}^{4}}\left[\vcenter{\hbox{\shortstack[l]{$\displaystyle  - 17 r^{4} - 3 r^{2} + 12 z^{2} - 120 r^{3} z^{3}$\\
$\displaystyle  - 48 r z^{3} - 32 r^{3} z - 24 r z + 4 \sqrt{3} r$\\
$\displaystyle  + 12 \sqrt{3} r^{5} + 30 \sqrt{3} r^{3} + 48 r^{2} z^{4}$\\
$\displaystyle  + 69 r^{4} z^{2} + 115 r^{2} z^{2} - 40 \sqrt{3} r^{4} z$\\
$\displaystyle  - 24 \sqrt{3} r^{2} z - 12 \sqrt{3} r^{5} z^{2}$\\
$\displaystyle  - 8 \sqrt{3} r^{2} z^{3} + 4 \sqrt{3} r z^{2}$\\
$\displaystyle  + 10 \sqrt{3} r^{3} z^{2} + 24 \sqrt{3} r^{4} z^{3}$}}}\right]
\label{eq:cat-auxiliary-plus-0011-m2}
\end{equation}
\begin{equation}
u_{-1}^{(+),0011}=- \frac{4 i r^{2} \left(3 r - \sqrt{3}\right) \left(z - 1\right)^{2} \left(z + 1\right)^{2}}{q_{+}^{2}}
\label{eq:cat-auxiliary-plus-0011-m1}
\end{equation}
\paragraph{Coefficients $u_p^{(+),0022}$.}
\begin{equation}
u_{-4}^{(+),0022}=\frac{32 \left(z - 1\right) \left(z + 1\right) \left(3 r z - 1\right)^{2}}{q_{+}^{2}}
\label{eq:cat-auxiliary-plus-0022-m4}
\end{equation}
\begin{equation}
u_{-3}^{(+),0022}=\frac{32 i \left(z - 1\right) \left(z + 1\right) \left(3 r z - 1\right) \left(2 \sqrt{3} r z - r - \sqrt{3}\right)}{q_{+}^{2}}
\label{eq:cat-auxiliary-plus-0022-m3}
\end{equation}
\begin{equation}
u_{-2}^{(+),0022}=- \frac{8 \left(z - 1\right) \left(z + 1\right) \left(12 r^{2} z^{2} - 4 \sqrt{3} r^{2} z + r^{2} - 12 r z + 2 \sqrt{3} r + 3\right)}{q_{+}^{2}}
\label{eq:cat-auxiliary-plus-0022-m2}
\end{equation}
\paragraph{Coefficients $u_p^{(+),0101}$.}
\begin{equation}
u_{-4}^{(+),0101}=\frac{2 \left(z^{2} + 1\right) \left(- 18 r^{4} z + 15 r^{3} z^{2} - 5 r^{3} + 18 r^{2} z^{3} - 2 r^{2} z + 15 r z^{2} - 5 r - 18 z\right)}{q_{+}^{2} r}
\label{eq:cat-auxiliary-plus-0101-m4}
\end{equation}
\begin{equation}
u_{-3}^{(+),0101}=\frac{2 i r - 2 i}{q_{+}^{4}}\left[\vcenter{\hbox{\shortstack[l]{$\displaystyle  - 1 + \sqrt{3} - r^{4} - 27 z^{4} - 18 z - 18 z^{3} - 6 r$\\
$\displaystyle  - 6 r^{3} - 4 z^{2} - 2 r^{2} + \sqrt{3} r^{4}$\\
$\displaystyle  - 36 r^{2} z^{3} - 27 r^{4} z^{4} - 20 r z^{3}$\\
$\displaystyle  - 20 r^{3} z^{3} - 18 r^{4} z - 18 r^{4} z^{3}$\\
$\displaystyle  - 6 r z^{4} - 6 \sqrt{3} z - 6 \sqrt{3} z^{3} - 6 r^{3} z^{4}$\\
$\displaystyle  - 4 r^{4} z^{2} + 2 \sqrt{3} r^{2} + 2 r^{2} z^{4}$\\
$\displaystyle  + 3 \sqrt{3} z^{4} + 4 r z + 4 r^{3} z + 4 \sqrt{3} z^{2}$\\
$\displaystyle  + 18 r^{2} z + 18 r^{2} z^{6} + 24 r z^{5}$\\
$\displaystyle  + 24 r^{3} z^{5} + 30 r^{2} z^{2} + 36 r z^{2}$\\
$\displaystyle  + 36 r^{3} z^{2} + 42 r^{2} z^{5} - 28 \sqrt{3} r^{2} z^{3}$\\
$\displaystyle  - 20 \sqrt{3} r z^{3} - 20 \sqrt{3} r^{3} z^{3}$\\
$\displaystyle  - 16 \sqrt{3} r^{2} z^{5} - 12 \sqrt{3} r z^{5} - 12 \sqrt{3} r^{2} z$\\
$\displaystyle  - 12 \sqrt{3} r^{3} z^{5} - 8 \sqrt{3} r z - 8 \sqrt{3} r^{3} z$\\
$\displaystyle  - 6 \sqrt{3} r^{4} z - 6 \sqrt{3} r^{4} z^{3}$\\
$\displaystyle  + 3 \sqrt{3} r^{4} z^{4} + 4 \sqrt{3} r^{4} z^{2}$\\
$\displaystyle  + 12 \sqrt{3} r^{2} z^{6} + 20 \sqrt{3} r z^{2} + 20 \sqrt{3} r z^{4}$\\
$\displaystyle  + 20 \sqrt{3} r^{2} z^{2} + 20 \sqrt{3} r^{3} z^{2}$\\
$\displaystyle  + 20 \sqrt{3} r^{3} z^{4} + 30 \sqrt{3} r^{2} z^{4}$}}}\right]
\label{eq:cat-auxiliary-plus-0101-m3}
\end{equation}
\begin{equation}
u_{-2}^{(+),0101}=\frac{2 r}{q_{+}^{4}}\left[\vcenter{\hbox{\shortstack[l]{$\displaystyle 1 + r^{4} - \sqrt{3} - z^{3} - 6 z^{4} - 5 z^{2} - 3 z^{5}$\\
$\displaystyle  + 2 z + 2 r^{2} + \sqrt{3} z^{3} - \sqrt{3} r^{4} - r^{4} z^{3}$\\
$\displaystyle  - 32 r^{2} z^{4} - 18 r^{2} z^{2} - 12 r^{2} z^{6}$\\
$\displaystyle  - 10 r^{2} z^{5} - 6 r^{2} z^{3} - 6 r^{4} z^{4}$\\
$\displaystyle  - 5 r^{4} z^{2} - 3 \sqrt{3} z^{2} - 3 \sqrt{3} z^{5}$\\
$\displaystyle  - 3 r^{4} z^{5} - 2 r z^{2} - 2 \sqrt{3} r^{2}$\\
$\displaystyle  - 2 r^{3} z^{2} + 2 r z + 2 \sqrt{3} r + 2 r^{3} z$\\
$\displaystyle  + 2 r^{4} z + 2 \sqrt{3} r^{3} + 4 r z^{4} + 4 \sqrt{3} z$\\
$\displaystyle  + 4 r^{2} z + 4 r^{3} z^{4} + 6 r z^{6} + 6 \sqrt{3} z^{4}$\\
$\displaystyle  + 6 r^{3} z^{6} + 18 r z^{5} + 18 r^{3} z^{5}$\\
$\displaystyle  + 20 r z^{3} + 20 r^{3} z^{3} + \sqrt{3} r z^{5}$\\
$\displaystyle  + \sqrt{3} r^{3} z^{5} + \sqrt{3} r^{4} z^{3} - 12 \sqrt{3} r z^{2}$\\
$\displaystyle  - 12 \sqrt{3} r^{2} z^{5} - 12 \sqrt{3} r^{2} z^{6}$\\
$\displaystyle  - 12 \sqrt{3} r^{3} z^{2} - 10 \sqrt{3} r^{2} z - 7 \sqrt{3} r z$\\
$\displaystyle  - 7 \sqrt{3} r^{3} z - 3 \sqrt{3} r^{4} z^{2}$\\
$\displaystyle  - 3 \sqrt{3} r^{4} z^{5} - 2 \sqrt{3} r^{2} z^{2}$\\
$\displaystyle  + 4 \sqrt{3} r^{4} z + 4 \sqrt{3} r^{2} z^{4} + 6 \sqrt{3} r z^{6}$\\
$\displaystyle  + 6 \sqrt{3} r^{3} z^{6} + 6 \sqrt{3} r^{4} z^{4}$\\
$\displaystyle  + 8 \sqrt{3} r z^{4} + 8 \sqrt{3} r^{3} z^{4} + 10 \sqrt{3} r z^{3}$\\
$\displaystyle  + 10 \sqrt{3} r^{2} z^{3} + 10 \sqrt{3} r^{3} z^{3}$}}}\right]
\label{eq:cat-auxiliary-plus-0101-m2}
\end{equation}
\begin{equation}
u_{-1}^{(+),0101}=- \frac{2 i r^{2} z \left(r - 1\right) \left(z + 1\right) \left(z^{2} + 1\right) \left(\sqrt{3} z + 3 z - 1 + \sqrt{3}\right)}{q_{+}^{2}}
\label{eq:cat-auxiliary-plus-0101-m1}
\end{equation}
\paragraph{Coefficients $u_p^{(+),0111}$.}
\begin{equation}
u_{-4}^{(+),0111}=\frac{2 \sqrt{3} \left(z - 1\right) \left(z + 1\right) \left(z^{2} + 1\right) \left(- 15 r^{2} + 6 r z - 7\right)}{q_{+}^{2}}
\label{eq:cat-auxiliary-plus-0111-m4}
\end{equation}
\begin{equation}
u_{-3}^{(+),0111}=\frac{- 2 i z^{2} + 2 i}{q_{+}^{4}}\left[\vcenter{\hbox{\shortstack[l]{$\displaystyle 3 + \sqrt{3} + 3 z^{2} + 15 r^{4} + 18 r^{2} - 36 r^{3} z$\\
$\displaystyle  - 36 r^{3} z^{3} - 34 \sqrt{3} r^{3} - 14 \sqrt{3} r - 12 r z$\\
$\displaystyle  - 12 r z^{3} - 6 \sqrt{3} r^{5} + 2 \sqrt{3} r^{2} + 9 \sqrt{3} z^{2}$\\
$\displaystyle  + 12 r^{2} z^{4} + 15 \sqrt{3} r^{4} + 15 r^{4} z^{2}$\\
$\displaystyle  + 30 r^{2} z^{2} - 72 \sqrt{3} r^{3} z^{3} - 40 \sqrt{3} r z^{3}$\\
$\displaystyle  - 28 \sqrt{3} r^{3} z^{2} - 18 \sqrt{3} r^{4} z^{3} - 8 \sqrt{3} r z$\\
$\displaystyle  - 8 \sqrt{3} r^{3} z - 6 \sqrt{3} r z^{2} + 2 \sqrt{3} r^{2} z^{3}$\\
$\displaystyle  + 6 \sqrt{3} r^{3} z^{4} + 18 \sqrt{3} r^{5} z^{2}$\\
$\displaystyle  + 30 \sqrt{3} r^{2} z^{4} + 32 \sqrt{3} r^{2} z^{2}$\\
$\displaystyle  + 34 \sqrt{3} r^{2} z + 39 \sqrt{3} r^{4} z^{2} + 46 \sqrt{3} r^{4} z$}}}\right]
\label{eq:cat-auxiliary-plus-0111-m3}
\end{equation}
\begin{equation}
u_{-2}^{(+),0111}=\frac{2 z^{2} - 2}{q_{+}^{4}}\left[\vcenter{\hbox{\shortstack[l]{$\displaystyle  - 33 r^{3} - 6 r - 6 r^{2} - 6 r^{5} + 12 z^{2}$\\
$\displaystyle  + 15 r^{4} - 72 r^{3} z^{3} - 39 r z^{3} - 18 r^{4} z^{3}$\\
$\displaystyle  - 15 \sqrt{3} r^{3} - 15 r^{2} z^{3} - 15 r^{3} z^{2}$\\
$\displaystyle  - 12 r^{3} z - 6 \sqrt{3} r^{5} - 3 r z - 2 \sqrt{3} r$\\
$\displaystyle  + 6 r z^{2} + 6 r^{3} z^{4} + 7 \sqrt{3} r^{2}$\\
$\displaystyle  + 18 r^{5} z^{2} + 21 r^{2} z + 30 r^{2} z^{4}$\\
$\displaystyle  + 36 r^{2} z^{2} + 39 r^{4} z^{2} + 42 r^{4} z - \sqrt{3} r z$\\
$\displaystyle  - \sqrt{3} r z^{3} - 22 \sqrt{3} r^{3} z^{3} - 18 \sqrt{3} r^{3} z$\\
$\displaystyle  - 10 \sqrt{3} r^{5} z - 9 \sqrt{3} r^{3} z^{2}$\\
$\displaystyle  - 6 \sqrt{3} r^{3} z^{4} - 6 \sqrt{3} r^{5} z^{2}$\\
$\displaystyle  - 6 \sqrt{3} r^{5} z^{3} + 2 \sqrt{3} r z^{2}$\\
$\displaystyle  + 2 \sqrt{3} r^{2} z^{4} + 6 \sqrt{3} r^{2} z^{3}$\\
$\displaystyle  + 10 \sqrt{3} r^{2} z + 11 \sqrt{3} r^{4} z + 12 \sqrt{3} r^{4} z^{4}$\\
$\displaystyle  + 13 \sqrt{3} r^{2} z^{2} + 15 \sqrt{3} r^{4} z^{3}$\\
$\displaystyle  + 24 \sqrt{3} r^{4} z^{2}$}}}\right]
\label{eq:cat-auxiliary-plus-0111-m2}
\end{equation}
\begin{equation}
u_{-1}^{(+),0111}=- \frac{2 i r^{2} \left(z - 1\right) \left(z + 1\right)^{2} \left(z^{2} + 1\right) \left(6 r - 3 - \sqrt{3}\right)}{q_{+}^{2}}
\label{eq:cat-auxiliary-plus-0111-m1}
\end{equation}
\paragraph{Coefficients $u_p^{(+),0122}$.}
\begin{equation}
u_{-3}^{(+),0122}=\frac{16 \sqrt{3} i \left(r + 1\right) \left(z - 1\right) \left(z + 1\right) \left(r z + 1\right) \left(3 r z - 1\right)}{q_{+}^{2}}
\label{eq:cat-auxiliary-plus-0122-m3}
\end{equation}
\begin{equation}
u_{-2}^{(+),0122}=- \frac{8 \left(r + 1\right) \left(z - 1\right) \left(z + 1\right) \left(r z + 1\right) \left(6 r z - \sqrt{3} r - 3\right)}{q_{+}^{2}}
\label{eq:cat-auxiliary-plus-0122-m2}
\end{equation}
\paragraph{Coefficients $u_p^{(+),0202}$.}
\begin{equation}
u_{-4}^{(+),0202}=- \frac{4 z \left(- 18 r^{4} z + 15 r^{3} z^{2} - 5 r^{3} + 18 r^{2} z^{3} - 2 r^{2} z + 15 r z^{2} - 5 r - 18 z\right)}{q_{+}^{2} r}
\label{eq:cat-auxiliary-plus-0202-m4}
\end{equation}
\begin{equation}
u_{-3}^{(+),0202}=\frac{- 4 i r z + 4 i z}{q_{+}^{4}}\left[\vcenter{\hbox{\shortstack[l]{$\displaystyle 5 + \sqrt{3} - 21 z^{2} - 18 z + 5 r^{4} + 6 r + 6 r^{3}$\\
$\displaystyle  + 10 r^{2} + \sqrt{3} r^{4} - 21 r^{4} z^{2} - 18 r^{4} z$\\
$\displaystyle  - 8 r z - 8 r^{3} z - 6 \sqrt{3} z - 6 r^{2} z$\\
$\displaystyle  - 4 r^{2} z^{2} + 2 \sqrt{3} r^{2} + 3 \sqrt{3} z^{2} + 6 r z^{2}$\\
$\displaystyle  + 6 r^{3} z^{2} + 12 r z^{3} + 12 r^{3} z^{3}$\\
$\displaystyle  + 18 r^{2} z^{3} + 18 r^{2} z^{4} - 16 \sqrt{3} r^{2} z^{3}$\\
$\displaystyle  - 12 \sqrt{3} r z^{3} - 12 \sqrt{3} r^{2} z - 12 \sqrt{3} r^{3} z^{3}$\\
$\displaystyle  - 8 \sqrt{3} r z - 8 \sqrt{3} r^{3} z - 6 \sqrt{3} r^{4} z$\\
$\displaystyle  + 3 \sqrt{3} r^{4} z^{2} + 12 \sqrt{3} r^{2} z^{4}$\\
$\displaystyle  + 18 \sqrt{3} r^{2} z^{2} + 20 \sqrt{3} r z^{2}$\\
$\displaystyle  + 20 \sqrt{3} r^{3} z^{2}$}}}\right]
\label{eq:cat-auxiliary-plus-0202-m3}
\end{equation}
\begin{equation}
u_{-2}^{(+),0202}=- \frac{4 r z}{q_{+}^{4}}\left[\vcenter{\hbox{\shortstack[l]{$\displaystyle 1 + r^{4} - 6 z^{2} - 3 \sqrt{3} - 3 z^{3} + 2 z + 2 r^{2}$\\
$\displaystyle  - 20 r^{2} z^{2} - 12 r^{2} z^{4} - 10 r^{2} z^{3}$\\
$\displaystyle  - 6 \sqrt{3} r^{2} - 6 r^{4} z^{2} - 3 \sqrt{3} r^{4}$\\
$\displaystyle  - 3 \sqrt{3} z^{3} - 3 r^{4} z^{3} - 2 \sqrt{3} r - 2 r z^{2}$\\
$\displaystyle  - 2 \sqrt{3} r^{3} - 2 r^{3} z^{2} + 2 r z + 2 r^{3} z$\\
$\displaystyle  + 2 r^{4} z + 4 \sqrt{3} z + 4 r^{2} z + 4 \sqrt{3} z^{2}$\\
$\displaystyle  + 6 r z^{4} + 6 r^{3} z^{4} + 18 r z^{3} + 18 r^{3} z^{3}$\\
$\displaystyle  - 12 \sqrt{3} r^{2} z^{4} - 4 \sqrt{3} r^{2} z^{3} - 3 \sqrt{3} r z$\\
$\displaystyle  - 3 \sqrt{3} r^{3} z - 3 \sqrt{3} r^{4} z^{3} - 2 \sqrt{3} r z^{2}$\\
$\displaystyle  - 2 \sqrt{3} r^{2} z - 2 \sqrt{3} r^{3} z^{2} + 4 \sqrt{3} r^{4} z$\\
$\displaystyle  + 4 \sqrt{3} r^{4} z^{2} + 5 \sqrt{3} r z^{3}$\\
$\displaystyle  + 5 \sqrt{3} r^{3} z^{3} + 6 \sqrt{3} r z^{4}$\\
$\displaystyle  + 6 \sqrt{3} r^{3} z^{4} + 12 \sqrt{3} r^{2} z^{2}$}}}\right]
\label{eq:cat-auxiliary-plus-0202-m2}
\end{equation}
\begin{equation}
u_{-1}^{(+),0202}=\frac{4 i r^{2} z^{2} \left(r - 1\right) \left(z + 1\right) \left(\sqrt{3} z + 3 z - 1 + \sqrt{3}\right)}{q_{+}^{2}}
\label{eq:cat-auxiliary-plus-0202-m1}
\end{equation}
\paragraph{Coefficients $u_p^{(+),0212}$.}
\begin{equation}
u_{-4}^{(+),0212}=- \frac{4 \sqrt{3} z \left(z - 1\right) \left(z + 1\right) \left(- 15 r^{2} + 6 r z - 7\right)}{q_{+}^{2}}
\label{eq:cat-auxiliary-plus-0212-m4}
\end{equation}
\begin{equation}
u_{-3}^{(+),0212}=\frac{- 4 i z^{2} + 4 i}{q_{+}^{4}}\left[\vcenter{\hbox{\shortstack[l]{$\displaystyle  - 3 z - 18 r^{2} z - 15 r^{4} z - 12 r^{2} z^{3}$\\
$\displaystyle  - 5 \sqrt{3} z + 4 \sqrt{3} r + 4 \sqrt{3} r^{2} + 4 \sqrt{3} r^{3}$\\
$\displaystyle  + 4 \sqrt{3} r^{4} + 12 r z^{2} + 36 r^{3} z^{2}$\\
$\displaystyle  - 27 \sqrt{3} r^{4} z - 26 \sqrt{3} r^{2} z - 22 \sqrt{3} r^{2} z^{2}$\\
$\displaystyle  - 18 \sqrt{3} r^{4} z^{2} - 6 \sqrt{3} r^{5} z$\\
$\displaystyle  - 6 \sqrt{3} r^{2} z^{3} + 10 \sqrt{3} r z + 10 \sqrt{3} r^{3} z$\\
$\displaystyle  + 18 \sqrt{3} r^{3} z^{3} + 20 \sqrt{3} r z^{2}$\\
$\displaystyle  + 36 \sqrt{3} r^{3} z^{2}$}}}\right]
\label{eq:cat-auxiliary-plus-0212-m3}
\end{equation}
\begin{equation}
u_{-2}^{(+),0212}=\frac{- 4 z^{2} + 4}{q_{+}^{4}}\left[\vcenter{\hbox{\shortstack[l]{$\displaystyle  - 6 r - 6 r^{2} - 6 r^{3} - 6 r^{4} + 6 z$\\
$\displaystyle  - 36 r^{3} z^{2} - 18 r^{3} z^{3} - 15 r z^{2}$\\
$\displaystyle  - 3 r^{3} z - 2 \sqrt{3} r^{2} - 2 \sqrt{3} r^{3} - 2 \sqrt{3} r^{4}$\\
$\displaystyle  - 2 \sqrt{3} r^{5} + 6 r^{5} z + 6 r^{2} z^{3}$\\
$\displaystyle  + 9 r^{2} z^{2} + 18 r^{4} z^{2} + 24 r^{2} z$\\
$\displaystyle  + 27 r^{4} z - \sqrt{3} r z^{2} - 18 \sqrt{3} r^{3} z^{2}$\\
$\displaystyle  - 9 \sqrt{3} r^{3} z - 6 \sqrt{3} r^{5} z - 6 \sqrt{3} r^{3} z^{3}$\\
$\displaystyle  - 6 \sqrt{3} r^{5} z^{2} + 2 \sqrt{3} r^{2} z^{3}$\\
$\displaystyle  + 6 \sqrt{3} r^{4} z + 9 \sqrt{3} r^{2} z + 10 \sqrt{3} r^{2} z^{2}$\\
$\displaystyle  + 12 \sqrt{3} r^{4} z^{3} + 15 \sqrt{3} r^{4} z^{2}$}}}\right]
\label{eq:cat-auxiliary-plus-0212-m2}
\end{equation}
\begin{equation}
u_{-1}^{(+),0212}=\frac{4 i r^{2} z \left(z - 1\right) \left(z + 1\right)^{2} \left(6 r - 3 - \sqrt{3}\right)}{q_{+}^{2}}
\label{eq:cat-auxiliary-plus-0212-m1}
\end{equation}
\paragraph{Coefficients $u_p^{(+),1111}$.}
\begin{equation}
u_{-4}^{(+),1111}=- 30 \left(z^{2} + 1\right)^{2}
\label{eq:cat-auxiliary-plus-1111-m4}
\end{equation}
\begin{equation}
u_{-3}^{(+),1111}=- \frac{12 i \left(r - 1\right) \left(z^{2} + 1\right)^{2} \left(- r^{2} + 6 r z + 4 r - 1\right)}{q_{+}^{2}}
\label{eq:cat-auxiliary-plus-1111-m3}
\end{equation}
\begin{equation}
u_{-2}^{(+),1111}=\frac{- 6 z - 6}{q_{+}^{4}}\left[\vcenter{\hbox{\shortstack[l]{$\displaystyle  - 4 z^{3} - 3 r^{2} - 3 r^{4} + 4 r + 4 r^{5} + 4 z^{2}$\\
$\displaystyle  + 22 r^{3} - 50 r^{2} z^{3} - 50 r^{4} z^{3}$\\
$\displaystyle  - 26 r^{2} z^{2} - 26 r^{4} z^{2} - 19 r^{2} z^{5}$\\
$\displaystyle  - 19 r^{4} z^{5} - 16 r z^{3} - 16 r^{5} z^{3}$\\
$\displaystyle  - 4 r^{6} z^{3} - 3 r^{2} z - 3 r^{4} z + 4 r^{6} z^{2}$\\
$\displaystyle  + 6 r^{3} z^{5} + 8 r z + 8 r z^{2} + 8 r^{5} z$\\
$\displaystyle  + 8 r^{3} z^{6} + 8 r^{5} z^{2} + 12 r z^{4}$\\
$\displaystyle  + 12 r^{3} z^{3} + 12 r^{5} z^{4} + 13 r^{2} z^{4}$\\
$\displaystyle  + 13 r^{4} z^{4} + 20 r^{3} z^{2} + 22 r^{3} z$\\
$\displaystyle  + 54 r^{3} z^{4}$}}}\right]
\label{eq:cat-auxiliary-plus-1111-m2}
\end{equation}
\begin{equation}
u_{-1}^{(+),1111}=- \frac{12 i r^{2} \left(r - 1\right) \left(z + 1\right)^{2} \left(z^{2} + 1\right)^{2}}{q_{+}^{2}}
\label{eq:cat-auxiliary-plus-1111-m1}
\end{equation}
\paragraph{Coefficients $u_p^{(+),1122}$.}
\begin{equation}
u_{-2}^{(+),1122}=- \frac{24 \left(r + 1\right)^{2} \left(z - 1\right) \left(z + 1\right) \left(r z + 1\right)^{2}}{q_{+}^{2}}
\label{eq:cat-auxiliary-plus-1122-m2}
\end{equation}
\paragraph{Coefficients $u_p^{(+),1212}$.}
\begin{equation}
u_{-4}^{(+),1212}=60 z \left(z^{2} + 1\right)
\label{eq:cat-auxiliary-plus-1212-m4}
\end{equation}
\begin{equation}
u_{-3}^{(+),1212}=\frac{24 i z \left(r - 1\right) \left(z^{2} + 1\right) \left(- r^{2} + 6 r z + 4 r - 1\right)}{q_{+}^{2}}
\label{eq:cat-auxiliary-plus-1212-m3}
\end{equation}
\begin{equation}
u_{-2}^{(+),1212}=\frac{12 z + 12}{q_{+}^{4}}\left[\vcenter{\hbox{\shortstack[l]{$\displaystyle  - 2 z^{2} + 2 r + 2 z + 2 r^{5} + 4 r^{2} + 4 r^{3}$\\
$\displaystyle  + 4 r^{4} - 25 r^{2} z^{2} - 25 r^{4} z^{2}$\\
$\displaystyle  - 15 r^{2} z^{4} - 15 r^{4} z^{4} - 6 r z^{2}$\\
$\displaystyle  - 6 r^{5} z^{2} - 5 r^{2} z - 5 r^{4} z - 3 r^{2} z^{3}$\\
$\displaystyle  - 3 r^{4} z^{3} - 2 r^{6} z^{2} + 2 r^{6} z$\\
$\displaystyle  + 2 r^{3} z^{2} + 6 r z + 6 r z^{3} + 6 r^{5} z$\\
$\displaystyle  + 6 r^{5} z^{3} + 8 r^{3} z^{5} + 14 r^{3} z^{4}$\\
$\displaystyle  + 18 r^{3} z + 26 r^{3} z^{3}$}}}\right]
\label{eq:cat-auxiliary-plus-1212-m2}
\end{equation}
\begin{equation}
u_{-1}^{(+),1212}=\frac{24 i r^{2} z \left(r - 1\right) \left(z + 1\right)^{2} \left(z^{2} + 1\right)}{q_{+}^{2}}
\label{eq:cat-auxiliary-plus-1212-m1}
\end{equation}
\paragraph{Coefficients $u_p^{(+),2222}$.}
\begin{equation}
u_{-4}^{(+),2222}=- 120 z^{2}
\label{eq:cat-auxiliary-plus-2222-m4}
\end{equation}
\begin{equation}
u_{-3}^{(+),2222}=- \frac{48 i z^{2} \left(r - 1\right) \left(- r^{2} + 6 r z + 4 r - 1\right)}{q_{+}^{2}}
\label{eq:cat-auxiliary-plus-2222-m3}
\end{equation}
\begin{equation}
u_{-2}^{(+),2222}=\frac{- 24 z - 24}{q_{+}^{4}}\left[\vcenter{\hbox{\shortstack[l]{$\displaystyle 1 + r^{6} - z + 2 r + 2 r^{5} + 3 r^{2} + 3 r^{4}$\\
$\displaystyle  + 4 r^{3} - r^{6} z - 15 r^{2} z^{3} - 15 r^{4} z^{3}$\\
$\displaystyle  - 7 r^{2} z^{2} - 7 r^{4} z^{2} - 4 r^{3} z - 3 r^{2} z$\\
$\displaystyle  - 3 r^{4} z - 2 r z - 2 r^{5} z + 4 r z^{2}$\\
$\displaystyle  + 4 r^{5} z^{2} + 8 r^{3} z^{4} + 14 r^{3} z^{2}$\\
$\displaystyle  + 14 r^{3} z^{3}$}}}\right]
\label{eq:cat-auxiliary-plus-2222-m2}
\end{equation}
\begin{equation}
u_{-1}^{(+),2222}=- \frac{48 i r^{2} z^{2} \left(r - 1\right) \left(z + 1\right)^{2}}{q_{+}^{2}}
\label{eq:cat-auxiliary-plus-2222-m1}
\end{equation}

\subsection{All 27 partial derivatives and five local vertices}
\label{sec:cat-local-vertices}
This subsection applies to a general potential $ U\in C^6 $ on the regular local branch $ \chi<0 $, $ \Delta_\Phi\ne0 $, not only to the radiation-like background. The function $ \mathcal W(\Phi,s) $ is obtained by algebraically eliminating $ L $ from the local density (162) of the main article. The equation for $ L $ is $ sL+L^5U'(L)+4\Phi=0 $. Here $ A_L=s+d(L^5U')/dL=L^3\Delta_\Phi $, $ f_m=d^m(L^5U')/dL^m $ for $ m=2,3,4,5 $, as in Eq. (417) of the main article.
% Main article: eq:scalar-local-action; eq:local-sixth-finite-jets.

We write $ \mathcal W_{ab}=\partial_\Phi^a\partial_s^b\mathcal W $, where $ a,b $ are nonnegative integer indices, $ 1\leq a+b\leq6 $. When differentiating with respect to one argument, the other is held fixed. After differentiation, all coefficients are evaluated on the background. In local inertial coordinates, set $ \phi=\delta\Phi $, $ s=\dot{\bar\Phi}>0 $, $ u=\dot\phi $, and $ G_\phi=\delta^{ij}\partial_i\phi\partial_j\phi $. A dot denotes differentiation with respect to local time. The notation $ u $ and $ G_\phi $ applies only to this subsection. The local densities $ \mathcal L_n $ with $ n=2,\ldots,6 $ contain all terms of the corresponding order in the perturbations and are defined by Eqs. (416) and (423) of the main article.
% Main article: eq:v34-local-vertex-generator; eq:local-sixth-complete-vertices.

As a check, for $ U=0 $ we have $ f_m=0 $, $ A_L=s $, and $ \mathcal W=-s^4/(768\Phi^3) $. These local vertices do not replace the elimination of the gravitational constraints required to obtain the full fifth- and sixth-order action.
\begin{equation}
\mathcal W_{10}=\frac{1}{L^{4}}
\label{eq:cat-local-W-1-0}
\end{equation}
\begin{equation}
\mathcal W_{01}=\frac{1}{3 L^{3}}
\label{eq:cat-local-W-0-1}
\end{equation}
\begin{equation}
\mathcal W_{20}=\frac{16}{A_{L} L^{5}}
\label{eq:cat-local-W-2-0}
\end{equation}
\begin{equation}
\mathcal W_{11}=\frac{4}{A_{L} L^{4}}
\label{eq:cat-local-W-1-1}
\end{equation}
\begin{equation}
\mathcal W_{02}=\frac{1}{A_{L} L^{3}}
\label{eq:cat-local-W-0-2}
\end{equation}
\begin{equation}
\mathcal W_{30}=\frac{320 A_{L} + 64 L f_{2}}{A_{L}^{3} L^{6}}
\label{eq:cat-local-W-3-0}
\end{equation}
\begin{equation}
\mathcal W_{21}=\frac{64 A_{L} + 16 L f_{2}}{A_{L}^{3} L^{5}}
\label{eq:cat-local-W-2-1}
\end{equation}
\begin{equation}
\mathcal W_{12}=\frac{12 A_{L} + 4 L f_{2}}{A_{L}^{3} L^{4}}
\label{eq:cat-local-W-1-2}
\end{equation}
\begin{equation}
\mathcal W_{03}=\frac{2 A_{L} + L f_{2}}{A_{L}^{3} L^{3}}
\label{eq:cat-local-W-0-3}
\end{equation}
\begin{equation}
\mathcal W_{40}=\frac{7680 A_{L}^{2} - 256 A_{L} L^{2} f_{3} + 3840 A_{L} L f_{2} + 768 L^{2} f_{2}^{2}}{A_{L}^{5} L^{7}}
\label{eq:cat-local-W-4-0}
\end{equation}
\begin{equation}
\mathcal W_{31}=\frac{1280 A_{L}^{2} - 64 A_{L} L^{2} f_{3} + 768 A_{L} L f_{2} + 192 L^{2} f_{2}^{2}}{A_{L}^{5} L^{6}}
\label{eq:cat-local-W-3-1}
\end{equation}
\begin{equation}
\mathcal W_{22}=\frac{192 A_{L}^{2} - 16 A_{L} L^{2} f_{3} + 144 A_{L} L f_{2} + 48 L^{2} f_{2}^{2}}{A_{L}^{5} L^{5}}
\label{eq:cat-local-W-2-2}
\end{equation}
\begin{equation}
\mathcal W_{13}=\frac{24 A_{L}^{2} - 4 A_{L} L^{2} f_{3} + 24 A_{L} L f_{2} + 12 L^{2} f_{2}^{2}}{A_{L}^{5} L^{4}}
\label{eq:cat-local-W-1-3}
\end{equation}
\begin{equation}
\mathcal W_{04}=\frac{2 A_{L}^{2} - A_{L} L^{2} f_{3} + 3 A_{L} L f_{2} + 3 L^{2} f_{2}^{2}}{A_{L}^{5} L^{3}}
\label{eq:cat-local-W-0-4}
\end{equation}
\begin{equation}
\mathcal W_{50}=\frac{1}{A_{L}^{7} L^{8}}\left(\vcenter{\hbox{\shortstack[l]{$\displaystyle 215040 A_{L}^{3} + 15360 L^{3} f_{2}^{3} - 20480 A_{L}^{2} L^{2} f_{3}$\\
$\displaystyle + 1024 A_{L}^{2} L^{3} f_{4} + 76800 A_{L} L^{2} f_{2}^{2} + 184320 A_{L}^{2} L f_{2}$\\
$\displaystyle - 10240 A_{L} L^{3} f_{2} f_{3}$}}}\right)
\label{eq:cat-local-W-5-0}
\end{equation}
\begin{equation}
\mathcal W_{41}=\frac{1}{A_{L}^{7} L^{7}}\left(\vcenter{\hbox{\shortstack[l]{$\displaystyle 30720 A_{L}^{3} + 3840 L^{3} f_{2}^{3} - 4096 A_{L}^{2} L^{2} f_{3}$\\
$\displaystyle + 256 A_{L}^{2} L^{3} f_{4} + 15360 A_{L} L^{2} f_{2}^{2} + 30720 A_{L}^{2} L f_{2}$\\
$\displaystyle - 2560 A_{L} L^{3} f_{2} f_{3}$}}}\right)
\label{eq:cat-local-W-4-1}
\end{equation}
\begin{equation}
\mathcal W_{32}=\frac{1}{A_{L}^{7} L^{6}}\left(\vcenter{\hbox{\shortstack[l]{$\displaystyle 3840 A_{L}^{3} + 960 L^{3} f_{2}^{3} - 768 A_{L}^{2} L^{2} f_{3}$\\
$\displaystyle + 64 A_{L}^{2} L^{3} f_{4} + 2880 A_{L} L^{2} f_{2}^{2} + 4608 A_{L}^{2} L f_{2}$\\
$\displaystyle - 640 A_{L} L^{3} f_{2} f_{3}$}}}\right)
\label{eq:cat-local-W-3-2}
\end{equation}
\begin{equation}
\mathcal W_{23}=\frac{1}{A_{L}^{7} L^{5}}\left(\vcenter{\hbox{\shortstack[l]{$\displaystyle 384 A_{L}^{3} + 240 L^{3} f_{2}^{3} - 128 A_{L}^{2} L^{2} f_{3}$\\
$\displaystyle + 16 A_{L}^{2} L^{3} f_{4} + 480 A_{L} L^{2} f_{2}^{2} + 576 A_{L}^{2} L f_{2}$\\
$\displaystyle - 160 A_{L} L^{3} f_{2} f_{3}$}}}\right)
\label{eq:cat-local-W-2-3}
\end{equation}
\begin{equation}
\mathcal W_{14}=\frac{1}{A_{L}^{7} L^{4}}\left(\vcenter{\hbox{\shortstack[l]{$\displaystyle 24 A_{L}^{3} + 60 L^{3} f_{2}^{3} - 16 A_{L}^{2} L^{2} f_{3}$\\
$\displaystyle + 4 A_{L}^{2} L^{3} f_{4} + 48 A_{L}^{2} L f_{2} + 60 A_{L} L^{2} f_{2}^{2}$\\
$\displaystyle - 40 A_{L} L^{3} f_{2} f_{3}$}}}\right)
\label{eq:cat-local-W-1-4}
\end{equation}
\begin{equation}
\mathcal W_{05}=\frac{A_{L}^{2} f_{4} - 10 A_{L} f_{2} f_{3} + 15 f_{2}^{3}}{A_{L}^{7}}
\label{eq:cat-local-W-0-5}
\end{equation}
\begin{equation}
\mathcal W_{60}=\frac{1}{A_{L}^{9} L^{9}}\left(\vcenter{\hbox{\shortstack[l]{$\displaystyle 6881280 A_{L}^{4} + 430080 L^{4} f_{2}^{4} - 1228800 A_{L}^{3} L^{2} f_{3}$\\
$\displaystyle - 4096 A_{L}^{3} L^{4} f_{5} + 40960 A_{L}^{2} L^{4} f_{3}^{2} + 102400 A_{L}^{3} L^{3} f_{4}$\\
$\displaystyle + 2150400 A_{L} L^{3} f_{2}^{3} + 5529600 A_{L}^{2} L^{2} f_{2}^{2} + 8601600 A_{L}^{3} L f_{2}$\\
$\displaystyle - 1228800 A_{L}^{2} L^{3} f_{2} f_{3} - 430080 A_{L} L^{4} f_{2}^{2} f_{3} + 61440 A_{L}^{2} L^{4} f_{2} f_{4}$}}}\right)
\label{eq:cat-local-W-6-0}
\end{equation}
\begin{equation}
\mathcal W_{51}=\frac{1}{A_{L}^{9} L^{8}}\left(\vcenter{\hbox{\shortstack[l]{$\displaystyle 860160 A_{L}^{4} + 107520 L^{4} f_{2}^{4} - 204800 A_{L}^{3} L^{2} f_{3}$\\
$\displaystyle - 1024 A_{L}^{3} L^{4} f_{5} + 10240 A_{L}^{2} L^{4} f_{3}^{2} + 20480 A_{L}^{3} L^{3} f_{4}$\\
$\displaystyle + 430080 A_{L} L^{3} f_{2}^{3} + 921600 A_{L}^{2} L^{2} f_{2}^{2} + 1228800 A_{L}^{3} L f_{2}$\\
$\displaystyle - 245760 A_{L}^{2} L^{3} f_{2} f_{3} - 107520 A_{L} L^{4} f_{2}^{2} f_{3} + 15360 A_{L}^{2} L^{4} f_{2} f_{4}$}}}\right)
\label{eq:cat-local-W-5-1}
\end{equation}
\begin{equation}
\mathcal W_{42}=\frac{1}{A_{L}^{9} L^{7}}\left(\vcenter{\hbox{\shortstack[l]{$\displaystyle 92160 A_{L}^{4} + 26880 L^{4} f_{2}^{4} - 30720 A_{L}^{3} L^{2} f_{3}$\\
$\displaystyle - 256 A_{L}^{3} L^{4} f_{5} + 2560 A_{L}^{2} L^{4} f_{3}^{2} + 3840 A_{L}^{3} L^{3} f_{4}$\\
$\displaystyle + 80640 A_{L} L^{3} f_{2}^{3} + 138240 A_{L}^{2} L^{2} f_{2}^{2} + 153600 A_{L}^{3} L f_{2}$\\
$\displaystyle - 46080 A_{L}^{2} L^{3} f_{2} f_{3} - 26880 A_{L} L^{4} f_{2}^{2} f_{3} + 3840 A_{L}^{2} L^{4} f_{2} f_{4}$}}}\right)
\label{eq:cat-local-W-4-2}
\end{equation}
\begin{equation}
\mathcal W_{33}=\frac{1}{A_{L}^{9} L^{6}}\left(\vcenter{\hbox{\shortstack[l]{$\displaystyle 7680 A_{L}^{4} + 6720 L^{4} f_{2}^{4} - 3840 A_{L}^{3} L^{2} f_{3}$\\
$\displaystyle - 64 A_{L}^{3} L^{4} f_{5} + 640 A_{L}^{3} L^{3} f_{4} + 640 A_{L}^{2} L^{4} f_{3}^{2}$\\
$\displaystyle + 13440 A_{L} L^{3} f_{2}^{3} + 15360 A_{L}^{3} L f_{2} + 17280 A_{L}^{2} L^{2} f_{2}^{2}$\\
$\displaystyle - 7680 A_{L}^{2} L^{3} f_{2} f_{3} - 6720 A_{L} L^{4} f_{2}^{2} f_{3} + 960 A_{L}^{2} L^{4} f_{2} f_{4}$}}}\right)
\label{eq:cat-local-W-3-3}
\end{equation}
\begin{equation}
\mathcal W_{24}=\frac{1}{A_{L}^{9} L^{5}}\left(\vcenter{\hbox{\shortstack[l]{$\displaystyle 384 A_{L}^{4} + 1680 L^{4} f_{2}^{4} - 320 A_{L}^{3} L^{2} f_{3}$\\
$\displaystyle - 16 A_{L}^{3} L^{4} f_{5} + 80 A_{L}^{3} L^{3} f_{4} + 160 A_{L}^{2} L^{4} f_{3}^{2}$\\
$\displaystyle + 960 A_{L}^{3} L f_{2} + 1440 A_{L}^{2} L^{2} f_{2}^{2} + 1680 A_{L} L^{3} f_{2}^{3}$\\
$\displaystyle - 1680 A_{L} L^{4} f_{2}^{2} f_{3} - 960 A_{L}^{2} L^{3} f_{2} f_{3} + 240 A_{L}^{2} L^{4} f_{2} f_{4}$}}}\right)
\label{eq:cat-local-W-2-4}
\end{equation}
\begin{equation}
\mathcal W_{15}=\frac{- 4 A_{L}^{3} f_{5} + 60 A_{L}^{2} f_{2} f_{4} + 40 A_{L}^{2} f_{3}^{2} - 420 A_{L} f_{2}^{2} f_{3} + 420 f_{2}^{4}}{A_{L}^{9}}
\label{eq:cat-local-W-1-5}
\end{equation}
\begin{equation}
\mathcal W_{06}=\frac{1}{A_{L}^{9}}\left(\vcenter{\hbox{\shortstack[l]{$\displaystyle - 105 A_{L} f_{2}^{3} - 5 A_{L}^{3} f_{4} + 105 L f_{2}^{4}$\\
$\displaystyle - A_{L}^{3} L f_{5} + 10 A_{L}^{2} L f_{3}^{2} + 60 A_{L}^{2} f_{2} f_{3}$\\
$\displaystyle - 105 A_{L} L f_{2}^{2} f_{3} + 15 A_{L}^{2} L f_{2} f_{4}$}}}\right)
\label{eq:cat-local-W-0-6}
\end{equation}
\begin{equation}
\mathcal L_{2}=\left(\vcenter{\hbox{\shortstack[l]{$\displaystyle \frac{\mathcal W_{02} u^{2}}{2} + \frac{\mathcal W_{20} \phi^{2}}{2} + \mathcal W_{11} \phi u$\\
$\displaystyle - \frac{G_{\phi} \mathcal W_{01}}{2 s}$}}}\right)
\label{eq:cat-local-L-2}
\end{equation}
\begin{equation}
\mathcal L_{3}=\left(\vcenter{\hbox{\shortstack[l]{$\displaystyle \frac{\mathcal W_{03} u^{3}}{6} + \frac{\mathcal W_{30} \phi^{3}}{6} + \frac{\mathcal W_{12} \phi u^{2}}{2}$\\
$\displaystyle + \frac{\mathcal W_{21} \phi^{2} u}{2} + \frac{G_{\phi} \mathcal W_{01} u}{2 s^{2}} - \frac{G_{\phi} \mathcal W_{02} u}{2 s}$\\
$\displaystyle - \frac{G_{\phi} \mathcal W_{11} \phi}{2 s}$}}}\right)
\label{eq:cat-local-L-3}
\end{equation}
\begin{equation}
\mathcal L_{4}=\left(\vcenter{\hbox{\shortstack[l]{$\displaystyle \frac{\mathcal W_{04} u^{4}}{24} + \frac{\mathcal W_{40} \phi^{4}}{24} - \frac{G_{\phi}^{2} \mathcal W_{01}}{8 s^{3}}$\\
$\displaystyle + \frac{\mathcal W_{22} \phi^{2} u^{2}}{4} + \frac{\mathcal W_{13} \phi u^{3}}{6} + \frac{\mathcal W_{31} \phi^{3} u}{6}$\\
$\displaystyle + \frac{G_{\phi}^{2} \mathcal W_{02}}{8 s^{2}} + \frac{G_{\phi} \mathcal W_{02} u^{2}}{2 s^{2}} - \frac{G_{\phi} \mathcal W_{01} u^{2}}{2 s^{3}}$\\
$\displaystyle - \frac{G_{\phi} \mathcal W_{03} u^{2}}{4 s} - \frac{G_{\phi} \mathcal W_{21} \phi^{2}}{4 s} + \frac{G_{\phi} \mathcal W_{11} \phi u}{2 s^{2}}$\\
$\displaystyle - \frac{G_{\phi} \mathcal W_{12} \phi u}{2 s}$}}}\right)
\label{eq:cat-local-L-4}
\end{equation}
\begin{equation}
\mathcal L_{5}=\left(\vcenter{\hbox{\shortstack[l]{$\displaystyle \frac{\mathcal W_{05} u^{5}}{120} + \frac{\mathcal W_{50} \phi^{5}}{120} + \frac{\mathcal W_{23} \phi^{2} u^{3}}{12}$\\
$\displaystyle + \frac{\mathcal W_{32} \phi^{3} u^{2}}{12} + \frac{\mathcal W_{14} \phi u^{4}}{24} + \frac{\mathcal W_{41} \phi^{4} u}{24}$\\
$\displaystyle + \frac{G_{\phi} \mathcal W_{01} u^{3}}{2 s^{4}} - \frac{3 G_{\phi}^{2} \mathcal W_{02} u}{8 s^{3}} - \frac{G_{\phi} \mathcal W_{02} u^{3}}{2 s^{3}}$\\
$\displaystyle - \frac{G_{\phi}^{2} \mathcal W_{11} \phi}{8 s^{3}} - \frac{G_{\phi} \mathcal W_{04} u^{3}}{12 s} - \frac{G_{\phi} \mathcal W_{31} \phi^{3}}{12 s}$\\
$\displaystyle + \frac{G_{\phi} \mathcal W_{03} u^{3}}{4 s^{2}} + \frac{G_{\phi}^{2} \mathcal W_{03} u}{8 s^{2}} + \frac{G_{\phi}^{2} \mathcal W_{12} \phi}{8 s^{2}}$\\
$\displaystyle + \frac{3 G_{\phi}^{2} \mathcal W_{01} u}{8 s^{4}} + \frac{G_{\phi} \mathcal W_{12} \phi u^{2}}{2 s^{2}} - \frac{G_{\phi} \mathcal W_{11} \phi u^{2}}{2 s^{3}}$\\
$\displaystyle - \frac{G_{\phi} \mathcal W_{13} \phi u^{2}}{4 s} - \frac{G_{\phi} \mathcal W_{22} \phi^{2} u}{4 s} + \frac{G_{\phi} \mathcal W_{21} \phi^{2} u}{4 s^{2}}$}}}\right)
\label{eq:cat-local-L-5}
\end{equation}
\begin{equation}
\mathcal L_{6}=\left(\vcenter{\hbox{\shortstack[l]{$\displaystyle \frac{\mathcal W_{06} u^{6}}{720} + \frac{\mathcal W_{60} \phi^{6}}{720} - \frac{G_{\phi}^{3} \mathcal W_{01}}{16 s^{5}}$\\
$\displaystyle - \frac{G_{\phi}^{3} \mathcal W_{03}}{48 s^{3}} + \frac{G_{\phi}^{3} \mathcal W_{02}}{16 s^{4}} + \frac{\mathcal W_{33} \phi^{3} u^{3}}{36}$\\
$\displaystyle + \frac{\mathcal W_{24} \phi^{2} u^{4}}{48} + \frac{\mathcal W_{42} \phi^{4} u^{2}}{48} + \frac{\mathcal W_{15} \phi u^{5}}{120}$\\
$\displaystyle + \frac{\mathcal W_{51} \phi^{5} u}{120} + \frac{G_{\phi} \mathcal W_{02} u^{4}}{2 s^{4}} - \frac{5 G_{\phi}^{2} \mathcal W_{03} u^{2}}{16 s^{3}}$\\
$\displaystyle - \frac{3 G_{\phi}^{2} \mathcal W_{01} u^{2}}{4 s^{5}} - \frac{G_{\phi} \mathcal W_{01} u^{4}}{2 s^{5}} - \frac{G_{\phi} \mathcal W_{03} u^{4}}{4 s^{3}}$\\
$\displaystyle - \frac{G_{\phi}^{2} \mathcal W_{21} \phi^{2}}{16 s^{3}} - \frac{G_{\phi} \mathcal W_{05} u^{4}}{48 s} - \frac{G_{\phi} \mathcal W_{41} \phi^{4}}{48 s}$\\
$\displaystyle + \frac{G_{\phi} \mathcal W_{04} u^{4}}{12 s^{2}} + \frac{G_{\phi}^{2} \mathcal W_{04} u^{2}}{16 s^{2}} + \frac{G_{\phi}^{2} \mathcal W_{22} \phi^{2}}{16 s^{2}}$\\
$\displaystyle + \frac{3 G_{\phi}^{2} \mathcal W_{02} u^{2}}{4 s^{4}} + \frac{G_{\phi} \mathcal W_{11} \phi u^{3}}{2 s^{4}} - \frac{3 G_{\phi}^{2} \mathcal W_{12} \phi u}{8 s^{3}}$\\
$\displaystyle - \frac{G_{\phi} \mathcal W_{12} \phi u^{3}}{2 s^{3}} - \frac{G_{\phi} \mathcal W_{21} \phi^{2} u^{2}}{4 s^{3}} - \frac{G_{\phi} \mathcal W_{23} \phi^{2} u^{2}}{8 s}$\\
$\displaystyle - \frac{G_{\phi} \mathcal W_{14} \phi u^{3}}{12 s} - \frac{G_{\phi} \mathcal W_{32} \phi^{3} u}{12 s} + \frac{G_{\phi} \mathcal W_{13} \phi u^{3}}{4 s^{2}}$\\
$\displaystyle + \frac{G_{\phi} \mathcal W_{22} \phi^{2} u^{2}}{4 s^{2}} + \frac{G_{\phi}^{2} \mathcal W_{13} \phi u}{8 s^{2}} + \frac{G_{\phi} \mathcal W_{31} \phi^{3} u}{12 s^{2}}$\\
$\displaystyle + \frac{3 G_{\phi}^{2} \mathcal W_{11} \phi u}{8 s^{4}}$}}}\right)
\label{eq:cat-local-L-6}
\end{equation}
\subsection{Explicit algebraic functions}
\label{sec:cat-rationals}
The functions $ R_s(r,z) $ given below determine the coefficients on the radiation-like background in Subsections~\ref{sec:cat-contacts}--\ref{sec:cat-angular}. Their arguments satisfy the conditions $ r>0 $ and $ -1<z<1 $ of Subsection~\ref{sec:cat-scope}. The joint asymptotics in the soft momentum-transfer limit are defined for the combined transition kernel, not for an individual function $ R_s $.

For the chosen crossed partition, set $ q=q_+ $ at $ \varsigma=+1 $ and $ q=q_- $ at $ \varsigma=-1 $, where the momentum transfers are defined in \eqref{eq:cat-partition-transfers}. The limit is specified by the conditions $ q\to0^+ $, $ r=1+\mu q $, and $ z=\varsigma(1+r^2-q^2)/(2r) $. The parameter $ \mu\in[-1,1] $, the scale $ k>0 $, the chosen external fields, and the finite time interval $ 0<\eta_i<\eta_f $ are held fixed. Hence $ x_i=k\eta_i $ and $ x_f=k\eta_f $ are also fixed. In this limit, $ k'=kr\to k $ and $ z\to\varsigma $. The values $ |\mu|=1 $ are understood as boundary values approached from $ -1<z<1 $. For the momentum transfer $ q_+ $, this is the setup preceding Eq. (195) of the main article, where the catalogue partition $ 02|13 $ is denoted by $ 13|24 $.
% Main article: eq:v34r-soft-scalar-vertex; eq:v33-soft-singularity; eq:v34-transfer-variable.

Before taking the limit, the contact contribution with coefficients $ b_p,u_p^{(\pm)},v_p^{(\pm)} $ and the scalar and tensor exchange contributions for both time orderings are calculated using Eq. (192) of the main article. They are added to all five boundary corrections in Eq. (199) of the main article. The explicit corrections for each time-ordered partition are given in Eqs. (385)--(387) of the main article. The reverse time ordering is added separately. The resulting expression is $ \widetilde{\mathcal F}_{\ne0} $ in Eq. (395) of the main article. In the boundary transformation, the upper limit $ Q_* $ of the dimensionful soft band is held fixed with $ 2Q_*<k,k' $. Its lower limit $ \epsilon $ serves as a soft regulator: the contributions are first summed at $ \epsilon>0 $, and the regulator is then removed in their sum. The contact terms $ u_p^{(0)},v_p^{(0)} $ and the exactly homogeneous exchange are retained separately with the normalization (403) of the main article and are not determined by the $ q\to0 $ limit of the nonzero-transfer contributions.
% Main article: eq:v33-radial-kernel; eq:v33-boundary-order-two; eq:v34r-boundary-complete-correction; eq:v34r-boundary-mixed-terms; eq:v34r-boundary-square-terms; eq:v34r-dressed-tree-bound; eq:v34-isolated-zero-distribution.
\begin{equation}
R_{1}(r,z)=\frac{N_{1}(r,z)}{D_{1}(r,z)}.
\label{eq:cat-rational-1}
\end{equation}
\begin{equation}
N_{1}(r,z)=\shortstack[l]{$\displaystyle  - 368 - 736 r^{2} - 368 r^{4} - 288 z^{2} + 72 z^{4}$\\
$\displaystyle  - 288 r^{4} z^{2} - 48 r^{2} z^{4} + 72 r^{4} z^{4}$\\
$\displaystyle  + 1952 r^{2} z^{2}$}
\label{eq:cat-numerator-1}
\end{equation}
\begin{equation}
D_{1}(r,z)=\left(- r^{2} + 2 r z - 1\right) \left(r^{2} + 2 r z + 1\right)
\label{eq:cat-denominator-1}
\end{equation}
\begin{equation}
R_{2}(r,z)=\frac{N_{2}(r,z)}{D_{2}(r,z)}.
\label{eq:cat-rational-2}
\end{equation}
\begin{equation}
N_{2}(r,z)=\shortstack[l]{$\displaystyle 104 \sqrt{3} i - 648 \sqrt{3} i r^{5} - 440 \sqrt{3} i r^{3}$\\
$\displaystyle  - 424 \sqrt{3} i r^{7} - 112 \sqrt{3} i r - 104 \sqrt{3} i r^{9}$\\
$\displaystyle  - 24 \sqrt{3} i z^{4} + 32 \sqrt{3} i z^{2} + 112 \sqrt{3} i r^{8}$\\
$\displaystyle  + 424 \sqrt{3} i r^{2} + 440 \sqrt{3} i r^{6} + 648 \sqrt{3} i r^{4}$\\
$\displaystyle  - 2312 \sqrt{3} i r^{5} z^{4} - 1840 \sqrt{3} i r^{4} z^{2}$\\
$\displaystyle  - 960 \sqrt{3} i r^{2} z^{2} - 704 \sqrt{3} i r^{6} z^{2}$\\
$\displaystyle  - 664 \sqrt{3} i r^{6} z^{4} - 192 \sqrt{3} i r^{3} z^{6}$\\
$\displaystyle  - 160 \sqrt{3} i r^{4} z^{6} - 144 \sqrt{3} i r z^{2}$\\
$\displaystyle  - 104 \sqrt{3} i r^{2} z^{4} - 96 \sqrt{3} i r^{7} z^{6}$\\
$\displaystyle  - 48 \sqrt{3} i r^{8} z^{4} - 32 \sqrt{3} i r^{9} z^{2}$\\
$\displaystyle  + 24 \sqrt{3} i r^{9} z^{4} + 48 \sqrt{3} i r z^{4}$\\
$\displaystyle  + 96 \sqrt{3} i r^{2} z^{6} + 104 \sqrt{3} i r^{7} z^{4}$\\
$\displaystyle  + 144 \sqrt{3} i r^{8} z^{2} + 160 \sqrt{3} i r^{5} z^{6}$\\
$\displaystyle  + 192 \sqrt{3} i r^{6} z^{6} + 664 \sqrt{3} i r^{3} z^{4}$\\
$\displaystyle  + 704 \sqrt{3} i r^{3} z^{2} + 960 \sqrt{3} i r^{7} z^{2}$\\
$\displaystyle  + 1840 \sqrt{3} i r^{5} z^{2} + 2312 \sqrt{3} i r^{4} z^{4}$}
\label{eq:cat-numerator-2}
\end{equation}
\begin{equation}
D_{2}(r,z)=\left(- r^{2} + 2 r z - 1\right)^{2} \left(r^{2} + 2 r z + 1\right)^{2}
\label{eq:cat-denominator-2}
\end{equation}
\begin{equation}
R_{3}(r,z)=\frac{N_{3}(r,z)}{D_{3}(r,z)}.
\label{eq:cat-rational-3}
\end{equation}
\begin{equation}
N_{3}(r,z)=\shortstack[l]{$\displaystyle  - 80 - 836 r^{4} - 836 r^{6} - 412 r^{2} - 412 r^{8}$\\
$\displaystyle  - 80 r^{10} + 256 r + 256 r^{9} + 1048 r^{3} + 1048 r^{7}$\\
$\displaystyle  + 1584 r^{5} - 4624 r^{5} z^{2} - 2216 r^{3} z^{2}$\\
$\displaystyle  - 2216 r^{7} z^{2} - 1280 r^{4} z^{4} - 1280 r^{6} z^{4}$\\
$\displaystyle  - 768 r^{5} z^{6} - 576 r^{3} z^{4} - 576 r^{7} z^{4}$\\
$\displaystyle  - 144 r z^{4} - 144 r^{9} z^{4} - 48 r^{2} z^{4}$\\
$\displaystyle  - 48 r^{8} z^{4} + 144 r z^{2} + 144 r^{9} z^{2}$\\
$\displaystyle  + 192 r^{4} z^{6} + 192 r^{6} z^{6} + 576 r^{3} z^{6}$\\
$\displaystyle  + 576 r^{7} z^{6} + 628 r^{2} z^{2} + 628 r^{8} z^{2}$\\
$\displaystyle  + 1836 r^{4} z^{2} + 1836 r^{6} z^{2} + 5632 r^{5} z^{4}$}
\label{eq:cat-numerator-3}
\end{equation}
\begin{equation}
D_{3}(r,z)=3 \left(- r^{2} + 2 r z - 1\right)^{2} \left(r^{2} + 2 r z + 1\right)^{2}
\label{eq:cat-denominator-3}
\end{equation}
\begin{equation}
R_{4}(r,z)=\frac{N_{4}(r,z)}{D_{4}(r,z)}.
\label{eq:cat-rational-4}
\end{equation}
\begin{equation}
N_{4}(r,z)=\shortstack[l]{$\displaystyle  - 128 \sqrt{3} i r^{3} - 64 \sqrt{3} i r - 64 \sqrt{3} i r^{5}$\\
$\displaystyle  + 64 \sqrt{3} i r^{2} + 64 \sqrt{3} i r^{6} + 128 \sqrt{3} i r^{4}$\\
$\displaystyle  - 328 \sqrt{3} i r^{4} z^{2} - 120 \sqrt{3} i r^{3} z^{4}$\\
$\displaystyle  - 72 \sqrt{3} i r^{2} z^{4} - 24 \sqrt{3} i r^{5} z^{2}$\\
$\displaystyle  + 24 \sqrt{3} i r^{2} z^{2} + 72 \sqrt{3} i r^{5} z^{4}$\\
$\displaystyle  + 120 \sqrt{3} i r^{4} z^{4} + 328 \sqrt{3} i r^{3} z^{2}$}
\label{eq:cat-numerator-4}
\end{equation}
\begin{equation}
D_{4}(r,z)=9 \left(- r^{2} + 2 r z - 1\right) \left(r^{2} + 2 r z + 1\right)
\label{eq:cat-denominator-4}
\end{equation}
\begin{equation}
R_{5}(r,z)=- \frac{4 r^{2} \left(z^{2} - 3\right)}{9}
\label{eq:cat-rational-5}
\end{equation}
\begin{equation}
R_{6}(r,z)=- \frac{24 \left(3 r^{4} z^{2} + 10 r^{2} z^{2} + 26 r^{2} + 3 z^{2}\right)}{r^{2}}
\label{eq:cat-rational-6}
\end{equation}
\begin{equation}
R_{7}(r,z)=- \frac{16 \sqrt{3} i \left(r - 1\right) \left(9 r^{2} z^{2} - 6 r z^{2} - 38 r + 9 z^{2}\right)}{3 r}
\label{eq:cat-rational-7}
\end{equation}
\begin{equation}
R_{8}(r,z)=40 r^{2} z^{2} + \frac{104 r^{2}}{3} - 80 r z^{2} - 176 r + 40 z^{2} + \frac{104}{3}
\label{eq:cat-rational-8}
\end{equation}
\begin{equation}
R_{9}(r,z)=\frac{16 \sqrt{3} i r \left(r - 1\right) \left(3 z^{2} + 5\right)}{9}
\label{eq:cat-rational-9}
\end{equation}
\begin{equation}
R_{10}(r,z)=- \frac{4 r^{2} \left(z^{2} + 1\right)}{3}
\label{eq:cat-rational-10}
\end{equation}
\begin{equation}
R_{11}(r,z)=\frac{N_{11}(r,z)}{D_{11}(r,z)}.
\label{eq:cat-rational-11}
\end{equation}
\begin{equation}
N_{11}(r,z)=\shortstack[l]{$\displaystyle 668 - 444 z^{4} + 312 z^{2} + 668 r^{8} + 2672 r^{2}$\\
$\displaystyle  + 2672 r^{6} + 4008 r^{4} - 9104 r^{4} z^{2}$\\
$\displaystyle  - 4240 r^{2} z^{2} - 4240 r^{6} z^{2} - 1680 r^{2} z^{4}$\\
$\displaystyle  - 1680 r^{6} z^{4} - 444 r^{8} z^{4} + 312 r^{8} z^{2}$\\
$\displaystyle  + 1104 r^{2} z^{6} + 1104 r^{6} z^{6} + 2208 r^{4} z^{6}$\\
$\displaystyle  + 6104 r^{4} z^{4}$}
\label{eq:cat-numerator-11}
\end{equation}
\begin{equation}
D_{11}(r,z)=\left(- r^{2} + 2 r z - 1\right)^{2} \left(r^{2} + 2 r z + 1\right)^{2}
\label{eq:cat-denominator-11}
\end{equation}
\begin{equation}
R_{12}(r,z)=\frac{N_{12}(r,z)}{D_{12}(r,z)}.
\label{eq:cat-rational-12}
\end{equation}
\begin{equation}
N_{12}(r,z)=\shortstack[l]{$\displaystyle 728 \sqrt{3} i - 4368 \sqrt{3} i r^{5} - 2912 \sqrt{3} i r^{3}$\\
$\displaystyle  - 2912 \sqrt{3} i r^{7} - 728 \sqrt{3} i r - 728 \sqrt{3} i r^{9}$\\
$\displaystyle  - 72 \sqrt{3} i z^{4} + 144 \sqrt{3} i z^{2} + 728 \sqrt{3} i r^{8}$\\
$\displaystyle  + 2912 \sqrt{3} i r^{2} + 2912 \sqrt{3} i r^{6} + 4368 \sqrt{3} i r^{4}$\\
$\displaystyle  - 10064 \sqrt{3} i r^{4} z^{2} - 7136 \sqrt{3} i r^{5} z^{4}$\\
$\displaystyle  - 5056 \sqrt{3} i r^{2} z^{2} - 4384 \sqrt{3} i r^{6} z^{2}$\\
$\displaystyle  - 2496 \sqrt{3} i r^{6} z^{4} - 2208 \sqrt{3} i r^{3} z^{6}$\\
$\displaystyle  - 1920 \sqrt{3} i r^{5} z^{6} - 888 \sqrt{3} i r^{8} z^{4}$\\
$\displaystyle  - 864 \sqrt{3} i r^{2} z^{4} - 480 \sqrt{3} i r z^{2}$\\
$\displaystyle  - 288 \sqrt{3} i r^{7} z^{6} - 144 \sqrt{3} i r^{9} z^{2}$\\
$\displaystyle  + 72 \sqrt{3} i r^{9} z^{4} + 288 \sqrt{3} i r^{2} z^{6}$\\
$\displaystyle  + 480 \sqrt{3} i r^{8} z^{2} + 864 \sqrt{3} i r^{7} z^{4}$\\
$\displaystyle  + 888 \sqrt{3} i r z^{4} + 1920 \sqrt{3} i r^{4} z^{6}$\\
$\displaystyle  + 2208 \sqrt{3} i r^{6} z^{6} + 2496 \sqrt{3} i r^{3} z^{4}$\\
$\displaystyle  + 4384 \sqrt{3} i r^{3} z^{2} + 5056 \sqrt{3} i r^{7} z^{2}$\\
$\displaystyle  + 7136 \sqrt{3} i r^{4} z^{4} + 10064 \sqrt{3} i r^{5} z^{2}$}
\label{eq:cat-numerator-12}
\end{equation}
\begin{equation}
D_{12}(r,z)=3 \left(- r^{2} + 2 r z - 1\right)^{2} \left(r^{2} + 2 r z + 1\right)^{2}
\label{eq:cat-denominator-12}
\end{equation}
\begin{equation}
R_{13}(r,z)=\frac{N_{13}(r,z)}{D_{13}(r,z)}.
\label{eq:cat-rational-13}
\end{equation}
\begin{equation}
N_{13}(r,z)=\shortstack[l]{$\displaystyle  - 144 - 1440 r^{4} - 1440 r^{6} - 720 r^{2} - 720 r^{8}$\\
$\displaystyle  - 144 r^{10} + 800 r + 800 r^{9} + 3200 r^{3} + 3200 r^{7}$\\
$\displaystyle  + 4800 r^{5} - 12080 r^{5} z^{2} - 5896 r^{3} z^{2}$\\
$\displaystyle  - 5896 r^{7} z^{2} - 1008 r^{4} z^{6} - 1008 r^{6} z^{6}$\\
$\displaystyle  - 936 r^{3} z^{4} - 936 r^{7} z^{4} - 324 r^{4} z^{4}$\\
$\displaystyle  - 324 r^{6} z^{4} - 144 r z^{4} - 144 r^{9} z^{4}$\\
$\displaystyle  + 144 r z^{2} + 144 r^{9} z^{2} + 420 r^{2} z^{4}$\\
$\displaystyle  + 420 r^{8} z^{4} + 576 r^{3} z^{6} + 576 r^{7} z^{6}$\\
$\displaystyle  + 804 r^{2} z^{2} + 804 r^{8} z^{2} + 1632 r^{5} z^{6}$\\
$\displaystyle  + 2412 r^{4} z^{2} + 2412 r^{6} z^{2} + 10160 r^{5} z^{4}$}
\label{eq:cat-numerator-13}
\end{equation}
\begin{equation}
D_{13}(r,z)=3 \left(- r^{2} + 2 r z - 1\right)^{2} \left(r^{2} + 2 r z + 1\right)^{2}
\label{eq:cat-denominator-13}
\end{equation}
\begin{equation}
R_{14}(r,z)=\frac{N_{14}(r,z)}{D_{14}(r,z)}.
\label{eq:cat-rational-14}
\end{equation}
\begin{equation}
N_{14}(r,z)=\shortstack[l]{$\displaystyle  - 320 \sqrt{3} i r^{3} - 160 \sqrt{3} i r - 160 \sqrt{3} i r^{5}$\\
$\displaystyle  + 160 \sqrt{3} i r^{2} + 160 \sqrt{3} i r^{6} + 320 \sqrt{3} i r^{4}$\\
$\displaystyle  - 616 \sqrt{3} i r^{4} z^{2} - 72 \sqrt{3} i r^{2} z^{4}$\\
$\displaystyle  - 24 \sqrt{3} i r^{5} z^{2} - 24 \sqrt{3} i r^{3} z^{4}$\\
$\displaystyle  + 24 \sqrt{3} i r^{2} z^{2} + 24 \sqrt{3} i r^{4} z^{4}$\\
$\displaystyle  + 72 \sqrt{3} i r^{5} z^{4} + 616 \sqrt{3} i r^{3} z^{2}$}
\label{eq:cat-numerator-14}
\end{equation}
\begin{equation}
D_{14}(r,z)=9 \left(- r^{2} + 2 r z - 1\right) \left(r^{2} + 2 r z + 1\right)
\label{eq:cat-denominator-14}
\end{equation}
\begin{equation}
R_{15}(r,z)=\frac{32 r^{2}}{9}
\label{eq:cat-rational-15}
\end{equation}
\begin{equation}
R_{16}(r,z)=\frac{N_{16}(r,z)}{D_{16}(r,z)}.
\label{eq:cat-rational-16}
\end{equation}
\begin{equation}
N_{16}(r,z)=\shortstack[l]{$\displaystyle 72 z^{2} + 324 r^{2} + 324 r^{10} + 1296 r^{4} + 1296 r^{8}$\\
$\displaystyle  + 1944 r^{6} - 5136 r^{6} z^{2} - 2640 r^{4} z^{4}$\\
$\displaystyle  - 2640 r^{8} z^{4} - 2136 r^{4} z^{2} - 2136 r^{8} z^{2}$\\
$\displaystyle  - 204 r^{2} z^{4} - 204 r^{10} z^{4} + 72 r^{12} z^{2}$\\
$\displaystyle  + 336 r^{4} z^{6} + 336 r^{8} z^{6} + 504 r^{2} z^{2}$\\
$\displaystyle  + 504 r^{10} z^{2} + 1440 r^{6} z^{6} + 6648 r^{6} z^{4}$}
\label{eq:cat-numerator-16}
\end{equation}
\begin{equation}
D_{16}(r,z)=r^{2} \left(- r^{2} + 2 r z - 1\right)^{2} \left(r^{2} + 2 r z + 1\right)^{2}
\label{eq:cat-denominator-16}
\end{equation}
\begin{equation}
R_{17}(r,z)=\frac{N_{17}(r,z)}{D_{17}(r,z)}.
\label{eq:cat-rational-17}
\end{equation}
\begin{equation}
N_{17}(r,z)=\shortstack[l]{$\displaystyle  - 408 \sqrt{3} i r^{6} - 280 \sqrt{3} i r^{4} - 264 \sqrt{3} i r^{8}$\\
$\displaystyle  - 72 \sqrt{3} i r^{2} - 64 \sqrt{3} i r^{10} - 48 \sqrt{3} i z^{2}$\\
$\displaystyle  + 64 \sqrt{3} i r + 72 \sqrt{3} i r^{9} + 264 \sqrt{3} i r^{3}$\\
$\displaystyle  + 280 \sqrt{3} i r^{7} + 408 \sqrt{3} i r^{5}$\\
$\displaystyle  - 1512 \sqrt{3} i r^{6} z^{4} - 1240 \sqrt{3} i r^{7} z^{4}$\\
$\displaystyle  - 1056 \sqrt{3} i r^{5} z^{2} - 528 \sqrt{3} i r^{3} z^{2}$\\
$\displaystyle  - 480 \sqrt{3} i r^{6} z^{6} - 456 \sqrt{3} i r^{3} z^{4}$\\
$\displaystyle  - 256 \sqrt{3} i r^{2} z^{2} - 256 \sqrt{3} i r^{7} z^{2}$\\
$\displaystyle  - 224 \sqrt{3} i r^{4} z^{6} - 136 \sqrt{3} i r^{9} z^{4}$\\
$\displaystyle  - 64 \sqrt{3} i r^{10} z^{2} + 48 \sqrt{3} i r^{11} z^{2}$\\
$\displaystyle  + 64 \sqrt{3} i r z^{2} + 136 \sqrt{3} i r^{2} z^{4}$\\
$\displaystyle  + 224 \sqrt{3} i r^{7} z^{6} + 256 \sqrt{3} i r^{4} z^{2}$\\
$\displaystyle  + 256 \sqrt{3} i r^{9} z^{2} + 456 \sqrt{3} i r^{8} z^{4}$\\
$\displaystyle  + 480 \sqrt{3} i r^{5} z^{6} + 528 \sqrt{3} i r^{8} z^{2}$\\
$\displaystyle  + 1056 \sqrt{3} i r^{6} z^{2} + 1240 \sqrt{3} i r^{4} z^{4}$\\
$\displaystyle  + 1512 \sqrt{3} i r^{5} z^{4}$}
\label{eq:cat-numerator-17}
\end{equation}
\begin{equation}
D_{17}(r,z)=r \left(- r^{2} + 2 r z - 1\right)^{2} \left(r^{2} + 2 r z + 1\right)^{2}
\label{eq:cat-denominator-17}
\end{equation}
\begin{equation}
R_{18}(r,z)=\frac{N_{18}(r,z)}{D_{18}(r,z)}.
\label{eq:cat-rational-18}
\end{equation}
\begin{equation}
N_{18}(r,z)=\shortstack[l]{$\displaystyle  - 40 - 436 r^{4} - 436 r^{6} - 212 r^{2} - 212 r^{8}$\\
$\displaystyle  - 120 z^{2} - 48 r^{5} - 40 r^{3} - 40 r^{7} - 40 r^{10}$\\
$\displaystyle  - 16 r - 16 r^{9} - 1560 r^{3} z^{4} - 1560 r^{7} z^{4}$\\
$\displaystyle  - 720 r^{4} z^{6} - 720 r^{6} z^{6} - 120 r^{10} z^{2}$\\
$\displaystyle  + 56 r^{2} z^{2} + 56 r^{8} z^{2} + 80 r^{5} z^{4}$\\
$\displaystyle  + 240 r z^{2} + 240 r^{9} z^{2} + 260 r^{4} z^{4}$\\
$\displaystyle  + 260 r^{6} z^{4} + 416 r^{3} z^{2} + 416 r^{7} z^{2}$\\
$\displaystyle  + 448 r^{5} z^{2} + 492 r^{2} z^{4} + 492 r^{8} z^{4}$\\
$\displaystyle  + 720 r^{4} z^{2} + 720 r^{6} z^{2} + 1440 r^{5} z^{6}$}
\label{eq:cat-numerator-18}
\end{equation}
\begin{equation}
D_{18}(r,z)=3 \left(- r^{2} + 2 r z - 1\right)^{2} \left(r^{2} + 2 r z + 1\right)^{2}
\label{eq:cat-denominator-18}
\end{equation}
\begin{equation}
R_{19}(r,z)=\frac{N_{19}(r,z)}{D_{19}(r,z)}.
\label{eq:cat-rational-19}
\end{equation}
\begin{equation}
N_{19}(r,z)=\shortstack[l]{$\displaystyle  - 32 \sqrt{3} i r^{4} - 16 \sqrt{3} i r^{2} - 16 \sqrt{3} i r^{6}$\\
$\displaystyle  + 16 \sqrt{3} i r + 16 \sqrt{3} i r^{5} + 32 \sqrt{3} i r^{3}$\\
$\displaystyle  - 96 \sqrt{3} i r^{4} z^{4} - 64 \sqrt{3} i r^{3} z^{2}$\\
$\displaystyle  - 48 \sqrt{3} i r z^{2} - 48 \sqrt{3} i r^{5} z^{2}$\\
$\displaystyle  + 48 \sqrt{3} i r^{2} z^{2} + 48 \sqrt{3} i r^{6} z^{2}$\\
$\displaystyle  + 64 \sqrt{3} i r^{4} z^{2} + 96 \sqrt{3} i r^{3} z^{4}$}
\label{eq:cat-numerator-19}
\end{equation}
\begin{equation}
D_{19}(r,z)=9 \left(- r^{2} + 2 r z - 1\right) \left(r^{2} + 2 r z + 1\right)
\label{eq:cat-denominator-19}
\end{equation}
\begin{equation}
R_{20}(r,z)=\frac{8 r^{2} \left(z - 1\right) \left(z + 1\right)}{9}
\label{eq:cat-rational-20}
\end{equation}
\begin{equation}
R_{21}(r,z)=\frac{N_{21}(r,z)}{D_{21}(r,z)}.
\label{eq:cat-rational-21}
\end{equation}
\begin{equation}
N_{21}(r,z)=\shortstack[l]{$\displaystyle  - 48 \sqrt{3} - 588 \sqrt{3} r^{4} - 492 \sqrt{3} r^{6} - 292 \sqrt{3} r^{2}$\\
$\displaystyle  - 148 \sqrt{3} r^{8} + 24 \sqrt{3} z^{2} + 24 \sqrt{3} z^{4}$\\
$\displaystyle  - 1420 \sqrt{3} r^{4} z^{4} - 956 \sqrt{3} r^{6} z^{4}$\\
$\displaystyle  - 244 \sqrt{3} r^{2} z^{4} - 60 \sqrt{3} r^{8} z^{4}$\\
$\displaystyle  - 48 \sqrt{3} r^{2} z^{6} + 192 \sqrt{3} r^{6} z^{6}$\\
$\displaystyle  + 208 \sqrt{3} r^{8} z^{2} + 400 \sqrt{3} r^{4} z^{6}$\\
$\displaystyle  + 584 \sqrt{3} r^{2} z^{2} + 1256 \sqrt{3} r^{6} z^{2}$\\
$\displaystyle  + 1608 \sqrt{3} r^{4} z^{2}$}
\label{eq:cat-numerator-21}
\end{equation}
\begin{equation}
D_{21}(r,z)=\left(- r^{2} + 2 r z - 1\right)^{2} \left(r^{2} + 2 r z + 1\right)^{2}
\label{eq:cat-denominator-21}
\end{equation}
\begin{equation}
R_{22}(r,z)=\frac{N_{22}(r,z)}{D_{22}(r,z)}.
\label{eq:cat-rational-22}
\end{equation}
\begin{equation}
N_{22}(r,z)=\shortstack[l]{$\displaystyle  - 16 i - 664 i r^{4} - 636 i r^{6} - 252 i r^{2}$\\
$\displaystyle  - 208 i r^{8} - 8 i z^{2} + 24 i r + 24 i z^{4} + 60 i r^{9}$\\
$\displaystyle  + 136 i r^{3} + 208 i r^{7} + 260 i r^{5}$\\
$\displaystyle  - 2072 i r^{4} z^{4} - 1580 i r^{6} z^{4} - 704 i r^{5} z^{2}$\\
$\displaystyle  - 520 i r^{7} z^{2} - 296 i r^{3} z^{2} - 172 i r^{2} z^{4}$\\
$\displaystyle  - 144 i r^{5} z^{6} - 144 i r^{7} z^{6} - 120 i r^{8} z^{4}$\\
$\displaystyle  - 96 i r^{9} z^{2} - 48 i r^{2} z^{6} - 24 i r z^{4}$\\
$\displaystyle  + 4 \sqrt{3} i r^{9} + 36 i r^{9} z^{4} + 48 \sqrt{3} i r$\\
$\displaystyle  + 48 i r^{3} z^{6} + 60 \sqrt{3} i r^{7} + 112 i r^{3} z^{4}$\\
$\displaystyle  + 148 \sqrt{3} i r^{3} + 156 \sqrt{3} i r^{5} + 328 i r^{8} z^{2}$\\
$\displaystyle  + 384 i r^{6} z^{6} + 456 i r^{7} z^{4} + 472 i r^{2} z^{2}$\\
$\displaystyle  + 588 i r^{5} z^{4} + 736 i r^{4} z^{6} + 1832 i r^{6} z^{2}$\\
$\displaystyle  + 2000 i r^{4} z^{2} - 408 \sqrt{3} i r^{5} z^{2}$\\
$\displaystyle  - 392 \sqrt{3} i r^{3} z^{2} - 104 \sqrt{3} i r^{7} z^{2}$\\
$\displaystyle  - 96 \sqrt{3} i r^{3} z^{6} - 72 \sqrt{3} i r z^{2}$\\
$\displaystyle  - 48 \sqrt{3} i r^{7} z^{6} - 16 \sqrt{3} i r^{9} z^{2}$\\
$\displaystyle  - 16 \sqrt{3} i r^{5} z^{6} + 12 \sqrt{3} i r^{9} z^{4}$\\
$\displaystyle  + 24 \sqrt{3} i r z^{4} + 92 \sqrt{3} i r^{7} z^{4}$\\
$\displaystyle  + 268 \sqrt{3} i r^{5} z^{4} + 340 \sqrt{3} i r^{3} z^{4}$}
\label{eq:cat-numerator-22}
\end{equation}
\begin{equation}
D_{22}(r,z)=\left(- r^{2} + 2 r z - 1\right)^{2} \left(r^{2} + 2 r z + 1\right)^{2}
\label{eq:cat-denominator-22}
\end{equation}
\begin{equation}
R_{23}(r,z)=\frac{N_{23}(r,z)}{D_{23}(r,z)}.
\label{eq:cat-rational-23}
\end{equation}
\begin{equation}
N_{23}(r,z)=\shortstack[l]{$\displaystyle  - 156 r^{5} - 144 r^{3} - 72 r^{7} - 48 r - 12 r^{9} - 8 \sqrt{3}$\\
$\displaystyle  + 72 r^{2} + 84 r^{8} + 228 r^{4} + 240 r^{6}$\\
$\displaystyle  - 708 r^{4} z^{2} - 600 r^{6} z^{2} - 552 r^{3} z^{4}$\\
$\displaystyle  - 456 r^{7} z^{4} - 288 r^{4} z^{6} - 276 \sqrt{3} r^{5}$\\
$\displaystyle  - 260 \sqrt{3} r^{7} - 240 r^{5} z^{4} - 144 r^{2} z^{2}$\\
$\displaystyle  - 108 \sqrt{3} r^{3} - 96 r^{5} z^{6} - 84 \sqrt{3} r^{9}$\\
$\displaystyle  - 84 r^{8} z^{2} - 72 r z^{4} - 72 r^{9} z^{4} - 8 \sqrt{3} r$\\
$\displaystyle  + 8 \sqrt{3} z^{2} + 28 \sqrt{3} r^{2} + 52 \sqrt{3} r^{8}$\\
$\displaystyle  + 72 r^{2} z^{4} + 84 r^{9} z^{2} + 120 r z^{2}$\\
$\displaystyle  + 132 \sqrt{3} r^{4} + 148 \sqrt{3} r^{6} + 240 r^{7} z^{2}$\\
$\displaystyle  + 288 r^{3} z^{6} + 288 r^{7} z^{6} + 360 r^{6} z^{4}$\\
$\displaystyle  + 408 r^{3} z^{2} + 492 r^{5} z^{2} + 768 r^{4} z^{4}$\\
$\displaystyle  - 872 \sqrt{3} r^{5} z^{4} - 804 \sqrt{3} r^{7} z^{4}$\\
$\displaystyle  - 460 \sqrt{3} r^{6} z^{2} - 408 \sqrt{3} r^{4} z^{2}$\\
$\displaystyle  - 176 \sqrt{3} r^{4} z^{6} - 96 \sqrt{3} r^{6} z^{6}$\\
$\displaystyle  - 88 \sqrt{3} r^{8} z^{2} - 72 \sqrt{3} r^{9} z^{4}$\\
$\displaystyle  - 68 \sqrt{3} r^{3} z^{4} - 48 \sqrt{3} r^{3} z^{6}$\\
$\displaystyle  - 28 \sqrt{3} r^{2} z^{2} - 16 \sqrt{3} r z^{2} + 24 \sqrt{3} r z^{4}$\\
$\displaystyle  + 36 \sqrt{3} r^{8} z^{4} + 156 \sqrt{3} r^{9} z^{2}$\\
$\displaystyle  + 224 \sqrt{3} r^{3} z^{2} + 288 \sqrt{3} r^{5} z^{6}$\\
$\displaystyle  + 288 \sqrt{3} r^{7} z^{6} + 408 \sqrt{3} r^{6} z^{4}$\\
$\displaystyle  + 452 \sqrt{3} r^{4} z^{4} + 776 \sqrt{3} r^{7} z^{2}$\\
$\displaystyle  + 860 \sqrt{3} r^{5} z^{2}$}
\label{eq:cat-numerator-23}
\end{equation}
\begin{equation}
D_{23}(r,z)=3 \left(- r^{2} + 2 r z - 1\right)^{2} \left(r^{2} + 2 r z + 1\right)^{2}
\label{eq:cat-denominator-23}
\end{equation}
\begin{equation}
R_{24}(r,z)=\frac{N_{24}(r,z)}{D_{24}(r,z)}.
\label{eq:cat-rational-24}
\end{equation}
\begin{equation}
N_{24}(r,z)=\shortstack[l]{$\displaystyle  - 8 i r + 8 i r^{3} + 16 i r^{5} - 52 i r^{5} z^{2}$\\
$\displaystyle  - 24 i r^{4} z^{2} - 12 i r^{3} z^{2} - 8 \sqrt{3} i r^{2}$\\
$\displaystyle  - 8 \sqrt{3} i r^{4} + 4 i r^{3} z^{4} + 8 i r z^{2}$\\
$\displaystyle  + 24 i r^{4} z^{4} + 36 i r^{5} z^{4} - 24 \sqrt{3} i r^{2} z^{4}$\\
$\displaystyle  - 20 \sqrt{3} i r^{3} z^{4} - 12 \sqrt{3} i r^{5} z^{2}$\\
$\displaystyle  + 8 \sqrt{3} i r^{4} z^{2} + 12 \sqrt{3} i r^{5} z^{4}$\\
$\displaystyle  + 20 \sqrt{3} i r^{3} z^{2} + 32 \sqrt{3} i r^{2} z^{2}$}
\label{eq:cat-numerator-24}
\end{equation}
\begin{equation}
D_{24}(r,z)=3 \left(- r^{2} + 2 r z - 1\right) \left(r^{2} + 2 r z + 1\right)
\label{eq:cat-denominator-24}
\end{equation}
\begin{equation}
R_{25}(r,z)=- 40 \sqrt{3} \left(z - 1\right) \left(z + 1\right)
\label{eq:cat-rational-25}
\end{equation}
\begin{equation}
R_{26}(r,z)=16 \sqrt{3} i r \left(z - 1\right) \left(z + 1\right) + 8 i \left(3 r - 5\right) \left(z - 1\right) \left(z + 1\right)
\label{eq:cat-rational-26}
\end{equation}
\begin{equation}
R_{27}(r,z)=- 16 r \left(z - 1\right) \left(z + 1\right) - \frac{8 \sqrt{3} \left(3 r - 1\right) \left(z - 1\right) \left(z + 1\right)}{3}
\label{eq:cat-rational-27}
\end{equation}
\begin{equation}
R_{28}(r,z)=- \frac{8 i r \left(z - 1\right) \left(z + 1\right)}{3}
\label{eq:cat-rational-28}
\end{equation}
\begin{equation}
R_{29}(r,z)=\frac{N_{29}(r,z)}{D_{29}(r,z)}.
\label{eq:cat-rational-29}
\end{equation}
\begin{equation}
N_{29}(r,z)=\shortstack[l]{$\displaystyle 4 \sqrt{3} - 52 \sqrt{3} z^{4} + 20 \sqrt{3} r^{8} + 32 \sqrt{3} r^{2}$\\
$\displaystyle  + 48 \sqrt{3} z^{2} + 64 \sqrt{3} r^{6} + 72 \sqrt{3} r^{4}$\\
$\displaystyle  - 1320 \sqrt{3} r^{4} z^{4} - 992 \sqrt{3} r^{6} z^{4}$\\
$\displaystyle  - 324 \sqrt{3} r^{8} z^{4} - 320 \sqrt{3} r^{2} z^{4}$\\
$\displaystyle  - 16 \sqrt{3} r^{2} z^{6} + 256 \sqrt{3} r^{4} z^{2}$\\
$\displaystyle  + 304 \sqrt{3} r^{2} z^{2} + 304 \sqrt{3} r^{6} z^{2}$\\
$\displaystyle  + 304 \sqrt{3} r^{8} z^{2} + 624 \sqrt{3} r^{6} z^{6}$\\
$\displaystyle  + 992 \sqrt{3} r^{4} z^{6}$}
\label{eq:cat-numerator-29}
\end{equation}
\begin{equation}
D_{29}(r,z)=\left(- r^{2} + 2 r z - 1\right)^{2} \left(r^{2} + 2 r z + 1\right)^{2}
\label{eq:cat-denominator-29}
\end{equation}
\begin{equation}
R_{30}(r,z)=\frac{N_{30}(r,z)}{D_{30}(r,z)}.
\label{eq:cat-rational-30}
\end{equation}
\begin{equation}
N_{30}(r,z)=\shortstack[l]{$\displaystyle  - 48 i r^{5} - 40 i r^{7} - 24 i z^{2} - 24 i r^{3}$\\
$\displaystyle  - 12 i r^{9} - 4 i r + 24 i z^{4} + 32 i r^{2} + 32 i r^{8}$\\
$\displaystyle  + 96 i r^{4} + 96 i r^{6} - 1952 i r^{6} z^{4}$\\
$\displaystyle  - 1744 i r^{4} z^{4} - 648 i r^{8} z^{4} - 576 i r^{5} z^{6}$\\
$\displaystyle  - 288 i r^{2} z^{4} - 264 i r^{3} z^{2} - 152 i r^{5} z^{2}$\\
$\displaystyle  - 144 i r^{7} z^{6} - 96 i r^{2} z^{6} - 48 i r z^{2}$\\
$\displaystyle  - 24 i r^{9} z^{2} - 8 \sqrt{3} i r^{3} - 4 \sqrt{3} i r$\\
$\displaystyle  + 4 \sqrt{3} i r^{9} + 8 \sqrt{3} i r^{7} + 16 i r^{3} z^{6}$\\
$\displaystyle  + 36 i r^{9} z^{4} + 40 i r^{7} z^{2} + 52 i r z^{4}$\\
$\displaystyle  + 144 i r^{7} z^{4} + 272 i r^{3} z^{4} + 352 i r^{2} z^{2}$\\
$\displaystyle  + 368 i r^{4} z^{2} + 608 i r^{6} z^{2} + 616 i r^{8} z^{2}$\\
$\displaystyle  + 776 i r^{5} z^{4} + 1248 i r^{6} z^{6} + 1280 i r^{4} z^{6}$\\
$\displaystyle  - 224 \sqrt{3} i r^{5} z^{6} - 176 \sqrt{3} i r^{3} z^{6}$\\
$\displaystyle  - 160 \sqrt{3} i r^{7} z^{2} - 112 \sqrt{3} i r^{5} z^{2}$\\
$\displaystyle  - 96 \sqrt{3} i r z^{2} - 64 \sqrt{3} i r^{3} z^{2}$\\
$\displaystyle  - 48 \sqrt{3} i r^{7} z^{6} - 16 \sqrt{3} i r^{9} z^{2}$\\
$\displaystyle  + 12 \sqrt{3} i r^{9} z^{4} + 100 \sqrt{3} i r z^{4}$\\
$\displaystyle  + 200 \sqrt{3} i r^{7} z^{4} + 248 \sqrt{3} i r^{3} z^{4}$\\
$\displaystyle  + 336 \sqrt{3} i r^{5} z^{4}$}
\label{eq:cat-numerator-30}
\end{equation}
\begin{equation}
D_{30}(r,z)=\left(- r^{2} + 2 r z - 1\right)^{2} \left(r^{2} + 2 r z + 1\right)^{2}
\label{eq:cat-denominator-30}
\end{equation}
\begin{equation}
R_{31}(r,z)=\frac{N_{31}(r,z)}{D_{31}(r,z)}.
\label{eq:cat-rational-31}
\end{equation}
\begin{equation}
N_{31}(r,z)=\shortstack[l]{$\displaystyle  - 36 r^{4} - 36 r^{5} - 36 r^{6} - 36 r^{7} - 12 r^{2}$\\
$\displaystyle  - 12 r^{3} - 12 r^{8} - 12 r^{9} - 984 r^{5} z^{4}$\\
$\displaystyle  - 780 r^{7} z^{4} - 528 r^{4} z^{6} - 396 r^{3} z^{4}$\\
$\displaystyle  - 288 r^{2} z^{2} - 180 r^{4} z^{2} - 84 r^{6} z^{4}$\\
$\displaystyle  - 72 r z^{4} - 72 r^{9} z^{4} - 36 \sqrt{3} r^{4} - 36 \sqrt{3} r^{6}$\\
$\displaystyle  - 12 \sqrt{3} r^{2} - 12 \sqrt{3} r^{8} + 12 \sqrt{3} r^{3}$\\
$\displaystyle  + 12 \sqrt{3} r^{9} + 12 r^{8} z^{2} + 36 \sqrt{3} r^{5}$\\
$\displaystyle  + 36 \sqrt{3} r^{7} + 72 r z^{2} + 84 r^{9} z^{2}$\\
$\displaystyle  + 120 r^{3} z^{2} + 120 r^{6} z^{2} + 288 r^{3} z^{6}$\\
$\displaystyle  + 288 r^{7} z^{6} + 300 r^{2} z^{4} + 492 r^{5} z^{2}$\\
$\displaystyle  + 528 r^{7} z^{2} + 528 r^{5} z^{6} + 744 r^{4} z^{4}$\\
$\displaystyle  - 1080 \sqrt{3} r^{5} z^{4} - 720 \sqrt{3} r^{6} z^{6}$\\
$\displaystyle  - 396 \sqrt{3} r^{7} z^{4} - 356 \sqrt{3} r^{6} z^{2}$\\
$\displaystyle  - 336 \sqrt{3} r^{8} z^{2} - 320 \sqrt{3} r^{4} z^{6}$\\
$\displaystyle  - 204 \sqrt{3} r^{3} z^{4} - 96 \sqrt{3} r^{3} z^{6}$\\
$\displaystyle  - 88 \sqrt{3} r^{4} z^{2} - 72 \sqrt{3} r^{9} z^{4}$\\
$\displaystyle  - 68 \sqrt{3} r^{2} z^{2} - 24 \sqrt{3} r z^{2} + 24 \sqrt{3} r z^{4}$\\
$\displaystyle  + 60 \sqrt{3} r^{9} z^{2} + 72 \sqrt{3} r^{7} z^{2}$\\
$\displaystyle  + 80 \sqrt{3} r^{2} z^{4} + 288 \sqrt{3} r^{3} z^{2}$\\
$\displaystyle  + 288 \sqrt{3} r^{7} z^{6} + 324 \sqrt{3} r^{5} z^{2}$\\
$\displaystyle  + 348 \sqrt{3} r^{8} z^{4} + 444 \sqrt{3} r^{4} z^{4}$\\
$\displaystyle  + 720 \sqrt{3} r^{5} z^{6} + 1112 \sqrt{3} r^{6} z^{4}$}
\label{eq:cat-numerator-31}
\end{equation}
\begin{equation}
D_{31}(r,z)=3 \left(- r^{2} + 2 r z - 1\right)^{2} \left(r^{2} + 2 r z + 1\right)^{2}
\label{eq:cat-denominator-31}
\end{equation}
\begin{equation}
R_{32}(r,z)=\frac{N_{32}(r,z)}{D_{32}(r,z)}.
\label{eq:cat-rational-32}
\end{equation}
\begin{equation}
N_{32}(r,z)=\shortstack[l]{$\displaystyle  - 36 i r^{3} z^{2} - 36 i r^{5} z^{2} - 24 i r^{4} z^{4}$\\
$\displaystyle  + 24 i r^{4} z^{2} + 36 i r^{3} z^{4} + 36 i r^{5} z^{4}$\\
$\displaystyle  - 24 \sqrt{3} i r^{2} z^{4} - 12 \sqrt{3} i r^{5} z^{2}$\\
$\displaystyle  - 4 \sqrt{3} i r^{3} z^{4} + 4 \sqrt{3} i r^{3} z^{2}$\\
$\displaystyle  + 12 \sqrt{3} i r^{5} z^{4} + 24 \sqrt{3} i r^{2} z^{2}$}
\label{eq:cat-numerator-32}
\end{equation}
\begin{equation}
D_{32}(r,z)=3 \left(- r^{2} + 2 r z - 1\right) \left(r^{2} + 2 r z + 1\right)
\label{eq:cat-denominator-32}
\end{equation}
\begin{equation}
R_{33}(r,z)=\frac{N_{33}(r,z)}{D_{33}(r,z)}.
\label{eq:cat-rational-33}
\end{equation}
\begin{equation}
N_{33}(r,z)=\shortstack[l]{$\displaystyle  - 92 \sqrt{3} - 900 \sqrt{3} r^{4} - 716 \sqrt{3} r^{6} - 484 \sqrt{3} r^{2}$\\
$\displaystyle  - 208 \sqrt{3} r^{8} + 16 \sqrt{3} z^{2} + 76 \sqrt{3} z^{4}$\\
$\displaystyle  - 1380 \sqrt{3} r^{4} z^{4} - 432 \sqrt{3} r^{6} z^{6}$\\
$\displaystyle  - 284 \sqrt{3} r^{6} z^{4} - 244 \sqrt{3} r^{2} z^{4}$\\
$\displaystyle  - 56 \sqrt{3} r^{8} z^{2} - 32 \sqrt{3} r^{2} z^{6}$\\
$\displaystyle  + 48 \sqrt{3} r^{4} z^{6} + 264 \sqrt{3} r^{8} z^{4}$\\
$\displaystyle  + 760 \sqrt{3} r^{2} z^{2} + 1432 \sqrt{3} r^{6} z^{2}$\\
$\displaystyle  + 2232 \sqrt{3} r^{4} z^{2}$}
\label{eq:cat-numerator-33}
\end{equation}
\begin{equation}
D_{33}(r,z)=\left(- r^{2} + 2 r z - 1\right)^{2} \left(r^{2} + 2 r z + 1\right)^{2}
\label{eq:cat-denominator-33}
\end{equation}
\begin{equation}
R_{34}(r,z)=\frac{N_{34}(r,z)}{D_{34}(r,z)}.
\label{eq:cat-rational-34}
\end{equation}
\begin{equation}
N_{34}(r,z)=\shortstack[l]{$\displaystyle  - 56 i - 1000 i r^{4} - 892 i r^{6} - 444 i r^{2}$\\
$\displaystyle  - 280 i r^{8} + 52 i r + 56 i z^{2} + 96 i r^{9}$\\
$\displaystyle  + 256 i r^{3} + 344 i r^{7} + 452 i r^{5}$\\
$\displaystyle  - 1608 i r^{4} z^{4} - 1080 i r^{5} z^{2} - 864 i r^{6} z^{6}$\\
$\displaystyle  - 848 i r^{7} z^{2} - 320 i r^{3} z^{2} - 248 i r^{8} z^{2}$\\
$\displaystyle  - 204 i r^{2} z^{4} - 96 i r^{9} z^{2} - 76 i r z^{4}$\\
$\displaystyle  + 16 \sqrt{3} i r^{9} + 24 i r z^{2} + 32 i r^{3} z^{4}$\\
$\displaystyle  + 32 i r^{3} z^{6} + 48 i r^{2} z^{6} + 48 i r^{5} z^{6}$\\
$\displaystyle  + 52 i r^{6} z^{4} + 68 \sqrt{3} i r + 96 i r^{4} z^{6}$\\
$\displaystyle  + 116 \sqrt{3} i r^{7} + 220 \sqrt{3} i r^{3} + 252 \sqrt{3} i r^{5}$\\
$\displaystyle  + 504 i r^{7} z^{4} + 528 i r^{8} z^{4} + 580 i r^{5} z^{4}$\\
$\displaystyle  + 600 i r^{2} z^{2} + 1704 i r^{6} z^{2} + 2512 i r^{4} z^{2}$\\
$\displaystyle  - 648 \sqrt{3} i r^{5} z^{2} - 520 \sqrt{3} i r^{3} z^{2}$\\
$\displaystyle  - 136 \sqrt{3} i r^{7} z^{2} - 76 \sqrt{3} i r z^{4}$\\
$\displaystyle  - 48 \sqrt{3} i r^{5} z^{6} - 16 \sqrt{3} i r^{9} z^{2}$\\
$\displaystyle  + 8 \sqrt{3} i r z^{2} + 20 \sqrt{3} i r^{7} z^{4}$\\
$\displaystyle  + 80 \sqrt{3} i r^{3} z^{6} + 220 \sqrt{3} i r^{3} z^{4}$\\
$\displaystyle  + 444 \sqrt{3} i r^{5} z^{4}$}
\label{eq:cat-numerator-34}
\end{equation}
\begin{equation}
D_{34}(r,z)=\left(- r^{2} + 2 r z - 1\right)^{2} \left(r^{2} + 2 r z + 1\right)^{2}
\label{eq:cat-denominator-34}
\end{equation}
\begin{equation}
R_{35}(r,z)=\frac{N_{35}(r,z)}{D_{35}(r,z)}.
\label{eq:cat-rational-35}
\end{equation}
\begin{equation}
N_{35}(r,z)=\shortstack[l]{$\displaystyle  - 408 r^{5} - 324 r^{3} - 228 r^{7} - 96 r - 48 r^{9}$\\
$\displaystyle  + 84 r^{2} + 96 r^{8} + 264 r^{4} + 276 r^{6}$\\
$\displaystyle  - 792 r^{5} z^{4} - 720 r^{6} z^{2} - 540 r^{3} z^{4}$\\
$\displaystyle  - 528 r^{4} z^{2} - 456 \sqrt{3} r^{5} - 392 \sqrt{3} r^{7}$\\
$\displaystyle  - 228 r^{2} z^{4} - 216 \sqrt{3} r^{3} - 120 \sqrt{3} r^{9}$\\
$\displaystyle  - 96 r^{8} z^{2} - 60 r^{7} z^{4} - 32 \sqrt{3} r$\\
$\displaystyle  + 24 r^{4} z^{4} + 48 r^{9} z^{2} + 72 \sqrt{3} r^{2}$\\
$\displaystyle  + 72 \sqrt{3} r^{8} + 96 r z^{2} + 144 r^{2} z^{2}$\\
$\displaystyle  + 144 r^{5} z^{6} + 216 \sqrt{3} r^{4} + 216 \sqrt{3} r^{6}$\\
$\displaystyle  + 240 r^{4} z^{6} + 288 r^{7} z^{2} + 444 r^{6} z^{4}$\\
$\displaystyle  + 864 r^{3} z^{2} + 1056 r^{5} z^{2} - 640 \sqrt{3} r^{6} z^{4}$\\
$\displaystyle  - 600 \sqrt{3} r^{7} z^{4} - 560 \sqrt{3} r^{5} z^{4}$\\
$\displaystyle  - 496 \sqrt{3} r^{4} z^{2} - 312 \sqrt{3} r^{8} z^{4}$\\
$\displaystyle  - 200 \sqrt{3} r^{6} z^{2} - 56 \sqrt{3} r^{2} z^{2}$\\
$\displaystyle  - 56 \sqrt{3} r^{3} z^{4} - 48 \sqrt{3} r^{5} z^{6}$\\
$\displaystyle  - 16 \sqrt{3} r^{2} z^{4} + 16 \sqrt{3} r^{4} z^{6} + 32 \sqrt{3} r z^{2}$\\
$\displaystyle  + 48 \sqrt{3} r^{3} z^{6} + 120 \sqrt{3} r^{9} z^{2}$\\
$\displaystyle  + 224 \sqrt{3} r^{3} z^{2} + 240 \sqrt{3} r^{8} z^{2}$\\
$\displaystyle  + 264 \sqrt{3} r^{4} z^{4} + 624 \sqrt{3} r^{6} z^{6}$\\
$\displaystyle  + 992 \sqrt{3} r^{7} z^{2} + 1064 \sqrt{3} r^{5} z^{2}$}
\label{eq:cat-numerator-35}
\end{equation}
\begin{equation}
D_{35}(r,z)=3 \left(- r^{2} + 2 r z - 1\right)^{2} \left(r^{2} + 2 r z + 1\right)^{2}
\label{eq:cat-denominator-35}
\end{equation}
\begin{equation}
R_{36}(r,z)=\frac{N_{36}(r,z)}{D_{36}(r,z)}.
\label{eq:cat-rational-36}
\end{equation}
\begin{equation}
N_{36}(r,z)=\shortstack[l]{$\displaystyle 24 i r^{3} + 24 i r^{5} - 48 i r^{4} z^{2} - 24 i r^{3} z^{2}$\\
$\displaystyle  - 24 i r^{5} z^{2} - 8 \sqrt{3} i r^{2} - 8 \sqrt{3} i r^{4}$\\
$\displaystyle  + 48 i r^{4} z^{4} - 16 \sqrt{3} i r^{3} z^{4}$\\
$\displaystyle  + 8 \sqrt{3} i r^{2} z^{2} + 8 \sqrt{3} i r^{4} z^{2}$\\
$\displaystyle  + 16 \sqrt{3} i r^{3} z^{2}$}
\label{eq:cat-numerator-36}
\end{equation}
\begin{equation}
D_{36}(r,z)=3 \left(- r^{2} + 2 r z - 1\right) \left(r^{2} + 2 r z + 1\right)
\label{eq:cat-denominator-36}
\end{equation}
\begin{equation}
R_{37}(r,z)=\frac{N_{37}(r,z)}{D_{37}(r,z)}.
\label{eq:cat-rational-37}
\end{equation}
\begin{equation}
N_{37}(r,z)=\shortstack[l]{$\displaystyle 16 - 48 z^{2} + 72 z^{4} + 76 r^{8} + 136 r^{2} + 256 r^{6}$\\
$\displaystyle  + 300 r^{4} - 1176 r^{4} z^{6} - 1072 r^{4} z^{2}$\\
$\displaystyle  - 648 r^{6} z^{4} - 524 r^{6} z^{2} - 428 r^{2} z^{2}$\\
$\displaystyle  - 252 r^{8} z^{4} - 156 r^{2} z^{6} + 144 r^{4} z^{8}$\\
$\displaystyle  + 216 r^{8} z^{2} + 288 r^{2} z^{4} + 756 r^{6} z^{6}$\\
$\displaystyle  + 2044 r^{4} z^{4}$}
\label{eq:cat-numerator-37}
\end{equation}
\begin{equation}
D_{37}(r,z)=\left(- r^{2} + 2 r z - 1\right)^{2} \left(r^{2} + 2 r z + 1\right)^{2}
\label{eq:cat-denominator-37}
\end{equation}
\begin{equation}
R_{38}(r,z)=\frac{N_{38}(r,z)}{D_{38}(r,z)}.
\label{eq:cat-rational-38}
\end{equation}
\begin{equation}
N_{38}(r,z)=\shortstack[l]{$\displaystyle  - 240 i r^{3} - 216 i r^{5} - 96 i r - 96 i r^{7}$\\
$\displaystyle  - 24 i r^{9} - 8 \sqrt{3} i - 2592 i r^{5} z^{6}$\\
$\displaystyle  - 960 i r^{7} z^{2} - 864 i r^{7} z^{6} - 432 i r^{9} z^{2}$\\
$\displaystyle  - 336 i r^{5} z^{2} - 144 i r z^{4} - 96 i r^{3} z^{2}$\\
$\displaystyle  - 48 \sqrt{3} i z^{2} + 72 \sqrt{3} i z^{4} + 88 \sqrt{3} i r^{8}$\\
$\displaystyle  + 112 \sqrt{3} i r^{2} + 216 i r^{9} z^{4} + 288 i r^{3} z^{6}$\\
$\displaystyle  + 304 \sqrt{3} i r^{6} + 336 \sqrt{3} i r^{4} + 1008 i r^{3} z^{4}$\\
$\displaystyle  + 1704 i r^{5} z^{4} + 2880 i r^{7} z^{4}$\\
$\displaystyle  - 2832 \sqrt{3} i r^{4} z^{6} - 1552 \sqrt{3} i r^{4} z^{2}$\\
$\displaystyle  - 1152 \sqrt{3} i r^{6} z^{4} - 728 \sqrt{3} i r^{6} z^{2}$\\
$\displaystyle  - 632 \sqrt{3} i r^{2} z^{2} - 504 \sqrt{3} i r^{8} z^{4}$\\
$\displaystyle  - 120 \sqrt{3} i r^{2} z^{6} + 288 \sqrt{3} i r^{4} z^{8}$\\
$\displaystyle  + 432 \sqrt{3} i r^{8} z^{2} + 576 \sqrt{3} i r^{2} z^{4}$\\
$\displaystyle  + 1512 \sqrt{3} i r^{6} z^{6} + 3856 \sqrt{3} i r^{4} z^{4}$}
\label{eq:cat-numerator-38}
\end{equation}
\begin{equation}
D_{38}(r,z)=3 \left(- r^{2} + 2 r z - 1\right)^{2} \left(r^{2} + 2 r z + 1\right)^{2}
\label{eq:cat-denominator-38}
\end{equation}
\begin{equation}
R_{39}(r,z)=\frac{N_{39}(r,z)}{D_{39}(r,z)}.
\label{eq:cat-rational-39}
\end{equation}
\begin{equation}
N_{39}(r,z)=\shortstack[l]{$\displaystyle 48 - 48 z^{2} + 84 r^{8} + 192 r^{2} + 264 r^{6} + 324 r^{4}$\\
$\displaystyle  - 1044 r^{6} z^{6} - 912 r^{4} z^{2} - 900 r^{6} z^{2}$\\
$\displaystyle  - 432 \sqrt{3} r^{5} - 408 r^{8} z^{2} - 376 \sqrt{3} r^{7}$\\
$\displaystyle  - 324 r^{2} z^{2} - 184 \sqrt{3} r^{3} - 144 r^{4} z^{8}$\\
$\displaystyle  - 112 \sqrt{3} r^{9} - 36 r^{2} z^{6} - 16 \sqrt{3} r$\\
$\displaystyle  + 132 r^{4} z^{4} + 168 r^{2} z^{4} + 324 r^{8} z^{4}$\\
$\displaystyle  + 600 r^{4} z^{6} + 1680 r^{6} z^{4} - 2128 \sqrt{3} r^{5} z^{4}$\\
$\displaystyle  - 1752 \sqrt{3} r^{7} z^{4} - 408 \sqrt{3} r^{3} z^{4}$\\
$\displaystyle  - 144 \sqrt{3} r^{9} z^{4} - 96 \sqrt{3} r^{3} z^{6}$\\
$\displaystyle  + 48 \sqrt{3} r z^{4} + 288 \sqrt{3} r^{9} z^{2}$\\
$\displaystyle  + 560 \sqrt{3} r^{3} z^{2} + 576 \sqrt{3} r^{7} z^{6}$\\
$\displaystyle  + 1152 \sqrt{3} r^{5} z^{6} + 1424 \sqrt{3} r^{7} z^{2}$\\
$\displaystyle  + 1600 \sqrt{3} r^{5} z^{2}$}
\label{eq:cat-numerator-39}
\end{equation}
\begin{equation}
D_{39}(r,z)=3 \left(- r^{2} + 2 r z - 1\right)^{2} \left(r^{2} + 2 r z + 1\right)^{2}
\label{eq:cat-denominator-39}
\end{equation}
\begin{equation}
R_{40}(r,z)=\frac{N_{40}(r,z)}{D_{40}(r,z)}.
\label{eq:cat-rational-40}
\end{equation}
\begin{equation}
N_{40}(r,z)=\shortstack[l]{$\displaystyle 24 i r^{3} + 24 i r^{5} - 48 i r^{3} z^{2} - 48 i r^{5} z^{2}$\\
$\displaystyle  - 8 \sqrt{3} i r^{2} - 8 \sqrt{3} i r^{4} + 24 i r^{3} z^{4}$\\
$\displaystyle  + 24 i r^{5} z^{4} - 8 \sqrt{3} i r^{2} z^{4}$\\
$\displaystyle  - 8 \sqrt{3} i r^{4} z^{4} + 16 \sqrt{3} i r^{2} z^{2}$\\
$\displaystyle  + 16 \sqrt{3} i r^{4} z^{2}$}
\label{eq:cat-numerator-40}
\end{equation}
\begin{equation}
D_{40}(r,z)=\left(- r^{2} + 2 r z - 1\right) \left(r^{2} + 2 r z + 1\right)
\label{eq:cat-denominator-40}
\end{equation}
\begin{equation}
R_{41}(r,z)=- 144 r^{2} z^{2} + 144 z^{4} - 144 z^{2} - 16
\label{eq:cat-rational-41}
\end{equation}
\begin{equation}
R_{42}(r,z)=32 i r - 96 \sqrt{3} i z^{2} \left(r^{2} - z^{2} + 1\right)
\label{eq:cat-rational-42}
\end{equation}
\begin{equation}
R_{43}(r,z)=16 \left(3 z^{2} + 1\right) \left(r^{2} - z^{2} + 1\right)
\label{eq:cat-rational-43}
\end{equation}
\begin{equation}
R_{44}(r,z)=\frac{N_{44}(r,z)}{D_{44}(r,z)}.
\label{eq:cat-rational-44}
\end{equation}
\begin{equation}
N_{44}(r,z)=\shortstack[l]{$\displaystyle 36 - 176 r^{6} - 108 r^{8} - 100 z^{4} + 8 r^{4} + 112 r^{2}$\\
$\displaystyle  + 136 z^{2} - 2352 r^{6} z^{4} - 2304 r^{4} z^{8}$\\
$\displaystyle  - 1896 r^{4} z^{4} - 1424 r^{2} z^{4} - 756 r^{8} z^{4}$\\
$\displaystyle  + 272 r^{2} z^{2} + 304 r^{4} z^{2} + 656 r^{6} z^{2}$\\
$\displaystyle  + 752 r^{2} z^{6} + 936 r^{8} z^{2} + 1584 r^{6} z^{6}$\\
$\displaystyle  + 4320 r^{4} z^{6}$}
\label{eq:cat-numerator-44}
\end{equation}
\begin{equation}
D_{44}(r,z)=\left(- r^{2} + 2 r z - 1\right)^{2} \left(r^{2} + 2 r z + 1\right)^{2}
\label{eq:cat-denominator-44}
\end{equation}
\begin{equation}
R_{45}(r,z)=\frac{N_{45}(r,z)}{D_{45}(r,z)}.
\label{eq:cat-rational-45}
\end{equation}
\begin{equation}
N_{45}(r,z)=\shortstack[l]{$\displaystyle  - 624 i r^{5} - 528 i r^{3} - 528 i r^{7} - 216 i r$\\
$\displaystyle  - 216 i r^{9} + 56 \sqrt{3} i - 2880 i r^{5} z^{6}$\\
$\displaystyle  - 1296 i r^{5} z^{4} - 864 i r^{7} z^{6} - 768 i r^{7} z^{2}$\\
$\displaystyle  - 448 \sqrt{3} i r^{6} - 432 i r^{9} z^{2} - 280 \sqrt{3} i r^{8}$\\
$\displaystyle  - 240 i r z^{2} - 216 \sqrt{3} i z^{4} + 24 i r z^{4}$\\
$\displaystyle  + 48 i r^{3} z^{4} + 216 i r^{9} z^{4} + 224 \sqrt{3} i r^{2}$\\
$\displaystyle  + 240 \sqrt{3} i z^{2} + 960 i r^{3} z^{2} + 1248 i r^{3} z^{6}$\\
$\displaystyle  + 2208 i r^{5} z^{2} + 3888 i r^{7} z^{4}$\\
$\displaystyle  - 4608 \sqrt{3} i r^{4} z^{8} - 4224 \sqrt{3} i r^{6} z^{4}$\\
$\displaystyle  - 4000 \sqrt{3} i r^{4} z^{4} - 3168 \sqrt{3} i r^{2} z^{4}$\\
$\displaystyle  - 1512 \sqrt{3} i r^{8} z^{4} + 608 \sqrt{3} i r^{2} z^{2}$\\
$\displaystyle  + 1024 \sqrt{3} i r^{4} z^{2} + 1184 \sqrt{3} i r^{6} z^{2}$\\
$\displaystyle  + 1872 \sqrt{3} i r^{8} z^{2} + 2016 \sqrt{3} i r^{2} z^{6}$\\
$\displaystyle  + 3168 \sqrt{3} i r^{6} z^{6} + 8064 \sqrt{3} i r^{4} z^{6}$}
\label{eq:cat-numerator-45}
\end{equation}
\begin{equation}
D_{45}(r,z)=3 \left(- r^{2} + 2 r z - 1\right)^{2} \left(r^{2} + 2 r z + 1\right)^{2}
\label{eq:cat-denominator-45}
\end{equation}
\begin{equation}
R_{46}(r,z)=\frac{N_{46}(r,z)}{D_{46}(r,z)}.
\label{eq:cat-rational-46}
\end{equation}
\begin{equation}
N_{46}(r,z)=\shortstack[l]{$\displaystyle  - 16 - 144 z^{2} - 28 r^{2} + 144 z^{4} + 180 r^{4}$\\
$\displaystyle  + 188 r^{8} + 380 r^{6} - 3696 r^{4} z^{6} - 1872 r^{6} z^{6}$\\
$\displaystyle  - 1152 r^{2} z^{6} - 1032 r^{8} z^{2} - 1000 r^{6} z^{2}$\\
$\displaystyle  - 920 r^{4} z^{2} - 424 r^{2} z^{2} - 8 \sqrt{3} r^{7}$\\
$\displaystyle  + 16 \sqrt{3} r^{9} + 48 \sqrt{3} r^{5} + 112 \sqrt{3} r + 184 \sqrt{3} r^{3}$\\
$\displaystyle  + 828 r^{8} z^{4} + 1668 r^{2} z^{4} + 2036 r^{4} z^{4}$\\
$\displaystyle  + 2304 r^{4} z^{8} + 2556 r^{6} z^{4} - 2088 \sqrt{3} r^{7} z^{4}$\\
$\displaystyle  - 552 \sqrt{3} r^{3} z^{4} - 192 \sqrt{3} r^{3} z^{6}$\\
$\displaystyle  - 176 \sqrt{3} r^{5} z^{4} - 160 \sqrt{3} r^{5} z^{2}$\\
$\displaystyle  - 144 \sqrt{3} r^{9} z^{4} - 80 \sqrt{3} r^{3} z^{2}$\\
$\displaystyle  + 48 \sqrt{3} r z^{4} + 288 \sqrt{3} r^{9} z^{2}$\\
$\displaystyle  + 576 \sqrt{3} r^{7} z^{6} + 880 \sqrt{3} r^{7} z^{2}$\\
$\displaystyle  + 1248 \sqrt{3} r^{5} z^{6}$}
\label{eq:cat-numerator-46}
\end{equation}
\begin{equation}
D_{46}(r,z)=3 \left(- r^{2} + 2 r z - 1\right)^{2} \left(r^{2} + 2 r z + 1\right)^{2}
\label{eq:cat-denominator-46}
\end{equation}
\begin{equation}
R_{47}(r,z)=\frac{N_{47}(r,z)}{D_{47}(r,z)}.
\label{eq:cat-rational-47}
\end{equation}
\begin{equation}
N_{47}(r,z)=\shortstack[l]{$\displaystyle  - 32 i r + 8 i r^{3} + 40 i r^{5} - 144 i r^{5} z^{2}$\\
$\displaystyle  - 24 \sqrt{3} i r^{2} - 24 \sqrt{3} i r^{4} - 16 i r^{3} z^{2}$\\
$\displaystyle  + 72 i r^{3} z^{4} + 72 i r^{5} z^{4} - 24 \sqrt{3} i r^{2} z^{4}$\\
$\displaystyle  - 24 \sqrt{3} i r^{4} z^{4} + 48 \sqrt{3} i r^{2} z^{2}$\\
$\displaystyle  + 48 \sqrt{3} i r^{4} z^{2}$}
\label{eq:cat-numerator-47}
\end{equation}
\begin{equation}
D_{47}(r,z)=3 \left(- r^{2} + 2 r z - 1\right) \left(r^{2} + 2 r z + 1\right)
\label{eq:cat-denominator-47}
\end{equation}
\begin{equation}
R_{48}(r,z)=\frac{N_{48}(r,z)}{D_{48}(r,z)}.
\label{eq:cat-rational-48}
\end{equation}
\begin{equation}
N_{48}(r,z)=\shortstack[l]{$\displaystyle  - 4 - 40 z^{2} + 28 z^{4} + 88 r^{2} + 200 r^{8} + 388 r^{4}$\\
$\displaystyle  + 496 r^{6} - 2328 r^{6} z^{4} - 888 r^{4} z^{6}$\\
$\displaystyle  - 792 r^{8} z^{4} - 192 r^{4} z^{2} - 124 r^{4} z^{4}$\\
$\displaystyle  - 108 r^{2} z^{2} - 16 r^{2} z^{4} + 132 r^{6} z^{2}$\\
$\displaystyle  + 144 r^{10} z^{2} + 144 r^{4} z^{8} + 244 r^{2} z^{6}$\\
$\displaystyle  + 2628 r^{6} z^{6}$}
\label{eq:cat-numerator-48}
\end{equation}
\begin{equation}
D_{48}(r,z)=\left(- r^{2} + 2 r z - 1\right)^{2} \left(r^{2} + 2 r z + 1\right)^{2}
\label{eq:cat-denominator-48}
\end{equation}
\begin{equation}
R_{49}(r,z)=\frac{N_{49}(r,z)}{D_{49}(r,z)}.
\label{eq:cat-rational-49}
\end{equation}
\begin{equation}
N_{49}(r,z)=\shortstack[l]{$\displaystyle  - 168 i r^{5} - 96 i r^{3} - 64 \sqrt{3} i + 24 i r + 48 i r^{7}$\\
$\displaystyle  + 96 i r^{9} - 1008 i r^{5} z^{2} - 1008 i r^{7} z^{4}$\\
$\displaystyle  - 960 i r^{3} z^{6} - 288 i r^{3} z^{2} - 168 i r z^{4}$\\
$\displaystyle  - 112 \sqrt{3} i r^{2} + 240 i r z^{2} + 288 i r^{5} z^{6}$\\
$\displaystyle  + 336 \sqrt{3} i r^{4} + 368 \sqrt{3} i r^{8} + 576 i r^{7} z^{2}$\\
$\displaystyle  + 752 \sqrt{3} i r^{6} + 960 i r^{3} z^{4} + 1464 i r^{5} z^{4}$\\
$\displaystyle  - 4992 \sqrt{3} i r^{6} z^{4} - 1680 \sqrt{3} i r^{4} z^{6}$\\
$\displaystyle  - 1584 \sqrt{3} i r^{8} z^{4} - 784 \sqrt{3} i r^{4} z^{4}$\\
$\displaystyle  + 168 \sqrt{3} i r^{2} z^{6} + 200 \sqrt{3} i r^{2} z^{2}$\\
$\displaystyle  + 288 \sqrt{3} i r^{10} z^{2} + 288 \sqrt{3} i r^{2} z^{4}$\\
$\displaystyle  + 288 \sqrt{3} i r^{4} z^{8} + 304 \sqrt{3} i r^{4} z^{2}$\\
$\displaystyle  + 968 \sqrt{3} i r^{6} z^{2} + 5256 \sqrt{3} i r^{6} z^{6}$}
\label{eq:cat-numerator-49}
\end{equation}
\begin{equation}
D_{49}(r,z)=3 \left(- r^{2} + 2 r z - 1\right)^{2} \left(r^{2} + 2 r z + 1\right)^{2}
\label{eq:cat-denominator-49}
\end{equation}
\begin{equation}
R_{50}(r,z)=\frac{N_{50}(r,z)}{D_{50}(r,z)}.
\label{eq:cat-rational-50}
\end{equation}
\begin{equation}
N_{50}(r,z)=\shortstack[l]{$\displaystyle 16 - 596 r^{6} - 344 r^{8} - 336 r^{4} - 48 r^{10} - 20 r^{2}$\\
$\displaystyle  - 2628 r^{6} z^{6} - 480 \sqrt{3} r^{5} - 368 \sqrt{3} r^{3}$\\
$\displaystyle  - 368 \sqrt{3} r^{7} - 156 r^{2} z^{4} - 144 r^{10} z^{2}$\\
$\displaystyle  - 144 r^{4} z^{8} - 128 \sqrt{3} r - 128 \sqrt{3} r^{9}$\\
$\displaystyle  - 44 r^{2} z^{2} - 36 r^{2} z^{6} + 4 r^{6} z^{2}$\\
$\displaystyle  + 8 r^{4} z^{2} + 336 r^{8} z^{2} + 456 r^{4} z^{6}$\\
$\displaystyle  + 792 r^{8} z^{4} + 880 r^{4} z^{4} + 2004 r^{6} z^{4}$\\
$\displaystyle  - 1952 \sqrt{3} r^{5} z^{4} - 96 \sqrt{3} r^{5} z^{6}$\\
$\displaystyle  + 96 \sqrt{3} r^{3} z^{6} + 144 \sqrt{3} r^{3} z^{4}$\\
$\displaystyle  + 336 \sqrt{3} r^{7} z^{4} + 544 \sqrt{3} r^{7} z^{2}$\\
$\displaystyle  + 640 \sqrt{3} r^{3} z^{2} + 1760 \sqrt{3} r^{5} z^{2}$}
\label{eq:cat-numerator-50}
\end{equation}
\begin{equation}
D_{50}(r,z)=3 \left(- r^{2} + 2 r z - 1\right)^{2} \left(r^{2} + 2 r z + 1\right)^{2}
\label{eq:cat-denominator-50}
\end{equation}
\begin{equation}
R_{51}(r,z)=- \frac{32 i r}{3}
\label{eq:cat-rational-51}
\end{equation}
\begin{equation}
R_{52}(r,z)=\frac{N_{52}(r,z)}{D_{52}(r,z)}.
\label{eq:cat-rational-52}
\end{equation}
\begin{equation}
N_{52}(r,z)=\shortstack[l]{$\displaystyle  - 56 - 160 r^{2} - 104 r^{4} + 96 z^{2} - 432 r^{4} z^{4}$\\
$\displaystyle  - 48 r^{2} z^{4} + 128 r^{2} z^{2} + 576 r^{4} z^{2}$}
\label{eq:cat-numerator-52}
\end{equation}
\begin{equation}
D_{52}(r,z)=\left(- r^{2} + 2 r z - 1\right) \left(r^{2} + 2 r z + 1\right)
\label{eq:cat-denominator-52}
\end{equation}
\begin{equation}
R_{53}(r,z)=\frac{N_{53}(r,z)}{D_{53}(r,z)}.
\label{eq:cat-rational-53}
\end{equation}
\begin{equation}
N_{53}(r,z)=\shortstack[l]{$\displaystyle  - 240 i r^{5} - 192 i r^{3} - 80 \sqrt{3} i + 48 i r$\\
$\displaystyle  - 352 \sqrt{3} i r^{2} - 288 i r z^{2} - 272 \sqrt{3} i r^{4}$\\
$\displaystyle  + 96 \sqrt{3} i z^{2} + 96 i r^{3} z^{2} + 576 i r^{3} z^{4}$\\
$\displaystyle  - 864 \sqrt{3} i r^{4} z^{4} + 96 \sqrt{3} i r^{2} z^{4}$\\
$\displaystyle  + 224 \sqrt{3} i r^{2} z^{2} + 1152 \sqrt{3} i r^{4} z^{2}$}
\label{eq:cat-numerator-53}
\end{equation}
\begin{equation}
D_{53}(r,z)=3 \left(- r^{2} + 2 r z - 1\right) \left(r^{2} + 2 r z + 1\right)
\label{eq:cat-denominator-53}
\end{equation}
\begin{equation}
R_{54}(r,z)=\frac{N_{54}(r,z)}{D_{54}(r,z)}.
\label{eq:cat-rational-54}
\end{equation}
\begin{equation}
N_{54}(r,z)=\shortstack[l]{$\displaystyle 48 - 48 z^{2} + 192 r^{4} + 240 r^{2} - 624 r^{4} z^{2}$\\
$\displaystyle  - 288 r^{2} z^{2} - 64 \sqrt{3} r - 32 \sqrt{3} r^{3} + 32 \sqrt{3} r^{5}$\\
$\displaystyle  + 48 r^{2} z^{4} + 432 r^{4} z^{4} - 192 \sqrt{3} r^{3} z^{4}$\\
$\displaystyle  + 96 \sqrt{3} r z^{2} + 160 \sqrt{3} r^{3} z^{2}$}
\label{eq:cat-numerator-54}
\end{equation}
\begin{equation}
D_{54}(r,z)=3 \left(- r^{2} + 2 r z - 1\right) \left(r^{2} + 2 r z + 1\right)
\label{eq:cat-denominator-54}
\end{equation}
\begin{equation}
R_{55}(r,z)=144 r^{2} z^{2} - 144 z^{2} + 160
\label{eq:cat-rational-55}
\end{equation}
\begin{equation}
R_{56}(r,z)=- 32 i r + 96 \sqrt{3} i \left(r^{2} z^{2} - z^{2} + 1\right)
\label{eq:cat-rational-56}
\end{equation}
\begin{equation}
R_{57}(r,z)=- 48 r^{2} z^{2} - 16 r^{2} + 64 z^{2} - 64
\label{eq:cat-rational-57}
\end{equation}
\begin{equation}
R_{58}(r,z)=\frac{N_{58}(r,z)}{D_{58}(r,z)}.
\label{eq:cat-rational-58}
\end{equation}
\begin{equation}
N_{58}(r,z)=\shortstack[l]{$\displaystyle 136 - 64 z^{2} + 72 r^{4} + 208 r^{2} - 576 r^{4} z^{4}$\\
$\displaystyle  - 544 r^{2} z^{2} + 192 r^{2} z^{4} + 576 r^{4} z^{2}$}
\label{eq:cat-numerator-58}
\end{equation}
\begin{equation}
D_{58}(r,z)=\left(- r^{2} + 2 r z - 1\right) \left(r^{2} + 2 r z + 1\right)
\label{eq:cat-denominator-58}
\end{equation}
\begin{equation}
R_{59}(r,z)=\frac{N_{59}(r,z)}{D_{59}(r,z)}.
\label{eq:cat-rational-59}
\end{equation}
\begin{equation}
N_{59}(r,z)=\shortstack[l]{$\displaystyle  - 672 i r^{3} - 432 i r^{5} - 240 i r + 272 \sqrt{3} i$\\
$\displaystyle  - 192 i r z^{2} - 192 \sqrt{3} i z^{2} + 80 \sqrt{3} i r^{4}$\\
$\displaystyle  + 352 \sqrt{3} i r^{2} + 384 i r^{3} z^{2} + 1152 i r^{3} z^{4}$\\
$\displaystyle  - 1280 \sqrt{3} i r^{2} z^{2} - 1152 \sqrt{3} i r^{4} z^{4}$\\
$\displaystyle  + 768 \sqrt{3} i r^{2} z^{4} + 1152 \sqrt{3} i r^{4} z^{2}$}
\label{eq:cat-numerator-59}
\end{equation}
\begin{equation}
D_{59}(r,z)=3 \left(- r^{2} + 2 r z - 1\right) \left(r^{2} + 2 r z + 1\right)
\label{eq:cat-denominator-59}
\end{equation}
\begin{equation}
R_{60}(r,z)=\frac{N_{60}(r,z)}{D_{60}(r,z)}.
\label{eq:cat-rational-60}
\end{equation}
\begin{equation}
N_{60}(r,z)=\shortstack[l]{$\displaystyle  - 160 - 224 r^{2} - 64 r^{4} + 144 z^{2} - 576 r^{2} z^{4}$\\
$\displaystyle  - 528 r^{4} z^{2} + 64 \sqrt{3} r + 160 \sqrt{3} r^{5}$\\
$\displaystyle  + 224 \sqrt{3} r^{3} + 576 r^{4} z^{4} + 832 r^{2} z^{2}$\\
$\displaystyle  - 384 \sqrt{3} r^{3} z^{4} - 160 \sqrt{3} r^{3} z^{2}$\\
$\displaystyle  + 96 \sqrt{3} r z^{2}$}
\label{eq:cat-numerator-60}
\end{equation}
\begin{equation}
D_{60}(r,z)=3 \left(- r^{2} + 2 r z - 1\right) \left(r^{2} + 2 r z + 1\right)
\label{eq:cat-denominator-60}
\end{equation}
\begin{equation}
R_{61}(r,z)=\frac{N_{61}(r,z)}{D_{61}(r,z)}.
\label{eq:cat-rational-61}
\end{equation}
\begin{equation}
N_{61}(r,z)=\shortstack[l]{$\displaystyle  - 32 - 48 r^{2} - 16 r^{4} + 16 z^{2} - 432 r^{4} z^{4}$\\
$\displaystyle  - 112 r^{2} z^{2} + 144 r^{4} z^{2} + 144 r^{6} z^{2}$\\
$\displaystyle  + 336 r^{2} z^{4}$}
\label{eq:cat-numerator-61}
\end{equation}
\begin{equation}
D_{61}(r,z)=\left(- r^{2} + 2 r z - 1\right) \left(r^{2} + 2 r z + 1\right)
\label{eq:cat-denominator-61}
\end{equation}
\begin{equation}
R_{62}(r,z)=\frac{N_{62}(r,z)}{D_{62}(r,z)}.
\label{eq:cat-rational-62}
\end{equation}
\begin{equation}
N_{62}(r,z)=\shortstack[l]{$\displaystyle  - 64 \sqrt{3} i + 96 i r^{5} + 192 i r + 288 i r^{3}$\\
$\displaystyle  - 576 i r^{3} z^{4} - 128 \sqrt{3} i r^{2} - 96 i r z^{2}$\\
$\displaystyle  - 64 \sqrt{3} i r^{4} + 96 i r^{3} z^{2} - 864 \sqrt{3} i r^{4} z^{4}$\\
$\displaystyle  + 64 \sqrt{3} i r^{2} z^{2} + 288 \sqrt{3} i r^{4} z^{2}$\\
$\displaystyle  + 288 \sqrt{3} i r^{6} z^{2} + 480 \sqrt{3} i r^{2} z^{4}$}
\label{eq:cat-numerator-62}
\end{equation}
\begin{equation}
D_{62}(r,z)=3 \left(- r^{2} + 2 r z - 1\right) \left(r^{2} + 2 r z + 1\right)
\label{eq:cat-denominator-62}
\end{equation}
\begin{equation}
R_{63}(r,z)=\frac{N_{63}(r,z)}{D_{63}(r,z)}.
\label{eq:cat-rational-63}
\end{equation}
\begin{equation}
N_{63}(r,z)=\shortstack[l]{$\displaystyle 16 - 48 r^{6} - 32 r^{4} + 32 r^{2} - 256 \sqrt{3} r^{3}$\\
$\displaystyle  - 144 r^{6} z^{2} - 144 r^{2} z^{4} - 128 \sqrt{3} r$\\
$\displaystyle  - 128 \sqrt{3} r^{5} - 112 r^{2} z^{2} + 432 r^{4} z^{4}$\\
$\displaystyle  + 192 \sqrt{3} r^{3} z^{4} + 320 \sqrt{3} r^{3} z^{2}$}
\label{eq:cat-numerator-63}
\end{equation}
\begin{equation}
D_{63}(r,z)=3 \left(- r^{2} + 2 r z - 1\right) \left(r^{2} + 2 r z + 1\right)
\label{eq:cat-denominator-63}
\end{equation}
\begin{equation}
R_{64}(r,z)=\frac{32 i r}{3}
\label{eq:cat-rational-64}
\end{equation}
\begin{equation}
R_{65}(r,z)=\frac{N_{65}(r,z)}{D_{65}(r,z)}.
\label{eq:cat-rational-65}
\end{equation}
\begin{equation}
N_{65}(r,z)=\shortstack[l]{$\displaystyle  - 68 - 198 r^{2} - 198 r^{4} - 68 r^{6} - 36 z^{4} - 32 z^{2}$\\
$\displaystyle  - 652 r^{3} z^{3} - 396 r^{3} z^{7} - 364 r z^{3}$\\
$\displaystyle  - 364 r^{5} z^{3} - 36 r^{6} z^{4} - 32 r^{6} z^{2} + 38 r z$\\
$\displaystyle  + 38 r^{5} z + 38 r^{2} z^{4} + 38 r^{4} z^{4}$\\
$\displaystyle  + 42 r^{2} z^{6} + 42 r^{4} z^{6} + 54 r z^{5}$\\
$\displaystyle  + 54 r^{5} z^{5} + 64 r^{3} z + 254 r^{2} z^{2}$\\
$\displaystyle  + 254 r^{4} z^{2} + 1528 r^{3} z^{5}$}
\label{eq:cat-numerator-65}
\end{equation}
\begin{equation}
D_{65}(r,z)=\left(- r^{2} + 2 r z - 1\right) \left(r^{2} + 2 r z + 1\right)^{2}
\label{eq:cat-denominator-65}
\end{equation}
\begin{equation}
R_{66}(r,z)=\frac{N_{66}(r,z)}{D_{66}(r,z)}.
\label{eq:cat-rational-66}
\end{equation}
\begin{equation}
N_{66}(r,z)=\shortstack[l]{$\displaystyle  - 4 i - 206 i r^{3} - 204 i r^{5} - 68 i r - 62 i r^{7}$\\
$\displaystyle  - 16 i z^{2} + 4 \sqrt{3} i + 4 i r^{9} + 12 i z^{4} + 62 i r^{2}$\\
$\displaystyle  + 68 i r^{8} + 204 i r^{4} + 206 i r^{6} - 810 i r^{6} z^{5}$\\
$\displaystyle  - 706 i r^{6} z^{4} - 584 i r^{5} z^{6} - 456 i r^{5} z^{3}$\\
$\displaystyle  - 368 i r^{4} z^{2} - 240 i r^{3} z^{7} - 228 i r^{3} z^{6}$\\
$\displaystyle  - 168 i r z^{3} - 156 i r^{4} z - 152 i r^{2} z^{2}$\\
$\displaystyle  - 144 i r^{4} z^{8} - 132 i r^{3} z^{3} - 128 i r z^{2}$\\
$\displaystyle  - 128 i r^{6} z^{2} - 114 i r^{2} z - 108 i r^{5} z^{4}$\\
$\displaystyle  - 72 i r^{8} z^{5} - 72 i r^{5} z^{7} - 60 i r^{8} z^{4}$\\
$\displaystyle  - 60 i r^{5} z^{5} - 48 i r z - 34 i r^{2} z^{4}$\\
$\displaystyle  - 28 \sqrt{3} i r^{5} - 24 \sqrt{3} i r^{3} - 18 i r^{2} z^{5}$\\
$\displaystyle  - 16 \sqrt{3} i r^{7} - 12 i r^{2} z^{3} - 12 i r^{9} z^{4}$\\
$\displaystyle  - 12 i r^{7} z^{6} - 8 \sqrt{3} i r - 8 \sqrt{3} i z^{2}$\\
$\displaystyle  - 6 i r^{3} z - 4 \sqrt{3} i r^{9} + 6 i r^{6} z + 8 \sqrt{3} i r^{8}$\\
$\displaystyle  + 12 \sqrt{3} i z^{4} + 12 i r^{7} z^{3} + 12 i r^{2} z^{6}$\\
$\displaystyle  + 16 \sqrt{3} i r^{2} + 16 i r^{9} z^{2} + 18 i r^{7} z^{5}$\\
$\displaystyle  + 24 \sqrt{3} i r^{6} + 28 \sqrt{3} i r^{4} + 34 i r^{7} z^{4}$\\
$\displaystyle  + 48 i r^{8} z + 60 i r z^{4} + 60 i r^{4} z^{5}$\\
$\displaystyle  + 72 i r z^{5} + 72 i r^{4} z^{7} + 108 i r^{4} z^{4}$\\
$\displaystyle  + 114 i r^{7} z + 128 i r^{3} z^{2} + 128 i r^{8} z^{2}$\\
$\displaystyle  + 132 i r^{6} z^{3} + 144 i r^{5} z^{8} + 152 i r^{7} z^{2}$\\
$\displaystyle  + 156 i r^{5} z + 168 i r^{8} z^{3} + 228 i r^{6} z^{6}$\\
$\displaystyle  + 240 i r^{6} z^{7} + 368 i r^{5} z^{2} + 456 i r^{4} z^{3}$\\
$\displaystyle  + 584 i r^{4} z^{6} + 706 i r^{3} z^{4} + 810 i r^{3} z^{5}$\\
$\displaystyle  - 788 \sqrt{3} i r^{5} z^{6} - 560 \sqrt{3} i r^{4} z^{4}$\\
$\displaystyle  - 560 \sqrt{3} i r^{5} z^{5} - 540 \sqrt{3} i r^{4} z^{3}$\\
$\displaystyle  - 530 \sqrt{3} i r^{3} z^{5} - 496 \sqrt{3} i r^{6} z^{3}$\\
$\displaystyle  - 264 \sqrt{3} i r^{4} z^{8} - 216 \sqrt{3} i r^{6} z^{7}$\\
$\displaystyle  - 164 \sqrt{3} i r^{5} z - 160 \sqrt{3} i r^{4} z^{7}$\\
$\displaystyle  - 158 \sqrt{3} i r^{3} z - 136 \sqrt{3} i r^{2} z^{4}$\\
$\displaystyle  - 136 \sqrt{3} i r^{6} z^{4} - 112 \sqrt{3} i r^{2} z^{3}$\\
$\displaystyle  - 108 \sqrt{3} i r^{8} z^{3} - 80 \sqrt{3} i r^{5} z^{2}$\\
$\displaystyle  - 62 \sqrt{3} i r^{7} z - 60 \sqrt{3} i r z^{5} - 56 \sqrt{3} i r z$\\
$\displaystyle  - 50 \sqrt{3} i r^{7} z^{2} - 42 \sqrt{3} i r^{7} z^{5}$\\
$\displaystyle  - 34 \sqrt{3} i r^{3} z^{2} - 30 \sqrt{3} i r^{7} z^{6}$\\
$\displaystyle  - 22 \sqrt{3} i r^{3} z^{6} - 12 \sqrt{3} i r z^{4}$\\
$\displaystyle  - 12 \sqrt{3} i r^{9} z^{4} - 4 \sqrt{3} i r^{8} z^{2}$\\
$\displaystyle  + 4 \sqrt{3} i r z^{2} + 8 \sqrt{3} i r^{9} z^{2}$\\
$\displaystyle  + 12 \sqrt{3} i r^{8} z^{4} + 22 \sqrt{3} i r^{6} z^{6}$\\
$\displaystyle  + 30 \sqrt{3} i r^{2} z^{6} + 34 \sqrt{3} i r^{6} z^{2}$\\
$\displaystyle  + 42 \sqrt{3} i r^{2} z^{5} + 50 \sqrt{3} i r^{2} z^{2}$\\
$\displaystyle  + 56 \sqrt{3} i r^{8} z + 60 \sqrt{3} i r^{8} z^{5}$\\
$\displaystyle  + 62 \sqrt{3} i r^{2} z + 80 \sqrt{3} i r^{4} z^{2}$\\
$\displaystyle  + 108 \sqrt{3} i r z^{3} + 112 \sqrt{3} i r^{7} z^{3}$\\
$\displaystyle  + 136 \sqrt{3} i r^{3} z^{4} + 136 \sqrt{3} i r^{7} z^{4}$\\
$\displaystyle  + 158 \sqrt{3} i r^{6} z + 160 \sqrt{3} i r^{5} z^{7}$\\
$\displaystyle  + 164 \sqrt{3} i r^{4} z + 216 \sqrt{3} i r^{3} z^{7}$\\
$\displaystyle  + 264 \sqrt{3} i r^{5} z^{8} + 496 \sqrt{3} i r^{3} z^{3}$\\
$\displaystyle  + 530 \sqrt{3} i r^{6} z^{5} + 540 \sqrt{3} i r^{5} z^{3}$\\
$\displaystyle  + 560 \sqrt{3} i r^{5} z^{4} + 560 \sqrt{3} i r^{4} z^{5}$\\
$\displaystyle  + 788 \sqrt{3} i r^{4} z^{6}$}
\label{eq:cat-numerator-66}
\end{equation}
\begin{equation}
D_{66}(r,z)=\left(- r^{2} + 2 r z - 1\right)^{2} \left(r^{2} + 2 r z + 1\right)^{2}
\label{eq:cat-denominator-66}
\end{equation}
\begin{equation}
R_{67}(r,z)=\frac{N_{67}(r,z)}{D_{67}(r,z)}.
\label{eq:cat-rational-67}
\end{equation}
\begin{equation}
N_{67}(r,z)=\shortstack[l]{$\displaystyle  - 48 r^{5} - 30 r^{3} - 30 r^{7} - 6 r - 6 r^{9}$\\
$\displaystyle  - 828 r^{5} z^{4} - 492 r^{4} z^{5} - 492 r^{6} z^{5}$\\
$\displaystyle  - 292 r^{5} z^{5} - 264 r^{5} z^{8} - 162 r^{3} z^{4}$\\
$\displaystyle  - 162 r^{7} z^{4} - 84 r^{4} z - 84 r^{6} z - 60 r^{2} z^{5}$\\
$\displaystyle  - 60 r^{8} z^{5} - 54 \sqrt{3} r^{4} - 54 \sqrt{3} r^{6} - 48 r^{5} z$\\
$\displaystyle  - 48 r^{4} z^{4} - 48 r^{6} z^{4} - 40 r^{3} z^{5}$\\
$\displaystyle  - 40 r^{7} z^{5} - 36 r^{2} z - 36 r^{8} z - 36 r^{4} z^{6}$\\
$\displaystyle  - 36 r^{6} z^{6} - 32 r^{3} z - 32 r^{7} z - 18 \sqrt{3} r^{2}$\\
$\displaystyle  - 18 \sqrt{3} r^{8} - 12 r^{2} z^{6} - 12 r^{8} z^{6} - 8 r z$\\
$\displaystyle  - 8 r^{9} z - 6 r z^{2} - 6 r z^{5} - 6 r^{9} z^{2}$\\
$\displaystyle  - 6 r^{9} z^{5} + 2 \sqrt{3} r + 2 r z^{3} + 2 \sqrt{3} r^{9}$\\
$\displaystyle  + 2 r^{9} z^{3} + 8 \sqrt{3} r^{3} + 8 \sqrt{3} r^{7} + 12 \sqrt{3} r^{5}$\\
$\displaystyle  + 12 r^{2} z^{2} + 12 r^{8} z^{2} + 24 r^{3} z^{7}$\\
$\displaystyle  + 24 r^{7} z^{7} + 36 r^{4} z^{2} + 36 r^{6} z^{2}$\\
$\displaystyle  + 48 r^{4} z^{8} + 48 r^{6} z^{8} + 78 r^{3} z^{6}$\\
$\displaystyle  + 78 r^{7} z^{6} + 80 r^{5} z^{7} + 96 r^{2} z^{3}$\\
$\displaystyle  + 96 r^{3} z^{3} + 96 r^{7} z^{3} + 96 r^{8} z^{3}$\\
$\displaystyle  + 162 r^{3} z^{2} + 162 r^{7} z^{2} + 188 r^{5} z^{3}$\\
$\displaystyle  + 216 r^{4} z^{7} + 216 r^{6} z^{7} + 312 r^{5} z^{2}$\\
$\displaystyle  + 360 r^{4} z^{3} + 360 r^{6} z^{3} + 756 r^{5} z^{6}$\\
$\displaystyle  - 220 \sqrt{3} r^{4} z^{3} - 220 \sqrt{3} r^{6} z^{3}$\\
$\displaystyle  - 124 \sqrt{3} r^{4} z^{6} - 124 \sqrt{3} r^{6} z^{6}$\\
$\displaystyle  - 104 \sqrt{3} r^{3} z^{4} - 104 \sqrt{3} r^{7} z^{4}$\\
$\displaystyle  - 104 \sqrt{3} r^{4} z^{7} - 104 \sqrt{3} r^{6} z^{7}$\\
$\displaystyle  - 96 \sqrt{3} r^{5} z - 96 \sqrt{3} r^{5} z^{8} - 84 \sqrt{3} r^{5} z^{2}$\\
$\displaystyle  - 60 \sqrt{3} r^{2} z^{3} - 60 \sqrt{3} r^{8} z^{3}$\\
$\displaystyle  - 56 \sqrt{3} r^{3} z^{5} - 56 \sqrt{3} r^{7} z^{5} - 48 \sqrt{3} r^{3} z$\\
$\displaystyle  - 48 \sqrt{3} r^{7} z - 32 \sqrt{3} r^{3} z^{2} - 32 \sqrt{3} r^{7} z^{2}$\\
$\displaystyle  - 20 \sqrt{3} r^{2} z^{2} - 20 \sqrt{3} r^{8} z^{2}$\\
$\displaystyle  - 12 \sqrt{3} r^{2} z^{6} - 12 \sqrt{3} r^{8} z^{6} - 8 \sqrt{3} r z^{4}$\\
$\displaystyle  - 8 \sqrt{3} r^{9} z^{4} - 6 \sqrt{3} r z^{5} - 6 \sqrt{3} r^{9} z^{5}$\\
$\displaystyle  - 4 \sqrt{3} r^{4} z^{2} - 4 \sqrt{3} r^{6} z^{2}$\\
$\displaystyle  - 4 \sqrt{3} r^{5} z^{5} + 6 \sqrt{3} r^{2} z + 6 \sqrt{3} r^{8} z$\\
$\displaystyle  + 10 \sqrt{3} r z^{2} + 10 \sqrt{3} r z^{3} + 10 \sqrt{3} r^{9} z^{2}$\\
$\displaystyle  + 10 \sqrt{3} r^{9} z^{3} + 16 \sqrt{3} r^{3} z^{3}$\\
$\displaystyle  + 16 \sqrt{3} r^{7} z^{3} + 18 \sqrt{3} r^{4} z + 18 \sqrt{3} r^{6} z$\\
$\displaystyle  + 24 \sqrt{3} r^{3} z^{7} + 24 \sqrt{3} r^{7} z^{7}$\\
$\displaystyle  + 26 \sqrt{3} r^{2} z^{4} + 26 \sqrt{3} r^{8} z^{4}$\\
$\displaystyle  + 30 \sqrt{3} r^{2} z^{5} + 30 \sqrt{3} r^{8} z^{5}$\\
$\displaystyle  + 32 \sqrt{3} r^{5} z^{4} + 48 \sqrt{3} r^{5} z^{7}$\\
$\displaystyle  + 48 \sqrt{3} r^{4} z^{8} + 48 \sqrt{3} r^{6} z^{8}$\\
$\displaystyle  + 64 \sqrt{3} r^{3} z^{6} + 64 \sqrt{3} r^{7} z^{6}$\\
$\displaystyle  + 158 \sqrt{3} r^{4} z^{4} + 158 \sqrt{3} r^{6} z^{4}$\\
$\displaystyle  + 172 \sqrt{3} r^{5} z^{3} + 256 \sqrt{3} r^{5} z^{6}$\\
$\displaystyle  + 330 \sqrt{3} r^{4} z^{5} + 330 \sqrt{3} r^{6} z^{5}$}
\label{eq:cat-numerator-67}
\end{equation}
\begin{equation}
D_{67}(r,z)=\left(- r^{2} + 2 r z - 1\right)^{2} \left(r^{2} + 2 r z + 1\right)^{2}
\label{eq:cat-denominator-67}
\end{equation}
\begin{equation}
R_{68}(r,z)=\frac{N_{68}(r,z)}{D_{68}(r,z)}.
\label{eq:cat-rational-68}
\end{equation}
\begin{equation}
N_{68}(r,z)=\shortstack[l]{$\displaystyle  - 2 i r^{2} - 2 i r^{4} + 2 i r^{3} + 2 i r^{5}$\\
$\displaystyle  - 28 i r^{3} z^{3} - 18 i r^{4} z - 14 i r^{2} z^{4}$\\
$\displaystyle  - 12 i r^{3} z^{6} - 8 i r^{5} z^{2} - 6 i r^{4} z^{4}$\\
$\displaystyle  - 6 i r^{2} z^{5} - 4 i r^{3} z^{2} - 4 i r^{2} z^{3}$\\
$\displaystyle  - 2 i r^{5} z - 2 \sqrt{3} i r^{2} - 2 \sqrt{3} i r^{4}$\\
$\displaystyle  - 2 i r^{4} z^{5} + 2 i r^{2} z + 2 \sqrt{3} i r^{3}$\\
$\displaystyle  + 2 \sqrt{3} i r^{5} + 2 i r^{3} z^{5} + 4 i r^{4} z^{2}$\\
$\displaystyle  + 4 i r^{5} z^{3} + 6 i r^{3} z^{4} + 6 i r^{5} z^{5}$\\
$\displaystyle  + 8 i r^{2} z^{2} + 12 i r^{4} z^{6} + 14 i r^{5} z^{4}$\\
$\displaystyle  + 18 i r^{3} z + 28 i r^{4} z^{3} - 12 \sqrt{3} i r^{3} z^{3}$\\
$\displaystyle  - 6 \sqrt{3} i r^{4} z - 6 \sqrt{3} i r^{2} z^{4}$\\
$\displaystyle  - 4 \sqrt{3} i r^{3} z^{2} - 4 \sqrt{3} i r^{2} z^{3}$\\
$\displaystyle  - 4 \sqrt{3} i r^{3} z^{6} - 2 \sqrt{3} i r^{2} z$\\
$\displaystyle  - 2 \sqrt{3} i r^{3} z^{4} - 2 \sqrt{3} i r^{2} z^{5}$\\
$\displaystyle  - 2 \sqrt{3} i r^{3} z^{5} + 2 \sqrt{3} i r^{5} z$\\
$\displaystyle  + 2 \sqrt{3} i r^{4} z^{4} + 2 \sqrt{3} i r^{4} z^{5}$\\
$\displaystyle  + 2 \sqrt{3} i r^{5} z^{5} + 4 \sqrt{3} i r^{4} z^{2}$\\
$\displaystyle  + 4 \sqrt{3} i r^{5} z^{3} + 4 \sqrt{3} i r^{4} z^{6}$\\
$\displaystyle  + 6 \sqrt{3} i r^{3} z + 6 \sqrt{3} i r^{5} z^{4}$\\
$\displaystyle  + 12 \sqrt{3} i r^{4} z^{3}$}
\label{eq:cat-numerator-68}
\end{equation}
\begin{equation}
D_{68}(r,z)=\left(- r^{2} + 2 r z - 1\right) \left(r^{2} + 2 r z + 1\right)
\label{eq:cat-denominator-68}
\end{equation}
\begin{equation}
R_{69}(r,z)=\frac{4 \left(3 r^{2} z^{3} + 3 r^{2} z - 27 r z^{4} - 8 r z^{2} - 17 r + 3 z^{3} + 3 z\right)}{r}
\label{eq:cat-rational-69}
\end{equation}
\begin{equation}
R_{70}(r,z)=4 i \left(r - 1\right) \left(3 z + 5\right) \left(z^{2} + 1\right) + 4 \sqrt{3} i \left(r - 1\right) \left(9 z^{4} + z^{3} + z^{2} + z + 4\right)
\label{eq:cat-rational-70}
\end{equation}
\begin{equation}
R_{71}(r,z)=- 8 \sqrt{3} r \left(z + 1\right) \left(z^{2} + 1\right) - 4 r \left(9 z^{4} - 2 z^{3} + 2 z^{2} + 2 z + 5\right)
\label{eq:cat-rational-71}
\end{equation}
\begin{equation}
R_{72}(r,z)=\frac{N_{72}(r,z)}{D_{72}(r,z)}.
\label{eq:cat-rational-72}
\end{equation}
\begin{equation}
N_{72}(r,z)=\shortstack[l]{$\displaystyle 36 z + 36 z^{3} + 52 r + 52 r^{7} + 212 r^{3} + 212 r^{5}$\\
$\displaystyle  - 880 r^{4} z^{5} - 644 r^{3} z^{4} - 644 r^{5} z^{4}$\\
$\displaystyle  - 344 r^{4} z^{3} - 240 r^{2} z^{5} - 240 r^{6} z^{5}$\\
$\displaystyle  + 36 r^{8} z + 36 r^{8} z^{3} + 60 r z^{4} + 60 r^{7} z^{4}$\\
$\displaystyle  + 64 r^{3} z^{2} + 64 r^{5} z^{2} + 120 r^{3} z^{6}$\\
$\displaystyle  + 120 r^{5} z^{6} + 136 r z^{2} + 136 r^{7} z^{2}$\\
$\displaystyle  + 208 r^{2} z + 208 r^{6} z + 232 r^{4} z + 240 r^{2} z^{3}$\\
$\displaystyle  + 240 r^{6} z^{3} + 432 r^{4} z^{7}$}
\label{eq:cat-numerator-72}
\end{equation}
\begin{equation}
D_{72}(r,z)=r \left(- r^{2} + 2 r z - 1\right) \left(r^{2} + 2 r z + 1\right)^{2}
\label{eq:cat-denominator-72}
\end{equation}
\begin{equation}
R_{73}(r,z)=\frac{N_{73}(r,z)}{D_{73}(r,z)}.
\label{eq:cat-rational-73}
\end{equation}
\begin{align}
N_{73}(r,z)={}& - 4 i - 132 i r^{6} - 112 i r^{4} - 52 i r^{8} - 36 i r^{2}\notag\\
& - 8 \sqrt{3} i + 4 i r^{9} + 20 i z^{2} + 36 i z + 36 i z^{3}\notag\\
& + 36 i r^{7} + 48 i z^{4} + 52 i r + 112 i r^{5}\notag\\
& + 132 i r^{3} - 876 i r^{4} z^{5} - 536 i r^{3} z^{3}\notag\\
& - 444 i r^{6} z^{5} - 392 i r^{5} z^{3} - 384 i r^{3} z^{7}\notag\\
& - 384 i r^{5} z^{7} - 376 i r^{7} z^{3} - 324 i r^{2} z^{5}\notag\\
& - 300 i r^{3} z^{6} - 252 i r^{3} z^{2} - 240 i r^{5} z^{2}\notag\\
& - 192 i r^{8} z^{4} - 188 i r^{4} z - 180 i r^{8} z^{5}\notag\\
& - 168 i r^{4} z^{4} - 152 i r^{4} z^{6} - 144 i r^{4} z^{8}\notag\\
& - 128 \sqrt{3} i r^{4} - 100 i r^{2} z - 84 \sqrt{3} i r^{6}\notag\\
& - 68 i r z^{2} - 68 \sqrt{3} i r^{2} - 48 i r^{9} z^{4}\notag\\
& - 44 i r^{6} z - 44 \sqrt{3} i z^{2} - 44 i r^{3} z^{4}\notag\\
& - 44 i r^{7} z^{4} - 36 i r^{9} z - 36 i r^{2} z^{2}\notag\\
& - 36 i r^{9} z^{3} - 20 i r^{9} z^{2} - 16 \sqrt{3} i r^{8}\notag\\
& - 12 i r^{7} z^{6} - 8 i r z - 4 i r z^{3} + 4 i r^{8} z^{3}\notag\\
& + 8 i r^{8} z + 8 \sqrt{3} i r^{9} + 12 \sqrt{3} i z\notag\\
& + 12 \sqrt{3} i z^{3} + 12 i r^{2} z^{6} + 16 \sqrt{3} i r\notag\\
& + 36 \sqrt{3} i z^{4} + 36 i r^{7} z^{2} + 44 i r^{3} z\notag\\
& + 44 i r^{2} z^{4} + 44 i r^{6} z^{4} + 68 \sqrt{3} i r^{7}\notag\\
& + 68 i r^{8} z^{2} + 84 \sqrt{3} i r^{3} + 100 i r^{7} z\notag\\
& + 128 \sqrt{3} i r^{5} + 144 i r^{5} z^{8} + 152 i r^{5} z^{6}\notag\\
& + 168 i r^{5} z^{4} + 180 i r z^{5} + 188 i r^{5} z\notag\\
& + 192 i r z^{4} + 240 i r^{4} z^{2} + 252 i r^{6} z^{2}\notag\\
& + 300 i r^{6} z^{6} + 324 i r^{7} z^{5} + 376 i r^{2} z^{3}\notag\\
& + 384 i r^{4} z^{7} + 384 i r^{6} z^{7} + 392 i r^{4} z^{3}\notag\\
& + 444 i r^{3} z^{5} + 536 i r^{6} z^{3} + 876 i r^{5} z^{5}\notag\\
& - 592 \sqrt{3} i r^{4} z^{6} - 340 \sqrt{3} i r^{3} z^{4}\notag\\
& - 288 \sqrt{3} i r^{5} z^{8} - 276 \sqrt{3} i r^{7} z^{4}\notag\\
& - 224 \sqrt{3} i r^{3} z^{5} - 216 \sqrt{3} i r^{2} z^{6}\notag\\
& - 200 \sqrt{3} i r^{6} z^{6} - 192 \sqrt{3} i r^{5} z^{4}\notag\\
& - 164 \sqrt{3} i r^{5} z^{5} - 152 \sqrt{3} i r^{6} z^{3}\notag\\
& - 144 \sqrt{3} i r^{6} z^{7} - 136 \sqrt{3} i r^{4} z^{3}\notag\\
& - 120 \sqrt{3} i r^{2} z^{5} - 96 \sqrt{3} i r^{5} z^{2}\notag\\
& - 88 \sqrt{3} i r^{7} z^{2} - 84 \sqrt{3} i r^{8} z^{3}\notag\\
& - 64 \sqrt{3} i r^{5} z^{7} - 56 \sqrt{3} i r z\notag\\
& - 56 \sqrt{3} i r^{3} z^{2} - 48 \sqrt{3} i r^{7} z\notag\\
& - 36 \sqrt{3} i r z^{5} - 36 \sqrt{3} i r^{9} z^{4}\notag\\
& - 24 \sqrt{3} i r^{3} z - 20 \sqrt{3} i r^{8} z^{2}\notag\\
& - 12 \sqrt{3} i r^{9} z - 12 \sqrt{3} i r^{9} z^{3}\notag\\
& - 8 \sqrt{3} i r^{2} z^{3} - 4 \sqrt{3} i r^{5} z - 4 \sqrt{3} i r z^{4}\notag\\
& + 4 \sqrt{3} i r^{4} z + 4 \sqrt{3} i r^{8} z^{4}\notag\\
& + 8 \sqrt{3} i r^{7} z^{3} + 20 \sqrt{3} i r z^{2}\notag\\
& + 24 \sqrt{3} i r^{6} z + 36 \sqrt{3} i r^{8} z^{5}\notag\\
& + 44 \sqrt{3} i r^{9} z^{2} + 48 \sqrt{3} i r^{2} z\notag\\
& + 56 \sqrt{3} i r^{8} z + 56 \sqrt{3} i r^{6} z^{2}\notag\\
& + 64 \sqrt{3} i r^{4} z^{7} + 84 \sqrt{3} i r z^{3}\notag\\
& + 88 \sqrt{3} i r^{2} z^{2} + 96 \sqrt{3} i r^{4} z^{2}\notag\\
& + 120 \sqrt{3} i r^{7} z^{5} + 136 \sqrt{3} i r^{5} z^{3}\notag\\
& + 144 \sqrt{3} i r^{3} z^{7} + 152 \sqrt{3} i r^{3} z^{3}\notag\\
& + 164 \sqrt{3} i r^{4} z^{5} + 192 \sqrt{3} i r^{4} z^{4}\notag\\
& + 200 \sqrt{3} i r^{3} z^{6} + 216 \sqrt{3} i r^{7} z^{6}\notag\\
& + 224 \sqrt{3} i r^{6} z^{5} + 276 \sqrt{3} i r^{2} z^{4}\notag\\
& + 288 \sqrt{3} i r^{4} z^{8} + 340 \sqrt{3} i r^{6} z^{4}\notag\\
& + 592 \sqrt{3} i r^{5} z^{6}
\label{eq:cat-numerator-73}
\end{align}
\begin{equation}
D_{73}(r,z)=\left(- r^{2} + 2 r z - 1\right)^{2} \left(r^{2} + 2 r z + 1\right)^{2}
\label{eq:cat-denominator-73}
\end{equation}
\begin{equation}
R_{74}(r,z)=\frac{N_{74}(r,z)}{D_{74}(r,z)}.
\label{eq:cat-rational-74}
\end{equation}
\begin{equation}
N_{74}(r,z)=\shortstack[l]{$\displaystyle  - 172 r^{5} - 96 r^{3} - 96 r^{7} - 10 r - 10 r^{9}$\\
$\displaystyle  - 464 r^{5} z^{6} - 204 r^{4} z^{5} - 204 r^{6} z^{5}$\\
$\displaystyle  - 168 r^{3} z^{6} - 168 r^{7} z^{6} - 60 \sqrt{3} r^{5}$\\
$\displaystyle  - 48 r^{4} z^{4} - 48 r^{6} z^{4} - 44 r^{2} z - 44 r^{8} z$\\
$\displaystyle  - 40 \sqrt{3} r^{3} - 40 \sqrt{3} r^{7} - 36 r^{2} z^{5}$\\
$\displaystyle  - 36 r^{5} z^{5} - 36 r^{8} z^{5} - 36 r^{4} z^{6}$\\
$\displaystyle  - 36 r^{6} z^{6} - 34 r z^{2} - 34 r^{9} z^{2}$\\
$\displaystyle  - 32 r^{4} z^{3} - 32 r^{6} z^{3} - 28 r^{5} z^{3}$\\
$\displaystyle  - 16 r^{3} z^{3} - 16 r^{7} z^{3} - 12 r^{2} z^{6}$\\
$\displaystyle  - 12 r^{8} z^{6} - 10 \sqrt{3} r - 10 \sqrt{3} r^{9} - 8 r^{3} z^{5}$\\
$\displaystyle  - 8 r^{7} z^{5} - 6 r z^{5} - 6 r^{9} z^{5} - 2 r z^{3}$\\
$\displaystyle  - 2 r^{9} z^{3} + 4 r z + 4 r^{9} z + 12 r^{2} z^{2}$\\
$\displaystyle  + 12 r^{8} z^{2} + 16 r^{3} z + 16 r^{7} z + 16 r^{5} z^{7}$\\
$\displaystyle  + 18 \sqrt{3} r^{2} + 18 \sqrt{3} r^{8} + 24 r^{5} z + 24 r z^{4}$\\
$\displaystyle  + 24 r^{9} z^{4} + 24 r^{3} z^{7} + 24 r^{7} z^{7}$\\
$\displaystyle  + 36 r^{4} z^{2} + 36 r^{6} z^{2} + 48 r^{4} z^{8}$\\
$\displaystyle  + 48 r^{6} z^{8} + 54 \sqrt{3} r^{4} + 54 \sqrt{3} r^{6}$\\
$\displaystyle  + 80 r^{2} z^{3} + 80 r^{8} z^{3} + 92 r^{4} z + 92 r^{6} z$\\
$\displaystyle  + 144 r^{4} z^{7} + 144 r^{6} z^{7} + 152 r^{3} z^{4}$\\
$\displaystyle  + 152 r^{7} z^{4} + 192 r^{3} z^{2} + 192 r^{7} z^{2}$\\
$\displaystyle  + 228 r^{5} z^{2} + 288 r^{5} z^{8} - 412 \sqrt{3} r^{5} z^{5}$\\
$\displaystyle  - 196 \sqrt{3} r^{4} z^{3} - 196 \sqrt{3} r^{6} z^{3}$\\
$\displaystyle  - 180 \sqrt{3} r^{4} z^{2} - 180 \sqrt{3} r^{6} z^{2}$\\
$\displaystyle  - 152 \sqrt{3} r^{5} z - 152 \sqrt{3} r^{4} z^{7}$\\
$\displaystyle  - 152 \sqrt{3} r^{6} z^{7} - 116 \sqrt{3} r^{3} z^{5}$\\
$\displaystyle  - 116 \sqrt{3} r^{7} z^{5} - 96 \sqrt{3} r^{5} z^{8}$\\
$\displaystyle  - 76 \sqrt{3} r^{3} z - 76 \sqrt{3} r^{7} z - 36 \sqrt{3} r^{2} z^{2}$\\
$\displaystyle  - 36 \sqrt{3} r^{8} z^{2} - 36 \sqrt{3} r^{2} z^{3}$\\
$\displaystyle  - 36 \sqrt{3} r^{8} z^{3} - 28 \sqrt{3} r^{4} z^{6}$\\
$\displaystyle  - 28 \sqrt{3} r^{6} z^{6} - 18 \sqrt{3} r^{4} z - 18 \sqrt{3} r^{6} z$\\
$\displaystyle  - 12 \sqrt{3} r^{2} z^{6} - 12 \sqrt{3} r^{8} z^{6}$\\
$\displaystyle  - 8 \sqrt{3} r^{3} z^{4} - 8 \sqrt{3} r^{7} z^{4} - 6 \sqrt{3} r^{2} z$\\
$\displaystyle  - 6 \sqrt{3} r^{8} z - 6 \sqrt{3} r z^{5} - 6 \sqrt{3} r^{9} z^{5}$\\
$\displaystyle  + 4 \sqrt{3} r z^{4} + 4 \sqrt{3} r^{9} z^{4} + 8 \sqrt{3} r^{5} z^{4}$\\
$\displaystyle  + 10 \sqrt{3} r z^{2} + 10 \sqrt{3} r z^{3} + 10 \sqrt{3} r^{9} z^{2}$\\
$\displaystyle  + 10 \sqrt{3} r^{9} z^{3} + 12 \sqrt{3} r^{5} z^{2}$\\
$\displaystyle  + 16 \sqrt{3} r^{3} z^{2} + 16 \sqrt{3} r^{7} z^{2}$\\
$\displaystyle  + 24 \sqrt{3} r^{3} z^{7} + 24 \sqrt{3} r^{7} z^{7}$\\
$\displaystyle  + 48 \sqrt{3} r^{4} z^{8} + 48 \sqrt{3} r^{6} z^{8}$\\
$\displaystyle  + 54 \sqrt{3} r^{2} z^{4} + 54 \sqrt{3} r^{8} z^{4}$\\
$\displaystyle  + 64 \sqrt{3} r^{3} z^{6} + 64 \sqrt{3} r^{5} z^{6}$\\
$\displaystyle  + 64 \sqrt{3} r^{7} z^{6} + 66 \sqrt{3} r^{2} z^{5}$\\
$\displaystyle  + 66 \sqrt{3} r^{8} z^{5} + 82 \sqrt{3} r^{4} z^{4}$\\
$\displaystyle  + 82 \sqrt{3} r^{6} z^{4} + 200 \sqrt{3} r^{3} z^{3}$\\
$\displaystyle  + 200 \sqrt{3} r^{7} z^{3} + 236 \sqrt{3} r^{5} z^{3}$\\
$\displaystyle  + 256 \sqrt{3} r^{5} z^{7} + 342 \sqrt{3} r^{4} z^{5}$\\
$\displaystyle  + 342 \sqrt{3} r^{6} z^{5}$}
\label{eq:cat-numerator-74}
\end{equation}
\begin{equation}
D_{74}(r,z)=\left(- r^{2} + 2 r z - 1\right)^{2} \left(r^{2} + 2 r z + 1\right)^{2}
\label{eq:cat-denominator-74}
\end{equation}
\begin{equation}
R_{75}(r,z)=\frac{N_{75}(r,z)}{D_{75}(r,z)}.
\label{eq:cat-rational-75}
\end{equation}
\begin{equation}
N_{75}(r,z)=\shortstack[l]{$\displaystyle  - 6 i r^{2} - 6 i r^{4} + 6 i r^{3} + 6 i r^{5}$\\
$\displaystyle  - 28 i r^{3} z^{3} - 12 i r^{3} z^{6} - 10 i r^{4} z$\\
$\displaystyle  - 10 i r^{2} z^{4} - 10 i r^{4} z^{5} - 8 i r^{5} z^{2}$\\
$\displaystyle  - 6 i r^{2} z^{5} - 4 i r^{3} z^{2} - 4 i r^{2} z^{3}$\\
$\displaystyle  - 2 i r^{5} z - 2 \sqrt{3} i r^{2} - 2 \sqrt{3} i r^{4}$\\
$\displaystyle  - 2 i r^{4} z^{4} + 2 i r^{2} z + 2 \sqrt{3} i r^{3}$\\
$\displaystyle  + 2 \sqrt{3} i r^{5} + 2 i r^{3} z^{4} + 4 i r^{4} z^{2}$\\
$\displaystyle  + 4 i r^{5} z^{3} + 6 i r^{5} z^{5} + 8 i r^{2} z^{2}$\\
$\displaystyle  + 10 i r^{3} z + 10 i r^{5} z^{4} + 10 i r^{3} z^{5}$\\
$\displaystyle  + 12 i r^{4} z^{6} + 28 i r^{4} z^{3} - 12 \sqrt{3} i r^{3} z^{3}$\\
$\displaystyle  - 6 \sqrt{3} i r^{4} z - 6 \sqrt{3} i r^{2} z^{4}$\\
$\displaystyle  - 4 \sqrt{3} i r^{3} z^{2} - 4 \sqrt{3} i r^{2} z^{3}$\\
$\displaystyle  - 4 \sqrt{3} i r^{3} z^{6} - 2 \sqrt{3} i r^{2} z$\\
$\displaystyle  - 2 \sqrt{3} i r^{3} z^{4} - 2 \sqrt{3} i r^{2} z^{5}$\\
$\displaystyle  - 2 \sqrt{3} i r^{3} z^{5} + 2 \sqrt{3} i r^{5} z$\\
$\displaystyle  + 2 \sqrt{3} i r^{4} z^{4} + 2 \sqrt{3} i r^{4} z^{5}$\\
$\displaystyle  + 2 \sqrt{3} i r^{5} z^{5} + 4 \sqrt{3} i r^{4} z^{2}$\\
$\displaystyle  + 4 \sqrt{3} i r^{5} z^{3} + 4 \sqrt{3} i r^{4} z^{6}$\\
$\displaystyle  + 6 \sqrt{3} i r^{3} z + 6 \sqrt{3} i r^{5} z^{4}$\\
$\displaystyle  + 12 \sqrt{3} i r^{4} z^{3}$}
\label{eq:cat-numerator-75}
\end{equation}
\begin{equation}
D_{75}(r,z)=\left(- r^{2} + 2 r z - 1\right) \left(r^{2} + 2 r z + 1\right)
\label{eq:cat-denominator-75}
\end{equation}
\begin{equation}
R_{76}(r,z)=\frac{N_{76}(r,z)}{D_{76}(r,z)}.
\label{eq:cat-rational-76}
\end{equation}
\begin{equation}
N_{76}(r,z)=\shortstack[l]{$\displaystyle  - 614 r^{3} - 614 r^{5} - 188 r - 188 r^{7} - 24 z - 24 z^{3}$\\
$\displaystyle  - 668 r^{2} z^{3} - 668 r^{6} z^{3} - 368 r^{4} z$\\
$\displaystyle  - 258 r^{2} z - 258 r^{6} z - 180 r z^{4} - 180 r^{7} z^{4}$\\
$\displaystyle  - 176 r z^{2} - 176 r^{7} z^{2} - 24 r^{8} z - 24 r^{8} z^{3}$\\
$\displaystyle  + 30 r^{2} z^{5} + 30 r^{6} z^{5} + 36 r^{4} z^{7}$\\
$\displaystyle  + 84 r^{4} z^{3} + 258 r^{3} z^{6} + 258 r^{5} z^{6}$\\
$\displaystyle  + 438 r^{3} z^{2} + 438 r^{5} z^{2} + 462 r^{3} z^{4}$\\
$\displaystyle  + 462 r^{5} z^{4} + 2136 r^{4} z^{5}$}
\label{eq:cat-numerator-76}
\end{equation}
\begin{equation}
D_{76}(r,z)=r \left(- r^{2} + 2 r z - 1\right) \left(r^{2} + 2 r z + 1\right)^{2}
\label{eq:cat-denominator-76}
\end{equation}
\begin{equation}
R_{77}(r,z)=\frac{N_{77}(r,z)}{D_{77}(r,z)}.
\label{eq:cat-rational-77}
\end{equation}
\begin{equation}
N_{77}(r,z)=\shortstack[l]{$\displaystyle 20 i - 436 i r^{5} - 418 i r^{3} - 178 i r^{7} - 140 i r$\\
$\displaystyle  - 36 i z^{4} - 24 i z - 24 i z^{3} - 20 i r^{9} - 16 i z^{2}$\\
$\displaystyle  + 28 \sqrt{3} i + 140 i r^{8} + 178 i r^{2} + 418 i r^{6}$\\
$\displaystyle  + 436 i r^{4} - 1056 i r^{5} z^{6} - 936 i r^{5} z^{5}$\\
$\displaystyle  - 910 i r^{6} z^{4} - 808 i r^{4} z^{2} - 462 i r^{6} z^{5}$\\
$\displaystyle  - 460 i r^{6} z^{2} - 452 i r^{6} z^{3} - 436 i r^{2} z^{3}$\\
$\displaystyle  - 276 i r^{5} z^{4} - 252 \sqrt{3} i r^{5} - 238 i r^{2} z^{4}$\\
$\displaystyle  - 210 i r^{7} z^{5} - 196 i r^{2} z^{2} - 176 i r z^{3}$\\
$\displaystyle  - 172 \sqrt{3} i r^{3} - 148 \sqrt{3} i r^{7} - 144 i r^{6} z^{7}$\\
$\displaystyle  - 132 i r z^{4} - 120 i r^{4} z^{7} - 108 i r z^{5}$\\
$\displaystyle  - 104 i r^{5} z - 98 i r^{3} z - 80 i r z^{2}$\\
$\displaystyle  - 72 i r^{6} z^{6} - 56 i r^{4} z^{3} - 52 i r z$\\
$\displaystyle  - 40 \sqrt{3} i r - 34 i r^{7} z - 28 \sqrt{3} i r^{9} - 8 \sqrt{3} i z$\\
$\displaystyle  - 8 \sqrt{3} i z^{3} + 12 \sqrt{3} i z^{4} + 16 i r^{9} z^{2}$\\
$\displaystyle  + 24 i r^{9} z + 24 i r^{9} z^{3} + 34 i r^{2} z$\\
$\displaystyle  + 36 i r^{9} z^{4} + 40 \sqrt{3} i z^{2} + 40 \sqrt{3} i r^{8}$\\
$\displaystyle  + 52 i r^{8} z + 56 i r^{5} z^{3} + 72 i r^{3} z^{6}$\\
$\displaystyle  + 80 i r^{8} z^{2} + 98 i r^{6} z + 104 i r^{4} z$\\
$\displaystyle  + 108 i r^{8} z^{5} + 120 i r^{5} z^{7} + 132 i r^{8} z^{4}$\\
$\displaystyle  + 144 i r^{3} z^{7} + 148 \sqrt{3} i r^{2} + 172 \sqrt{3} i r^{6}$\\
$\displaystyle  + 176 i r^{8} z^{3} + 196 i r^{7} z^{2} + 210 i r^{2} z^{5}$\\
$\displaystyle  + 238 i r^{7} z^{4} + 252 \sqrt{3} i r^{4} + 276 i r^{4} z^{4}$\\
$\displaystyle  + 436 i r^{7} z^{3} + 452 i r^{3} z^{3} + 460 i r^{3} z^{2}$\\
$\displaystyle  + 462 i r^{3} z^{5} + 808 i r^{5} z^{2} + 910 i r^{3} z^{4}$\\
$\displaystyle  + 936 i r^{4} z^{5} + 1056 i r^{4} z^{6}$\\
$\displaystyle  - 868 \sqrt{3} i r^{5} z^{6} - 444 \sqrt{3} i r^{4} z^{3}$\\
$\displaystyle  - 396 \sqrt{3} i r^{5} z^{5} - 364 \sqrt{3} i r^{6} z^{4}$\\
$\displaystyle  - 360 \sqrt{3} i r^{6} z^{3} - 344 \sqrt{3} i r^{4} z^{4}$\\
$\displaystyle  - 300 \sqrt{3} i r^{2} z^{4} - 274 \sqrt{3} i r^{3} z^{5}$\\
$\displaystyle  - 248 \sqrt{3} i r^{4} z^{2} - 184 \sqrt{3} i r^{5} z$\\
$\displaystyle  - 160 \sqrt{3} i r^{4} z^{7} - 150 \sqrt{3} i r^{3} z$\\
$\displaystyle  - 150 \sqrt{3} i r^{2} z^{2} - 134 \sqrt{3} i r^{6} z^{2}$\\
$\displaystyle  - 130 \sqrt{3} i r^{7} z^{5} - 120 \sqrt{3} i r^{2} z^{3}$\\
$\displaystyle  - 72 \sqrt{3} i r^{6} z^{7} - 66 \sqrt{3} i r^{6} z^{6}$\\
$\displaystyle  - 44 \sqrt{3} i r z^{4} - 42 \sqrt{3} i r^{2} z^{6}$\\
$\displaystyle  - 40 \sqrt{3} i r^{9} z^{2} - 30 \sqrt{3} i r^{7} z$\\
$\displaystyle  - 24 \sqrt{3} i r z^{5} - 24 \sqrt{3} i r^{5} z^{8}$\\
$\displaystyle  - 20 \sqrt{3} i r z^{2} - 20 \sqrt{3} i r^{8} z^{3}$\\
$\displaystyle  - 12 \sqrt{3} i r^{9} z^{4} - 4 \sqrt{3} i r z + 4 \sqrt{3} i r^{8} z$\\
$\displaystyle  + 8 \sqrt{3} i r^{9} z + 8 \sqrt{3} i r^{9} z^{3} + 20 \sqrt{3} i r z^{3}$\\
$\displaystyle  + 20 \sqrt{3} i r^{8} z^{2} + 24 \sqrt{3} i r^{8} z^{5}$\\
$\displaystyle  + 24 \sqrt{3} i r^{4} z^{8} + 30 \sqrt{3} i r^{2} z$\\
$\displaystyle  + 42 \sqrt{3} i r^{7} z^{6} + 44 \sqrt{3} i r^{8} z^{4}$\\
$\displaystyle  + 66 \sqrt{3} i r^{3} z^{6} + 72 \sqrt{3} i r^{3} z^{7}$\\
$\displaystyle  + 120 \sqrt{3} i r^{7} z^{3} + 130 \sqrt{3} i r^{2} z^{5}$\\
$\displaystyle  + 134 \sqrt{3} i r^{3} z^{2} + 150 \sqrt{3} i r^{6} z$\\
$\displaystyle  + 150 \sqrt{3} i r^{7} z^{2} + 160 \sqrt{3} i r^{5} z^{7}$\\
$\displaystyle  + 184 \sqrt{3} i r^{4} z + 248 \sqrt{3} i r^{5} z^{2}$\\
$\displaystyle  + 274 \sqrt{3} i r^{6} z^{5} + 300 \sqrt{3} i r^{7} z^{4}$\\
$\displaystyle  + 344 \sqrt{3} i r^{5} z^{4} + 360 \sqrt{3} i r^{3} z^{3}$\\
$\displaystyle  + 364 \sqrt{3} i r^{3} z^{4} + 396 \sqrt{3} i r^{4} z^{5}$\\
$\displaystyle  + 444 \sqrt{3} i r^{5} z^{3} + 868 \sqrt{3} i r^{4} z^{6}$}
\label{eq:cat-numerator-77}
\end{equation}
\begin{equation}
D_{77}(r,z)=\left(- r^{2} + 2 r z - 1\right)^{2} \left(r^{2} + 2 r z + 1\right)^{2}
\label{eq:cat-denominator-77}
\end{equation}
\begin{equation}
R_{78}(r,z)=\frac{N_{78}(r,z)}{D_{78}(r,z)}.
\label{eq:cat-rational-78}
\end{equation}
\begin{equation}
N_{78}(r,z)=\shortstack[l]{$\displaystyle 24 r + 24 r^{9} + 146 r^{3} + 146 r^{7} + 244 r^{5}$\\
$\displaystyle  - 420 r^{5} z^{4} - 288 r^{4} z^{5} - 288 r^{6} z^{5}$\\
$\displaystyle  - 234 r^{3} z^{4} - 234 r^{7} z^{4} - 188 r^{5} z^{2}$\\
$\displaystyle  - 176 r^{4} z - 176 r^{6} z - 158 r^{3} z^{2}$\\
$\displaystyle  - 158 r^{7} z^{2} - 108 \sqrt{3} r^{4} - 108 \sqrt{3} r^{6}$\\
$\displaystyle  - 64 r^{5} z^{7} - 42 r^{3} z^{6} - 42 r^{7} z^{6}$\\
$\displaystyle  - 36 \sqrt{3} r^{2} - 36 \sqrt{3} r^{8} - 24 r^{5} z - 24 r^{2} z^{5}$\\
$\displaystyle  - 24 r^{8} z^{5} - 16 r^{3} z - 16 r^{7} z - 4 r z$\\
$\displaystyle  - 4 r^{9} z - 4 r z^{3} - 4 r^{9} z^{3} + 8 r^{2} z$\\
$\displaystyle  + 8 r^{8} z + 12 r z^{4} + 12 r^{9} z^{4} + 16 r^{2} z^{3}$\\
$\displaystyle  + 16 r^{3} z^{3} + 16 r^{7} z^{3} + 16 r^{8} z^{3}$\\
$\displaystyle  + 20 \sqrt{3} r + 20 \sqrt{3} r^{9} + 24 r^{5} z^{8} + 32 r^{3} z^{5}$\\
$\displaystyle  + 32 r^{7} z^{5} + 36 r z^{2} + 36 r^{9} z^{2}$\\
$\displaystyle  + 40 r^{5} z^{3} + 72 r^{4} z^{7} + 72 r^{6} z^{7}$\\
$\displaystyle  + 80 \sqrt{3} r^{3} + 80 \sqrt{3} r^{7} + 120 \sqrt{3} r^{5}$\\
$\displaystyle  + 392 r^{4} z^{3} + 392 r^{6} z^{3} + 772 r^{5} z^{6}$\\
$\displaystyle  - 216 \sqrt{3} r^{3} z^{3} - 216 \sqrt{3} r^{7} z^{3}$\\
$\displaystyle  - 176 \sqrt{3} r^{5} z^{2} - 160 \sqrt{3} r^{3} z^{4}$\\
$\displaystyle  - 160 \sqrt{3} r^{7} z^{4} - 144 \sqrt{3} r^{5} z^{3}$\\
$\displaystyle  - 96 \sqrt{3} r^{4} z^{6} - 96 \sqrt{3} r^{6} z^{6}$\\
$\displaystyle  - 80 \sqrt{3} r^{3} z^{2} - 80 \sqrt{3} r^{7} z^{2}$\\
$\displaystyle  - 80 \sqrt{3} r^{5} z^{7} - 36 \sqrt{3} r^{2} z^{5}$\\
$\displaystyle  - 36 \sqrt{3} r^{8} z^{5} - 28 \sqrt{3} r^{2} z^{4}$\\
$\displaystyle  - 28 \sqrt{3} r^{8} z^{4} - 24 \sqrt{3} r^{2} z^{3}$\\
$\displaystyle  - 24 \sqrt{3} r^{4} z^{3} - 24 \sqrt{3} r^{6} z^{3}$\\
$\displaystyle  - 24 \sqrt{3} r^{8} z^{3} - 12 \sqrt{3} r z^{4} - 12 \sqrt{3} r^{9} z^{4}$\\
$\displaystyle  - 12 \sqrt{3} r^{4} z^{5} - 12 \sqrt{3} r^{6} z^{5}$\\
$\displaystyle  - 4 \sqrt{3} r^{3} z^{5} - 4 \sqrt{3} r^{7} z^{5} + 8 \sqrt{3} r z$\\
$\displaystyle  + 8 \sqrt{3} r^{9} z + 8 \sqrt{3} r z^{2} + 8 \sqrt{3} r z^{3}$\\
$\displaystyle  + 8 \sqrt{3} r^{9} z^{2} + 8 \sqrt{3} r^{9} z^{3} + 12 \sqrt{3} r^{2} z$\\
$\displaystyle  + 12 \sqrt{3} r^{8} z + 16 \sqrt{3} r^{2} z^{2} + 16 \sqrt{3} r^{8} z^{2}$\\
$\displaystyle  + 24 \sqrt{3} r^{5} z^{4} + 36 \sqrt{3} r^{4} z + 36 \sqrt{3} r^{6} z$\\
$\displaystyle  + 48 \sqrt{3} r^{4} z^{7} + 48 \sqrt{3} r^{6} z^{7} + 60 \sqrt{3} r^{3} z$\\
$\displaystyle  + 60 \sqrt{3} r^{7} z + 76 \sqrt{3} r^{4} z^{4} + 76 \sqrt{3} r^{6} z^{4}$\\
$\displaystyle  + 104 \sqrt{3} r^{5} z + 176 \sqrt{3} r^{4} z^{2}$\\
$\displaystyle  + 176 \sqrt{3} r^{6} z^{2} + 320 \sqrt{3} r^{5} z^{6}$\\
$\displaystyle  + 408 \sqrt{3} r^{5} z^{5}$}
\label{eq:cat-numerator-78}
\end{equation}
\begin{equation}
D_{78}(r,z)=\left(- r^{2} + 2 r z - 1\right)^{2} \left(r^{2} + 2 r z + 1\right)^{2}
\label{eq:cat-denominator-78}
\end{equation}
\begin{equation}
R_{79}(r,z)=- \frac{4 i r^{2} \left(r - 1\right) \left(z - 1\right) \left(z + 1\right) \left(z^{2} + 1\right)}{r^{2} - 2 r z + 1}
\label{eq:cat-rational-79}
\end{equation}
\begin{equation}
R_{80}(r,z)=\frac{N_{80}(r,z)}{D_{80}(r,z)}.
\label{eq:cat-rational-80}
\end{equation}
\begin{equation}
N_{80}(r,z)=\shortstack[l]{$\displaystyle  - 4 \sqrt{3} + 4 \sqrt{3} z^{4} + 8 \sqrt{3} r^{2} + 12 \sqrt{3} r^{4}$\\
$\displaystyle  - 48 \sqrt{3} r^{2} z^{6} - 12 \sqrt{3} r^{4} z^{4}$\\
$\displaystyle  - 8 \sqrt{3} r^{2} z^{4} + 48 \sqrt{3} r^{2} z^{2}$}
\label{eq:cat-numerator-80}
\end{equation}
\begin{equation}
D_{80}(r,z)=\left(- r^{2} + 2 r z - 1\right) \left(r^{2} + 2 r z + 1\right)
\label{eq:cat-denominator-80}
\end{equation}
\begin{equation}
R_{81}(r,z)=\frac{N_{81}(r,z)}{D_{81}(r,z)}.
\label{eq:cat-rational-81}
\end{equation}
\begin{equation}
N_{81}(r,z)=\shortstack[l]{$\displaystyle 12 i - 24 i r^{6} - 12 i z^{4} - 12 i r^{8} + 4 \sqrt{3} i$\\
$\displaystyle  + 24 i r^{2} - 192 i r^{4} z^{8} - 112 \sqrt{3} i r^{5}$\\
$\displaystyle  - 96 i r^{2} z^{2} - 96 i r^{4} z^{2} - 88 \sqrt{3} i r^{7}$\\
$\displaystyle  - 56 \sqrt{3} i r^{3} - 40 \sqrt{3} i r^{4} - 40 \sqrt{3} i r^{6}$\\
$\displaystyle  - 24 \sqrt{3} i r^{9} - 24 i r^{2} z^{4} - 12 \sqrt{3} i r^{8}$\\
$\displaystyle  - 8 \sqrt{3} i r - 8 \sqrt{3} i r^{2} - 4 \sqrt{3} i z^{4}$\\
$\displaystyle  + 12 i r^{8} z^{4} + 24 i r^{6} z^{4} + 96 i r^{2} z^{6}$\\
$\displaystyle  + 96 i r^{4} z^{6} + 192 i r^{4} z^{4}$\\
$\displaystyle  - 752 \sqrt{3} i r^{5} z^{4} - 440 \sqrt{3} i r^{4} z^{4}$\\
$\displaystyle  - 408 \sqrt{3} i r^{6} z^{4} - 360 \sqrt{3} i r^{7} z^{4}$\\
$\displaystyle  - 120 \sqrt{3} i r^{3} z^{6} - 80 \sqrt{3} i r^{5} z^{6}$\\
$\displaystyle  - 56 \sqrt{3} i r^{2} z^{4} - 52 \sqrt{3} i r^{8} z^{4}$\\
$\displaystyle  - 40 \sqrt{3} i r^{9} z^{4} - 8 \sqrt{3} i r^{3} z^{4}$\\
$\displaystyle  + 8 \sqrt{3} i r z^{4} + 8 \sqrt{3} i r^{7} z^{6}$\\
$\displaystyle  + 24 \sqrt{3} i r^{2} z^{6} + 32 \sqrt{3} i r^{4} z^{8}$\\
$\displaystyle  + 40 \sqrt{3} i r^{2} z^{2} + 64 \sqrt{3} i r^{8} z^{2}$\\
$\displaystyle  + 64 \sqrt{3} i r^{9} z^{2} + 152 \sqrt{3} i r^{6} z^{6}$\\
$\displaystyle  + 184 \sqrt{3} i r^{3} z^{2} + 208 \sqrt{3} i r^{4} z^{6}$\\
$\displaystyle  + 240 \sqrt{3} i r^{4} z^{2} + 296 \sqrt{3} i r^{6} z^{2}$\\
$\displaystyle  + 416 \sqrt{3} i r^{5} z^{8} + 440 \sqrt{3} i r^{7} z^{2}$\\
$\displaystyle  + 528 \sqrt{3} i r^{5} z^{2}$}
\label{eq:cat-numerator-81}
\end{equation}
\begin{equation}
D_{81}(r,z)=\left(- r^{2} + 2 r z - 1\right)^{2} \left(r^{2} + 2 r z + 1\right)^{2}
\label{eq:cat-denominator-81}
\end{equation}
\begin{equation}
R_{82}(r,z)=\frac{N_{82}(r,z)}{D_{82}(r,z)}.
\label{eq:cat-rational-82}
\end{equation}
\begin{equation}
N_{82}(r,z)=\shortstack[l]{$\displaystyle  - 48 z^{2} + 36 r^{8} + 48 z^{4} + 72 r + 72 r^{2} + 72 r^{9}$\\
$\displaystyle  + 144 r^{6} + 180 r^{4} + 324 r^{3} + 324 r^{7} + 504 r^{5}$\\
$\displaystyle  - 1728 r^{5} z^{2} - 1272 r^{7} z^{2} - 1248 r^{5} z^{8}$\\
$\displaystyle  - 840 r^{3} z^{2} - 840 r^{6} z^{2} - 720 r^{4} z^{2}$\\
$\displaystyle  - 456 r^{6} z^{6} - 384 r^{4} z^{6} - 264 r^{2} z^{2}$\\
$\displaystyle  - 192 r^{8} z^{2} - 192 r^{9} z^{2} - 168 r^{2} z^{6}$\\
$\displaystyle  - 96 r^{4} z^{8} - 48 r z^{2} - 24 r z^{4} - 24 r^{7} z^{6}$\\
$\displaystyle  - 16 \sqrt{3} z^{2} + 16 \sqrt{3} z^{4} + 24 \sqrt{3} r + 36 \sqrt{3} r^{2}$\\
$\displaystyle  + 48 \sqrt{3} r^{10} + 72 \sqrt{3} r^{9} + 108 r^{3} z^{4}$\\
$\displaystyle  + 120 r^{9} z^{4} + 156 \sqrt{3} r^{3} + 156 r^{8} z^{4}$\\
$\displaystyle  + 168 \sqrt{3} r^{4} + 192 \sqrt{3} r^{8} + 252 \sqrt{3} r^{7}$\\
$\displaystyle  + 276 \sqrt{3} r^{6} + 312 \sqrt{3} r^{5} + 360 r^{2} z^{4}$\\
$\displaystyle  + 408 r^{3} z^{6} + 624 r^{5} z^{6} + 972 r^{7} z^{4}$\\
$\displaystyle  + 1020 r^{4} z^{4} + 1152 r^{6} z^{4} + 1848 r^{5} z^{4}$\\
$\displaystyle  - 1064 \sqrt{3} r^{6} z^{6} - 896 \sqrt{3} r^{4} z^{6}$\\
$\displaystyle  - 712 \sqrt{3} r^{6} z^{2} - 528 \sqrt{3} r^{4} z^{2}$\\
$\displaystyle  - 480 \sqrt{3} r^{5} z^{2} - 400 \sqrt{3} r^{8} z^{2}$\\
$\displaystyle  - 312 \sqrt{3} r^{7} z^{2} - 300 \sqrt{3} r^{7} z^{4}$\\
$\displaystyle  - 288 \sqrt{3} r^{5} z^{8} - 264 \sqrt{3} r^{3} z^{2}$\\
$\displaystyle  - 264 \sqrt{3} r^{5} z^{4} - 152 \sqrt{3} r^{2} z^{6}$\\
$\displaystyle  - 144 \sqrt{3} r^{8} z^{6} - 136 \sqrt{3} r^{2} z^{2}$\\
$\displaystyle  - 72 \sqrt{3} r^{9} z^{4} - 48 \sqrt{3} r z^{2}$\\
$\displaystyle  - 48 \sqrt{3} r^{10} z^{2} + 24 \sqrt{3} r z^{4}$\\
$\displaystyle  + 24 \sqrt{3} r^{3} z^{6} + 84 \sqrt{3} r^{3} z^{4}$\\
$\displaystyle  + 252 \sqrt{3} r^{2} z^{4} + 352 \sqrt{3} r^{8} z^{4}$\\
$\displaystyle  + 352 \sqrt{3} r^{4} z^{8} + 360 \sqrt{3} r^{7} z^{6}$\\
$\displaystyle  + 576 \sqrt{3} r^{6} z^{8} + 720 \sqrt{3} r^{5} z^{6}$\\
$\displaystyle  + 904 \sqrt{3} r^{4} z^{4} + 924 \sqrt{3} r^{6} z^{4}$}
\label{eq:cat-numerator-82}
\end{equation}
\begin{equation}
D_{82}(r,z)=3 \left(- r^{2} + 2 r z - 1\right)^{2} \left(r^{2} + 2 r z + 1\right)^{2}
\label{eq:cat-denominator-82}
\end{equation}
\begin{equation}
R_{83}(r,z)=\frac{N_{83}(r,z)}{D_{83}(r,z)}.
\label{eq:cat-rational-83}
\end{equation}
\begin{equation}
N_{83}(r,z)=\shortstack[l]{$\displaystyle  - 24 i r^{3} - 24 i r^{5} + 12 i r^{2} + 12 i r^{4}$\\
$\displaystyle  - 48 i r^{4} z^{2} - 24 i r^{3} z^{6} - 12 i r^{2} z^{4}$\\
$\displaystyle  - 12 i r^{4} z^{4} + 4 \sqrt{3} i r^{2} + 4 \sqrt{3} i r^{4}$\\
$\displaystyle  + 24 i r^{3} z^{2} + 24 i r^{3} z^{4} + 24 i r^{5} z^{4}$\\
$\displaystyle  + 48 i r^{4} z^{6} - 8 \sqrt{3} i r^{3} z^{6}$\\
$\displaystyle  - 4 \sqrt{3} i r^{2} z^{4} - 4 \sqrt{3} i r^{4} z^{4}$\\
$\displaystyle  + 8 \sqrt{3} i r^{3} z^{2}$}
\label{eq:cat-numerator-83}
\end{equation}
\begin{equation}
D_{83}(r,z)=\left(- r^{2} + 2 r z - 1\right) \left(r^{2} + 2 r z + 1\right)
\label{eq:cat-denominator-83}
\end{equation}
\begin{equation}
R_{84}(r,z)=32 \sqrt{3} i z^{2} \left(r + 1\right) \left(z - 1\right) \left(z + 1\right)
\label{eq:cat-rational-84}
\end{equation}
\begin{equation}
R_{85}(r,z)=- 32 z^{2} \left(r + 1\right) \left(z - 1\right) \left(z + 1\right) - \frac{16 \sqrt{3} \left(3 r^{2} - z^{2}\right) \left(z - 1\right) \left(z + 1\right)}{3}
\label{eq:cat-rational-85}
\end{equation}
\begin{equation}
R_{86}(r,z)=\frac{N_{86}(r,z)}{D_{86}(r,z)}.
\label{eq:cat-rational-86}
\end{equation}
\begin{equation}
N_{86}(r,z)=\shortstack[l]{$\displaystyle  - 28 \sqrt{3} - 88 \sqrt{3} r^{2} - 60 \sqrt{3} r^{4} + 28 \sqrt{3} z^{4}$\\
$\displaystyle  - 48 \sqrt{3} r^{2} z^{6} + 48 \sqrt{3} r^{2} z^{2}$\\
$\displaystyle  + 60 \sqrt{3} r^{4} z^{4} + 88 \sqrt{3} r^{2} z^{4}$}
\label{eq:cat-numerator-86}
\end{equation}
\begin{equation}
D_{86}(r,z)=\left(- r^{2} + 2 r z - 1\right) \left(r^{2} + 2 r z + 1\right)
\label{eq:cat-denominator-86}
\end{equation}
\begin{equation}
R_{87}(r,z)=\frac{N_{87}(r,z)}{D_{87}(r,z)}.
\label{eq:cat-rational-87}
\end{equation}
\begin{equation}
N_{87}(r,z)=\shortstack[l]{$\displaystyle 12 i - 12 i z^{4} + 4 \sqrt{3} i + 60 i r^{8} + 96 i r^{2}$\\
$\displaystyle  + 192 i r^{6} + 216 i r^{4} - 384 i r^{4} z^{2}$\\
$\displaystyle  - 352 \sqrt{3} i r^{5} - 288 i r^{6} z^{2} - 248 \sqrt{3} i r^{3}$\\
$\displaystyle  - 192 i r^{6} z^{4} - 192 i r^{4} z^{8} - 184 \sqrt{3} i r^{7}$\\
$\displaystyle  - 96 i r^{2} z^{2} - 96 i r^{2} z^{4} - 60 i r^{8} z^{4}$\\
$\displaystyle  - 56 \sqrt{3} i r - 36 \sqrt{3} i z^{4} - 24 \sqrt{3} i r^{9}$\\
$\displaystyle  - 24 i r^{4} z^{4} + 16 \sqrt{3} i r^{2} + 32 \sqrt{3} i z^{2}$\\
$\displaystyle  + 60 \sqrt{3} i r^{8} + 80 \sqrt{3} i r^{4} + 96 i r^{2} z^{6}$\\
$\displaystyle  + 128 \sqrt{3} i r^{6} + 288 i r^{6} z^{6} + 384 i r^{4} z^{6}$\\
$\displaystyle  - 1264 \sqrt{3} i r^{4} z^{4} - 1216 \sqrt{3} i r^{5} z^{4}$\\
$\displaystyle  - 960 \sqrt{3} i r^{6} z^{4} - 648 \sqrt{3} i r^{7} z^{4}$\\
$\displaystyle  - 480 \sqrt{3} i r^{4} z^{8} - 464 \sqrt{3} i r^{2} z^{4}$\\
$\displaystyle  - 200 \sqrt{3} i r^{3} z^{4} - 156 \sqrt{3} i r^{8} z^{4}$\\
$\displaystyle  - 96 \sqrt{3} i r^{5} z^{8} - 72 \sqrt{3} i r^{9} z^{4}$\\
$\displaystyle  - 24 \sqrt{3} i r^{7} z^{6} + 24 \sqrt{3} i r z^{4}$\\
$\displaystyle  + 32 \sqrt{3} i r z^{2} + 40 \sqrt{3} i r^{3} z^{6}$\\
$\displaystyle  + 72 \sqrt{3} i r^{2} z^{2} + 96 \sqrt{3} i r^{8} z^{2}$\\
$\displaystyle  + 96 \sqrt{3} i r^{9} z^{2} + 144 \sqrt{3} i r^{4} z^{2}$\\
$\displaystyle  + 376 \sqrt{3} i r^{2} z^{6} + 408 \sqrt{3} i r^{3} z^{2}$\\
$\displaystyle  + 408 \sqrt{3} i r^{6} z^{6} + 424 \sqrt{3} i r^{6} z^{2}$\\
$\displaystyle  + 752 \sqrt{3} i r^{5} z^{6} + 856 \sqrt{3} i r^{7} z^{2}$\\
$\displaystyle  + 912 \sqrt{3} i r^{5} z^{2} + 1520 \sqrt{3} i r^{4} z^{6}$}
\label{eq:cat-numerator-87}
\end{equation}
\begin{equation}
D_{87}(r,z)=\left(- r^{2} + 2 r z - 1\right)^{2} \left(r^{2} + 2 r z + 1\right)^{2}
\label{eq:cat-denominator-87}
\end{equation}
\begin{equation}
R_{88}(r,z)=\frac{N_{88}(r,z)}{D_{88}(r,z)}.
\label{eq:cat-rational-88}
\end{equation}
\begin{equation}
N_{88}(r,z)=\shortstack[l]{$\displaystyle  - 96 r^{6} - 60 r^{8} - 48 z^{2} - 12 r^{4} + 24 r + 24 r^{2}$\\
$\displaystyle  + 24 r^{9} + 48 z^{4} + 180 r^{3} + 180 r^{7} + 312 r^{5}$\\
$\displaystyle  - 1440 r^{4} z^{6} - 840 r^{7} z^{2} - 816 r^{5} z^{6}$\\
$\displaystyle  - 768 r^{5} z^{2} - 408 r^{3} z^{2} - 408 r^{6} z^{2}$\\
$\displaystyle  - 408 r^{6} z^{6} - 312 r^{2} z^{6} - 216 r^{2} z^{2}$\\
$\displaystyle  - 144 r^{4} z^{2} - 120 r^{3} z^{6} - 96 r^{8} z^{2}$\\
$\displaystyle  - 96 r^{9} z^{2} - 56 \sqrt{3} r^{4} - 48 r z^{2} - 28 \sqrt{3} r^{2}$\\
$\displaystyle  - 28 \sqrt{3} r^{6} + 8 \sqrt{3} r + 24 r z^{4} + 24 \sqrt{3} r^{9}$\\
$\displaystyle  + 24 r^{7} z^{6} + 72 r^{9} z^{4} + 76 \sqrt{3} r^{3}$\\
$\displaystyle  + 96 r^{5} z^{8} + 108 \sqrt{3} r^{7} + 152 \sqrt{3} r^{5}$\\
$\displaystyle  + 156 r^{8} z^{4} + 348 r^{3} z^{4} + 480 r^{4} z^{8}$\\
$\displaystyle  + 504 r^{2} z^{4} + 636 r^{7} z^{4} + 912 r^{6} z^{4}$\\
$\displaystyle  + 1116 r^{4} z^{4} + 1176 r^{5} z^{4} - 244 \sqrt{3} r^{6} z^{4}$\\
$\displaystyle  - 184 \sqrt{3} r^{3} z^{2} - 160 \sqrt{3} r^{5} z^{2}$\\
$\displaystyle  - 152 \sqrt{3} r^{6} z^{6} - 136 \sqrt{3} r^{5} z^{4}$\\
$\displaystyle  - 128 \sqrt{3} r^{4} z^{6} - 124 \sqrt{3} r^{7} z^{4}$\\
$\displaystyle  - 104 \sqrt{3} r^{7} z^{2} - 96 \sqrt{3} r^{5} z^{8}$\\
$\displaystyle  - 48 \sqrt{3} r^{8} z^{6} - 24 \sqrt{3} r^{9} z^{4} - 16 \sqrt{3} r z^{2}$\\
$\displaystyle  - 16 \sqrt{3} r^{8} z^{4} - 8 \sqrt{3} r^{2} z^{2}$\\
$\displaystyle  - 8 \sqrt{3} r^{2} z^{6} + 4 \sqrt{3} r^{3} z^{4} + 8 \sqrt{3} r z^{4}$\\
$\displaystyle  + 32 \sqrt{3} r^{4} z^{8} + 44 \sqrt{3} r^{2} z^{4}$\\
$\displaystyle  + 48 \sqrt{3} r^{4} z^{2} + 64 \sqrt{3} r^{8} z^{2}$\\
$\displaystyle  + 104 \sqrt{3} r^{4} z^{4} + 104 \sqrt{3} r^{3} z^{6}$\\
$\displaystyle  + 120 \sqrt{3} r^{7} z^{6} + 192 \sqrt{3} r^{6} z^{8}$\\
$\displaystyle  + 232 \sqrt{3} r^{6} z^{2} + 240 \sqrt{3} r^{5} z^{6}$}
\label{eq:cat-numerator-88}
\end{equation}
\begin{equation}
D_{88}(r,z)=\left(- r^{2} + 2 r z - 1\right)^{2} \left(r^{2} + 2 r z + 1\right)^{2}
\label{eq:cat-denominator-88}
\end{equation}
\begin{equation}
R_{89}(r,z)=\frac{N_{89}(r,z)}{D_{89}(r,z)}.
\label{eq:cat-rational-89}
\end{equation}
\begin{equation}
N_{89}(r,z)=\shortstack[l]{$\displaystyle 24 \sqrt{3} - 24 \sqrt{3} z^{4} + 72 \sqrt{3} r^{4} + 96 \sqrt{3} r^{2}$\\
$\displaystyle  - 96 \sqrt{3} r^{2} z^{4} - 72 \sqrt{3} r^{4} z^{4}$}
\label{eq:cat-numerator-89}
\end{equation}
\begin{equation}
D_{89}(r,z)=\left(- r^{2} + 2 r z - 1\right) \left(r^{2} + 2 r z + 1\right)
\label{eq:cat-denominator-89}
\end{equation}
\begin{equation}
R_{90}(r,z)=\frac{N_{90}(r,z)}{D_{90}(r,z)}.
\label{eq:cat-rational-90}
\end{equation}
\begin{equation}
N_{90}(r,z)=\shortstack[l]{$\displaystyle  - 216 i r^{4} - 216 i r^{6} - 72 i r^{2} - 72 i r^{8}$\\
$\displaystyle  - 288 i r^{4} z^{6} - 288 i r^{6} z^{6} - 168 \sqrt{3} i r^{6}$\\
$\displaystyle  - 120 \sqrt{3} i r^{4} - 72 \sqrt{3} i r^{8} - 24 \sqrt{3} i r^{2}$\\
$\displaystyle  + 48 \sqrt{3} i r + 72 i r^{2} z^{4} + 72 i r^{8} z^{4}$\\
$\displaystyle  + 96 \sqrt{3} i r^{7} + 192 \sqrt{3} i r^{3} + 216 i r^{4} z^{4}$\\
$\displaystyle  + 216 i r^{6} z^{4} + 240 \sqrt{3} i r^{5} + 288 i r^{4} z^{2}$\\
$\displaystyle  + 288 i r^{6} z^{2} - 288 \sqrt{3} i r^{7} z^{2}$\\
$\displaystyle  - 288 \sqrt{3} i r^{4} z^{6} - 240 \sqrt{3} i r^{5} z^{4}$\\
$\displaystyle  - 192 \sqrt{3} i r^{5} z^{2} - 192 \sqrt{3} i r^{3} z^{4}$\\
$\displaystyle  - 96 \sqrt{3} i r^{3} z^{2} - 96 \sqrt{3} i r^{7} z^{4}$\\
$\displaystyle  - 96 \sqrt{3} i r^{2} z^{6} - 48 \sqrt{3} i r z^{4}$\\
$\displaystyle  + 24 \sqrt{3} i r^{2} z^{4} + 72 \sqrt{3} i r^{8} z^{4}$\\
$\displaystyle  + 96 \sqrt{3} i r^{2} z^{2} + 96 \sqrt{3} i r^{3} z^{6}$\\
$\displaystyle  + 120 \sqrt{3} i r^{4} z^{4} + 168 \sqrt{3} i r^{6} z^{4}$\\
$\displaystyle  + 192 \sqrt{3} i r^{5} z^{6} + 288 \sqrt{3} i r^{4} z^{2}$\\
$\displaystyle  + 288 \sqrt{3} i r^{7} z^{6}$}
\label{eq:cat-numerator-90}
\end{equation}
\begin{equation}
D_{90}(r,z)=\left(- r^{2} + 2 r z - 1\right)^{2} \left(r^{2} + 2 r z + 1\right)^{2}
\label{eq:cat-denominator-90}
\end{equation}
\begin{equation}
R_{91}(r,z)=\frac{N_{91}(r,z)}{D_{91}(r,z)}.
\label{eq:cat-rational-91}
\end{equation}
\begin{equation}
N_{91}(r,z)=\shortstack[l]{$\displaystyle  - 144 r^{5} - 72 r^{3} - 72 r^{7} + 72 r^{4} + 72 r^{8}$\\
$\displaystyle  + 144 r^{6} - 288 r^{4} z^{2} - 288 r^{7} z^{6}$\\
$\displaystyle  - 144 r^{6} z^{4} - 72 r^{4} z^{4} - 72 r^{8} z^{4}$\\
$\displaystyle  - 48 \sqrt{3} r^{5} - 24 \sqrt{3} r^{3} - 24 \sqrt{3} r^{7}$\\
$\displaystyle  + 24 \sqrt{3} r^{2} + 24 \sqrt{3} r^{6} + 48 \sqrt{3} r^{4}$\\
$\displaystyle  + 72 r^{3} z^{4} + 72 r^{7} z^{4} + 144 r^{5} z^{4}$\\
$\displaystyle  + 288 r^{7} z^{2} + 288 r^{4} z^{6} - 96 \sqrt{3} r^{6} z^{2}$\\
$\displaystyle  - 96 \sqrt{3} r^{3} z^{6} - 48 \sqrt{3} r^{4} z^{4}$\\
$\displaystyle  - 24 \sqrt{3} r^{2} z^{4} - 24 \sqrt{3} r^{6} z^{4}$\\
$\displaystyle  + 24 \sqrt{3} r^{3} z^{4} + 24 \sqrt{3} r^{7} z^{4}$\\
$\displaystyle  + 48 \sqrt{3} r^{5} z^{4} + 96 \sqrt{3} r^{3} z^{2}$\\
$\displaystyle  + 96 \sqrt{3} r^{6} z^{6}$}
\label{eq:cat-numerator-91}
\end{equation}
\begin{equation}
D_{91}(r,z)=\left(- r^{2} + 2 r z - 1\right)^{2} \left(r^{2} + 2 r z + 1\right)^{2}
\label{eq:cat-denominator-91}
\end{equation}
\begin{equation}
R_{92}(r,z)=\frac{N_{92}(r,z)}{D_{92}(r,z)}.
\label{eq:cat-rational-92}
\end{equation}
\begin{equation}
N_{92}(r,z)=\shortstack[l]{$\displaystyle 32 \sqrt{3} i r^{2} + 32 \sqrt{3} i r^{3} + 32 \sqrt{3} i r^{4}$\\
$\displaystyle  + 32 \sqrt{3} i r^{5} - 96 \sqrt{3} i r^{2} z^{4}$\\
$\displaystyle  - 96 \sqrt{3} i r^{3} z^{4} - 96 \sqrt{3} i r^{4} z^{4}$\\
$\displaystyle  - 96 \sqrt{3} i r^{5} z^{4} + 64 \sqrt{3} i r^{2} z^{2}$\\
$\displaystyle  + 64 \sqrt{3} i r^{3} z^{2} + 64 \sqrt{3} i r^{4} z^{2}$\\
$\displaystyle  + 64 \sqrt{3} i r^{5} z^{2}$}
\label{eq:cat-numerator-92}
\end{equation}
\begin{equation}
D_{92}(r,z)=\left(- r^{2} + 2 r z - 1\right) \left(r^{2} + 2 r z + 1\right)
\label{eq:cat-denominator-92}
\end{equation}
\begin{equation}
R_{93}(r,z)=\frac{N_{93}(r,z)}{D_{93}(r,z)}.
\label{eq:cat-rational-93}
\end{equation}
\begin{equation}
N_{93}(r,z)=\shortstack[l]{$\displaystyle 48 - 96 r^{4} - 96 r^{5} - 48 z^{2} - 48 r^{2} - 48 r^{3}$\\
$\displaystyle  + 16 \sqrt{3} + 48 r - 192 r^{4} z^{2} - 192 r^{5} z^{2}$\\
$\displaystyle  - 144 r^{2} z^{2} - 144 r^{3} z^{2} - 48 r z^{2}$\\
$\displaystyle  - 16 \sqrt{3} z^{2} + 48 \sqrt{3} r + 48 \sqrt{3} r^{3} + 48 \sqrt{3} r^{6}$\\
$\displaystyle  + 128 \sqrt{3} r^{2} + 160 \sqrt{3} r^{4} + 192 r^{2} z^{4}$\\
$\displaystyle  + 192 r^{3} z^{4} + 288 r^{4} z^{4} + 288 r^{5} z^{4}$\\
$\displaystyle  - 256 \sqrt{3} r^{4} z^{2} - 192 \sqrt{3} r^{2} z^{2}$\\
$\displaystyle  - 96 \sqrt{3} r^{3} z^{4} - 48 \sqrt{3} r z^{2} - 48 \sqrt{3} r^{6} z^{2}$\\
$\displaystyle  + 48 \sqrt{3} r^{3} z^{2} + 64 \sqrt{3} r^{2} z^{4}$\\
$\displaystyle  + 96 \sqrt{3} r^{4} z^{4}$}
\label{eq:cat-numerator-93}
\end{equation}
\begin{equation}
D_{93}(r,z)=3 \left(- r^{2} + 2 r z - 1\right) \left(r^{2} + 2 r z + 1\right)
\label{eq:cat-denominator-93}
\end{equation}
\begin{equation}
R_{94}(r,z)=32 \sqrt{3} i \left(r + 1\right) \left(z - 1\right) \left(z + 1\right)
\label{eq:cat-rational-94}
\end{equation}
\begin{equation}
R_{95}(r,z)=- 32 \left(r + 1\right) \left(z - 1\right) \left(z + 1\right) + \frac{16 \sqrt{3} \left(3 r^{2} + 1\right) \left(z - 1\right) \left(z + 1\right)}{3}
\label{eq:cat-rational-95}
\end{equation}
\begin{equation}
R_{96}(r,z)=\frac{N_{96}(r,z)}{D_{96}(r,z)}.
\label{eq:cat-rational-96}
\end{equation}
\begin{equation}
N_{96}(r,z)=\shortstack[l]{$\displaystyle  - 32 \sqrt{3} i - 32 \sqrt{3} i r - 32 \sqrt{3} i r^{2} - 32 \sqrt{3} i r^{3}$\\
$\displaystyle  + 32 \sqrt{3} i z^{2} - 224 \sqrt{3} i r^{2} z^{4}$\\
$\displaystyle  - 224 \sqrt{3} i r^{3} z^{4} - 96 \sqrt{3} i r^{4} z^{4}$\\
$\displaystyle  - 96 \sqrt{3} i r^{5} z^{4} + 32 \sqrt{3} i r z^{2}$\\
$\displaystyle  + 96 \sqrt{3} i r^{4} z^{2} + 96 \sqrt{3} i r^{5} z^{2}$\\
$\displaystyle  + 256 \sqrt{3} i r^{2} z^{2} + 256 \sqrt{3} i r^{3} z^{2}$}
\label{eq:cat-numerator-96}
\end{equation}
\begin{equation}
D_{96}(r,z)=\left(- r^{2} + 2 r z - 1\right) \left(r^{2} + 2 r z + 1\right)
\label{eq:cat-denominator-96}
\end{equation}
\begin{equation}
R_{97}(r,z)=\frac{N_{97}(r,z)}{D_{97}(r,z)}.
\label{eq:cat-rational-97}
\end{equation}
\begin{equation}
N_{97}(r,z)=\shortstack[l]{$\displaystyle 48 - 48 z^{2} + 48 r + 48 r^{2} + 48 r^{3} - 240 r^{2} z^{2}$\\
$\displaystyle  - 240 r^{3} z^{2} - 96 r^{4} z^{2} - 96 r^{5} z^{2}$\\
$\displaystyle  - 48 r z^{2} + 16 \sqrt{3} r + 16 \sqrt{3} r^{2} + 16 \sqrt{3} r^{3}$\\
$\displaystyle  + 16 \sqrt{3} r^{4} + 96 r^{4} z^{4} + 96 r^{5} z^{4}$\\
$\displaystyle  + 192 r^{2} z^{4} + 192 r^{3} z^{4} - 32 \sqrt{3} r^{3} z^{4}$\\
$\displaystyle  - 32 \sqrt{3} r^{4} z^{4} - 16 \sqrt{3} r z^{2} - 16 \sqrt{3} r^{2} z^{2}$\\
$\displaystyle  + 16 \sqrt{3} r^{3} z^{2} + 16 \sqrt{3} r^{4} z^{2}$}
\label{eq:cat-numerator-97}
\end{equation}
\begin{equation}
D_{97}(r,z)=\left(- r^{2} + 2 r z - 1\right) \left(r^{2} + 2 r z + 1\right)
\label{eq:cat-denominator-97}
\end{equation}
\begin{equation}
R_{98}(r,z)=\frac{N_{98}(r,z)}{D_{98}(r,z)}.
\label{eq:cat-rational-98}
\end{equation}
\begin{equation}
N_{98}(r,z)=\shortstack[l]{$\displaystyle  - 24 r - 24 r^{3} + 36 z^{3} + 100 z - 520 r^{2} z^{3}$\\
$\displaystyle  + 24 r z^{4} + 24 r^{3} z^{4} + 36 r^{4} z^{3} + 100 r^{4} z$\\
$\displaystyle  + 248 r^{2} z$}
\label{eq:cat-numerator-98}
\end{equation}
\begin{equation}
D_{98}(r,z)=\left(- r^{2} + 2 r z - 1\right) \left(r^{2} + 2 r z + 1\right)
\label{eq:cat-denominator-98}
\end{equation}
\begin{equation}
R_{99}(r,z)=\frac{N_{99}(r,z)}{D_{99}(r,z)}.
\label{eq:cat-rational-99}
\end{equation}
\begin{equation}
N_{99}(r,z)=\shortstack[l]{$\displaystyle  - 48 i r^{3} - 24 i r - 24 i r^{5} - 12 i z^{3} + 20 i z$\\
$\displaystyle  + 24 i r^{2} + 24 i r^{6} + 48 i r^{4} - 728 i r^{3} z^{4}$\\
$\displaystyle  - 580 i r^{3} z^{3} - 376 i r^{4} z^{2} - 276 i r^{4} z$\\
$\displaystyle  - 272 i r^{2} z^{2} - 200 i r^{6} z^{2} - 196 i r^{6} z$\\
$\displaystyle  - 144 i r^{4} z^{6} - 120 i r^{5} z^{4} - 116 i r^{2} z^{3}$\\
$\displaystyle  - 60 i r^{2} z - 60 i r z^{3} - 48 i r z^{4}$\\
$\displaystyle  - 32 \sqrt{3} i r^{4} - 24 i r^{4} z^{5} - 24 i r^{5} z^{5}$\\
$\displaystyle  - 20 i r^{7} z - 16 \sqrt{3} i r^{2} - 16 \sqrt{3} i r^{6}$\\
$\displaystyle  - 12 \sqrt{3} i z^{3} + 4 \sqrt{3} i z + 12 i r^{7} z^{3}$\\
$\displaystyle  + 16 \sqrt{3} i r + 16 \sqrt{3} i r^{5} + 24 i r^{2} z^{5}$\\
$\displaystyle  + 24 i r^{3} z^{5} + 32 \sqrt{3} i r^{3} + 48 i r^{6} z^{4}$\\
$\displaystyle  + 60 i r^{5} z + 60 i r^{6} z^{3} + 116 i r^{5} z^{3}$\\
$\displaystyle  + 120 i r^{2} z^{4} + 144 i r^{3} z^{6} + 196 i r z$\\
$\displaystyle  + 200 i r z^{2} + 272 i r^{5} z^{2} + 276 i r^{3} z$\\
$\displaystyle  + 376 i r^{3} z^{2} + 580 i r^{4} z^{3} + 728 i r^{4} z^{4}$\\
$\displaystyle  - 108 \sqrt{3} i r^{5} z^{3} - 96 \sqrt{3} i r^{3} z^{4}$\\
$\displaystyle  - 76 \sqrt{3} i r^{3} z^{3} - 68 \sqrt{3} i r^{4} z$\\
$\displaystyle  - 60 \sqrt{3} i r^{2} z - 48 \sqrt{3} i r^{6} z^{4}$\\
$\displaystyle  - 48 \sqrt{3} i r^{2} z^{5} - 40 \sqrt{3} i r z^{2}$\\
$\displaystyle  - 16 \sqrt{3} i r^{4} z^{2} - 16 \sqrt{3} i r^{2} z^{4}$\\
$\displaystyle  - 16 \sqrt{3} i r^{3} z^{5} - 12 \sqrt{3} i r^{6} z^{3}$\\
$\displaystyle  - 8 \sqrt{3} i r^{5} z^{2} - 4 \sqrt{3} i r^{6} z - 4 \sqrt{3} i r^{7} z$\\
$\displaystyle  + 4 \sqrt{3} i r z + 8 \sqrt{3} i r^{2} z^{2} + 12 \sqrt{3} i r z^{3}$\\
$\displaystyle  + 12 \sqrt{3} i r^{7} z^{3} + 16 \sqrt{3} i r^{3} z^{2}$\\
$\displaystyle  + 16 \sqrt{3} i r^{5} z^{4} + 16 \sqrt{3} i r^{4} z^{5}$\\
$\displaystyle  + 40 \sqrt{3} i r^{6} z^{2} + 48 \sqrt{3} i r z^{4}$\\
$\displaystyle  + 48 \sqrt{3} i r^{5} z^{5} + 60 \sqrt{3} i r^{5} z$\\
$\displaystyle  + 68 \sqrt{3} i r^{3} z + 76 \sqrt{3} i r^{4} z^{3}$\\
$\displaystyle  + 96 \sqrt{3} i r^{4} z^{4} + 108 \sqrt{3} i r^{2} z^{3}$}
\label{eq:cat-numerator-99}
\end{equation}
\begin{equation}
D_{99}(r,z)=\left(- r^{2} + 2 r z - 1\right)^{2} \left(r^{2} + 2 r z + 1\right)
\label{eq:cat-denominator-99}
\end{equation}
\begin{equation}
R_{100}(r,z)=\frac{N_{100}(r,z)}{D_{100}(r,z)}.
\label{eq:cat-rational-100}
\end{equation}
\begin{equation}
N_{100}(r,z)=\shortstack[l]{$\displaystyle 8 r + 8 r^{7} + 24 r^{2} + 24 r^{3} + 24 r^{5} + 24 r^{6}$\\
$\displaystyle  + 48 r^{4} - 96 r^{2} z^{2} - 96 r^{4} z^{2} - 96 r^{6} z^{2}$\\
$\displaystyle  - 96 r^{3} z^{5} - 96 r^{5} z^{5} - 60 r^{3} z - 60 r^{5} z$\\
$\displaystyle  - 56 r^{3} z^{2} - 56 r^{5} z^{2} - 48 r^{3} z^{6}$\\
$\displaystyle  - 48 r^{5} z^{6} - 32 r^{4} z - 16 r^{2} z - 16 r^{6} z$\\
$\displaystyle  - 8 r z^{2} - 8 r^{7} z^{2} - 8 r^{2} z^{3} - 8 r^{6} z^{3}$\\
$\displaystyle  + 8 \sqrt{3} r + 8 \sqrt{3} r^{7} + 12 r z + 12 r^{7} z + 12 r z^{4}$\\
$\displaystyle  + 12 r^{7} z^{4} + 24 \sqrt{3} r^{3} + 24 \sqrt{3} r^{5}$\\
$\displaystyle  + 32 r^{4} z^{5} + 48 r^{4} z^{3} + 48 r^{2} z^{4}$\\
$\displaystyle  + 48 r^{6} z^{4} + 68 r^{3} z^{4} + 68 r^{5} z^{4}$\\
$\displaystyle  + 96 r^{4} z^{4} + 144 r^{3} z^{3} + 144 r^{5} z^{3}$\\
$\displaystyle  - 280 \sqrt{3} r^{4} z^{4} - 160 \sqrt{3} r^{4} z^{3}$\\
$\displaystyle  - 96 \sqrt{3} r^{3} z^{2} - 96 \sqrt{3} r^{5} z^{2}$\\
$\displaystyle  - 48 \sqrt{3} r^{3} z^{6} - 48 \sqrt{3} r^{5} z^{6}$\\
$\displaystyle  - 40 \sqrt{3} r^{3} z^{5} - 40 \sqrt{3} r^{5} z^{5} - 36 \sqrt{3} r^{3} z$\\
$\displaystyle  - 36 \sqrt{3} r^{5} z - 28 \sqrt{3} r^{2} z^{4} - 28 \sqrt{3} r^{6} z^{4}$\\
$\displaystyle  - 24 \sqrt{3} r z^{2} - 24 \sqrt{3} r^{7} z^{2} - 12 \sqrt{3} r z$\\
$\displaystyle  - 12 \sqrt{3} r^{7} z - 8 \sqrt{3} r^{2} z^{3} - 8 \sqrt{3} r^{6} z^{3}$\\
$\displaystyle  + 8 \sqrt{3} r z^{3} + 8 \sqrt{3} r^{7} z^{3} + 12 \sqrt{3} r z^{4}$\\
$\displaystyle  + 12 \sqrt{3} r^{7} z^{4} + 16 \sqrt{3} r^{4} z^{5} + 40 \sqrt{3} r^{2} z$\\
$\displaystyle  + 40 \sqrt{3} r^{6} z + 60 \sqrt{3} r^{2} z^{2} + 60 \sqrt{3} r^{6} z^{2}$\\
$\displaystyle  + 80 \sqrt{3} r^{4} z + 80 \sqrt{3} r^{3} z^{3} + 80 \sqrt{3} r^{5} z^{3}$\\
$\displaystyle  + 96 \sqrt{3} r^{4} z^{6} + 120 \sqrt{3} r^{4} z^{2}$\\
$\displaystyle  + 124 \sqrt{3} r^{3} z^{4} + 124 \sqrt{3} r^{5} z^{4}$}
\label{eq:cat-numerator-100}
\end{equation}
\begin{equation}
D_{100}(r,z)=\left(- r^{2} + 2 r z - 1\right)^{2} \left(r^{2} + 2 r z + 1\right)
\label{eq:cat-denominator-100}
\end{equation}
\begin{equation}
R_{101}(r,z)=\frac{N_{101}(r,z)}{D_{101}(r,z)}.
\label{eq:cat-rational-101}
\end{equation}
\begin{equation}
N_{101}(r,z)=\shortstack[l]{$\displaystyle  - 16 i r^{3} z^{3} - 12 i r^{3} z^{4} - 8 i r^{2} z$\\
$\displaystyle  - 4 i r^{2} z^{2} + 4 i r^{3} z^{2} + 8 i r^{3} z$\\
$\displaystyle  + 12 i r^{2} z^{4} + 16 i r^{2} z^{3} - 8 \sqrt{3} i r^{3} z^{3}$\\
$\displaystyle  - 4 \sqrt{3} i r^{3} z^{2} - 4 \sqrt{3} i r^{3} z^{4}$\\
$\displaystyle  + 4 \sqrt{3} i r^{2} z^{2} + 4 \sqrt{3} i r^{2} z^{4}$\\
$\displaystyle  + 8 \sqrt{3} i r^{2} z^{3}$}
\label{eq:cat-numerator-101}
\end{equation}
\begin{equation}
D_{101}(r,z)=- r^{2} + 2 r z - 1
\label{eq:cat-denominator-101}
\end{equation}
\begin{equation}
R_{102}(r,z)=- \frac{8 z \left(3 r^{2} z - 18 r z^{2} - 8 r + 3 z\right)}{r}
\label{eq:cat-rational-102}
\end{equation}
\begin{equation}
R_{103}(r,z)=- 8 i z \left(r - 1\right) \left(3 z + 5\right) - 8 \sqrt{3} i z \left(r - 1\right) \left(6 z^{2} + z + 1\right)
\label{eq:cat-rational-103}
\end{equation}
\begin{equation}
R_{104}(r,z)=16 \sqrt{3} r z \left(z + 1\right) + 8 r \left(3 z + 1\right) \left(2 z^{2} - z + 1\right)
\label{eq:cat-rational-104}
\end{equation}
\begin{equation}
R_{105}(r,z)=\frac{N_{105}(r,z)}{D_{105}(r,z)}.
\label{eq:cat-rational-105}
\end{equation}
\begin{equation}
N_{105}(r,z)=\shortstack[l]{$\displaystyle  - 72 z^{2} - 32 r^{2} - 32 r^{4} - 432 r^{3} z^{5}$\\
$\displaystyle  - 184 r^{2} z^{2} - 184 r^{4} z^{2} - 168 r^{3} z - 116 r z$\\
$\displaystyle  - 116 r^{5} z - 72 r^{6} z^{2} + 12 r z^{3} + 12 r^{5} z^{3}$\\
$\displaystyle  + 288 r^{2} z^{4} + 288 r^{4} z^{4} + 808 r^{3} z^{3}$}
\label{eq:cat-numerator-105}
\end{equation}
\begin{equation}
D_{105}(r,z)=r \left(- r^{2} + 2 r z - 1\right) \left(r^{2} + 2 r z + 1\right)
\label{eq:cat-denominator-105}
\end{equation}
\begin{equation}
R_{106}(r,z)=\frac{N_{106}(r,z)}{D_{106}(r,z)}.
\label{eq:cat-rational-106}
\end{equation}
\begin{equation}
N_{106}(r,z)=\shortstack[l]{$\displaystyle  - 84 i z^{3} - 72 i z^{2} - 32 i r^{3} - 16 i r - 16 i r^{5}$\\
$\displaystyle  + 16 i r^{2} + 16 i r^{6} + 20 i z + 32 i r^{4}$\\
$\displaystyle  - 432 i r^{5} z^{4} - 324 i r^{4} z^{3} - 264 i r^{2} z^{2}$\\
$\displaystyle  - 224 i r^{4} z^{2} - 216 i r^{3} z^{5} - 184 i r^{3} z^{4}$\\
$\displaystyle  - 168 i r^{5} z^{5} - 144 i r^{4} z^{6} - 120 i r z^{4}$\\
$\displaystyle  - 96 i r^{6} z^{2} - 84 i r z^{3} - 60 i r^{2} z^{3}$\\
$\displaystyle  - 52 i r^{5} z - 52 i r^{6} z - 36 \sqrt{3} i z^{3}$\\
$\displaystyle  - 32 \sqrt{3} i r^{4} - 24 \sqrt{3} i z^{2} - 20 i r^{4} z$\\
$\displaystyle  - 20 i r^{7} z - 16 \sqrt{3} i r^{2} - 16 \sqrt{3} i r^{6}$\\
$\displaystyle  + 16 \sqrt{3} i r + 16 \sqrt{3} i r^{5} + 20 i r^{3} z$\\
$\displaystyle  + 32 \sqrt{3} i r^{3} + 52 i r z + 52 \sqrt{3} i z + 52 i r^{2} z$\\
$\displaystyle  + 60 i r^{5} z^{3} + 72 i r^{7} z^{2} + 84 i r^{6} z^{3}$\\
$\displaystyle  + 84 i r^{7} z^{3} + 96 i r z^{2} + 120 i r^{6} z^{4}$\\
$\displaystyle  + 144 i r^{3} z^{6} + 168 i r^{2} z^{5} + 184 i r^{4} z^{4}$\\
$\displaystyle  + 216 i r^{4} z^{5} + 224 i r^{3} z^{2} + 264 i r^{5} z^{2}$\\
$\displaystyle  + 324 i r^{3} z^{3} + 432 i r^{2} z^{4}$\\
$\displaystyle  - 400 \sqrt{3} i r^{4} z^{4} - 348 \sqrt{3} i r^{4} z^{3}$\\
$\displaystyle  - 336 \sqrt{3} i r^{3} z^{5} - 288 \sqrt{3} i r^{3} z^{6}$\\
$\displaystyle  - 268 \sqrt{3} i r^{2} z^{3} - 144 \sqrt{3} i r^{5} z^{5}$\\
$\displaystyle  - 112 \sqrt{3} i r z^{2} - 88 \sqrt{3} i r^{5} z^{4}$\\
$\displaystyle  - 76 \sqrt{3} i r^{5} z - 72 \sqrt{3} i r^{3} z^{2}$\\
$\displaystyle  - 72 \sqrt{3} i r^{6} z^{4} - 60 \sqrt{3} i r^{3} z$\\
$\displaystyle  - 52 \sqrt{3} i r^{7} z - 52 \sqrt{3} i r^{6} z^{3} - 36 \sqrt{3} i r z$\\
$\displaystyle  + 24 \sqrt{3} i r^{7} z^{2} + 36 \sqrt{3} i r^{6} z$\\
$\displaystyle  + 36 \sqrt{3} i r^{7} z^{3} + 52 \sqrt{3} i r z^{3}$\\
$\displaystyle  + 60 \sqrt{3} i r^{4} z + 72 \sqrt{3} i r z^{4}$\\
$\displaystyle  + 72 \sqrt{3} i r^{4} z^{2} + 76 \sqrt{3} i r^{2} z$\\
$\displaystyle  + 88 \sqrt{3} i r^{2} z^{4} + 112 \sqrt{3} i r^{6} z^{2}$\\
$\displaystyle  + 144 \sqrt{3} i r^{2} z^{5} + 268 \sqrt{3} i r^{5} z^{3}$\\
$\displaystyle  + 288 \sqrt{3} i r^{4} z^{6} + 336 \sqrt{3} i r^{4} z^{5}$\\
$\displaystyle  + 348 \sqrt{3} i r^{3} z^{3} + 400 \sqrt{3} i r^{3} z^{4}$}
\label{eq:cat-numerator-106}
\end{equation}
\begin{equation}
D_{106}(r,z)=\left(- r^{2} + 2 r z - 1\right)^{2} \left(r^{2} + 2 r z + 1\right)
\label{eq:cat-denominator-106}
\end{equation}
\begin{equation}
R_{107}(r,z)=\frac{N_{107}(r,z)}{D_{107}(r,z)}.
\label{eq:cat-rational-107}
\end{equation}
\begin{equation}
N_{107}(r,z)=\shortstack[l]{$\displaystyle 32 r^{2} + 32 r^{6} + 64 r^{4} - 288 r^{4} z^{6}$\\
$\displaystyle  - 160 r^{4} z^{2} - 144 r^{2} z^{2} - 144 r^{6} z^{2}$\\
$\displaystyle  - 120 r^{3} z^{3} - 120 r^{5} z^{3} - 48 r^{3} z^{6}$\\
$\displaystyle  - 48 r^{5} z^{6} - 24 r z^{3} - 24 r^{3} z^{2}$\\
$\displaystyle  - 24 r^{5} z^{2} - 24 r^{7} z^{3} - 16 r^{4} z^{3}$\\
$\displaystyle  - 8 r z^{2} - 8 r^{7} z^{2} - 8 r^{2} z^{3} - 8 r^{6} z^{3}$\\
$\displaystyle  + 4 r^{3} z + 4 r^{5} z + 8 \sqrt{3} r + 8 \sqrt{3} r^{7}$\\
$\displaystyle  + 12 r z^{4} + 12 r^{7} z^{4} + 24 \sqrt{3} r^{3} + 24 \sqrt{3} r^{5}$\\
$\displaystyle  + 32 r^{4} z^{5} + 44 r z + 44 r^{7} z + 68 r^{3} z^{4}$\\
$\displaystyle  + 68 r^{5} z^{4} + 72 r^{2} z^{4} + 72 r^{6} z^{4}$\\
$\displaystyle  + 96 r^{3} z^{5} + 96 r^{5} z^{5} + 464 r^{4} z^{4}$\\
$\displaystyle  - 144 \sqrt{3} r^{4} z^{5} - 136 \sqrt{3} r^{4} z^{4}$\\
$\displaystyle  - 112 \sqrt{3} r^{3} z^{2} - 112 \sqrt{3} r^{5} z^{2}$\\
$\displaystyle  - 52 \sqrt{3} r^{2} z^{4} - 52 \sqrt{3} r^{6} z^{4}$\\
$\displaystyle  - 48 \sqrt{3} r^{3} z^{6} - 48 \sqrt{3} r^{5} z^{6}$\\
$\displaystyle  - 40 \sqrt{3} r^{3} z^{3} - 40 \sqrt{3} r^{5} z^{3} - 24 \sqrt{3} r z^{2}$\\
$\displaystyle  - 24 \sqrt{3} r^{7} z^{2} - 24 \sqrt{3} r^{2} z^{3}$\\
$\displaystyle  - 24 \sqrt{3} r^{6} z^{3} - 16 \sqrt{3} r z^{3} - 16 \sqrt{3} r^{7} z^{3}$\\
$\displaystyle  + 8 \sqrt{3} r^{2} z + 8 \sqrt{3} r^{6} z + 8 \sqrt{3} r^{3} z^{5}$\\
$\displaystyle  + 8 \sqrt{3} r^{5} z^{5} + 12 \sqrt{3} r z + 12 \sqrt{3} r^{7} z$\\
$\displaystyle  + 12 \sqrt{3} r z^{4} + 12 \sqrt{3} r^{7} z^{4} + 16 \sqrt{3} r^{4} z$\\
$\displaystyle  + 36 \sqrt{3} r^{3} z + 36 \sqrt{3} r^{5} z + 36 \sqrt{3} r^{2} z^{2}$\\
$\displaystyle  + 36 \sqrt{3} r^{6} z^{2} + 72 \sqrt{3} r^{4} z^{2}$\\
$\displaystyle  + 96 \sqrt{3} r^{4} z^{6} + 140 \sqrt{3} r^{3} z^{4}$\\
$\displaystyle  + 140 \sqrt{3} r^{5} z^{4} + 160 \sqrt{3} r^{4} z^{3}$}
\label{eq:cat-numerator-107}
\end{equation}
\begin{equation}
D_{107}(r,z)=\left(- r^{2} + 2 r z - 1\right)^{2} \left(r^{2} + 2 r z + 1\right)
\label{eq:cat-denominator-107}
\end{equation}
\begin{equation}
R_{108}(r,z)=\frac{N_{108}(r,z)}{D_{108}(r,z)}.
\label{eq:cat-rational-108}
\end{equation}
\begin{equation}
N_{108}(r,z)=\shortstack[l]{$\displaystyle  - 12 i r^{3} z^{4} - 8 i r^{3} z^{3} - 4 i r^{2} z^{2}$\\
$\displaystyle  + 4 i r^{3} z^{2} + 8 i r^{2} z^{3} + 12 i r^{2} z^{4}$\\
$\displaystyle  - 8 \sqrt{3} i r^{3} z^{3} - 4 \sqrt{3} i r^{3} z^{2}$\\
$\displaystyle  - 4 \sqrt{3} i r^{3} z^{4} + 4 \sqrt{3} i r^{2} z^{2}$\\
$\displaystyle  + 4 \sqrt{3} i r^{2} z^{4} + 8 \sqrt{3} i r^{2} z^{3}$}
\label{eq:cat-numerator-108}
\end{equation}
\begin{equation}
D_{108}(r,z)=- r^{2} + 2 r z - 1
\label{eq:cat-denominator-108}
\end{equation}
\begin{equation}
R_{109}(r,z)=\frac{N_{109}(r,z)}{D_{109}(r,z)}.
\label{eq:cat-rational-109}
\end{equation}
\begin{equation}
N_{109}(r,z)=\shortstack[l]{$\displaystyle 8 r^{2} + 8 r^{4} + 48 z^{2} - 1296 r^{3} z^{3}$\\
$\displaystyle  - 168 r^{2} z^{4} - 168 r^{4} z^{4} - 144 r^{3} z^{5}$\\
$\displaystyle  + 48 r^{6} z^{2} + 112 r^{2} z^{2} + 112 r^{4} z^{2}$\\
$\displaystyle  + 168 r z^{3} + 168 r^{5} z^{3} + 280 r z + 280 r^{5} z$\\
$\displaystyle  + 544 r^{3} z$}
\label{eq:cat-numerator-109}
\end{equation}
\begin{equation}
D_{109}(r,z)=r \left(- r^{2} + 2 r z - 1\right) \left(r^{2} + 2 r z + 1\right)
\label{eq:cat-denominator-109}
\end{equation}
\begin{equation}
R_{110}(r,z)=\frac{N_{110}(r,z)}{D_{110}(r,z)}.
\label{eq:cat-rational-110}
\end{equation}
\begin{equation}
N_{110}(r,z)=\shortstack[l]{$\displaystyle  - 40 i z - 16 i r^{3} - 8 i r - 8 i r^{5} + 8 i r^{2}$\\
$\displaystyle  + 8 i r^{6} + 16 i r^{4} + 48 i z^{2} + 72 i z^{3}$\\
$\displaystyle  - 968 i r^{3} z^{3} - 960 i r^{3} z^{4} - 384 i r^{4} z^{2}$\\
$\displaystyle  - 376 i r^{4} z - 232 i r^{2} z - 216 i r^{2} z^{4}$\\
$\displaystyle  - 208 i r^{6} z^{2} - 184 i r^{6} z - 160 i r^{2} z^{2}$\\
$\displaystyle  - 144 i r^{2} z^{5} - 72 i r^{6} z^{3} - 72 i r^{7} z^{3}$\\
$\displaystyle  - 72 i r^{6} z^{4} - 56 \sqrt{3} i z - 56 i r^{5} z^{3}$\\
$\displaystyle  - 48 i r^{7} z^{2} - 48 i r^{4} z^{5} - 24 \sqrt{3} i z^{3}$\\
$\displaystyle  + 16 \sqrt{3} i z^{2} + 40 i r^{7} z + 48 i r^{3} z^{5}$\\
$\displaystyle  + 56 i r^{2} z^{3} + 72 i r z^{3} + 72 i r z^{4}$\\
$\displaystyle  + 144 i r^{5} z^{5} + 160 i r^{5} z^{2} + 184 i r z$\\
$\displaystyle  + 208 i r z^{2} + 216 i r^{5} z^{4} + 232 i r^{5} z$\\
$\displaystyle  + 376 i r^{3} z + 384 i r^{3} z^{2} + 960 i r^{4} z^{4}$\\
$\displaystyle  + 968 i r^{4} z^{3} - 400 \sqrt{3} i r^{3} z^{4}$\\
$\displaystyle  - 280 \sqrt{3} i r^{3} z^{3} - 248 \sqrt{3} i r^{5} z^{3}$\\
$\displaystyle  - 168 \sqrt{3} i r^{2} z^{4} - 160 \sqrt{3} i r^{2} z$\\
$\displaystyle  - 152 \sqrt{3} i r^{4} z - 144 \sqrt{3} i r^{4} z^{2}$\\
$\displaystyle  - 96 \sqrt{3} i r^{6} z^{2} - 96 \sqrt{3} i r^{3} z^{6}$\\
$\displaystyle  - 72 \sqrt{3} i r^{6} z^{4} - 64 \sqrt{3} i r^{4} z^{5}$\\
$\displaystyle  - 48 \sqrt{3} i r^{6} z - 32 \sqrt{3} i r^{2} z^{2}$\\
$\displaystyle  - 24 \sqrt{3} i r^{6} z^{3} - 16 \sqrt{3} i r^{7} z^{2}$\\
$\displaystyle  + 24 \sqrt{3} i r z^{3} + 24 \sqrt{3} i r^{7} z^{3}$\\
$\displaystyle  + 32 \sqrt{3} i r^{5} z^{2} + 48 \sqrt{3} i r z + 56 \sqrt{3} i r^{7} z$\\
$\displaystyle  + 64 \sqrt{3} i r^{3} z^{5} + 72 \sqrt{3} i r z^{4}$\\
$\displaystyle  + 96 \sqrt{3} i r z^{2} + 96 \sqrt{3} i r^{4} z^{6}$\\
$\displaystyle  + 144 \sqrt{3} i r^{3} z^{2} + 152 \sqrt{3} i r^{3} z$\\
$\displaystyle  + 160 \sqrt{3} i r^{5} z + 168 \sqrt{3} i r^{5} z^{4}$\\
$\displaystyle  + 248 \sqrt{3} i r^{2} z^{3} + 280 \sqrt{3} i r^{4} z^{3}$\\
$\displaystyle  + 400 \sqrt{3} i r^{4} z^{4}$}
\label{eq:cat-numerator-110}
\end{equation}
\begin{equation}
D_{110}(r,z)=\left(- r^{2} + 2 r z - 1\right)^{2} \left(r^{2} + 2 r z + 1\right)
\label{eq:cat-denominator-110}
\end{equation}
\begin{equation}
R_{111}(r,z)=\frac{N_{111}(r,z)}{D_{111}(r,z)}.
\label{eq:cat-rational-111}
\end{equation}
\begin{equation}
N_{111}(r,z)=\shortstack[l]{$\displaystyle  - 16 r^{4} - 8 r^{2} - 8 r^{6} - 304 r^{4} z^{4}$\\
$\displaystyle  - 112 r^{3} z - 112 r^{5} z - 96 r^{4} z^{6} - 48 r z$\\
$\displaystyle  - 48 r^{7} z - 32 r^{4} z^{3} - 32 r^{3} z^{4}$\\
$\displaystyle  - 32 r^{5} z^{4} - 24 r z^{3} - 24 r^{7} z^{3}$\\
$\displaystyle  - 16 r^{2} z^{3} - 16 r^{6} z^{3} + 8 r z^{2} + 8 r^{7} z^{2}$\\
$\displaystyle  + 24 r^{3} z^{2} + 24 r^{5} z^{2} + 64 r^{4} z^{5}$\\
$\displaystyle  + 72 r^{2} z^{4} + 72 r^{6} z^{4} + 80 r^{2} z^{2}$\\
$\displaystyle  + 80 r^{6} z^{2} + 128 r^{4} z^{2} + 184 r^{3} z^{3}$\\
$\displaystyle  + 184 r^{5} z^{3} - 272 \sqrt{3} r^{4} z^{4}$\\
$\displaystyle  - 256 \sqrt{3} r^{4} z^{3} - 120 \sqrt{3} r^{3} z - 120 \sqrt{3} r^{5} z$\\
$\displaystyle  - 48 \sqrt{3} r^{3} z^{5} - 48 \sqrt{3} r^{5} z^{5} - 40 \sqrt{3} r z$\\
$\displaystyle  - 40 \sqrt{3} r^{7} z - 32 \sqrt{3} r^{3} z^{2} - 32 \sqrt{3} r^{5} z^{2}$\\
$\displaystyle  - 16 \sqrt{3} r z^{2} - 16 \sqrt{3} r^{7} z^{2} + 24 \sqrt{3} r z^{3}$\\
$\displaystyle  + 24 \sqrt{3} r^{7} z^{3} + 24 \sqrt{3} r^{2} z^{4}$\\
$\displaystyle  + 24 \sqrt{3} r^{6} z^{4} + 32 \sqrt{3} r^{2} z + 32 \sqrt{3} r^{6} z$\\
$\displaystyle  + 32 \sqrt{3} r^{4} z^{5} + 48 \sqrt{3} r^{2} z^{3}$\\
$\displaystyle  + 48 \sqrt{3} r^{6} z^{3} + 48 \sqrt{3} r^{3} z^{4}$\\
$\displaystyle  + 48 \sqrt{3} r^{5} z^{4} + 56 \sqrt{3} r^{2} z^{2}$\\
$\displaystyle  + 56 \sqrt{3} r^{6} z^{2} + 64 \sqrt{3} r^{4} z$\\
$\displaystyle  + 112 \sqrt{3} r^{4} z^{2} + 184 \sqrt{3} r^{3} z^{3}$\\
$\displaystyle  + 184 \sqrt{3} r^{5} z^{3}$}
\label{eq:cat-numerator-111}
\end{equation}
\begin{equation}
D_{111}(r,z)=\left(- r^{2} + 2 r z - 1\right)^{2} \left(r^{2} + 2 r z + 1\right)
\label{eq:cat-denominator-111}
\end{equation}
\begin{equation}
R_{112}(r,z)=\frac{8 i r^{2} z \left(r - 1\right) \left(z - 1\right) \left(z + 1\right)}{r^{2} - 2 r z + 1}
\label{eq:cat-rational-112}
\end{equation}
\begin{equation}
R_{113}(r,z)=- \frac{4 \sqrt{3} z \left(z - 1\right) \left(z + 1\right) \left(3 r^{2} + 6 r z - 1\right)}{r^{2} - 2 r z + 1}
\label{eq:cat-rational-113}
\end{equation}
\begin{equation}
R_{114}(r,z)=\frac{N_{114}(r,z)}{D_{114}(r,z)}.
\label{eq:cat-rational-114}
\end{equation}
\begin{equation}
N_{114}(r,z)=\shortstack[l]{$\displaystyle  - 12 i z + 12 i z^{3} - 144 i r^{3} z^{4} - 60 i r^{4} z^{3}$\\
$\displaystyle  - 48 i r^{2} z^{5} - 24 i r z^{4} - 24 i r^{5} z^{4}$\\
$\displaystyle  - 12 i r^{2} z - 12 i r^{6} z^{3} - 4 \sqrt{3} i z$\\
$\displaystyle  + 4 \sqrt{3} i z^{3} + 12 i r^{4} z + 12 i r^{6} z + 24 i r z^{2}$\\
$\displaystyle  + 24 i r^{5} z^{2} + 32 \sqrt{3} i r + 32 \sqrt{3} i r^{2}$\\
$\displaystyle  + 32 \sqrt{3} i r^{5} + 32 \sqrt{3} i r^{6} + 48 i r^{3} z^{2}$\\
$\displaystyle  + 48 i r^{4} z^{5} + 60 i r^{2} z^{3} + 64 \sqrt{3} i r^{3}$\\
$\displaystyle  + 64 \sqrt{3} i r^{4} + 96 i r^{3} z^{6} - 480 \sqrt{3} i r^{4} z^{2}$\\
$\displaystyle  - 384 \sqrt{3} i r^{3} z^{2} - 384 \sqrt{3} i r^{6} z^{2}$\\
$\displaystyle  - 376 \sqrt{3} i r^{4} z^{5} - 336 \sqrt{3} i r^{5} z^{2}$\\
$\displaystyle  - 328 \sqrt{3} i r^{5} z^{5} - 208 \sqrt{3} i r^{4} z^{6}$\\
$\displaystyle  - 112 \sqrt{3} i r^{2} z^{2} - 88 \sqrt{3} i r^{2} z^{5}$\\
$\displaystyle  - 64 \sqrt{3} i r z^{2} - 56 \sqrt{3} i r^{7} z^{3}$\\
$\displaystyle  - 44 \sqrt{3} i r^{6} z^{3} - 40 \sqrt{3} i r^{3} z^{5}$\\
$\displaystyle  - 20 \sqrt{3} i r^{2} z - 16 \sqrt{3} i r^{3} z^{6}$\\
$\displaystyle  - 8 \sqrt{3} i r z^{3} + 8 \sqrt{3} i r z + 16 \sqrt{3} i r^{3} z$\\
$\displaystyle  + 24 \sqrt{3} i r^{3} z^{3} + 28 \sqrt{3} i r^{4} z$\\
$\displaystyle  + 32 \sqrt{3} i r z^{4} + 44 \sqrt{3} i r^{6} z + 56 \sqrt{3} i r^{7} z$\\
$\displaystyle  + 64 \sqrt{3} i r^{5} z + 80 \sqrt{3} i r^{2} z^{4}$\\
$\displaystyle  + 108 \sqrt{3} i r^{2} z^{3} + 264 \sqrt{3} i r^{5} z^{3}$\\
$\displaystyle  + 304 \sqrt{3} i r^{5} z^{4} + 336 \sqrt{3} i r^{3} z^{4}$\\
$\displaystyle  + 348 \sqrt{3} i r^{4} z^{3} + 352 \sqrt{3} i r^{6} z^{4}$\\
$\displaystyle  + 624 \sqrt{3} i r^{4} z^{4}$}
\label{eq:cat-numerator-114}
\end{equation}
\begin{equation}
D_{114}(r,z)=\left(- r^{2} + 2 r z - 1\right)^{2} \left(r^{2} + 2 r z + 1\right)
\label{eq:cat-denominator-114}
\end{equation}
\begin{equation}
R_{115}(r,z)=\frac{N_{115}(r,z)}{D_{115}(r,z)}.
\label{eq:cat-rational-115}
\end{equation}
\begin{equation}
N_{115}(r,z)=\shortstack[l]{$\displaystyle  - 288 r^{3} - 288 r^{4} - 144 r - 144 r^{2} - 144 r^{5}$\\
$\displaystyle  - 144 r^{6} - 48 z^{3} + 48 z - 1764 r^{4} z^{4}$\\
$\displaystyle  - 1164 r^{4} z^{3} - 1056 r^{6} z^{4} - 912 r^{5} z^{4}$\\
$\displaystyle  - 900 r^{3} z^{4} - 876 r^{5} z^{3} - 408 r^{2} z^{3}$\\
$\displaystyle  - 240 \sqrt{3} r^{5} - 192 \sqrt{3} r^{3} - 168 r^{7} z$\\
$\displaystyle  - 156 r^{2} z^{4} - 132 r^{6} z - 108 r^{5} z - 96 \sqrt{3} r^{4}$\\
$\displaystyle  - 96 \sqrt{3} r^{7} - 96 r^{3} z^{5} - 48 \sqrt{3} r - 48 \sqrt{3} r^{2}$\\
$\displaystyle  - 48 \sqrt{3} r^{6} - 24 r z - 16 \sqrt{3} z^{3} - 12 r z^{4}$\\
$\displaystyle  + 16 \sqrt{3} z + 24 r z^{3} + 36 r^{3} z + 36 r^{4} z$\\
$\displaystyle  + 48 r^{3} z^{6} + 60 r^{3} z^{3} + 132 r^{6} z^{3}$\\
$\displaystyle  + 156 r z^{2} + 168 r^{7} z^{3} + 192 r^{2} z^{5}$\\
$\displaystyle  + 216 r^{2} z + 300 r^{2} z^{2} + 624 r^{4} z^{6}$\\
$\displaystyle  + 984 r^{5} z^{5} + 1056 r^{5} z^{2} + 1128 r^{4} z^{5}$\\
$\displaystyle  + 1140 r^{3} z^{2} + 1200 r^{6} z^{2} + 1428 r^{4} z^{2}$\\
$\displaystyle  - 916 \sqrt{3} r^{4} z^{3} - 288 \sqrt{3} r^{5} z^{5}$\\
$\displaystyle  - 288 \sqrt{3} r^{5} z^{6} - 268 \sqrt{3} r^{2} z^{3}$\\
$\displaystyle  - 256 \sqrt{3} r^{6} z^{3} - 176 \sqrt{3} r^{3} z^{6}$\\
$\displaystyle  - 120 \sqrt{3} r^{3} z^{5} - 108 \sqrt{3} r^{4} z^{4}$\\
$\displaystyle  - 72 \sqrt{3} r^{7} z - 36 \sqrt{3} r^{5} z - 36 \sqrt{3} r^{6} z^{4}$\\
$\displaystyle  - 24 \sqrt{3} r z^{3} + 4 \sqrt{3} r z^{2} + 24 \sqrt{3} r z$\\
$\displaystyle  + 24 \sqrt{3} r^{7} z^{2} + 44 \sqrt{3} r z^{4} + 48 \sqrt{3} r^{2} z^{4}$\\
$\displaystyle  + 60 \sqrt{3} r^{3} z + 60 \sqrt{3} r^{4} z^{2} + 60 \sqrt{3} r^{3} z^{3}$\\
$\displaystyle  + 64 \sqrt{3} r^{2} z^{5} + 72 \sqrt{3} r^{7} z^{3}$\\
$\displaystyle  + 72 \sqrt{3} r^{7} z^{4} + 84 \sqrt{3} r^{6} z^{2}$\\
$\displaystyle  + 136 \sqrt{3} r^{5} z^{2} + 140 \sqrt{3} r^{3} z^{2}$\\
$\displaystyle  + 144 \sqrt{3} r^{4} z^{6} + 204 \sqrt{3} r^{2} z$\\
$\displaystyle  + 228 \sqrt{3} r^{3} z^{4} + 256 \sqrt{3} r^{6} z$\\
$\displaystyle  + 324 \sqrt{3} r^{5} z^{3} + 392 \sqrt{3} r^{5} z^{4}$\\
$\displaystyle  + 444 \sqrt{3} r^{4} z + 472 \sqrt{3} r^{4} z^{5}$}
\label{eq:cat-numerator-115}
\end{equation}
\begin{equation}
D_{115}(r,z)=3 \left(- r^{2} + 2 r z - 1\right)^{2} \left(r^{2} + 2 r z + 1\right)
\label{eq:cat-denominator-115}
\end{equation}
\begin{equation}
R_{116}(r,z)=\frac{N_{116}(r,z)}{D_{116}(r,z)}.
\label{eq:cat-rational-116}
\end{equation}
\begin{equation}
N_{116}(r,z)=\shortstack[l]{$\displaystyle  - 24 i r^{3} z^{3} - 24 i r^{3} z^{4} - 12 i r^{2} z$\\
$\displaystyle  - 12 i r^{2} z^{2} + 12 i r^{2} z^{3} + 12 i r^{2} z^{4}$\\
$\displaystyle  + 24 i r^{3} z + 24 i r^{3} z^{2} - 4 \sqrt{3} i r^{2} z$\\
$\displaystyle  - 4 \sqrt{3} i r^{2} z^{2} + 4 \sqrt{3} i r^{2} z^{3}$\\
$\displaystyle  + 4 \sqrt{3} i r^{2} z^{4}$}
\label{eq:cat-numerator-116}
\end{equation}
\begin{equation}
D_{116}(r,z)=- r^{2} + 2 r z - 1
\label{eq:cat-denominator-116}
\end{equation}
\begin{equation}
R_{117}(r,z)=- 32 \sqrt{3} i z \left(r + 1\right) \left(z - 1\right) \left(z + 1\right)
\label{eq:cat-rational-117}
\end{equation}
\begin{equation}
R_{118}(r,z)=32 z \left(r + 1\right) \left(z - 1\right) \left(z + 1\right) + \frac{16 \sqrt{3} \left(3 r - z\right) \left(z - 1\right) \left(z + 1\right)}{3}
\label{eq:cat-rational-118}
\end{equation}
\begin{equation}
R_{119}(r,z)=\frac{4 \sqrt{3} z \left(z - 1\right) \left(z + 1\right) \left(15 r^{2} - 6 r z + 7\right)}{r^{2} - 2 r z + 1}
\label{eq:cat-rational-119}
\end{equation}
\begin{equation}
R_{120}(r,z)=\frac{N_{120}(r,z)}{D_{120}(r,z)}.
\label{eq:cat-rational-120}
\end{equation}
\begin{equation}
N_{120}(r,z)=\shortstack[l]{$\displaystyle  - 12 i z + 12 i z^{3} - 240 i r^{4} z^{5} - 144 i r^{3} z^{4}$\\
$\displaystyle  - 132 i r^{4} z - 84 i r^{2} z - 60 i r^{6} z$\\
$\displaystyle  - 48 i r^{2} z^{5} - 36 \sqrt{3} i z - 24 i r z^{4}$\\
$\displaystyle  - 24 i r^{5} z^{4} + 24 i r z^{2} + 24 i r^{5} z^{2}$\\
$\displaystyle  + 32 \sqrt{3} i r + 32 \sqrt{3} i r^{2} + 32 \sqrt{3} i r^{5}$\\
$\displaystyle  + 32 \sqrt{3} i r^{6} + 36 \sqrt{3} i z^{3} + 48 i r^{3} z^{2}$\\
$\displaystyle  + 60 i r^{6} z^{3} + 64 \sqrt{3} i r^{3} + 64 \sqrt{3} i r^{4}$\\
$\displaystyle  + 96 i r^{3} z^{6} + 132 i r^{2} z^{3} + 372 i r^{4} z^{3}$\\
$\displaystyle  - 648 \sqrt{3} i r^{4} z^{5} - 392 \sqrt{3} i r^{6} z^{2}$\\
$\displaystyle  - 344 \sqrt{3} i r^{5} z^{2} - 312 \sqrt{3} i r^{5} z^{5}$\\
$\displaystyle  - 304 \sqrt{3} i r^{4} z^{2} - 264 \sqrt{3} i r^{2} z^{5}$\\
$\displaystyle  - 208 \sqrt{3} i r^{3} z^{2} - 164 \sqrt{3} i r^{4} z$\\
$\displaystyle  - 140 \sqrt{3} i r^{2} z - 96 \sqrt{3} i r^{3} z^{4}$\\
$\displaystyle  - 60 \sqrt{3} i r^{6} z - 56 \sqrt{3} i r z^{4}$\\
$\displaystyle  - 40 \sqrt{3} i r^{3} z^{3} - 24 \sqrt{3} i r z^{3}$\\
$\displaystyle  - 24 \sqrt{3} i r^{2} z^{2} - 24 \sqrt{3} i r^{7} z^{3}$\\
$\displaystyle  - 24 \sqrt{3} i r^{3} z^{5} - 8 \sqrt{3} i r^{2} z^{4}$\\
$\displaystyle  + 24 \sqrt{3} i r z + 24 \sqrt{3} i r^{7} z + 24 \sqrt{3} i r z^{2}$\\
$\displaystyle  + 48 \sqrt{3} i r^{4} z^{6} + 60 \sqrt{3} i r^{6} z^{3}$\\
$\displaystyle  + 64 \sqrt{3} i r^{3} z + 64 \sqrt{3} i r^{5} z$\\
$\displaystyle  + 192 \sqrt{3} i r^{4} z^{4} + 240 \sqrt{3} i r^{3} z^{6}$\\
$\displaystyle  + 248 \sqrt{3} i r^{5} z^{3} + 312 \sqrt{3} i r^{5} z^{4}$\\
$\displaystyle  + 360 \sqrt{3} i r^{6} z^{4} + 404 \sqrt{3} i r^{2} z^{3}$\\
$\displaystyle  + 812 \sqrt{3} i r^{4} z^{3}$}
\label{eq:cat-numerator-120}
\end{equation}
\begin{equation}
D_{120}(r,z)=\left(- r^{2} + 2 r z - 1\right)^{2} \left(r^{2} + 2 r z + 1\right)
\label{eq:cat-denominator-120}
\end{equation}
\begin{equation}
R_{121}(r,z)=\frac{N_{121}(r,z)}{D_{121}(r,z)}.
\label{eq:cat-rational-121}
\end{equation}
\begin{equation}
N_{121}(r,z)=\shortstack[l]{$\displaystyle  - 96 r^{3} - 96 r^{4} - 48 r - 48 z^{3} - 48 r^{2} - 48 r^{5}$\\
$\displaystyle  - 48 r^{6} + 48 z - 828 r^{4} z^{3} - 360 r^{2} z^{3}$\\
$\displaystyle  - 360 r^{6} z^{4} - 312 r^{5} z^{4} - 300 r^{5} z^{3}$\\
$\displaystyle  - 240 r^{3} z^{6} - 132 r^{3} z^{3} - 132 r^{4} z^{4}$\\
$\displaystyle  - 60 r^{6} z^{3} - 48 r^{4} z^{6} - 32 \sqrt{3} r^{4}$\\
$\displaystyle  - 32 \sqrt{3} r^{5} - 24 r^{7} z - 24 r z^{3} - 16 \sqrt{3} r^{2}$\\
$\displaystyle  - 16 \sqrt{3} r^{3} - 16 \sqrt{3} r^{6} - 16 \sqrt{3} r^{7} - 12 r^{5} z$\\
$\displaystyle  - 12 r z^{2} + 12 r^{2} z^{4} + 24 r z + 24 r^{7} z^{3}$\\
$\displaystyle  + 36 r^{3} z + 36 r^{2} z^{2} + 60 r^{6} z + 60 r z^{4}$\\
$\displaystyle  + 96 r^{3} z^{5} + 156 r^{3} z^{4} + 168 r^{2} z$\\
$\displaystyle  + 180 r^{4} z + 180 r^{3} z^{2} + 192 r^{2} z^{5}$\\
$\displaystyle  + 276 r^{4} z^{2} + 312 r^{5} z^{5} + 360 r^{5} z^{2}$\\
$\displaystyle  + 408 r^{6} z^{2} + 648 r^{4} z^{5} - 148 \sqrt{3} r^{4} z^{3}$\\
$\displaystyle  - 96 \sqrt{3} r^{5} z^{5} - 96 \sqrt{3} r^{5} z^{6}$\\
$\displaystyle  - 88 \sqrt{3} r^{3} z^{5} - 80 \sqrt{3} r^{5} z^{2}$\\
$\displaystyle  - 48 \sqrt{3} r^{6} z^{3} - 44 \sqrt{3} r^{2} z^{3} - 36 \sqrt{3} r^{5} z$\\
$\displaystyle  - 24 \sqrt{3} r^{7} z - 20 \sqrt{3} r^{3} z^{2} - 16 \sqrt{3} r^{3} z^{6}$\\
$\displaystyle  - 12 \sqrt{3} r^{4} z^{4} - 12 \sqrt{3} r^{6} z^{4} - 8 \sqrt{3} r z^{3}$\\
$\displaystyle  - 8 \sqrt{3} r^{7} z^{2} - 8 \sqrt{3} r^{2} z^{4} - 4 \sqrt{3} r^{3} z$\\
$\displaystyle  - 4 \sqrt{3} r z^{2} - 4 \sqrt{3} r^{4} z^{2} + 4 \sqrt{3} r z^{4}$\\
$\displaystyle  + 8 \sqrt{3} r z + 24 \sqrt{3} r^{2} z^{2} + 24 \sqrt{3} r^{7} z^{3}$\\
$\displaystyle  + 24 \sqrt{3} r^{7} z^{4} + 28 \sqrt{3} r^{6} z^{2} + 44 \sqrt{3} r^{2} z$\\
$\displaystyle  + 48 \sqrt{3} r^{6} z + 48 \sqrt{3} r^{4} z^{6} + 52 \sqrt{3} r^{3} z^{4}$\\
$\displaystyle  + 56 \sqrt{3} r^{4} z^{5} + 92 \sqrt{3} r^{4} z + 92 \sqrt{3} r^{3} z^{3}$\\
$\displaystyle  + 132 \sqrt{3} r^{5} z^{3} + 208 \sqrt{3} r^{5} z^{4}$}
\label{eq:cat-numerator-121}
\end{equation}
\begin{equation}
D_{121}(r,z)=\left(- r^{2} + 2 r z - 1\right)^{2} \left(r^{2} + 2 r z + 1\right)
\label{eq:cat-denominator-121}
\end{equation}
\begin{equation}
R_{122}(r,z)=- \frac{24 \sqrt{3} z \left(3 r^{2} + 1\right) \left(z - 1\right) \left(z + 1\right)}{r^{2} - 2 r z + 1}
\label{eq:cat-rational-122}
\end{equation}
\begin{equation}
R_{123}(r,z)=\frac{N_{123}(r,z)}{D_{123}(r,z)}.
\label{eq:cat-rational-123}
\end{equation}
\begin{equation}
N_{123}(r,z)=\shortstack[l]{$\displaystyle  - 144 i r^{3} z^{2} - 72 i r^{2} z^{3} - 72 i r^{4} z^{3}$\\
$\displaystyle  + 72 i r^{2} z + 72 i r^{4} z + 144 i r^{3} z^{4}$\\
$\displaystyle  - 96 \sqrt{3} i r^{3} z - 72 \sqrt{3} i r^{3} z^{2}$\\
$\displaystyle  - 72 \sqrt{3} i r^{4} z^{3} - 72 \sqrt{3} i r^{2} z^{4}$\\
$\displaystyle  - 72 \sqrt{3} i r^{4} z^{4} - 48 \sqrt{3} i r z - 24 \sqrt{3} i r z^{2}$\\
$\displaystyle  - 24 \sqrt{3} i r^{2} z^{3} + 24 \sqrt{3} i r^{2} z$\\
$\displaystyle  + 24 \sqrt{3} i r z^{4} + 48 \sqrt{3} i r z^{3} + 72 \sqrt{3} i r^{4} z$\\
$\displaystyle  + 72 \sqrt{3} i r^{2} z^{2} + 72 \sqrt{3} i r^{4} z^{2}$\\
$\displaystyle  + 72 \sqrt{3} i r^{3} z^{4} + 96 \sqrt{3} i r^{3} z^{3}$}
\label{eq:cat-numerator-123}
\end{equation}
\begin{equation}
D_{123}(r,z)=\left(- r^{2} + 2 r z - 1\right)^{2}
\label{eq:cat-denominator-123}
\end{equation}
\begin{equation}
R_{124}(r,z)=\frac{N_{124}(r,z)}{D_{124}(r,z)}.
\label{eq:cat-rational-124}
\end{equation}
\begin{equation}
N_{124}(r,z)=\shortstack[l]{$\displaystyle  - 72 r^{4} z - 72 r^{4} z^{2} - 72 r^{3} z^{3}$\\
$\displaystyle  - 72 r^{3} z^{4} + 72 r^{3} z + 72 r^{3} z^{2}$\\
$\displaystyle  + 72 r^{4} z^{3} + 72 r^{4} z^{4} - 24 \sqrt{3} r^{2} z$\\
$\displaystyle  - 24 \sqrt{3} r^{2} z^{2} - 24 \sqrt{3} r^{3} z^{3}$\\
$\displaystyle  - 24 \sqrt{3} r^{3} z^{4} + 24 \sqrt{3} r^{3} z + 24 \sqrt{3} r^{3} z^{2}$\\
$\displaystyle  + 24 \sqrt{3} r^{2} z^{3} + 24 \sqrt{3} r^{2} z^{4}$}
\label{eq:cat-numerator-124}
\end{equation}
\begin{equation}
D_{124}(r,z)=\left(- r^{2} + 2 r z - 1\right)^{2}
\label{eq:cat-denominator-124}
\end{equation}
\begin{equation}
R_{125}(r,z)=12 z^{4} + 24 z^{2} + 4
\label{eq:cat-rational-125}
\end{equation}
\begin{equation}
R_{126}(r,z)=\frac{N_{126}(r,z)}{D_{126}(r,z)}.
\label{eq:cat-rational-126}
\end{equation}
\begin{equation}
N_{126}(r,z)=\shortstack[l]{$\displaystyle 40 i - 40 i r^{5} - 32 i r^{3} - 8 i r^{4} + 8 i r$\\
$\displaystyle  + 24 i z^{4} + 32 i r^{2} + 48 i z^{2} - 64 i r^{2} z^{2}$\\
$\displaystyle  - 48 i r^{4} z^{2} - 48 i r^{5} z^{2} - 24 i r^{4} z^{4}$\\
$\displaystyle  - 24 i r^{5} z^{4} + 24 i r z^{4} + 48 i r z^{2}$\\
$\displaystyle  + 64 i r^{3} z^{2}$}
\label{eq:cat-numerator-126}
\end{equation}
\begin{equation}
D_{126}(r,z)=\left(- r^{2} + 2 r z - 1\right) \left(r^{2} + 2 r z + 1\right)
\label{eq:cat-denominator-126}
\end{equation}
\begin{equation}
R_{127}(r,z)=\frac{N_{127}(r,z)}{D_{127}(r,z)}.
\label{eq:cat-rational-127}
\end{equation}
\begin{equation}
N_{127}(r,z)=\shortstack[l]{$\displaystyle 96 - 528 r^{5} - 344 r^{3} - 344 r^{7} - 160 z^{2} - 80 r$\\
$\displaystyle  - 80 r^{9} + 64 z^{4} + 96 r^{10} + 444 r^{2} + 444 r^{8}$\\
$\displaystyle  + 852 r^{4} + 852 r^{6} - 4756 r^{4} z^{2} - 4756 r^{6} z^{2}$\\
$\displaystyle  - 3100 r^{4} z^{6} - 3100 r^{6} z^{6} - 1868 r^{2} z^{2}$\\
$\displaystyle  - 1868 r^{8} z^{2} - 1040 r^{5} z^{6} - 880 r^{5} z^{2}$\\
$\displaystyle  - 608 r^{5} z^{8} - 552 r^{3} z^{2} - 552 r^{7} z^{2}$\\
$\displaystyle  - 356 r^{2} z^{6} - 356 r^{8} z^{6} - 160 r z^{2}$\\
$\displaystyle  - 160 r^{9} z^{2} - 160 r^{10} z^{2} + 16 r z^{4}$\\
$\displaystyle  + 16 r^{9} z^{4} + 64 r^{10} z^{4} + 136 r^{3} z^{6}$\\
$\displaystyle  + 136 r^{7} z^{6} + 496 r^{4} z^{8} + 496 r^{6} z^{8}$\\
$\displaystyle  + 1080 r^{3} z^{4} + 1080 r^{7} z^{4} + 1492 r^{2} z^{4}$\\
$\displaystyle  + 1492 r^{8} z^{4} + 2864 r^{5} z^{4} + 6796 r^{4} z^{4}$\\
$\displaystyle  + 6796 r^{6} z^{4}$}
\label{eq:cat-numerator-127}
\end{equation}
\begin{equation}
D_{127}(r,z)=\left(- r^{2} + 2 r z - 1\right)^{2} \left(r^{2} + 2 r z + 1\right)^{2}
\label{eq:cat-denominator-127}
\end{equation}
\begin{equation}
R_{128}(r,z)=\frac{N_{128}(r,z)}{D_{128}(r,z)}.
\label{eq:cat-rational-128}
\end{equation}
\begin{equation}
N_{128}(r,z)=\shortstack[l]{$\displaystyle  - 24 i r^{2} - 24 i r^{4} + 24 i r^{3} + 24 i r^{5}$\\
$\displaystyle  - 120 i r^{3} z^{4} - 72 i r^{2} z^{2} - 72 i r^{2} z^{4}$\\
$\displaystyle  - 72 i r^{3} z^{6} - 24 i r^{3} z^{2} - 24 i r^{2} z^{6}$\\
$\displaystyle  + 24 i r^{4} z^{2} + 24 i r^{5} z^{6} + 72 i r^{5} z^{2}$\\
$\displaystyle  + 72 i r^{5} z^{4} + 72 i r^{4} z^{6} + 120 i r^{4} z^{4}$}
\label{eq:cat-numerator-128}
\end{equation}
\begin{equation}
D_{128}(r,z)=\left(- r^{2} + 2 r z - 1\right) \left(r^{2} + 2 r z + 1\right)
\label{eq:cat-denominator-128}
\end{equation}
\begin{equation}
R_{129}(r,z)=16 \left(z - 1\right) \left(z + 1\right) \left(r^{2} z^{2} - 6 r^{2} - 2 r z^{2} + z^{2} - 6\right)
\label{eq:cat-rational-129}
\end{equation}
\begin{equation}
R_{130}(r,z)=- 60 z^{4} - 120 z^{2} - 52
\label{eq:cat-rational-130}
\end{equation}
\begin{equation}
R_{131}(r,z)=\frac{N_{131}(r,z)}{D_{131}(r,z)}.
\label{eq:cat-rational-131}
\end{equation}
\begin{equation}
N_{131}(r,z)=\shortstack[l]{$\displaystyle 8 i - 112 i r^{3} - 104 i r - 8 i r^{5} + 24 i z^{4}$\\
$\displaystyle  + 48 i z^{2} + 104 i r^{4} + 112 i r^{2} - 432 i r^{2} z^{4}$\\
$\displaystyle  - 288 i r^{2} z^{6} - 240 i r z^{2} - 120 i r z^{4}$\\
$\displaystyle  - 64 i r^{3} z^{2} - 48 i r^{5} z^{2} - 24 i r^{5} z^{4}$\\
$\displaystyle  + 64 i r^{2} z^{2} + 120 i r^{4} z^{4} + 240 i r^{4} z^{2}$\\
$\displaystyle  + 288 i r^{3} z^{6} + 432 i r^{3} z^{4}$}
\label{eq:cat-numerator-131}
\end{equation}
\begin{equation}
D_{131}(r,z)=\left(- r^{2} + 2 r z - 1\right) \left(r^{2} + 2 r z + 1\right)
\label{eq:cat-denominator-131}
\end{equation}
\begin{equation}
R_{132}(r,z)=\frac{N_{132}(r,z)}{D_{132}(r,z)}.
\label{eq:cat-rational-132}
\end{equation}
\begin{equation}
N_{132}(r,z)=\shortstack[l]{$\displaystyle  - 432 r^{5} - 232 r^{3} - 232 r^{7} - 48 z^{2} - 16 r$\\
$\displaystyle  - 16 r^{9} + 36 r^{2} + 36 r^{8} + 48 z^{4} + 108 r^{4}$\\
$\displaystyle  + 108 r^{6} - 2928 r^{5} z^{6} - 2384 r^{5} z^{2}$\\
$\displaystyle  - 1548 r^{4} z^{2} - 1548 r^{6} z^{2} - 1048 r^{3} z^{2}$\\
$\displaystyle  - 1048 r^{7} z^{2} - 996 r^{4} z^{6} - 996 r^{6} z^{6}$\\
$\displaystyle  - 672 r^{5} z^{8} - 324 r^{2} z^{2} - 324 r^{8} z^{2}$\\
$\displaystyle  - 192 r z^{2} - 192 r^{9} z^{2} - 156 r^{2} z^{6}$\\
$\displaystyle  - 156 r^{8} z^{6} - 48 r^{10} z^{2} + 48 r z^{4}$\\
$\displaystyle  + 48 r^{9} z^{4} + 48 r^{10} z^{4} + 312 r^{3} z^{6}$\\
$\displaystyle  + 312 r^{7} z^{6} + 528 r^{4} z^{8} + 528 r^{6} z^{8}$\\
$\displaystyle  + 732 r^{2} z^{4} + 732 r^{8} z^{4} + 1620 r^{4} z^{4}$\\
$\displaystyle  + 1620 r^{6} z^{4} + 2184 r^{3} z^{4} + 2184 r^{7} z^{4}$\\
$\displaystyle  + 4304 r^{5} z^{4}$}
\label{eq:cat-numerator-132}
\end{equation}
\begin{equation}
D_{132}(r,z)=\left(- r^{2} + 2 r z - 1\right)^{2} \left(r^{2} + 2 r z + 1\right)^{2}
\label{eq:cat-denominator-132}
\end{equation}
\begin{equation}
R_{133}(r,z)=72 z^{4} + 144 z^{2} + 56
\label{eq:cat-rational-133}
\end{equation}
\begin{equation}
R_{134}(r,z)=\frac{N_{134}(r,z)}{D_{134}(r,z)}.
\label{eq:cat-rational-134}
\end{equation}
\begin{equation}
N_{134}(r,z)=\shortstack[l]{$\displaystyle 32 i - 112 i r^{4} - 80 i r^{2} - 32 i r^{5} + 80 i r^{3}$\\
$\displaystyle  + 112 i r - 432 i r^{3} z^{4} - 288 i r^{4} z^{2}$\\
$\displaystyle  - 288 i r^{3} z^{6} - 144 i r^{4} z^{4} - 128 i r^{2} z^{2}$\\
$\displaystyle  + 128 i r^{3} z^{2} + 144 i r z^{4} + 288 i r z^{2}$\\
$\displaystyle  + 288 i r^{2} z^{6} + 432 i r^{2} z^{4}$}
\label{eq:cat-numerator-134}
\end{equation}
\begin{equation}
D_{134}(r,z)=\left(- r^{2} + 2 r z - 1\right) \left(r^{2} + 2 r z + 1\right)
\label{eq:cat-denominator-134}
\end{equation}
\begin{equation}
R_{135}(r,z)=\frac{N_{135}(r,z)}{D_{135}(r,z)}.
\label{eq:cat-rational-135}
\end{equation}
\begin{equation}
N_{135}(r,z)=\shortstack[l]{$\displaystyle  - 216 r^{4} - 216 r^{6} - 112 r^{3} - 112 r^{7} - 96 r^{5}$\\
$\displaystyle  - 72 r^{2} - 72 r^{8} - 64 r - 64 r^{9} - 736 r^{5} z^{4}$\\
$\displaystyle  - 720 r^{3} z^{4} - 720 r^{7} z^{4} - 432 r^{3} z^{6}$\\
$\displaystyle  - 432 r^{7} z^{6} - 288 r^{4} z^{8} - 288 r^{6} z^{8}$\\
$\displaystyle  - 216 r^{2} z^{2} - 216 r^{8} z^{2} - 216 r^{2} z^{4}$\\
$\displaystyle  - 216 r^{8} z^{4} - 72 r^{2} z^{6} - 72 r^{8} z^{6}$\\
$\displaystyle  + 72 r^{4} z^{6} + 72 r^{6} z^{6} + 216 r^{4} z^{2}$\\
$\displaystyle  + 216 r^{6} z^{2} + 368 r^{3} z^{2} + 368 r^{7} z^{2}$\\
$\displaystyle  + 576 r^{5} z^{8} + 792 r^{4} z^{4} + 792 r^{6} z^{4}$\\
$\displaystyle  + 864 r^{5} z^{6} + 1312 r^{5} z^{2}$}
\label{eq:cat-numerator-135}
\end{equation}
\begin{equation}
D_{135}(r,z)=\left(- r^{2} + 2 r z - 1\right)^{2} \left(r^{2} + 2 r z + 1\right)^{2}
\label{eq:cat-denominator-135}
\end{equation}
\begin{equation}
R_{136}(r,z)=8
\label{eq:cat-rational-136}
\end{equation}
\begin{equation}
R_{137}(r,z)=- 16 i \left(r - 1\right)
\label{eq:cat-rational-137}
\end{equation}
\begin{equation}
R_{138}(r,z)=\frac{N_{138}(r,z)}{D_{138}(r,z)}.
\label{eq:cat-rational-138}
\end{equation}
\begin{equation}
N_{138}(r,z)=\shortstack[l]{$\displaystyle 32 - 96 r - 96 r^{3} + 96 z^{2} + 128 r^{6} + 192 r^{2}$\\
$\displaystyle  + 288 r^{4} - 800 r^{4} z^{2} - 416 r^{2} z^{2}$\\
$\displaystyle  - 192 r^{3} z^{2} - 192 r^{4} z^{6} - 96 r^{6} z^{2}$\\
$\displaystyle  + 96 r^{2} z^{4} + 96 r^{3} z^{4} + 96 r^{5} z^{4}$\\
$\displaystyle  + 96 r^{6} z^{4} + 192 r z^{2} + 576 r^{4} z^{4}$}
\label{eq:cat-numerator-138}
\end{equation}
\begin{equation}
D_{138}(r,z)=\left(- r^{2} + 2 r z - 1\right) \left(r^{2} + 2 r z + 1\right)
\label{eq:cat-denominator-138}
\end{equation}
\begin{equation}
R_{139}(r,z)=- 48 r^{2} z^{4} + 48 r^{2} z^{2} - 128 r^{2} - 96 r z^{2} - 32 r - 48 z^{2} - 80
\label{eq:cat-rational-139}
\end{equation}
\begin{equation}
R_{140}(r,z)=-8
\label{eq:cat-rational-140}
\end{equation}
\begin{equation}
R_{141}(r,z)=16 i \left(r - 1\right)
\label{eq:cat-rational-141}
\end{equation}
\begin{equation}
R_{142}(r,z)=\frac{N_{142}(r,z)}{D_{142}(r,z)}.
\label{eq:cat-rational-142}
\end{equation}
\begin{equation}
N_{142}(r,z)=\shortstack[l]{$\displaystyle  - 48 - 96 r^{2} - 64 r - 48 r^{4} - 32 r^{3} + 32 r^{5}$\\
$\displaystyle  + 48 z^{2} - 512 r^{3} z^{2} - 240 r^{4} z^{2}$\\
$\displaystyle  - 144 r^{2} z^{2} - 96 r^{5} z^{2} - 48 r^{6} z^{2}$\\
$\displaystyle  + 48 r^{6} z^{4} + 96 r z^{2} + 96 r^{5} z^{4}$\\
$\displaystyle  + 240 r^{2} z^{4} + 288 r^{4} z^{4} + 480 r^{3} z^{4}$}
\label{eq:cat-numerator-142}
\end{equation}
\begin{equation}
D_{142}(r,z)=\left(- r^{2} + 2 r z - 1\right) \left(r^{2} + 2 r z + 1\right)
\label{eq:cat-denominator-142}
\end{equation}
\begin{equation}
R_{143}(r,z)=16
\label{eq:cat-rational-143}
\end{equation}
\begin{equation}
R_{144}(r,z)=- 32 i \left(r - 1\right)
\label{eq:cat-rational-144}
\end{equation}
\begin{equation}
R_{145}(r,z)=64 r
\label{eq:cat-rational-145}
\end{equation}
\begin{equation}
R_{146}(r,z)=- 12 z \left(z^{2} + 1\right)
\label{eq:cat-rational-146}
\end{equation}
\begin{equation}
R_{147}(r,z)=- \frac{24 i z \left(r - 1\right) \left(r + 1\right)^{2} \left(z^{2} + 1\right)}{r^{2} - 2 r z + 1}
\label{eq:cat-rational-147}
\end{equation}
\begin{equation}
R_{148}(r,z)=\frac{N_{148}(r,z)}{D_{148}(r,z)}.
\label{eq:cat-rational-148}
\end{equation}
\begin{equation}
N_{148}(r,z)=\shortstack[l]{$\displaystyle  - 288 r^{3} - 288 r^{5} - 96 r - 96 r^{7} + 32 z + 32 z^{3}$\\
$\displaystyle  + 96 r^{2} + 96 r^{6} + 192 r^{4} - 1744 r^{4} z^{3}$\\
$\displaystyle  - 560 r^{3} z^{4} - 560 r^{5} z^{4} - 464 r^{3} z^{3}$\\
$\displaystyle  - 464 r^{5} z^{3} - 400 r^{4} z^{4} - 376 r^{2} z^{4}$\\
$\displaystyle  - 376 r^{6} z^{4} - 216 r^{2} z^{3} - 216 r^{6} z^{3}$\\
$\displaystyle  - 208 r z^{4} - 208 r^{7} z^{4} - 192 r^{4} z^{7}$\\
$\displaystyle  + 32 r^{8} z + 32 r^{8} z^{3} + 72 r^{3} z + 72 r^{5} z$\\
$\displaystyle  + 80 r z + 80 r^{7} z + 80 r z^{2} + 80 r z^{3}$\\
$\displaystyle  + 80 r^{7} z^{2} + 80 r^{7} z^{3} + 104 r^{2} z^{2}$\\
$\displaystyle  + 104 r^{6} z^{2} + 196 r^{2} z^{5} + 196 r^{6} z^{5}$\\
$\displaystyle  + 232 r^{3} z^{5} + 232 r^{5} z^{5} + 232 r^{3} z^{6}$\\
$\displaystyle  + 232 r^{5} z^{6} + 256 r^{4} z^{2} + 304 r^{4} z^{6}$\\
$\displaystyle  + 356 r^{2} z + 356 r^{6} z + 488 r^{4} z^{5} + 648 r^{4} z$\\
$\displaystyle  + 840 r^{3} z^{2} + 840 r^{5} z^{2}$}
\label{eq:cat-numerator-148}
\end{equation}
\begin{equation}
D_{148}(r,z)=\left(- r^{2} + 2 r z - 1\right)^{2} \left(r^{2} + 2 r z + 1\right)
\label{eq:cat-denominator-148}
\end{equation}
\begin{equation}
R_{149}(r,z)=\frac{24 i r^{2} z \left(r - 1\right) \left(z + 1\right)^{2} \left(z^{2} + 1\right)}{r^{2} - 2 r z + 1}
\label{eq:cat-rational-149}
\end{equation}
\begin{equation}
R_{150}(r,z)=32 r^{2} z^{3} + 32 r^{2} z - 48 r z^{4} + 32 r z^{3} + 32 r z - 144 r + 32 z^{3} + 32 z
\label{eq:cat-rational-150}
\end{equation}
\begin{equation}
R_{151}(r,z)=60 z \left(z^{2} + 1\right)
\label{eq:cat-rational-151}
\end{equation}
\begin{equation}
R_{152}(r,z)=- \frac{24 i z \left(r - 1\right) \left(z^{2} + 1\right) \left(r^{2} - 6 r z - 4 r + 1\right)}{r^{2} - 2 r z + 1}
\label{eq:cat-rational-152}
\end{equation}
\begin{equation}
R_{153}(r,z)=\frac{N_{153}(r,z)}{D_{153}(r,z)}.
\label{eq:cat-rational-153}
\end{equation}
\begin{equation}
N_{153}(r,z)=\shortstack[l]{$\displaystyle 48 r + 48 r^{7} + 96 r^{2} + 96 r^{6} + 144 r^{3} + 144 r^{5}$\\
$\displaystyle  + 192 r^{4} - 480 r^{3} z^{4} - 480 r^{5} z^{4}$\\
$\displaystyle  - 456 r^{2} z^{4} - 456 r^{6} z^{4} - 432 r^{3} z^{3}$\\
$\displaystyle  - 432 r^{5} z^{3} - 360 r^{2} z^{3} - 360 r^{6} z^{3}$\\
$\displaystyle  - 360 r^{3} z^{6} - 360 r^{5} z^{6} - 264 r^{4} z$\\
$\displaystyle  - 240 r^{4} z^{3} - 240 r^{4} z^{4} - 132 r^{2} z$\\
$\displaystyle  - 132 r^{6} z - 96 r z^{4} - 96 r^{7} z^{4} + 24 r^{2} z^{2}$\\
$\displaystyle  + 24 r^{6} z^{2} + 48 r z + 48 r^{7} z + 48 r z^{3}$\\
$\displaystyle  + 48 r^{7} z^{3} + 120 r^{3} z + 120 r^{5} z + 144 r z^{2}$\\
$\displaystyle  + 144 r^{7} z^{2} + 156 r^{2} z^{5} + 156 r^{6} z^{5}$\\
$\displaystyle  + 192 r^{4} z^{7} + 216 r^{3} z^{5} + 216 r^{5} z^{5}$\\
$\displaystyle  + 336 r^{4} z^{6} + 384 r^{4} z^{2} + 600 r^{3} z^{2}$\\
$\displaystyle  + 600 r^{5} z^{2} + 984 r^{4} z^{5}$}
\label{eq:cat-numerator-153}
\end{equation}
\begin{equation}
D_{153}(r,z)=\left(- r^{2} + 2 r z - 1\right)^{2} \left(r^{2} + 2 r z + 1\right)
\label{eq:cat-denominator-153}
\end{equation}
\begin{equation}
R_{154}(r,z)=- 72 z \left(z^{2} + 1\right)
\label{eq:cat-rational-154}
\end{equation}
\begin{equation}
R_{155}(r,z)=- \frac{144 i r z \left(r - 1\right) \left(z + 1\right) \left(z^{2} + 1\right)}{r^{2} - 2 r z + 1}
\label{eq:cat-rational-155}
\end{equation}
\begin{equation}
R_{156}(r,z)=\frac{72 r^{2} z \left(r - 1\right)^{2} \left(z + 1\right)^{2} \left(z^{2} + 1\right)}{\left(r^{2} - 2 r z + 1\right)^{2}}
\label{eq:cat-rational-156}
\end{equation}
\begin{equation}
R_{157}(r,z)=48 z^{2} - 8
\label{eq:cat-rational-157}
\end{equation}
\begin{equation}
R_{158}(r,z)=\frac{N_{158}(r,z)}{D_{158}(r,z)}.
\label{eq:cat-rational-158}
\end{equation}
\begin{equation}
N_{158}(r,z)=\shortstack[l]{$\displaystyle 16 i - 32 i r^{3} - 16 i r - 16 i r^{5} + 16 i r^{4}$\\
$\displaystyle  + 32 i r^{2} + 96 i z^{2} - 96 i r^{4} z^{2}$\\
$\displaystyle  - 96 i r^{5} z^{2} - 64 i r^{2} z^{2} + 64 i r^{3} z^{2}$\\
$\displaystyle  + 96 i r z^{2}$}
\label{eq:cat-numerator-158}
\end{equation}
\begin{equation}
D_{158}(r,z)=\left(- r^{2} + 2 r z - 1\right) \left(r^{2} + 2 r z + 1\right)
\label{eq:cat-denominator-158}
\end{equation}
\begin{equation}
R_{159}(r,z)=\frac{N_{159}(r,z)}{D_{159}(r,z)}.
\label{eq:cat-rational-159}
\end{equation}
\begin{equation}
N_{159}(r,z)=\shortstack[l]{$\displaystyle 32 - 576 r^{5} - 384 r^{3} - 384 r^{7} - 96 r - 96 r^{9}$\\
$\displaystyle  - 32 z^{2} + 32 r^{10} + 160 r^{2} + 160 r^{8} + 320 r^{4}$\\
$\displaystyle  + 320 r^{6} - 2672 r^{4} z^{2} - 2672 r^{6} z^{2}$\\
$\displaystyle  - 2432 r^{5} z^{6} - 1088 r^{4} z^{6} - 1088 r^{6} z^{6}$\\
$\displaystyle  - 960 r^{5} z^{2} - 944 r^{2} z^{2} - 944 r^{8} z^{2}$\\
$\displaystyle  - 608 r^{3} z^{2} - 608 r^{7} z^{2} - 128 r z^{2}$\\
$\displaystyle  - 128 r^{9} z^{2} - 32 r^{10} z^{2} + 496 r^{2} z^{4}$\\
$\displaystyle  + 496 r^{8} z^{4} + 1312 r^{3} z^{4} + 1312 r^{7} z^{4}$\\
$\displaystyle  + 3728 r^{4} z^{4} + 3728 r^{6} z^{4} + 3776 r^{5} z^{4}$}
\label{eq:cat-numerator-159}
\end{equation}
\begin{equation}
D_{159}(r,z)=\left(- r^{2} + 2 r z - 1\right)^{2} \left(r^{2} + 2 r z + 1\right)^{2}
\label{eq:cat-denominator-159}
\end{equation}
\begin{equation}
R_{160}(r,z)=- \frac{96 i r^{2} z^{2} \left(r - 1\right) \left(r^{2} z^{2} + r^{2} + 4 r z^{2} + z^{2} + 1\right)}{\left(r^{2} - 2 r z + 1\right) \left(r^{2} + 2 r z + 1\right)}
\label{eq:cat-rational-160}
\end{equation}
\begin{equation}
R_{161}(r,z)=- 16 \left(z - 1\right) \left(z + 1\right) \left(5 r^{2} + 2 r + 5\right)
\label{eq:cat-rational-161}
\end{equation}
\begin{equation}
R_{162}(r,z)=- 240 z^{2} + 8
\label{eq:cat-rational-162}
\end{equation}
\begin{equation}
R_{163}(r,z)=\frac{N_{163}(r,z)}{D_{163}(r,z)}.
\label{eq:cat-rational-163}
\end{equation}
\begin{equation}
N_{163}(r,z)=\shortstack[l]{$\displaystyle  - 16 i - 32 i r^{2} - 16 i r^{4} + 16 i r + 16 i r^{5}$\\
$\displaystyle  + 32 i r^{3} + 96 i z^{2} - 1152 i r^{2} z^{4}$\\
$\displaystyle  - 640 i r^{3} z^{2} - 480 i r z^{2} - 96 i r^{5} z^{2}$\\
$\displaystyle  + 480 i r^{4} z^{2} + 640 i r^{2} z^{2} + 1152 i r^{3} z^{4}$}
\label{eq:cat-numerator-163}
\end{equation}
\begin{equation}
D_{163}(r,z)=\left(- r^{2} + 2 r z - 1\right) \left(r^{2} + 2 r z + 1\right)
\label{eq:cat-denominator-163}
\end{equation}
\begin{equation}
R_{164}(r,z)=\frac{N_{164}(r,z)}{D_{164}(r,z)}.
\label{eq:cat-rational-164}
\end{equation}
\begin{equation}
N_{164}(r,z)=\shortstack[l]{$\displaystyle  - 48 - 480 r^{4} - 480 r^{6} - 384 r^{5} - 256 r^{3}$\\
$\displaystyle  - 256 r^{7} - 240 r^{2} - 240 r^{8} - 64 r - 64 r^{9}$\\
$\displaystyle  - 48 r^{10} + 48 z^{2} - 4224 r^{5} z^{6} - 2432 r^{5} z^{2}$\\
$\displaystyle  - 2064 r^{4} z^{4} - 2064 r^{6} z^{4} - 1312 r^{3} z^{2}$\\
$\displaystyle  - 1312 r^{7} z^{2} - 96 r z^{2} - 96 r^{9} z^{2}$\\
$\displaystyle  + 48 r^{10} z^{2} + 144 r^{2} z^{4} + 144 r^{8} z^{4}$\\
$\displaystyle  + 384 r^{2} z^{2} + 384 r^{8} z^{2} + 912 r^{4} z^{2}$\\
$\displaystyle  + 912 r^{6} z^{2} + 1344 r^{4} z^{6} + 1344 r^{6} z^{6}$\\
$\displaystyle  + 2784 r^{3} z^{4} + 2784 r^{7} z^{4} + 4928 r^{5} z^{4}$}
\label{eq:cat-numerator-164}
\end{equation}
\begin{equation}
D_{164}(r,z)=\left(- r^{2} + 2 r z - 1\right)^{2} \left(r^{2} + 2 r z + 1\right)^{2}
\label{eq:cat-denominator-164}
\end{equation}
\begin{equation}
R_{165}(r,z)=288 z^{2} - 16
\label{eq:cat-rational-165}
\end{equation}
\begin{equation}
R_{166}(r,z)=\frac{N_{166}(r,z)}{D_{166}(r,z)}.
\label{eq:cat-rational-166}
\end{equation}
\begin{equation}
N_{166}(r,z)=\shortstack[l]{$\displaystyle 32 i - 64 i r^{3} - 32 i r - 32 i r^{5} + 32 i r^{4}$\\
$\displaystyle  + 64 i r^{2} - 1152 i r^{3} z^{4} - 704 i r^{2} z^{2}$\\
$\displaystyle  - 576 i r^{4} z^{2} + 576 i r z^{2} + 704 i r^{3} z^{2}$\\
$\displaystyle  + 1152 i r^{2} z^{4}$}
\label{eq:cat-numerator-166}
\end{equation}
\begin{equation}
D_{166}(r,z)=\left(- r^{2} + 2 r z - 1\right) \left(r^{2} + 2 r z + 1\right)
\label{eq:cat-denominator-166}
\end{equation}
\begin{equation}
R_{167}(r,z)=\frac{N_{167}(r,z)}{D_{167}(r,z)}.
\label{eq:cat-rational-167}
\end{equation}
\begin{equation}
N_{167}(r,z)=\shortstack[l]{$\displaystyle  - 384 r^{5} - 256 r^{3} - 256 r^{7} - 64 r - 64 r^{9}$\\
$\displaystyle  - 2176 r^{5} z^{4} - 1728 r^{3} z^{4} - 1728 r^{7} z^{4}$\\
$\displaystyle  - 1152 r^{4} z^{6} - 1152 r^{6} z^{6} - 864 r^{4} z^{2}$\\
$\displaystyle  - 864 r^{6} z^{2} - 288 r^{2} z^{2} - 288 r^{8} z^{2}$\\
$\displaystyle  - 288 r^{2} z^{4} - 288 r^{8} z^{4} + 1088 r^{3} z^{2}$\\
$\displaystyle  + 1088 r^{7} z^{2} + 2176 r^{5} z^{2} + 2304 r^{5} z^{6}$\\
$\displaystyle  + 2592 r^{4} z^{4} + 2592 r^{6} z^{4}$}
\label{eq:cat-numerator-167}
\end{equation}
\begin{equation}
D_{167}(r,z)=\left(- r^{2} + 2 r z - 1\right)^{2} \left(r^{2} + 2 r z + 1\right)^{2}
\label{eq:cat-denominator-167}
\end{equation}
\begin{equation}
R_{168}(r,z)=8 \sqrt{3}
\label{eq:cat-rational-168}
\end{equation}
\begin{equation}
R_{169}(r,z)=8 i
\label{eq:cat-rational-169}
\end{equation}
\begin{equation}
R_{170}(r,z)=- \frac{8 \sqrt{3}}{3}
\label{eq:cat-rational-170}
\end{equation}
\begin{equation}
R_{171}(r,z)=- \frac{16 \sqrt{3} i}{3}
\label{eq:cat-rational-171}
\end{equation}
\begin{equation}
R_{172}(r,z)=\frac{4}{3}
\label{eq:cat-rational-172}
\end{equation}
\begin{equation}
R_{173}(r,z)=- 8 \sqrt{3}
\label{eq:cat-rational-173}
\end{equation}
\begin{equation}
R_{174}(r,z)=- 8 i
\label{eq:cat-rational-174}
\end{equation}
\begin{equation}
R_{175}(r,z)=4 i + 4 \sqrt{3} i
\label{eq:cat-rational-175}
\end{equation}
\begin{equation}
R_{176}(r,z)=-16
\label{eq:cat-rational-176}
\end{equation}
\begin{equation}
R_{177}(r,z)=- 4 \sqrt{3} i - 4 i
\label{eq:cat-rational-177}
\end{equation}
\begin{equation}
R_{178}(r,z)=- 16 \sqrt{3}
\label{eq:cat-rational-178}
\end{equation}
\begin{equation}
R_{179}(r,z)=-4
\label{eq:cat-rational-179}
\end{equation}
\begin{equation}
R_{180}(r,z)=16 \sqrt{3}
\label{eq:cat-rational-180}
\end{equation}
\begin{equation}
R_{181}(r,z)=4
\label{eq:cat-rational-181}
\end{equation}
\begin{equation}
R_{182}(r,z)=- \frac{4 \left(3 r^{2} z^{2} + 7 r^{2} - 3 r z^{3} - 17 r z + 3 z^{2} + 7\right)}{r^{2} - 2 r z + 1}
\label{eq:cat-rational-182}
\end{equation}
\begin{equation}
R_{183}(r,z)=\frac{N_{183}(r,z)}{D_{183}(r,z)}.
\label{eq:cat-rational-183}
\end{equation}
\begin{equation}
N_{183}(r,z)=\shortstack[l]{$\displaystyle  - 20 \sqrt{3} i r^{2} - 20 \sqrt{3} i r^{4} - 6 \sqrt{3} i z$\\
$\displaystyle  + 20 \sqrt{3} i r + 20 \sqrt{3} i r^{3} - 52 \sqrt{3} i r^{2} z$\\
$\displaystyle  - 18 \sqrt{3} i r^{4} z^{2} - 12 \sqrt{3} i r^{2} z^{2}$\\
$\displaystyle  - 12 \sqrt{3} i r^{2} z^{3} + 6 \sqrt{3} i r^{5} z$\\
$\displaystyle  + 12 \sqrt{3} i r^{3} z^{2} + 12 \sqrt{3} i r^{3} z^{3}$\\
$\displaystyle  + 18 \sqrt{3} i r z^{2} + 52 \sqrt{3} i r^{3} z$}
\label{eq:cat-numerator-183}
\end{equation}
\begin{equation}
D_{183}(r,z)=3 r \left(- r^{2} + 2 r z - 1\right)
\label{eq:cat-denominator-183}
\end{equation}
\begin{equation}
R_{184}(r,z)=2 r^{2} z + \frac{2 r^{2}}{3} - 2 r z^{2} - \frac{2 r z}{3} - 4 r + 2 z + \frac{2}{3}
\label{eq:cat-rational-184}
\end{equation}
\begin{equation}
R_{185}(r,z)=\frac{2 \sqrt{3} i r \left(r - 1\right) \left(z + 1\right)}{9}
\label{eq:cat-rational-185}
\end{equation}
\begin{equation}
R_{186}(r,z)=- \frac{4 \sqrt{3} \left(z - 1\right) \left(z + 1\right) \left(2 r^{2} - 3 r z + 2\right)}{r^{2} - 2 r z + 1}
\label{eq:cat-rational-186}
\end{equation}
\begin{equation}
R_{187}(r,z)=\frac{2 i \left(r - 1\right) \left(z - 1\right) \left(z + 1\right) \left(3 r^{2} - 6 r z - r + 3\right)}{r^{2} - 2 r z + 1}
\label{eq:cat-rational-187}
\end{equation}
\begin{equation}
R_{188}(r,z)=- \frac{2 \sqrt{3} r \left(z - 1\right) \left(z + 1\right) \left(3 r^{2} - 6 r z - 2 r + 3\right)}{3 \left(r^{2} - 2 r z + 1\right)}
\label{eq:cat-rational-188}
\end{equation}
\begin{equation}
R_{189}(r,z)=\frac{2 \left(3 r^{4} z - 3 r^{3} z^{2} + 7 r^{3} - 14 r^{2} z - 3 r z^{2} + 7 r + 3 z\right)}{r \left(r^{2} - 2 r z + 1\right)}
\label{eq:cat-rational-189}
\end{equation}
\begin{equation}
R_{190}(r,z)=\frac{2 \sqrt{3} i \left(r - 1\right) \left(r^{2} z - r^{2} - r z^{2} + 3 r z + z - 1\right)}{r^{2} - 2 r z + 1}
\label{eq:cat-rational-190}
\end{equation}
\begin{equation}
R_{191}(r,z)=- \frac{2 r \left(z - 1\right)}{3}
\label{eq:cat-rational-191}
\end{equation}
\begin{equation}
R_{192}(r,z)=\frac{2 \sqrt{3} \left(r^{2} + 1\right) \left(z - 1\right) \left(z + 1\right)}{r^{2} - 2 r z + 1}
\label{eq:cat-rational-192}
\end{equation}
\begin{equation}
R_{193}(r,z)=\frac{2 i r \left(r - 1\right) \left(z - 1\right) \left(z + 1\right)}{r^{2} - 2 r z + 1}
\label{eq:cat-rational-193}
\end{equation}
\begin{equation}
R_{194}(r,z)=- \frac{4 \sqrt{3} \left(z - 1\right) \left(z + 1\right) \left(2 r^{2} - r z + 1\right)}{r^{2} - 2 r z + 1}
\label{eq:cat-rational-194}
\end{equation}
\begin{equation}
R_{195}(r,z)=\frac{N_{195}(r,z)}{D_{195}(r,z)}.
\label{eq:cat-rational-195}
\end{equation}
\begin{equation}
N_{195}(r,z)=\shortstack[l]{$\displaystyle  - 2 i - 8 i r^{2} + 2 i z^{2} - 4 i r z^{3} + 2 \sqrt{3} i r^{3}$\\
$\displaystyle  + 4 i r z + 4 \sqrt{3} i r + 8 i r^{2} z^{2} - 4 \sqrt{3} i r^{2} z$\\
$\displaystyle  - 4 \sqrt{3} i r z^{2} - 2 \sqrt{3} i r^{3} z^{2}$\\
$\displaystyle  + 4 \sqrt{3} i r^{2} z^{3}$}
\label{eq:cat-numerator-195}
\end{equation}
\begin{equation}
D_{195}(r,z)=- r^{2} + 2 r z - 1
\label{eq:cat-denominator-195}
\end{equation}
\begin{equation}
R_{196}(r,z)=- 2 r \left(z - 1\right) \left(z + 1\right) + \frac{4 \sqrt{3} \left(z - 1\right) \left(z + 1\right)}{3}
\label{eq:cat-rational-196}
\end{equation}
\begin{equation}
R_{197}(r,z)=- \frac{8 \left(2 r^{2} - 4 r z + z^{2} + 1\right)}{r^{2} - 2 r z + 1}
\label{eq:cat-rational-197}
\end{equation}
\begin{equation}
R_{198}(r,z)=- \frac{2 i r \left(3 r z - 4\right) \left(2 r^{2} - 4 r z + z^{2} + 1\right)}{r^{2} - 2 r z + 1} - \frac{2 \sqrt{3} i \left(2 r^{2} - 4 r z + z^{2} + 1\right)}{r^{2} - 2 r z + 1}
\label{eq:cat-rational-198}
\end{equation}
\begin{equation}
R_{199}(r,z)=\frac{N_{199}(r,z)}{D_{199}(r,z)}.
\label{eq:cat-rational-199}
\end{equation}
\begin{equation}
N_{199}(r,z)=\shortstack[l]{$\displaystyle 2 r^{2} + 4 r^{4} - 12 r^{3} z - 2 r z - 2 r z^{3} + 2 \sqrt{3} r$\\
$\displaystyle  + 4 \sqrt{3} r^{3} + 10 r^{2} z^{2} - 10 \sqrt{3} r^{2} z$\\
$\displaystyle  - 4 \sqrt{3} r^{4} z - 2 \sqrt{3} r^{2} z^{3} + 2 \sqrt{3} r z^{2}$\\
$\displaystyle  + 8 \sqrt{3} r^{3} z^{2}$}
\label{eq:cat-numerator-199}
\end{equation}
\begin{equation}
D_{199}(r,z)=- r^{2} + 2 r z - 1
\label{eq:cat-denominator-199}
\end{equation}
\begin{equation}
R_{200}(r,z)=\frac{2 \sqrt{3} \left(3 r^{2} + 1\right) \left(z - 1\right) \left(z + 1\right)}{r^{2} - 2 r z + 1}
\label{eq:cat-rational-200}
\end{equation}
\begin{equation}
R_{201}(r,z)=\frac{6 i r^{2} \left(z - 1\right) \left(z + 1\right)}{r^{2} - 2 r z + 1} - \frac{2 \sqrt{3} i r \left(z - 1\right) \left(z + 1\right)}{r^{2} - 2 r z + 1}
\label{eq:cat-rational-201}
\end{equation}
\begin{equation}
R_{202}(r,z)=\frac{2 \left(3 r z + 1\right) \left(2 r^{2} - 4 r z + z^{2} + 1\right)}{r^{2} - 2 r z + 1}
\label{eq:cat-rational-202}
\end{equation}
\begin{equation}
R_{203}(r,z)=\frac{2 \sqrt{3} i r z \left(2 r^{2} - 4 r z + z^{2} + 1\right)}{r^{2} - 2 r z + 1} - \frac{2 i r \left(2 r^{2} - 4 r z + z^{2} + 1\right)}{r^{2} - 2 r z + 1}
\label{eq:cat-rational-203}
\end{equation}
\begin{equation}
R_{204}(r,z)=- 16 r + 16 z
\label{eq:cat-rational-204}
\end{equation}
\begin{equation}
R_{205}(r,z)=- 4 i r \left(r - z\right) \left(3 r z - 4\right) - 4 \sqrt{3} i \left(r - z\right)
\label{eq:cat-rational-205}
\end{equation}
\begin{equation}
R_{206}(r,z)=- 4 r \left(r - z\right)^{2} + 4 \sqrt{3} r \left(r - z\right) \left(r z - 1\right)
\label{eq:cat-rational-206}
\end{equation}
\begin{equation}
R_{207}(r,z)=4 \left(r - z\right) \left(3 r z + 1\right)
\label{eq:cat-rational-207}
\end{equation}
\begin{equation}
R_{208}(r,z)=4 \sqrt{3} i r z \left(r - z\right) - 4 i r \left(r - z\right)
\label{eq:cat-rational-208}
\end{equation}
\begin{equation}
R_{209}(r,z)=- \frac{4 \sqrt{3} \left(z - 1\right) \left(z + 1\right) \left(r^{2} - r z + 2\right)}{r^{2} - 2 r z + 1}
\label{eq:cat-rational-209}
\end{equation}
\begin{equation}
R_{210}(r,z)=\frac{N_{210}(r,z)}{D_{210}(r,z)}.
\label{eq:cat-rational-210}
\end{equation}
\begin{equation}
N_{210}(r,z)=\shortstack[l]{$\displaystyle  - 2 \sqrt{3} i + 2 i r^{3} + 8 i r - 8 i r z^{2} - 4 i r^{2} z$\\
$\displaystyle  - 4 \sqrt{3} i r^{2} - 2 i r^{3} z^{2} + 2 \sqrt{3} i z^{2}$\\
$\displaystyle  + 4 i r^{2} z^{3} - 4 \sqrt{3} i r z^{3} + 4 \sqrt{3} i r z$\\
$\displaystyle  + 4 \sqrt{3} i r^{2} z^{2}$}
\label{eq:cat-numerator-210}
\end{equation}
\begin{equation}
D_{210}(r,z)=- r^{2} + 2 r z - 1
\label{eq:cat-denominator-210}
\end{equation}
\begin{equation}
R_{211}(r,z)=\frac{4 \sqrt{3} r^{2} \left(z - 1\right) \left(z + 1\right)}{3} - 2 r \left(z - 1\right) \left(z + 1\right)
\label{eq:cat-rational-211}
\end{equation}
\begin{equation}
R_{212}(r,z)=- \frac{8 \left(r^{2} z^{2} + r^{2} - 4 r z + 2\right)}{r^{2} - 2 r z + 1}
\label{eq:cat-rational-212}
\end{equation}
\begin{equation}
R_{213}(r,z)=\frac{N_{213}(r,z)}{D_{213}(r,z)}.
\label{eq:cat-rational-213}
\end{equation}
\begin{equation}
N_{213}(r,z)=\shortstack[l]{$\displaystyle  - 12 i z + 8 i r^{3} + 16 i r - 38 i r^{2} z$\\
$\displaystyle  - 6 i r^{2} z^{3} - 4 \sqrt{3} i r^{2} - 2 \sqrt{3} i r^{4}$\\
$\displaystyle  + 8 i r^{3} z^{2} + 24 i r z^{2} - 2 \sqrt{3} i r^{4} z^{2}$\\
$\displaystyle  + 8 \sqrt{3} i r^{3} z$}
\label{eq:cat-numerator-213}
\end{equation}
\begin{equation}
D_{213}(r,z)=r \left(- r^{2} + 2 r z - 1\right)
\label{eq:cat-denominator-213}
\end{equation}
\begin{equation}
R_{214}(r,z)=\frac{N_{214}(r,z)}{D_{214}(r,z)}.
\label{eq:cat-rational-214}
\end{equation}
\begin{equation}
N_{214}(r,z)=\shortstack[l]{$\displaystyle 4 + 2 r^{2} - 12 r z - 4 \sqrt{3} z - 2 r^{3} z - 2 r^{3} z^{3}$\\
$\displaystyle  + 2 \sqrt{3} r^{3} + 4 \sqrt{3} r + 10 r^{2} z^{2} - 10 \sqrt{3} r^{2} z$\\
$\displaystyle  - 2 \sqrt{3} r^{2} z^{3} + 2 \sqrt{3} r^{3} z^{2} + 8 \sqrt{3} r z^{2}$}
\label{eq:cat-numerator-214}
\end{equation}
\begin{equation}
D_{214}(r,z)=- r^{2} + 2 r z - 1
\label{eq:cat-denominator-214}
\end{equation}
\begin{equation}
R_{215}(r,z)=\frac{2 \sqrt{3} \left(r^{2} + 3\right) \left(z - 1\right) \left(z + 1\right)}{r^{2} - 2 r z + 1}
\label{eq:cat-rational-215}
\end{equation}
\begin{equation}
R_{216}(r,z)=\frac{2 \sqrt{3} i r^{2} \left(z - 1\right) \left(z + 1\right)}{r^{2} - 2 r z + 1} - \frac{6 i r \left(z - 1\right) \left(z + 1\right)}{r^{2} - 2 r z + 1}
\label{eq:cat-rational-216}
\end{equation}
\begin{equation}
R_{217}(r,z)=\frac{2 \left(r + 3 z\right) \left(r^{2} z^{2} + r^{2} - 4 r z + 2\right)}{r \left(r^{2} - 2 r z + 1\right)}
\label{eq:cat-rational-217}
\end{equation}
\begin{equation}
R_{218}(r,z)=- \frac{2 \sqrt{3} i z \left(r^{2} z^{2} + r^{2} - 4 r z + 2\right)}{r^{2} - 2 r z + 1} + \frac{2 i \left(r^{2} z^{2} + r^{2} - 4 r z + 2\right)}{r^{2} - 2 r z + 1}
\label{eq:cat-rational-218}
\end{equation}
\begin{equation}
R_{219}(r,z)=- 8 z^{2} - 8
\label{eq:cat-rational-219}
\end{equation}
\begin{equation}
R_{220}(r,z)=\frac{2 i \left(r - 1\right) \left(z^{2} + 1\right) \left(r^{2} - 5 r z - 3 r + 1\right)}{r^{2} - 2 r z + 1}
\label{eq:cat-rational-220}
\end{equation}
\begin{equation}
R_{221}(r,z)=- 2 r \left(z + 1\right) \left(z^{2} + 1\right)
\label{eq:cat-rational-221}
\end{equation}
\begin{equation}
R_{222}(r,z)=\frac{4 \sqrt{3} \left(z - 1\right) \left(z + 1\right) \left(2 r^{4} - 4 r^{3} z + r^{2} z^{2} + 5 r^{2} - 4 r z + 2\right)}{r^{2} - 2 r z + 1}
\label{eq:cat-rational-222}
\end{equation}
\begin{equation}
R_{223}(r,z)=6 z^{2} + 6
\label{eq:cat-rational-223}
\end{equation}
\begin{equation}
R_{224}(r,z)=\frac{6 i r \left(r - 1\right) \left(z + 1\right) \left(z^{2} + 1\right)}{r^{2} - 2 r z + 1}
\label{eq:cat-rational-224}
\end{equation}
\begin{equation}
R_{225}(r,z)=8 \sqrt{3} r \left(z - 1\right) \left(z + 1\right) \left(r^{2} - r z + 2\right)
\label{eq:cat-rational-225}
\end{equation}
\begin{equation}
R_{226}(r,z)=16 r z - 16
\label{eq:cat-rational-226}
\end{equation}
\begin{equation}
R_{227}(r,z)=- 4 \sqrt{3} i r \left(r z - 1\right) + \frac{4 i \left(4 r - 3 z\right) \left(r z - 1\right)}{r}
\label{eq:cat-rational-227}
\end{equation}
\begin{equation}
R_{228}(r,z)=4 \sqrt{3} \left(r - z\right) \left(r z - 1\right) - 4 \left(r z - 1\right)^{2}
\label{eq:cat-rational-228}
\end{equation}
\begin{equation}
R_{229}(r,z)=- \frac{4 \left(r + 3 z\right) \left(r z - 1\right)}{r}
\label{eq:cat-rational-229}
\end{equation}
\begin{equation}
R_{230}(r,z)=4 \sqrt{3} i z \left(r z - 1\right) - 4 i \left(r z - 1\right)
\label{eq:cat-rational-230}
\end{equation}
\begin{equation}
R_{231}(r,z)=8 \sqrt{3} \left(z - 1\right) \left(z + 1\right) \left(2 r^{2} - r z + 1\right)
\label{eq:cat-rational-231}
\end{equation}
\begin{equation}
R_{232}(r,z)=16 z
\label{eq:cat-rational-232}
\end{equation}
\begin{equation}
R_{233}(r,z)=- \frac{4 i z \left(r - 1\right) \left(r^{2} - 5 r z - 3 r + 1\right)}{r^{2} - 2 r z + 1}
\label{eq:cat-rational-233}
\end{equation}
\begin{equation}
R_{234}(r,z)=4 r z \left(z + 1\right)
\label{eq:cat-rational-234}
\end{equation}
\begin{equation}
R_{235}(r,z)=- \frac{8 \sqrt{3} r \left(z - 1\right) \left(z + 1\right) \left(2 r^{2} - 3 r z + 2\right)}{r^{2} - 2 r z + 1}
\label{eq:cat-rational-235}
\end{equation}
\begin{equation}
R_{236}(r,z)=- 12 z
\label{eq:cat-rational-236}
\end{equation}
\begin{equation}
R_{237}(r,z)=- \frac{12 i r z \left(r - 1\right) \left(z + 1\right)}{r^{2} - 2 r z + 1}
\label{eq:cat-rational-237}
\end{equation}
\begin{equation}
R_{238}(r,z)=- \frac{4 \left(3 r^{2} z^{2} + 7 r^{2} + 3 r z^{3} + 17 r z + 3 z^{2} + 7\right)}{r^{2} + 2 r z + 1}
\label{eq:cat-rational-238}
\end{equation}
\begin{equation}
R_{239}(r,z)=\frac{N_{239}(r,z)}{D_{239}(r,z)}.
\label{eq:cat-rational-239}
\end{equation}
\begin{equation}
N_{239}(r,z)=\shortstack[l]{$\displaystyle  - 20 \sqrt{3} i r - 20 \sqrt{3} i r^{3} - 6 \sqrt{3} i z$\\
$\displaystyle  + 20 \sqrt{3} i r^{2} + 20 \sqrt{3} i r^{4} - 52 \sqrt{3} i r^{2} z$\\
$\displaystyle  - 18 \sqrt{3} i r z^{2} - 12 \sqrt{3} i r^{3} z^{2}$\\
$\displaystyle  - 12 \sqrt{3} i r^{2} z^{3} + 6 \sqrt{3} i r^{5} z$\\
$\displaystyle  + 12 \sqrt{3} i r^{2} z^{2} + 12 \sqrt{3} i r^{3} z^{3}$\\
$\displaystyle  + 18 \sqrt{3} i r^{4} z^{2} + 52 \sqrt{3} i r^{3} z$}
\label{eq:cat-numerator-239}
\end{equation}
\begin{equation}
D_{239}(r,z)=3 r \left(r^{2} + 2 r z + 1\right)
\label{eq:cat-denominator-239}
\end{equation}
\begin{equation}
R_{240}(r,z)=- 2 r^{2} z + \frac{2 r^{2}}{3} - 2 r z^{2} + \frac{2 r z}{3} - 4 r - 2 z + \frac{2}{3}
\label{eq:cat-rational-240}
\end{equation}
\begin{equation}
R_{241}(r,z)=- \frac{2 \sqrt{3} i r \left(r - 1\right) \left(z - 1\right)}{9}
\label{eq:cat-rational-241}
\end{equation}
\begin{equation}
R_{242}(r,z)=- \frac{4 \sqrt{3} \left(z - 1\right) \left(z + 1\right) \left(2 r^{2} + 3 r z + 2\right)}{r^{2} + 2 r z + 1}
\label{eq:cat-rational-242}
\end{equation}
\begin{equation}
R_{243}(r,z)=\frac{2 i \left(r - 1\right) \left(z - 1\right) \left(z + 1\right) \left(3 r^{2} + 6 r z - r + 3\right)}{r^{2} + 2 r z + 1}
\label{eq:cat-rational-243}
\end{equation}
\begin{equation}
R_{244}(r,z)=- \frac{2 \sqrt{3} r \left(z - 1\right) \left(z + 1\right) \left(3 r^{2} + 6 r z - 2 r + 3\right)}{3 \left(r^{2} + 2 r z + 1\right)}
\label{eq:cat-rational-244}
\end{equation}
\begin{equation}
R_{245}(r,z)=- \frac{2 \left(3 r^{4} z + 3 r^{3} z^{2} - 7 r^{3} - 14 r^{2} z + 3 r z^{2} - 7 r + 3 z\right)}{r \left(r^{2} + 2 r z + 1\right)}
\label{eq:cat-rational-245}
\end{equation}
\begin{equation}
R_{246}(r,z)=- \frac{2 \sqrt{3} i \left(r - 1\right) \left(r^{2} z + r^{2} + r z^{2} + 3 r z + z + 1\right)}{r^{2} + 2 r z + 1}
\label{eq:cat-rational-246}
\end{equation}
\begin{equation}
R_{247}(r,z)=\frac{2 r \left(z + 1\right)}{3}
\label{eq:cat-rational-247}
\end{equation}
\begin{equation}
R_{248}(r,z)=\frac{2 \sqrt{3} \left(r^{2} + 1\right) \left(z - 1\right) \left(z + 1\right)}{r^{2} + 2 r z + 1}
\label{eq:cat-rational-248}
\end{equation}
\begin{equation}
R_{249}(r,z)=\frac{2 i r \left(r - 1\right) \left(z - 1\right) \left(z + 1\right)}{r^{2} + 2 r z + 1}
\label{eq:cat-rational-249}
\end{equation}
\begin{equation}
R_{250}(r,z)=- \frac{4 \sqrt{3} \left(z - 1\right) \left(z + 1\right) \left(2 r^{2} + r z + 1\right)}{r^{2} + 2 r z + 1}
\label{eq:cat-rational-250}
\end{equation}
\begin{equation}
R_{251}(r,z)=\frac{N_{251}(r,z)}{D_{251}(r,z)}.
\label{eq:cat-rational-251}
\end{equation}
\begin{equation}
N_{251}(r,z)=\shortstack[l]{$\displaystyle 2 i - 2 i z^{2} + 8 i r^{2} - 8 i r^{2} z^{2} - 4 \sqrt{3} i r$\\
$\displaystyle  - 4 i r z^{3} - 2 \sqrt{3} i r^{3} + 4 i r z - 4 \sqrt{3} i r^{2} z$\\
$\displaystyle  + 2 \sqrt{3} i r^{3} z^{2} + 4 \sqrt{3} i r z^{2}$\\
$\displaystyle  + 4 \sqrt{3} i r^{2} z^{3}$}
\label{eq:cat-numerator-251}
\end{equation}
\begin{equation}
D_{251}(r,z)=r^{2} + 2 r z + 1
\label{eq:cat-denominator-251}
\end{equation}
\begin{equation}
R_{252}(r,z)=- \frac{8 \left(2 r^{2} + 4 r z + z^{2} + 1\right)}{r^{2} + 2 r z + 1}
\label{eq:cat-rational-252}
\end{equation}
\begin{equation}
R_{253}(r,z)=\frac{2 i r \left(3 r z + 4\right) \left(2 r^{2} + 4 r z + z^{2} + 1\right)}{r^{2} + 2 r z + 1} - \frac{2 \sqrt{3} i \left(2 r^{2} + 4 r z + z^{2} + 1\right)}{r^{2} + 2 r z + 1}
\label{eq:cat-rational-253}
\end{equation}
\begin{equation}
R_{254}(r,z)=\frac{N_{254}(r,z)}{D_{254}(r,z)}.
\label{eq:cat-rational-254}
\end{equation}
\begin{equation}
N_{254}(r,z)=\shortstack[l]{$\displaystyle  - 4 r^{4} - 2 r^{2} - 12 r^{3} z - 10 r^{2} z^{2}$\\
$\displaystyle  - 4 \sqrt{3} r^{3} - 2 r z - 2 \sqrt{3} r - 2 r z^{3}$\\
$\displaystyle  - 10 \sqrt{3} r^{2} z - 8 \sqrt{3} r^{3} z^{2} - 4 \sqrt{3} r^{4} z$\\
$\displaystyle  - 2 \sqrt{3} r z^{2} - 2 \sqrt{3} r^{2} z^{3}$}
\label{eq:cat-numerator-254}
\end{equation}
\begin{equation}
D_{254}(r,z)=r^{2} + 2 r z + 1
\label{eq:cat-denominator-254}
\end{equation}
\begin{equation}
R_{255}(r,z)=\frac{2 \sqrt{3} \left(3 r^{2} + 1\right) \left(z - 1\right) \left(z + 1\right)}{r^{2} + 2 r z + 1}
\label{eq:cat-rational-255}
\end{equation}
\begin{equation}
R_{256}(r,z)=\frac{6 i r^{2} \left(z - 1\right) \left(z + 1\right)}{r^{2} + 2 r z + 1} - \frac{2 \sqrt{3} i r \left(z - 1\right) \left(z + 1\right)}{r^{2} + 2 r z + 1}
\label{eq:cat-rational-256}
\end{equation}
\begin{equation}
R_{257}(r,z)=- \frac{2 \left(3 r z - 1\right) \left(2 r^{2} + 4 r z + z^{2} + 1\right)}{r^{2} + 2 r z + 1}
\label{eq:cat-rational-257}
\end{equation}
\begin{equation}
R_{258}(r,z)=- \frac{2 \sqrt{3} i r z \left(2 r^{2} + 4 r z + z^{2} + 1\right)}{r^{2} + 2 r z + 1} - \frac{2 i r \left(2 r^{2} + 4 r z + z^{2} + 1\right)}{r^{2} + 2 r z + 1}
\label{eq:cat-rational-258}
\end{equation}
\begin{equation}
R_{259}(r,z)=- 16 r - 16 z
\label{eq:cat-rational-259}
\end{equation}
\begin{equation}
R_{260}(r,z)=4 i r \left(r + z\right) \left(3 r z + 4\right) - 4 \sqrt{3} i \left(r + z\right)
\label{eq:cat-rational-260}
\end{equation}
\begin{equation}
R_{261}(r,z)=- 4 r \left(r + z\right)^{2} - 4 \sqrt{3} r \left(r + z\right) \left(r z + 1\right)
\label{eq:cat-rational-261}
\end{equation}
\begin{equation}
R_{262}(r,z)=- 4 \left(r + z\right) \left(3 r z - 1\right)
\label{eq:cat-rational-262}
\end{equation}
\begin{equation}
R_{263}(r,z)=- 4 \sqrt{3} i r z \left(r + z\right) - 4 i r \left(r + z\right)
\label{eq:cat-rational-263}
\end{equation}
\begin{equation}
R_{264}(r,z)=- \frac{4 \sqrt{3} \left(z - 1\right) \left(z + 1\right) \left(r^{2} + r z + 2\right)}{r^{2} + 2 r z + 1}
\label{eq:cat-rational-264}
\end{equation}
\begin{equation}
R_{265}(r,z)=\frac{N_{265}(r,z)}{D_{265}(r,z)}.
\label{eq:cat-rational-265}
\end{equation}
\begin{equation}
N_{265}(r,z)=\shortstack[l]{$\displaystyle  - 8 i r - 2 i r^{3} + 2 \sqrt{3} i - 4 i r^{2} z - 2 \sqrt{3} i z^{2}$\\
$\displaystyle  + 2 i r^{3} z^{2} + 4 \sqrt{3} i r^{2} + 4 i r^{2} z^{3}$\\
$\displaystyle  + 8 i r z^{2} - 4 \sqrt{3} i r z^{3} - 4 \sqrt{3} i r^{2} z^{2}$\\
$\displaystyle  + 4 \sqrt{3} i r z$}
\label{eq:cat-numerator-265}
\end{equation}
\begin{equation}
D_{265}(r,z)=r^{2} + 2 r z + 1
\label{eq:cat-denominator-265}
\end{equation}
\begin{equation}
R_{266}(r,z)=- \frac{8 \left(r^{2} z^{2} + r^{2} + 4 r z + 2\right)}{r^{2} + 2 r z + 1}
\label{eq:cat-rational-266}
\end{equation}
\begin{equation}
R_{267}(r,z)=\frac{N_{267}(r,z)}{D_{267}(r,z)}.
\label{eq:cat-rational-267}
\end{equation}
\begin{equation}
N_{267}(r,z)=\shortstack[l]{$\displaystyle  - 16 i r - 12 i z - 8 i r^{3} - 38 i r^{2} z - 24 i r z^{2}$\\
$\displaystyle  - 8 i r^{3} z^{2} - 6 i r^{2} z^{3} + 2 \sqrt{3} i r^{4}$\\
$\displaystyle  + 4 \sqrt{3} i r^{2} + 2 \sqrt{3} i r^{4} z^{2} + 8 \sqrt{3} i r^{3} z$}
\label{eq:cat-numerator-267}
\end{equation}
\begin{equation}
D_{267}(r,z)=r \left(r^{2} + 2 r z + 1\right)
\label{eq:cat-denominator-267}
\end{equation}
\begin{equation}
R_{268}(r,z)=\frac{N_{268}(r,z)}{D_{268}(r,z)}.
\label{eq:cat-rational-268}
\end{equation}
\begin{equation}
N_{268}(r,z)=\shortstack[l]{$\displaystyle  - 4 - 2 r^{2} - 12 r z - 10 r^{2} z^{2} - 4 \sqrt{3} z - 4 \sqrt{3} r$\\
$\displaystyle  - 2 r^{3} z - 2 \sqrt{3} r^{3} - 2 r^{3} z^{3} - 10 \sqrt{3} r^{2} z$\\
$\displaystyle  - 8 \sqrt{3} r z^{2} - 2 \sqrt{3} r^{3} z^{2} - 2 \sqrt{3} r^{2} z^{3}$}
\label{eq:cat-numerator-268}
\end{equation}
\begin{equation}
D_{268}(r,z)=r^{2} + 2 r z + 1
\label{eq:cat-denominator-268}
\end{equation}
\begin{equation}
R_{269}(r,z)=\frac{2 \sqrt{3} \left(r^{2} + 3\right) \left(z - 1\right) \left(z + 1\right)}{r^{2} + 2 r z + 1}
\label{eq:cat-rational-269}
\end{equation}
\begin{equation}
R_{270}(r,z)=\frac{2 \sqrt{3} i r^{2} \left(z - 1\right) \left(z + 1\right)}{r^{2} + 2 r z + 1} - \frac{6 i r \left(z - 1\right) \left(z + 1\right)}{r^{2} + 2 r z + 1}
\label{eq:cat-rational-270}
\end{equation}
\begin{equation}
R_{271}(r,z)=\frac{2 \left(r - 3 z\right) \left(r^{2} z^{2} + r^{2} + 4 r z + 2\right)}{r \left(r^{2} + 2 r z + 1\right)}
\label{eq:cat-rational-271}
\end{equation}
\begin{equation}
R_{272}(r,z)=\frac{2 \sqrt{3} i z \left(r^{2} z^{2} + r^{2} + 4 r z + 2\right)}{r^{2} + 2 r z + 1} + \frac{2 i \left(r^{2} z^{2} + r^{2} + 4 r z + 2\right)}{r^{2} + 2 r z + 1}
\label{eq:cat-rational-272}
\end{equation}
\begin{equation}
R_{273}(r,z)=\frac{2 i \left(r - 1\right) \left(z^{2} + 1\right) \left(r^{2} + 5 r z - 3 r + 1\right)}{r^{2} + 2 r z + 1}
\label{eq:cat-rational-273}
\end{equation}
\begin{equation}
R_{274}(r,z)=2 r \left(z - 1\right) \left(z^{2} + 1\right)
\label{eq:cat-rational-274}
\end{equation}
\begin{equation}
R_{275}(r,z)=\frac{4 \sqrt{3} \left(z - 1\right) \left(z + 1\right) \left(2 r^{4} + 4 r^{3} z + r^{2} z^{2} + 5 r^{2} + 4 r z + 2\right)}{r^{2} + 2 r z + 1}
\label{eq:cat-rational-275}
\end{equation}
\begin{equation}
R_{276}(r,z)=- \frac{6 i r \left(r - 1\right) \left(z - 1\right) \left(z^{2} + 1\right)}{r^{2} + 2 r z + 1}
\label{eq:cat-rational-276}
\end{equation}
\begin{equation}
R_{277}(r,z)=8 \sqrt{3} r \left(z - 1\right) \left(z + 1\right) \left(r^{2} + r z + 2\right)
\label{eq:cat-rational-277}
\end{equation}
\begin{equation}
R_{278}(r,z)=- 16 r z - 16
\label{eq:cat-rational-278}
\end{equation}
\begin{equation}
R_{279}(r,z)=4 \sqrt{3} i r \left(r z + 1\right) - \frac{4 i \left(4 r + 3 z\right) \left(r z + 1\right)}{r}
\label{eq:cat-rational-279}
\end{equation}
\begin{equation}
R_{280}(r,z)=- 4 \sqrt{3} \left(r + z\right) \left(r z + 1\right) - 4 \left(r z + 1\right)^{2}
\label{eq:cat-rational-280}
\end{equation}
\begin{equation}
R_{281}(r,z)=\frac{4 \left(r - 3 z\right) \left(r z + 1\right)}{r}
\label{eq:cat-rational-281}
\end{equation}
\begin{equation}
R_{282}(r,z)=4 \sqrt{3} i z \left(r z + 1\right) + 4 i \left(r z + 1\right)
\label{eq:cat-rational-282}
\end{equation}
\begin{equation}
R_{283}(r,z)=8 \sqrt{3} \left(z - 1\right) \left(z + 1\right) \left(2 r^{2} + r z + 1\right)
\label{eq:cat-rational-283}
\end{equation}
\begin{equation}
R_{284}(r,z)=- 16 z
\label{eq:cat-rational-284}
\end{equation}
\begin{equation}
R_{285}(r,z)=\frac{4 i z \left(r - 1\right) \left(r^{2} + 5 r z - 3 r + 1\right)}{r^{2} + 2 r z + 1}
\label{eq:cat-rational-285}
\end{equation}
\begin{equation}
R_{286}(r,z)=4 r z \left(z - 1\right)
\label{eq:cat-rational-286}
\end{equation}
\begin{equation}
R_{287}(r,z)=- \frac{8 \sqrt{3} r \left(z - 1\right) \left(z + 1\right) \left(2 r^{2} + 3 r z + 2\right)}{r^{2} + 2 r z + 1}
\label{eq:cat-rational-287}
\end{equation}
\begin{equation}
R_{288}(r,z)=12 z
\label{eq:cat-rational-288}
\end{equation}
\begin{equation}
R_{289}(r,z)=- \frac{12 i r z \left(r - 1\right) \left(z - 1\right)}{r^{2} + 2 r z + 1}
\label{eq:cat-rational-289}
\end{equation}
\begin{equation}
R_{290}(r,z)=4 \sqrt{3} \left(z - 1\right) \left(z + 1\right)
\label{eq:cat-rational-290}
\end{equation}
\begin{equation}
R_{291}(r,z)=- 4 i r \left(z - 1\right) \left(z + 1\right)
\label{eq:cat-rational-291}
\end{equation}
\begin{equation}
R_{292}(r,z)=- \frac{4 \sqrt{3} r^{2} \left(z - 1\right) \left(z + 1\right)}{3}
\label{eq:cat-rational-292}
\end{equation}
\begin{equation}
R_{293}(r,z)=4 \sqrt{3} \left(3 z^{2} - 1\right)
\label{eq:cat-rational-293}
\end{equation}
\begin{equation}
R_{294}(r,z)=- 4 i r \left(3 z^{2} - 1\right)
\label{eq:cat-rational-294}
\end{equation}
\begin{equation}
R_{295}(r,z)=- \frac{4 \sqrt{3} r^{2} \left(3 z^{2} - 1\right)}{3}
\label{eq:cat-rational-295}
\end{equation}
\begin{equation}
R_{296}(r,z)=- 8 \sqrt{3} z \sqrt{- z^{2} + 1}
\label{eq:cat-rational-296}
\end{equation}
\begin{equation}
R_{297}(r,z)=8 i r z \sqrt{- z^{2} + 1}
\label{eq:cat-rational-297}
\end{equation}
\begin{equation}
R_{298}(r,z)=\frac{8 \sqrt{3} r^{2} z \sqrt{- z^{2} + 1}}{3}
\label{eq:cat-rational-298}
\end{equation}
\begin{equation}
R_{299}(r,z)=\frac{16 \sqrt{3} i r}{3}
\label{eq:cat-rational-299}
\end{equation}
\begin{equation}
R_{300}(r,z)=\frac{4 r^{2}}{3}
\label{eq:cat-rational-300}
\end{equation}
\begin{equation}
R_{301}(r,z)=- 4 \sqrt{3} \left(z - 1\right) \left(z + 1\right)
\label{eq:cat-rational-301}
\end{equation}
\begin{equation}
R_{302}(r,z)=4 i r \left(z - 1\right) \left(z + 1\right)
\label{eq:cat-rational-302}
\end{equation}
\begin{equation}
R_{303}(r,z)=- 4 \sqrt{3} \left(3 z^{2} - 1\right)
\label{eq:cat-rational-303}
\end{equation}
\begin{equation}
R_{304}(r,z)=4 i r \left(3 z^{2} - 1\right)
\label{eq:cat-rational-304}
\end{equation}
\begin{equation}
R_{305}(r,z)=8 \sqrt{3} z \sqrt{- z^{2} + 1}
\label{eq:cat-rational-305}
\end{equation}
\begin{equation}
R_{306}(r,z)=- 8 i r z \sqrt{- z^{2} + 1}
\label{eq:cat-rational-306}
\end{equation}
\begin{equation}
R_{307}(r,z)=8 z^{2} + 8
\label{eq:cat-rational-307}
\end{equation}
\begin{equation}
R_{308}(r,z)=- 2 \sqrt{3} i r \left(z^{2} + 1\right) - 2 i r \left(z^{2} + 1\right)
\label{eq:cat-rational-308}
\end{equation}
\begin{equation}
R_{309}(r,z)=24 \left(z - 1\right) \left(z + 1\right)
\label{eq:cat-rational-309}
\end{equation}
\begin{equation}
R_{310}(r,z)=- 6 \sqrt{3} i r \left(z - 1\right) \left(z + 1\right) - 6 i r \left(z - 1\right) \left(z + 1\right)
\label{eq:cat-rational-310}
\end{equation}
\begin{equation}
R_{311}(r,z)=- 16 z \sqrt{- z^{2} + 1}
\label{eq:cat-rational-311}
\end{equation}
\begin{equation}
R_{312}(r,z)=\sqrt{- z^{2} + 1} \left(4 i r z + 4 \sqrt{3} i r z\right)
\label{eq:cat-rational-312}
\end{equation}
\begin{equation}
R_{313}(r,z)=2 i r \left(z^{2} + 1\right) + 2 \sqrt{3} i r \left(z^{2} + 1\right)
\label{eq:cat-rational-313}
\end{equation}
\begin{equation}
R_{314}(r,z)=- 24 \left(z - 1\right) \left(z + 1\right)
\label{eq:cat-rational-314}
\end{equation}
\begin{equation}
R_{315}(r,z)=6 i r \left(z - 1\right) \left(z + 1\right) + 6 \sqrt{3} i r \left(z - 1\right) \left(z + 1\right)
\label{eq:cat-rational-315}
\end{equation}
\begin{equation}
R_{316}(r,z)=16 z \sqrt{- z^{2} + 1}
\label{eq:cat-rational-316}
\end{equation}
\begin{equation}
R_{317}(r,z)=\sqrt{- z^{2} + 1} \left(- 4 \sqrt{3} i r z - 4 i r z\right)
\label{eq:cat-rational-317}
\end{equation}
\begin{equation}
R_{318}(r,z)=- 4 \sqrt{3} i r z - 4 i r z
\label{eq:cat-rational-318}
\end{equation}
\begin{equation}
R_{319}(r,z)=- 16 \sqrt{- z^{2} + 1}
\label{eq:cat-rational-319}
\end{equation}
\begin{equation}
R_{320}(r,z)=\sqrt{- z^{2} + 1} \left(4 i r + 4 \sqrt{3} i r\right)
\label{eq:cat-rational-320}
\end{equation}
\begin{equation}
R_{321}(r,z)=4 i r z + 4 \sqrt{3} i r z
\label{eq:cat-rational-321}
\end{equation}
\begin{equation}
R_{322}(r,z)=16 \sqrt{- z^{2} + 1}
\label{eq:cat-rational-322}
\end{equation}
\begin{equation}
R_{323}(r,z)=\sqrt{- z^{2} + 1} \left(- 4 \sqrt{3} i r - 4 i r\right)
\label{eq:cat-rational-323}
\end{equation}
\begin{equation}
R_{324}(r,z)=- 8 \sqrt{3} r^{2} \left(z - 1\right) \left(z + 1\right)
\label{eq:cat-rational-324}
\end{equation}
\begin{equation}
R_{325}(r,z)=- 8 \sqrt{3} r^{2} \left(3 z^{2} - 1\right)
\label{eq:cat-rational-325}
\end{equation}
\begin{equation}
R_{326}(r,z)=16 \sqrt{3} r^{2} z \sqrt{- z^{2} + 1}
\label{eq:cat-rational-326}
\end{equation}
\begin{equation}
R_{327}(r,z)=8 i r
\label{eq:cat-rational-327}
\end{equation}
\begin{equation}
R_{328}(r,z)=8 \sqrt{3} r^{2} \left(z - 1\right) \left(z + 1\right)
\label{eq:cat-rational-328}
\end{equation}
\begin{equation}
R_{329}(r,z)=8 \sqrt{3} r^{2} \left(3 z^{2} - 1\right)
\label{eq:cat-rational-329}
\end{equation}
\begin{equation}
R_{330}(r,z)=- 16 \sqrt{3} r^{2} z \sqrt{- z^{2} + 1}
\label{eq:cat-rational-330}
\end{equation}
\begin{equation}
R_{331}(r,z)=- 8 i r
\label{eq:cat-rational-331}
\end{equation}

\clearpage
\phantomsection
\addcontentsline{toc}{section}{References}
\begin{thebibliography}{1}
\bibitem{YerokhinArticle}
D. Yerokhin, \textit{A minimal local covariant action for holographic dark energy:
constraints, perturbations, and nonlinear dynamics} (2026).
\end{thebibliography}
\end{document}
